\documentclass{book}
\input{preamble}

\begin{document}

\tableofcontents

% --- Begin differential-geometry.tex ---
\chapter{Differential Geometry}
\label{chapter-differential-geometry}




\section{Riemannian metrics}
\label{section-Riemannian-metrics}

\begin{definition}
\label{definition-product-metric}
Let $M_1$ and $M_2$ be Riemannian manifolds.
The {\it product metric} on
$M_1\times M_2$ is $\pi_1^*g_1+\pi_2^*g_2$.
\end{definition}

\begin{definition}
\label{definition-flat-torus}
$S^1\times\ldots\times S^1$ with the product
metric is called the {\it flat
torus}.
\end{definition}

\medskip\noindent
Let $G$ be a Lie group equipped with a
bi-invariant metric.
Then for any $U,V,X \in \mathfrak{g}$,
\begin{equation}
\label{equation-lie-group-bi-invar-metri}
\left\langle{}[U,X],V\right\rangle{}=-\left\langle{}U,[V,X]\right\rangle{}.
\end{equation}

\noindent
This follows from
$$
\begin{aligned}
\left\langle{}U,V\right\rangle{}&=
\left\langle{}dR_{\text{exp}(tX)}\circ
dL_{\text{exp}(tX)^{-1}}U,
dR_{\text{exp}(tX)}\circ
dL_{\text{exp}(tX)^{-1}}V\right\rangle{}\\
&=\left\langle{}dR_{\text{exp}(tX)}U,
dR_{\text{exp}(tX)}V\right\rangle{}
\end{aligned}
$$

\noindent
since both $U$ and $V$ are left invariant.
The key observation is that the map
$R_{\text{exp}(tX)}$
is exactly the flow of $X$.
Thus, differentiating and evaluating at $t=0$
we obtain
$$
0=\left\langle{}[U,X],V\right\rangle{}+\left\langle{}U,[V,X]\right\rangle{},
$$
where the $V$ and $U$ appear since evaluating
$R_{\text{exp}(tX)}$ at $t=0$ is the
identity.

A similar construction shows that the
derivative of the map
$\text{Ad}:G \to \text{End}(\mathfrak{g})$
is the map $\text{ad}:\mathfrak{g} \to
\text{End}(\mathfrak{g})$,
$X \mapsto [X,\cdot]$.
(Indeed, to compute this differential we use
the curve definition,
$\gamma(t)=\text{exp}(tX)$…)

\medskip\noindent

\begin{exercise}[Warped product]
\label{exercise-wraped-product}
See \cite{pet}. Para $(M,g_M)$ e $(N,g_N)$ e
$f:M \to \mathbb{R}_+$, definimos o
 {\it warped product} como sendo o produto
 cartesiano $M\times N$ com a métrica
 $g_M+f^2g_N$.

Mostre que se os fatores são completos, o
warped product é completo.
\end{exercise}

\begin{definition}
\label{definition-gradient}
Let  $f:M\to\mathbb{R}$. Define the {\it
gradient} of $f$ as the (unique) vector
field $\nabla f$ such that
\begin{equation}
\label{equation-gradient}
\left\langle{}\nabla f,X\right\rangle{}=Xf, \qquad \forall
X\in\mathfrak{X}(M).
\end{equation}
\end{definition}

Notice that for a frame $E_i$ with dual frame
$E^i$ we have
$$
\begin{aligned}
\left\langle{}\nabla f,X\right\rangle{}&=g_{ij}E^i\odot
E^j(\nabla f,X)
=g_{ij}(\nabla f)^iX^j\\
Xf&=df(X)=(df)_iE^i(X)=(df)_iX^i\\
\implies& g^{ij}(df)_iX^j=(\nabla f)^iX^i
\end{aligned}
$$
so substituting $X$ with a basic vector $E_i$
and $(df)_i$ with $\partial_if$ we
obtain
\begin{equation}
\label{equation-gradient-in-coordinates}
(\nabla f)^i=g^{ij}\partial_if \iff \nabla
f=g^{ij}\partial_ifE_j
\end{equation}

\begin{definition}
\label{definition-divergence}
For any $X\in \mathfrak{X}(M)$,
$$
\text{div}(X):=\text{tr}(v\mapsto \nabla_vX)
$$
\end{definition}

\begin{lemma}
\label{lemma-Laplacian-of-product}
\cite{doc}, Chapter III, Exercise 9(b).
\begin{equation}
\label{equation-Laplacian-of-product}
\Delta(fg)=f \Delta g + g \Delta
f+2\left\langle{}\nabla f,\nabla g\right\rangle{}
\end{equation}
\end{lemma}

\begin{proof}
Just compute.
\end{proof}

\begin{lemma}
\label{lemma-divergence-and-volume}
\begin{equation}
\label{equation-divergence-and-volume}
di_X\text{Vol}=\text{div}X\text{Vol},\qquad
\text{for all }X\in \mathfrak{X}(M).
\end{equation}
\end{lemma}

\begin{proof}
Pick a geodesic frame $E_i$ (see Lemma
\ref{lemma-geodesic-coordinates}) and
its dual coframe $\varepsilon^i$, i.e.
satisfying
$\varepsilon^i(E_j)=\delta_{ij}$. Then
$\operatorname{Vol}=\varepsilon^1
\wedge \ldots \wedge \varepsilon^n$. Then for
any $X=X^iE_i \in
\mathfrak{X}(U)$,

$$
\begin{aligned}
i_X\operatorname{Vol}
&=
\varepsilon^1\wedge\ldots\wedge
\varepsilon^n(X^iE_i,\cdot,\ldots,\cdot)
=X^i\varepsilon^1\wedge\ldots\wedge
\varepsilon^n(E_i,\cdot,\ldots,\cdot)
\end{aligned}
$$
How to compute that? Recall that for
top-forms we have
$$
\varepsilon^1\wedge\ldots\wedge
\varepsilon^n(Z_1,\ldots,Z_n)
=\det(\varepsilon^i(Z_j))
$$
so for example if $n=3$
\begin{align*}\varepsilon^1\wedge
\varepsilon^2\wedge\varepsilon^3(E_1,Z_2,Z_3)
&=\begin{vmatrix}
\varepsilon^1(E_1)&\varepsilon^1(Z_2)
&\varepsilon^1(Z_3)\\
\varepsilon^2(E_1)&\varepsilon^2(Z_2)&
\varepsilon^2(Z_3)\\
\varepsilon^3(E_1)&\varepsilon^3(Z_2)&
\varepsilon^3(Z_3)
\end{vmatrix}
=\begin{vmatrix}
1&\varepsilon^1(Z_2) &\varepsilon^1(Z_3)\\
0&\varepsilon^2(Z_2)&\varepsilon^2(Z_3)\\
0&\varepsilon^3(Z_2)&\varepsilon^3(Z_3)
\end{vmatrix}\\
&=\varepsilon^2(Z_2)\varepsilon^3(Z_3)-
\varepsilon^2(Z_3)\varepsilon^3(Z_2)
=\varepsilon^2\wedge\varepsilon^3(Z_2,Z_3),\\
\varepsilon^1\wedge\varepsilon^2\wedge
\varepsilon^3(E_2,Z_2,Z_3)
&=\begin{vmatrix}
0&\varepsilon^1(Z_2) &\varepsilon^1(Z_3)\\
1&\varepsilon^2(Z_2)&\varepsilon^2(Z_3)\\
0&\varepsilon^3(Z_2)&\varepsilon^3(Z_3)
\end{vmatrix}=-
\varepsilon^1\wedge\varepsilon^3(Z_2,Z_3)
\end{align*}
and so on. When we sum over all $i$, we get
$$
X^i\varepsilon^1\wedge\ldots\wedge
\varepsilon^n(E_i,\cdot,\ldots,\cdot)
=\sum_{i}(-1)^{i+1}
X^i\varepsilon^1\wedge\ldots\wedge\widehat{
\varepsilon^i}
\wedge\ldots\wedge\varepsilon^n.
$$
Now take exterior derivative of that, we get
$$
\begin{aligned}
d i_X\operatorname{Vol}
&=\sum_i(-1)^{i+1}(dX^i)
\wedge\varepsilon^1\wedge\ldots
\wedge\widehat{\varepsilon^i}
\wedge\ldots\wedge\varepsilon^n\\
&+\sum_{i}(-1)^{i+1}X^i\wedge
d(\varepsilon^1\wedge\ldots\wedge\widehat{
\varepsilon^i}\wedge\ldots
\wedge\varepsilon^n)\end{aligned}
$$
And then the first term actually is
$$
\begin{aligned}
\sum_i(-1)^{i+1}
E_jX^i\varepsilon^j\wedge\varepsilon^1
\wedge\ldots
\wedge\widehat{\varepsilon^i}
\wedge\ldots\wedge\varepsilon^n
&=E_iX^i\operatorname{Vol}
\end{aligned}
$$
while the second term vanishes because look,
$$
\begin{aligned}d \varepsilon^i(E_j,E_k)
&=E_j\varepsilon^i(E_k)-E_k\varepsilon^i(E_j)
-\varepsilon^i([E_j,E_k])\\
&=-\varepsilon^i(\nabla_{E_j}E_k-\nabla_{E_k}
E_j)\qquad
\text{torsion!}
\end{aligned}
$$
which vanishes at $p$ because we said that
this geodesic frame would have
 vanishing Christoffel symbols at $p$. So we
 conclude:
\[di_X
\operatorname{Vol}=E_iX^i\operatorname{Vol}\]

Now you just have to think what is
divergence:
$$
\begin{aligned}
\operatorname{div}X&=\sum_{i}
\left\langle{}\nabla_{E_i}X^jE_j,E_i\right\rangle{}
=\sum_{i}\left\langle{}E_iX^jE_j,E_i\right\rangle{}+
X^j\left\langle{}\nabla_{E_i}E_j,E_i\right\rangle{}\\
&=\sum_i
\left\langle{}E_iX^jE_j,E_i\right\rangle{}=
\sum_iE_iX^j\left\langle{}E_j,E_i\right\rangle{}=E_iX^i
\end{aligned}
$$
again using that we are a geodesic frame with
vanishing covariant
derivative at $p$.
\end{proof}

















\section{Connections}
\label{section-connections}

\begin{theorem}[Miracle]
\label{theorem-miracle}
There exists a unique metric torsion-free
connection on any Riemannian manifold.
\end{theorem}

\begin{proof}
$$
\begin{aligned}
2\left\langle{}\nabla_XY,Z\right\rangle{}
&=X\left\langle{}Y,Z\right\rangle{}+Y\left\langle{}X,Z\right\rangle{}-
Z\left\langle{}X,Y\right\rangle{}\\
&-\left\langle{}X,[Y,Z]\right\rangle{}+\left\langle{}Y,[Z,X]\right\rangle{}+
\left\langle{}Z,[X,Y]\right\rangle{}.
\end{aligned}
$$
\end{proof}

\begin{example}
\label{example-euclidean-connection}
The usual connection in $\mathbb{R}^n$ is
given by
$$
\nabla_XY=X(Y^i)\partial_i
$$
\end{example}

\begin{proposition}[Everything I know about
connections]
\label{proposition-exist-pullb-conne}
If $f:M \to N$ is a map of manifolds
and $E$ is a vector bundle on $N$
with a connection $\nabla$,
there exists a unique connection $\nabla^f$
on $f^*E$ such that
$$
\nabla^f_XY\circ f=\nabla_{f_*X}Y
$$
for $X \in TM$ and $Y \in \Gamma(E)$.
\end{proposition}

\begin{lemma}
\label{lemma-compo-pullb-conne}
\begin{reference}
\cite{aulas}, p.5
\end{reference}
If $g:M_1\to M_2$ and $f:M_2\to M_3$ are
smooth maps of
manifolds and $E$ is a vector bundle over
$M_3$, then
$$
g^* (f^* (E))=(f\circ g)^* (E)\qquad
\text{and}\qquad
(\nabla^f)^g=\nabla^{f\circ g}
$$
\end{lemma}

\begin{proposition}
\label{proposition-existence-of-levi-civita}
On any Riemann manifold
there exists a unique metric and torsion-free
connection called the {\it Levi-Civita
connection}.
\end{proposition}

\begin{proposition}[Symmetry and
compatibility]
\label{proposition-symme-compa}
Let $f:M \to N$ be a map of smooth manifolds
and suppose $M$ is Riemannian.
Then the pullback of the LC connection is
symmetric, i.e.
\begin{equation}
\label{equation-pullb-conne-symme}
\nabla^f_X f_*Y-\nabla^f_Yf_*X =
f_*[X,Y],\qquad
\forall X,Y \in \mathfrak{X}(N)
\end{equation}
and compatible with the pullback Riemannian
metric.
\end{proposition}

The following three lemmas are fundamental
properties of isometric
immersions and will be used in Section
\ref{section-isometric-immersions}. (Cf.
Lista 3, Exercise 2.)

\begin{lemma}
\label{lemma-tange-compo-pullb-conne-ii}
If $f:M \to \tilde{M}$ is an isometric
immersion, then
$f_*(\nabla_XY)=(\nabla^f_Xf_*Y)^\top$.
\end{lemma}

\begin{proof}
We have to check that $(\nabla^f_Xf_*Y)^\top$
is a torsion-free and compatible
connection on $TM$.
\end{proof}

\begin{lemma}
\label{lemma-isometric-immersion-along-curve}
Let $f:M \to \tilde{M}$ be an isometric
immersion, $c$ is a smooth curve in $M$
and $Y\in\mathfrak{X}_c$. Then
$$
f_*\nabla_{\frac{d}{dt}}^\gamma Y
=\left(\nabla^{f\circ\gamma}_{\frac{d}{dt}}
f_*Y\right)^\top
$$
\end{lemma}

\begin{proof}
This is just a mixture of Lemmas
\ref{lemma-compo-pullb-conne}
 and \ref{lemma-tange-compo-pullb-conne-ii}.
\end{proof}

The following result will be reminded when we
define totally geodesic
submanifolds, cf. Definition
\ref{definition-total-geode-subma},
addressing the case when also the normal
component vanishes:

\begin{lemma}
\label{lemma-geodesics-of-submanifold}
Let $f:M \to \tilde{M}$ be an isometric
immersion and $\gamma:I\to M$ a smooth
curve parametrized by arclength. Show that
$\gamma$ is a geodesic of M
if and only if its acceleration on
$\tilde{M}$ is perpendicular to $M$, that is,
$$
\tilde{\nabla}^{f\circ\gamma}_{\frac{d}
{dt}}(f\circ\gamma)'\perp
f_*(T_{\gamma(t)}M)
$$
\end{lemma}

\begin{proof}
Decompose the vector field
$\tilde{\nabla}^{f\circ\gamma}_{\frac{d}
{dt}}(f\circ\gamma)'$
in tangent and
normal part. Then apply Lemma
\ref{lemma-tange-compo-pullb-conne-ii}.
\end{proof}

\begin{exercise}[Parallel transport
determines connection]
\label{exercise-paral-trans-deter-conne}
\begin{reference}
\cite[Chapter II, Exercise 2]{doc}
\end{reference}
Let $X,Y$ be vector fields in a Riemannian
manifold $M$.
Let $p \in M$ and $c:I \to M$ an integral
curve of $X$ passing through $p$.
Show that the Riemannian connection of $M$ is
given by
$$
(\nabla_XY)(p)
=\frac{d}{dt}\Big|_{t=t_0}
(P^{-1}_{c;t_0;t}(Y(c(t)))
$$
where $P_{c;t_0;t}:T_{c(t_0)}M \to T_{c(t)}M$
is the parallel transport along $c$ from
$t_0$ to $t$.
\end{exercise}

\begin{proof}
Let $e_i$ be a basis of $T_{c(t_0)}M$.
Since parallel transport is a bijection
(its inverse is parallel transport
along the same curve but in reversed
direction;
+ uniqueness of ODE?),
we have that $P_t(e_i)$ is a basis of
$T_{c(t)}M$.
Thus we can write $P_tY$ in terms of this
basis, i.e.
$P_tY = Y^i_tP_te_i$.

Then
$$
\begin{aligned}
\nabla_XY&=\nabla_{c'}Y=\nabla_{c_*\frac{d}
{dt}}Y
=\nabla_{\frac{d}{dt}}^{c}Y \circ c\\
&=\nabla^c_{\frac{d}{dt}}Y^i_tP_te_i
=\frac{d}{dt}Y^i_tP_te_i+
Y^i_t\underbrace{\nabla_{\frac{d}{dt}}P_te_i}
_{=0}
=\frac{d}{dt}P_tY.
\end{aligned}
$$
\end{proof}

\begin{exercise}
\label{exercise-unique-connection-torsion}
Show that if $\nabla$ is any connection
on $M$, then there exists a unique
torsion-free connection on $M$
which has the same geodesics as $\nabla$.
\end{exercise}

\begin{proof}
Suppose that $\nabla'$  and $\nabla''$
are two such connections. Let
$A(X,Y)=\nabla'_XY-\nabla''_XY$.
Then $A$ is a symmetric tensor.
By tensoriality, the condition of
having the same geodesics gives $A(X,X)=0$.
Since $A$ is symmetric, we may polarize
and obtain that $A \equiv 0$.

For existence let
$\nabla'=\nabla-\frac{1}{2}T$
where $T$ is the torsion of $\nabla$.
\end{proof}

\begin{exercise}
\label{exercise-two-torsion-free-connections}
$\nabla$ and $\nabla'$ are two torsion
free connections on $M$ with the same
geodesics up to reparametrization
if and only if there exists a 1-form
$\omega$ such that
$\nabla_XY-\nabla'_XY
=\omega(X)Y-\omega(Y)X$.
\end{exercise}

\begin{proof}

\end{proof}

\begin{exercise}
\label{exercise-constant-curve-connection}
Let $M$ be a Riemannian manifold and $p$ a
point in $M$. Consider the constant curve
$f:I \to M$ given by $f(t)=p$. Let $V$ be a
differentiable vector field along $f$, that
is, a differentiable map from $I$ to $T_pM$.
Show that $\frac{DV}{dt}=\frac{dV}{dt}$, that
is, the covariant derivative coincides with
the usual derivative.
\end{exercise}

\begin{proof}
$$
\begin{aligned}
\frac{DV}{dt}&=\nabla^f_{\frac{d}{dt}}V
=\nabla^f_{\frac{d}{dt}}V^i(\partial_i\circ
f)
=\underbrace{\frac{d}{dt}V^i(t)
\partial_i\circ
f}_{\text{usual derivative}}
+V^i(t)\nabla_{\frac{d}{dt}}^f\partial_i
\circ f.
\end{aligned}
$$

\noindent
Note that the last term vanishes since
$$
\nabla^f_{\frac{d}{dt}}\partial_i \circ f
= \nabla_{f_*\frac{d}{dt}}\partial_i
=\nabla_0\frac{d}{dt}\partial_i=0
$$
since $f$ is constant.
\end{proof}

\begin{exercise}
\label{exercise-geodesics-tangent-bundle}
It's possible to introduce a Riemannian
metric on the tangent bundle
of a Riemannian manifold. Let $(p,v) \in TM$
and $V,W$ tangent vectors
of $TM$ at $(p,v)$. Let
$$
\alpha:t \to (p(t),v(t)),\qquad  \beta:s \to
(q(s),w(s))
$$
with $p(0)=q(0)=p$, $v(0)=w(0)=v$
and $V=\alpha'(0)$, $W=\beta'(0)$.
Define an inner product in $TM$ by
\begin{equation}
\label{equation-tangent-bundle-metric}
\left\langle{}V,W\right\rangle{}_{(p,v)}
=\left\langle{}\pi_*(V),\pi_*(W)\right\rangle{}_p
+\left\langle{}\nabla^p_{\frac{d}{dt}}v(0),
\nabla_{\frac{d}{ds}^q}w(0)\right\rangle{},
\end{equation}
where $\pi:TM \to M, \pi(p,v)=p$.

\begin{enumerate}
\item Show that this inner product is well
defined
and induces a Riemannian metric on $TM$.

\item A vector $(p,v) \in TM$ is orthogonal
(in the metric above) to the fiber
$\pi^{-1}(p) \cong TM$
is called a {\it horizontal vector}. A curve
$$
t \to (p(t),v(t))
$$
in $TM$ is called {\it horizontal} if its
tangent vector is
horizontal for all $t$.
Show that the curve $t \to (p(t),v(t))$
is horizontal if and only if the vector field
$v(t)$
is parallel along $p(t)$ in $M$.

\item Show that the geodesic field is a
horizontal vector field.

\item Show that the trajectories of the
geodesic field
are geodesics of $TM$ in the above metric.
\end{enumerate}
\end{exercise}

\begin{proof}
\begin{enumerate}
\item We just have to check that the right
term in
Equation \ref{equation-tangent-bundle-metric}
does not depend on the choice of $\alpha$ and
$\beta$.
In principle the covariant derivatives could
be different
for different choices of curves, but this is
not the
case by the hypothesis that $\alpha'(0)=V$
and $\beta'(0)=W$.
Indeed, the covariant derivative
$\nabla_{\frac{d}{dt}}^pv(0)$
determines the last $n$ coordinates of
$\alpha'(0)$.

\item We compute $\left\langle{}\alpha',V\right\rangle{}$ for
a curve $\alpha$ on $TM$
and a vertical vector field $V$, i.e. $V$ is
contained in $TM$.
In Equation
\ref{equation-tangent-bundle-metric},
the first term will always vanish by
verticality,
while the second depends on
$\nabla_{\frac{d}{dt}}^pv$.
This vanishes if and only if $v$ is parallel
along $p$.

\item This is immediate from the previous
item
since the velocity vector fields of geodesics
are parallel by definition.

\item Let $\alpha=(p,v)$ be a trajectory of
the geodesic field,
that is $v=p'$ and $v'=0$.
To show that $\alpha$ is a geodesic of $TM$
we must show that
$\nabla_{\alpha'}^{TM}\alpha'=0$.
Note that
$$
\begin{aligned}
2\left\langle{}\nabla_{\alpha'}^{TM}\alpha',
\alpha'\right\rangle{}
&=\frac{d}{dt}\left\langle{}\alpha',\alpha'\right\rangle{}\\
&=\frac{d}{dt}\left(\left\langle{}\pi_*\alpha',
\pi_*\alpha'\right\rangle{}
+\left\langle{}\nabla_{\frac{d}{dt}}^pv,
\nabla_{\frac{d}{dt}}^pv\right\rangle{}\right)\\
&=\frac{d}{dt}\left\langle{}p',p'\right\rangle{}=0
\end{aligned}
$$
since $p$ is a geodesic of $M$.
But $\alpha'$ cannot be zero because $p$ is a
geodesic.
\end{enumerate}
\end{proof}

\begin{exercise}
\label{exercise-lie-group-bi-invar-metri}
Let $G$ be a Lie group and $\mathfrak{g}$ its
Lie algebra.
Let $X \in \mathfrak{g}$. The trajectories of
$X$
determine a map
$\varphi:(-\varepsilon,\varepsilon) \to G$
with $\varphi(0)=e$ and
$\varphi'(t)=X(\varphi(t))$.

\begin{enumerate}
\item Show that $\varphi(t)$ is defined for
all $t \in \mathbb{R}$
and that $\varphi(t+s)=\varphi(t)\varphi(s)$.
$\varphi$ is called a {\it 1-parameter
subgroup} of $G$.
\end{enumerate}
\end{exercise}

\begin{proof}
\begin{enumerate}
\item Recall that $X \in \mathfrak{g}$
determines a vector field on $G$
by left translation, that is,
$X(g):=L_{g,*}X(e)$ where $X:=X(e)$.
To show that $\varphi$ is defined for all $t$
we may show that there is $\varepsilon>0$
such that for all $t$,
$\varphi(t+\varepsilon)$ is defined.
Suppose that $\varphi$ is a local solution of
$X$ (at $e$).
Pick any $t$ for which $\varphi$ is defined
and consider the translation
$L_{\varphi^{-1}(t)}$, which we use to
pushforward $X(\varphi(t))$
back at $e$.
But $X$ was constructed using left
translations,
so that the vector we obtain is again $X(e)$.
Integrating again we have a local solution up
to the same $t$.
Pushforward now using $L_{\varphi(t)}$.

To show that
$\varphi(t+s)=\varphi(t)\varphi(s)$
fix $t$. Then $\varphi(t+s)$ is an integral
curve of $X$
passing through $\varphi(t)$ at $s=0$.
But so is $\varphi(t)\varphi(s)$ since
$\varphi(0)=e$.
By uniqueness of solutions of ODEs we are
done.
\end{enumerate}
\end{proof}


\section{Riemannian submersions and
immersions}
\label{section-riema-subme-immer}

For any smooth manifold submersion
$\pi:\tilde{M}\to M$ there is a bundle
isomorphism I love:
\begin{equation}
\label{equation-subme-bundl-decom}
\pi^*TM \oplus \Ker \pi\cong
T\tilde{M} \end{equation}
 where $\Ker\pi$ is the {\it kernel bundle},
 whose
fiber is the kernel of $\pi_*$, and is also
called the {\it vertical bundle}.
Proving that $\Ker\pi$ is in fact a bundle
and confirming the isomorphism
\ref{equation-subme-bundl-decom} is a simple
exercise.

We call $\pi$ a {\it riemannian submersion}
only when $\pi^*TM$ is isometric to
$TM$, which makes sense only when $\tilde{M}$
and $M$ are Riemannian manifolds.

Following notation of
\cite{Milnor-Characteristic-Classes}, for any
subbundle
$\xi \subset\eta$ over a Riemannian manifold,
the {\it orthogonal complement
bundle} $\xi^\perp$ can be defined fiberwise
as the orthogonal complement of the
fiber and shown to be a bundle as in the case
of the kernel bundle. Then Milnor
defines the {\it normal bundle} when $\xi$ is
the tangent bundle of a
submanifold. Finally this construction is
noted to work just as well for any
isometric immersion $f:M\to \tilde{M}$, so
that
\begin{equation}
\label{equation-immer-bundl-decom}
f^*T\tilde{M}\cong TM\oplus T^\perp M
\end{equation}
Further, in problem 3-B we are asked to
define for $\xi \subset \eta$ the {\it
quotient bundle} and show it's a bundle. And
if $\eta$ has a Riemannian metric,
then $\xi^\perp\cong\eta/\xi$. So that's why
we sometimes say that we can define
the normal bundle using a Riemannian metric
but we don't need to.

An interesting thought is: given a subbundle,
is there a manifold whose tangent
bundle is that subbundle?

\section{Geodesics}
\label{section-geodesics}

\begin{example}[Geodesics of sphere]
\label{example-geodesics-of-Sn}
The geodesics of $S^n$ are
\begin{equation}
\label{equation-geodesics-of-Sn}
\gamma(t)=\cos t p+\sin t v
\end{equation}
for $p \in S^n$ and $v\in T^1_pS^n$.

Derivando como uma simples curva em
$\mathbb{R}^{n+1}$, vemos que
$\gamma''=-\gamma$, o que significa que
$(\gamma'')^{\top}=0$, i.e. $\gamma$ é
uma geodésica de $\mathbb{S}^n$.

Sendo essa uma geodésica partindo de um ponto
arbitrário numa direção
arbitrária, concluimos por unicidade das
geodésicas e \textit{rescaling lemma}
que todas as geodésicas de $\mathbb{S}^n$ são
como $\gamma$.
\end{example}

\begin{example}[Isometry group of sphere]
\label{example-isometry-group-of-Sn}
The isometry group of $S^n$ is
$$
\text{O}(n):=\{A\in\text{GL}(n+1):
AA^{\mathbf{T}}=\text{Id}\}
$$
It's immediate that $\mathcal{O}(n+1)\subset
\text{Isom}(S^n)$ since orthogonal
transformations preserve euclidean product:
$$
\begin{aligned}
A A^{\mathbf{T}}&=\text{Id}\iff \sum_k
A_{ik}A_{jk}=\delta_{ij}
\iff Ae_i\cdot Ae_j=\delta_{ij}\\
\iff Av\cdot Aw&= A\left( v^ie_i\right) \cdot
A\left( w^je_j\right)
=v^iw^j (Ae_i)\cdot (Ae_j)
=v^i\cdot w^j\delta_{ij}=v\cdot w
\end{aligned}
$$
For the other contention let $A:S^n \to S^n$
be an isometry.

Vamos mostrar que $A$ é a restrição de uma
função $\tilde{A} \in
\mathsf{O}(n+1)$. Defina
$$
\begin{aligned}
\tilde{A}: \mathbb{R}^{n+1} &\longrightarrow
\mathbb{R}^{n+1} \\
(r,\theta) &\longmapsto rA(1,\theta)\\
0&\longmapsto 0
\end{aligned}
$$
Se mostramos que $\tilde{A}$ é uma isometria
linear, é claro que ela é um
elemento de $\mathsf{O}(n+1)$ pela conta
anterior. De fato, basta mostrar que
$\tilde{A}$ é uma isometria, pois toda
isometria de espaços de Banach que fixa
a origem é linear (\cite{braga} Teo. 7.11).

Para ver que $\tilde{A}$ é uma isometria de
$\mathbb{R}^{n+1}$,
afirmo que a distância de $p$ a $q$ está
totalmente determinada
pelas normas  $\|p\|$ e $\|q\|$, e pela
distancia esférica entre
$\frac{p}{\|p\|}$ e $\frac{q}{\|q\|}$. Note
que essa afirmação é na verdade
um problema de geometria plana, pois todas
essas quantidades podem ser descritas
dentro do único plano que contém $0$, $p$ e
$q$.

Acabou que essa afirmação é simplesmente a
lei dos cosenos, já que a distância
esférica entre $\frac{p}{\|p\|}$ e
$\frac{q}{\|q\|}$ é exatamente o angulo
entre $p$ e $q$ (poque essa distância é um
segmento de círculo máximo!):
\[\text{lei dos cosenos:} \qquad
d(p,q)^2=\|p\|^2+\|q\|^2-2\|p\|\|q\| \cos
\angle(p,q) \]

Em fim, $\tilde{A}$ é uma isometria porque
$d_{\mathbb{R}^{n+1}}(p,q)=d_{\mathbb{R}^{n+
1}}(\tilde{A}p,\tilde{A}q)$
pelo
argumento anterior.
\end{example}

\begin{proposition}[Totally Normal
Neighbourhoods]
\label{proposition-total-norma-neigh}
For all $p \in M$ there is a neighbourhood $W
\ni p$ and a number $\delta>0$
 such that for every $q \in W$, $B_\delta(q)$
 is normal and
 $W \subset B_\delta(q)$.
\end{proposition}

\medskip\noindent
Now we discuss distance, geodesics and
exponential.
First notice that if $p$ and $q$ are two
points sufficiently near each other,
$d(p,q)=|\text{exp}^{-1}_p(q)|$.
Indeed, if $q=\gamma_w(1)$ for some geodesic
with
velocity $w \in T_pM$ at $p$ (this is what
we mean by ``sufficiently near''),
we have

$$
\begin{aligned}
d(p,q)=\int_0^1|\gamma'|=\int_0^1|w|=|w|.
\end{aligned}
$$

\begin{exercise}
\label{exercise-diffe-dista-funct}
Let $d_p:=d(p,\cdot)$. Compute $|\nabla d_p|$
and the integral
curves of $\nabla d_p$ inside
$B_\varepsilon(p)\setminus\{p\}$.
\end{exercise}

\begin{proof}
Suppose that $\gamma$ is a geodesic joining
$p$ and $q$
with $|\gamma'|=1$.
Let $\sigma$ be a curve contained in the
geodesic sphere $S_{d(p,q)}(p)$.
Then
$$
\left\langle{}\nabla d_p,\sigma'\right\rangle{}
=\sigma'(d_p\circ \sigma)
=\sigma'(\underbrace{d(p,q)}_{\text{cte}})=0.
$$
Thus $\nabla d_p=\lambda\gamma'$.
Now
$$
\left\langle{}\nabla d_p,\gamma'\right\rangle{}
=\lambda |\gamma'|=\lambda
$$
for some $\lambda \in \mathbb{R}$.
But
$$
\gamma'd_p=\frac{d}{dt}d_p \circ \gamma'
=\frac{d}{dt}t=1,
$$
so $\lambda=1$, and then $|\nabla d_p|=1$.
It is clear that the integral curves of
$\nabla d_p$
are unit speed geodesics.
\end{proof}

\begin{lemma}
\label{lemma-geodesic-coordinates}
At every point $p$ in $M$ there is an
orthonormal frame
$\{E_i\} \subset \mathfrak{X}(U)$ such that
\begin{equation}
\label{equation-norma-coord-vanis-chris}
\Big(\nabla_{E_i}E_j\Big)_p=0
\end{equation}
(This means the Christoffel symbols of the
connection vanish at the point $p$.)
\end{lemma}

\begin{proof}
This frame can be obtained by taking geodesic
coordinates at the point, an
orthonormal base $\{e_i\}$ of $T_pM$, and
taking parallel transport of the
vectors $e_i$ along radial geodesics
emanating from $p$. This immediately
ensures that $E_i$ is orthonormal since
parallel transport preserves angles.

To check that Christoffel symbols vanish at
$p$ we do as follows.
Take a random vector $v \in T_pM$ and its
geodesic
$\gamma_v(t)=\operatorname{exp}_p(tv)$. Then
$$
\begin{aligned}
0&=\nabla_{\frac{d}{dt}}^{\gamma}\gamma'
=\nabla^{\gamma}_{\frac{d}{dt}}v^i(E_i\circ
\gamma)
\end{aligned}
$$
where the $v=(v^1,\ldots,v^n)$. Indeed: this
is very silly but, since the
coordinate chart of geodesic coordinates is
$\operatorname{exp}_p^{-1}$,
the coordinate representation of $\gamma$ in
this chart is as simple as
$$
\hat{\gamma}(t)=(\underbrace{\varphi}
_{\text{chart}}
\circ \gamma)(t)
=\operatorname{exp}_p^{-1}
\operatorname{exp}_p(tv)=tv.
$$
And the composition $E_i \circ \gamma$ just
means that we take our local
frame \textit{along $\gamma$}. Continue:
$$
\begin{aligned}
&=v^i\nabla^{\gamma}_{\frac{d}{dt}}
E_i\circ\gamma
=v^i\nabla_{\gamma_{v,*} \frac{d}{dt}}E_i\\
&=v^i\nabla_{v^jE_j}E_i=v^iv^j\nabla_{E_j}
E_i\\&
=v^iv^j\Gamma_{ji}^kE_k
\end{aligned}
$$
along $\gamma$. Now choose $v=e_1$. You get
$\Gamma_{11}^k=0$ for all $k$ along
$\gamma_{e_1}$. Now choose $v=e_2$, then
$\Gamma_{22}^k=0$ along $\gamma_{e_2}$, so at
least at $p$ they both vanish. And now choose
$v=e_1+e_2$. You get
$$
\begin{aligned}
0=(v^1)^2 \underbrace{\Gamma_{1 1}^k}_{=0}
+v^1v^2 \Gamma_{1 2}^k+v^2v^1\Gamma_{2 1}^k
+(v^2)^2\underbrace{\Gamma_{2 2}^k}_{=0}
\end{aligned}
$$
So $\Gamma_{12}^k=0$ since Levi-Civita is
torsion-free, i.e. symmetric.
And so on. So the all Christoffel symbols
vanish at the same time at $p$.
\end{proof}

\begin{proposition}
\label{proposition-naturality-of-exponential}
If $f:M\to \tilde{M}$ is an isometry,
$$
\xymatrix{
B_\varepsilon(0)\subset T_pM \ar[r]^{d_p f}
\ar[d]_{\text{exp}_p} &
B_\varepsilon(0)\subset T_{f(p)}\tilde{M}
\ar[d]^{\text{exp}_{f(p)}} \\
B_\varepsilon(p)\subset M \ar[r]_{f} &
B_\varepsilon(f(p))\subset \tilde{M}
}
$$
\end{proposition}

\begin{proof}
Isometry maps geodesics to geodesics by Lemma
\ref{lemma-isome-maps-geode-geode},
then uniqueness of solutions
of ODEs.
\end{proof}

\begin{definition}
\label{definition-killing}
A {\it Killing vector field} is a vector
field
whose flow diffeomorphisms are isometries.
\end{definition}

\begin{exercise}
\label{exercise-killing}
$X$ is a Killing vector field if and only if
$$
\left\langle{}\nabla_YX,Z\right\rangle{} +
\left\langle{}\nabla_ZX,Y\right\rangle{}=0.
$$
\end{exercise}

\begin{exercise}
\label{exercise-killi-field-zero-order}
Let $X$ be a Killing field in a connected
Riemannian manifold $M$.
Suppose there is a point $q \in M$ such that
$X(q)=0$ and $\nabla_YX(q)=0$ for any $Y(q)
\in T_qM$.
Show that $X=0$.
\end{exercise}


\begin{exercise}
\label{exercise-maximum-principle}
\cite{doc}, Chapter III, Exercise 12. If $f$
is {\it subharmonic}, i.e.
$\Delta f\geq 0$, on $M$
compact connected $\implies$ $f$ constant.
\end{exercise}

\begin{proof}
To prove the theorem for subharmonic
functions first we show that in fact they
are harmonic via Eq.
\ref{equation-divergence-and-volume}
on $X:= \nabla f$ and Stokes:
$$
\int_M \Delta f \operatorname{Vol}=\int_M
\operatorname{div}X \operatorname{Vol}
=\int_M d(i_X
\operatorname{Vol})=\int_{\partial M}i_X
\operatorname{Vol}=0
$$
meaning that the non-negative function
$\Delta f$ is in fact 0, i.e. $f$ is
harmonic. Now we do it again for
$X:=\nabla(f^2/2)$:
$$
\int_M \Delta(f^2/2)\operatorname{Vol}
=\int_M d(i_X \operatorname{Vol})
=\int_{\partial M}i_X \operatorname{Vol}
=0
$$
And then apply Eq.
\ref{equation-Laplacian-of-product}:
$$
0=\int_M \Delta(f^2/2)\operatorname{Vol}
=\int_M f \Delta f \operatorname{Vol}
+\int_M \left\langle{}\nabla f,\nabla
f\right\rangle{}\operatorname{Vol}
$$
First one vanishes because $f$ is harmonic,
so second one is zero which says
 $f$ is constant!
\end{proof}


\section{Curvature}
\label{section-curvature}

The {\it curvature tensor} is
\begin{equation}
\label{equation-curvature-tensor}
R(X,Y)Z=\nabla_X\nabla_YZ-\nabla_Y\nabla_XZ-
\nabla_{[X,Y]}Z
\end{equation}

\begin{proposition}
\label{proposition-symme-curva-tenso}
The curvature tensor is a tensor and
\begin{enumerate}
\item It is antisymmetric in the first two
and last two entries,
\item
$\left\langle{}R(X,Y)Z,W\right\rangle{}=\left\langle{}R(Z,W)X,
Y\right\rangle{}$,
and
\item Bianchi identity.
\end{enumerate}
\end{proposition}

Given a dimension-2 vector subspace
$\sigma\subset T_pM$, the {\it sectional
curvature} of $M$ in $\sigma$ is
\begin{equation}
\label{equation-sectional-curvature}
K(\sigma)=\frac{\left\langle{}R(u,v)v,u\right\rangle{}}
{\left\langle{}u,u\right\rangle{}\left\langle{}v,v\right\rangle{}-\left\langle{}u,
v\right\rangle{}^2}
\end{equation}
for any base $\{u,v\}$ of $\sigma$.

\begin{lemma}
\label{lemma-secti-curva-well-defin}
Sectional curvature does not depend on the
choice of base.
\end{lemma}

\begin{proposition}
\label{proposition-curva-like-tenso-deter}
Let $R$ and $R'$ be two type $(0,4)$ tensors
on a vector space $\mathbb{V}$. If
their associated sectional curvatures $K$ and
$K'$ are equal, then so are $R$
and $R'$.
\end{proposition}

\begin{proof}
Polarize and make lots of computations, cf.
\cite{doc}.
\end{proof}

If a manifold has constant sectional
curvature $c$,
\begin{equation}
\label{equation-constant-curvature}
\left\langle{}R(X,Y)Z,W\right\rangle{}=
c(\left\langle{}X,Z\right\rangle{}\left\langle{}Y,W\right\rangle{}-\left\langle{}X,
W\right\rangle{}\left\langle{}Z,Y\right\rangle{})
\end{equation}
since the right-hand-side of Eq.
\ref{equation-constant-curvature} is a
so-called curvature tensor, i.e. a tensor
satisfying the symmetries of $R$, and
applying Proposition
\ref{proposition-curva-like-tenso-deter}.

The {\it Ricci tensor} is
\begin{equation}
\label{equation-Ricci-tensor}
\begin{aligned}
\text{Ric}: TM\times TM &\longrightarrow
\mathbb{R} \\
\text{Ric}(X,Y)&=\frac{1}{n-1}\text{tr}
(Z\mapsto
R(Z,X)Y)
\end{aligned}
\end{equation}
It is symmetric.

The {\it Ricci curvature} is
\begin{equation}
\label{equation-Ricci-curvature}
\begin{aligned}
\text{Ric}: T^1M &\longrightarrow \mathbb{R}
\\
\text{Ric}(X)=\text{Ric}(X,X)&=\frac{1}{n-1}
\sum_{i=1}^{n-1}K(X,E_i)
\end{aligned}
\end{equation}
where $\{X,E_1,\ldots,E_{n-1}\}$ is an
orthonormal basis.

The {\it scalar curvature} is
\begin{equation}
\label{equation-scalar-curvature}
\begin{aligned}
\text{scal}: M &\longrightarrow \mathbb{R} \\
\text{scal}(p)&=\frac{1}{n}\text{tr}
\text{Ric}
=\frac{1}{n}\sum_i\text{Ric}(E_i,E_i)
\end{aligned}
\end{equation}

\begin{lemma}
\label{lemma-scala-curva-mean-secti}
$$
\text{scal}(p)=\frac{1}{n}\sum_{i}\text{Ric}
(e_i,e_i)
=\frac{1}{n(n-1)}\sum_{i,j}K(e_i,e_j)
=\frac{1}{\text{Vol}(S^{n-1})}
\int_{\sigma\subset
T_pM}K(\sigma)
$$
\end{lemma}

\medskip\noindent
The following result is very important.

\begin{lemma}[Symmetry lemma for curvature]
\label{lemma-symmetry-lemma-curvature}
Let $f:U \subset \mathbb{R}^2 \to M$
be a map into a Riemannian manifold and $V
\in \mathfrak{X}^f$.
Then
$$
\nabla_{\partial_s}\nabla_{\partial_t}V
-\nabla_{\partial_t}\nabla_{\partial_s}V
=R(f_s,f_t)V
$$
where $f_s:=f_*\partial_s$,
$f_t:=f_*\partial_t$.
More generally, the pullback of $R$ is the
Riemannian curvature tensor of the pullback
vector bundle
for any affine vector bundle $(E,\nabla)$ on
$M$.
\end{lemma}


\begin{exercise}
\label{lemma-curvatures-of-product-metric}
Lista 5, Exercise 1. Let $(M_1,g_1)$ and
$(M_2,g_2)$ be Riemannian manifolds and
let $M_1\times M_2$ be equipped with the
product metric $g:=g_1\oplus g_2$. Show
that the curvature tensor, the Ricci
curvature and the scalar curvature of $g$
are given by
\begin{enumerate}
\item $R=\pi_1^*R_1+\pi_2^*R^2$,
\item
$\text{Ric}=\pi_1^*\text{Ric}_1+
\pi_2^*\text{Ric}_2$,
\item
$\text{scal}=\pi_1^*\text{scal}_1+
\pi_2^*\text{scal}_2$
\end{enumerate}
where $\pi_i:M_1\times M_2\to M_i$ are the
projections.
\end{exercise}

\begin{example}[Curvature of complex
projective space]
\label{example-curvature-of-CPn}
Defina uma métrica Riemanniana  em
$\mathbb{C}^{n+1}\setminus\{0\}$ do
seguinte modo. Se $Z \in
\mathbb{C}^{n+1}\setminus\{0\}$ e $V,W \in
T_Z
(\mathbb{C}^{n+1}\setminus\{0\})$,

$$
\left\langle{}V,W\right\rangle{}_Z=\frac{\text{Re}(V,W)}{(Z,
Z)}
$$
onde
$$
(Z,W)=z_0\overline{w}_0+\ldots+
z_n\overline{w}_n
$$

é o produto hermitiano em $\mathbb{C}^{n+1}$.
Observe que a métrica
$\left\langle{}\cdot,\cdot\right\rangle{}$ restrita a
$S^{2n+1}\subset
\mathbb{C}^{n+1}\setminus\{0\}$ coincide com
a métrica induzida por
$\mathbb{R}^{2n+2}$.
\begin{enumerate}

\item Mostre que, para todo $0\leq \theta
\leq 2\pi$, $e^{i\theta}:S^{2n+1}\to
S^{2n+1}$ é uma isometria, e que, portanto, é
possível definir uma métrica
Riemanniana em $\mathbb{P}^n(\mathbb{C})$ de
modo que a submersão $f$ seja
Riemanniana.

\item Mostre que, nesta métrica, a curvatura
seccional de
$\mathbb{P}^n(\mathbb{C})$ é dada por
$$
K(\sigma)=1+3\cos^2\varphi
$$
onde $\sigma$ é gerado pelo par ortonormal
$X$, $Y$, $\cos
\varphi=\left\langle{}\overline{X},i\overline{Y}
\right\rangle{}$,
e $\overline{X}$,
$\overline{Y}$ são os levantamentos
horizontais de $X$ e $Y$, respetivamente. Em
particular, $1 \leq  K(\sigma) \leq  4$.
\end{enumerate}
\end{example}

\section{Jacobi Fields}
\label{section-jacobi-fields}

\begin{proposition}[Everyday Jacobi field]
\label{proposition-everyday-jacobi-field}
\cite{doc} Chapter V, Proposition 2.4. Let
$\gamma:[0,a]\to M$ be a normalized
geodesic and let $J$ be a Jacobi field
through $\gamma$ with $J(0)=0$.
Let $p:=\gamma(0)$ $w:=J'(0)$ and
$v:=\gamma'(0)$.
Consider the variation
$$
f(s,t):=\text{exp}_{p}(t(v+sw))
$$
Then the associated Jacobi field is $J$. So
in particular the field
$$
J(t)=d_{tv}\text{exp}_p(tw)
$$
is a Jacobi field.

{\bf Remark:} the variation can also be taken
with respect to
any a curve $\sigma$ on $T_vT_{p}M$ passing
through $v$
at $s=0$ with velocity $w$.
\end{proposition}

\begin{example}[Jacobi Fields in constant
curvature]
\label{example-jacob-field-const-curva}
If $M$ has constant curvature $K$ and
$\gamma$ is a normalized geodesic, a
Jacobi field is the solution of the equation
\begin{equation}
\label{equation-jacob-field-const-curva}
J''+KJ=0
\end{equation}
This follows from Eq.
\ref{equation-constant-curvature} since it
implies that
for any $X\in\mathfrak{X}_\gamma$, which we
may take orthogonal to $\gamma'$,
$$
\left\langle{}X,R_{\gamma'}J\right\rangle{}
=K(\left\langle{}X,J\right\rangle{}\left\langle{}\gamma',
\gamma'\right\rangle{}
-\left\langle{}X,\gamma'\right\rangle{}\left\langle{}J,
\gamma'\right\rangle{})
$$
so that $R_{\gamma'}J=KJ$. The solutions of
Eq.
\ref{equation-jacob-field-const-curva} can be
computed taking a
parallel vector field $W$ along $\gamma$
orthogonal to $\gamma'$ and of unit
length. Then
\begin{equation}
\label{equation-jacob-field-const-curva-2}
J(t)=\begin{cases}
\sin(\sqrt{K}t)W(t)\qquad & K>0\\
tW(t)\qquad &K=0\\
\text{sinh}(\sqrt{-K}t)W(t)\qquad &K<0
\end{cases}
\end{equation}
provided that $J(0)=0$, $J'(t)\perp
\gamma'(t)$ and $J'(0)=W(0)$.
\end{example}

\begin{proposition}
\label{proposition-jacob-field-taylo-expan}
Differentiating a Jacobi field enough times,
we arrive to
\begin{equation}
\label{equation-jacob-field-taylo-expan}
\|J(t)\|^2=\|w\|^2t^2-\frac{1}{3}\left\langle{}R(w,v)
v,w\right\rangle{}t^4+O(t^4).
\end{equation}
\end{proposition}

Sometimes it's useful to work with

\begin{definition}
\label{definition-Jacobi-tensor}
The {\it Jacobi tensor} is the map
$$
\begin{aligned}
	\overline{J}_t:  &\longrightarrow
	T_{\gamma(t)}M \\
	 \overline{J}_t(w)&=d_{tv}\text{exp}_p(tw)
\end{aligned}
$$
\end{definition}

\begin{exercise}
\label{exercise-Jacobi-tensor-derivative}
Show that $\bar{J}'(0)=\text{Id}$.
\end{exercise}

\begin{proof}
Just notice that since
$\bar{J}_t=d_{tv}\text{exp}_p(tw)$ is linear
on
$tw$ (because it's a differential) we have
$$
\begin{aligned}
\overline{J}'(0)&=\frac{d}{dt}\Big|_{t=0}
(d_{tv}\text{exp}_p(tw))\\
&=\frac{d}{dt}\Big|_{t=0}t\cdot
(d_{tv}\text{exp}_p(w))\\
&=d_{tv}\text{exp}_p(w)|_{t=0}+
t\cdot(\text{something})|_{t=0}\\
&=\text{Id}_{T_pM}
\end{aligned}
$$
\end{proof}

\begin{remark}
\label{remark-dimen-space-jacob-field}
Lembre o seguinte fato geral mostrado em
aula:
\begin{equation}
\label{equation-dimen-space-ortho-jacob}
\dim\{J \in \mathfrak{X}^J_\gamma:J(0)=0, J
\perp \gamma\}=n-1
\end{equation}
Isso segue das seguintes observações:
\begin{enumerate}
\item $\dim\mathfrak{X}^J_\gamma=2n$ porque
são soluções de $n$ equações
diferenciais ordinárias de segunda ordem,
i.e. cada campo de Jacobi está
determinado pelas condições iniciais $J(0)$ e
$J'(0)$.
\item $\dim \{J \in
\mathfrak{X}^J_\gamma:J(0)=0\}=n$.
\item $\dim \{J \in
\mathfrak{X}^J_\gamma:J\perp \gamma'\}=2n-2$.
Para
confirmar isso note que pelas simetrias de
$R$ e a equação de Jacobi tem-se
que $\left\langle{}J,\gamma'\right\rangle{}''=0$, pelo que
$\left\langle{}J,\gamma'\right\rangle{}=a+bt$
para dois números reais $a,b$. Segue que
qualquer $J \in
\mathfrak{X}^J_\gamma$ se escreve como
$J(t)=a\gamma'(t)+bt\gamma'(t)+\hat{J}(t)$
para algum $\hat{J}$
perpendicular a  $\gamma'$. (Supondo que
$|\gamma'|=1$.) Então se $J \perp
\gamma$, temos que  $a=b=0$, então tiramos
dois números do $2n$ que
tínhamos.
\item Segue
\ref{equation-dimen-space-ortho-jacob}.
\end{enumerate}

\end{remark}


\section{Conjugate points}
\label{section-conjugate-points}

\begin{definition}
\label{definition-conjugate-points}
A point $p \in M$ is {\it conjugate} to $q
\in M$ along $\gamma$ if $\gamma$ is
a geodesic joining $p$ and $q$ and there is a
Jacobi field
$J\in\mathfrak{X}_\gamma^J$ such that
$J(0)=0$ and  $J(\ell)=0$.
\end{definition}

\begin{proposition}[Critical points of
exponential]
\label{proposition-critical-points-exp}
\cite{doc} Chapter V, Proposition 3.5. A
vector $v \in T_pM$ is a critical
point of $\text{exp}_p$ if and only if the
corresponding
$\gamma_v(1)$ is a conjugate point of $p$
along $\gamma_v$.
Moreover, in such points, the dimension of
the kernel of
$\text{exp}_p$ at this vector is the
multiplicity
of the conjugate point $\gamma_v(1)$
(i.e. the dimension of the space of Jacobi
fields vanishing
at the endpoints).
\end{proposition}

\begin{proof}
($\implies$) Suppose that $v$ is a critical
point of $\text{exp}_p$ and
let $w \in \ker d_v \text{exp}_p$ then a
Jacobi field along $\gamma$
vanishing at the endpoints is given by
$$
J(t):=d_{tv}\text{exp}_ptw
$$
($\impliedby$) If you have a critical point
along $\gamma$, then the
differential of $\text{exp}_p$ has kernel.

(Mulitplicity.) The proof every element of
the kernel gives a Jacobi
 field vanishing at the points. Linearly
 dependent elements will give
 linearly dependent Jacobi fields.
\end{proof}

\section{Isometric immersions}
\label{section-isometric-immersions}

Let $f:M \to \tilde{M}$ be an isometric
immersion of Riemannian manifolds. When
computing the covariant derivative along two
vector fields
$X,Y \in\mathfrak{X}(M)$ we get a vetor field
along $X$ that may be split into
its normal and tangent part thanks to
Equation
\ref{equation-immer-bundl-decom}. By Lemma
\ref{lemma-tange-compo-pullb-conne-ii} we see
that the
 tangent part is
in fact the Levi-Civita connection of $M$
since it is a symmetric and torsion
free connection on the tangent bundle. Since
the Levi-Civita connection of
$\tilde{M}$ is torsion free, the normal part
is a
symmetric tensor we call {\it second
fundamental form}:
\begin{equation}
\label{equation-second-fundamental-form}
\nabla^{f}_X f_*Y=f_*\nabla_XY+\alpha(X,Y)
\end{equation}

\medskip\noindent

Now let's differentiate a normal section
$\xi\in T^\perp M$.
We call the tangent part {\it shape
operator}. By Lemma
\ref{lemma-normal-connection} below, the
normal part is a connection
we call the {\it normal connection}.
\begin{equation}
\label{equation-shape-operator}
\nabla^f_X\xi=-A_\xi X+\nabla^\perp_X\xi
\end{equation}
\begin{lemma}
\label{lemma-normal-connection}
$\nabla^\perp$ is a connection on $T^\perp
M$.
\end{lemma}
Notice the adjointness relation
\begin{equation}
\label{equation-secon-funda-form-shape}
\left\langle{}\xi,\alpha(X,Y)\right\rangle{}=-\left\langle{}A_\xi
X,Y\right\rangle{}
\end{equation}
which follows from computing
$0=X\left\langle{}Y,\xi\right\rangle{}$.

The {\it Gauss equation}
$$
(R_{\nabla^f}(X,Y)Z)^\top=R_\nabla(X,Y)Z-
A_{\alpha(Y,Z)}X+A_{\alpha(X,Z)}Y
$$
may be derived by substituting Eq.
\ref{equation-second-fundamental-form}
into Equation
\ref{equation-curvature-tensor}. But most
likely we will use the
following form:
\begin{equation}
\label{equation-Gauss}
\left\langle{}\tilde{R}(X,Y)Z,W\right\rangle{}=\left\langle{}R(X,Y)Z,
W\right\rangle{}
-\left\langle{}\alpha(X,Z),\alpha(Y,W)\right\rangle{}+
\left\langle{}\alpha(X,W),\alpha(Y,Z)\right\rangle{}
\end{equation}
Substituting for $X=W$ and $Y=Z$ we obtain
the equivalent (by Proposition
\ref{proposition-curva-like-tenso-deter})
equation
\begin{equation}
\label{equation-Gauss-sectional-curvature}
K(X,Y)=\tilde{K}(X,Y)+\left\langle{}\alpha(X,Y),
\alpha(Y,Y)\right\rangle{}-\|\alpha(X,Y)\|^2
\end{equation}

Notice that we are seeing the determinant of
the second fundamental form, which
is the product of its eigenvalues, which may
be defined as the {\it principal
curvatures}. In particular, this shows that
the Gaussian curvature of a surface
in $\mathbb{R}^3$, defined as the determinant
of the second fundamental form,
coincides with sectional curvature when
substituting $\tilde{K}\equiv0$ for
 the flat metric on euclidean space. (!!!)

In the case of a hypersurface, the normal
vector $\alpha(X,Y)$ must be
proportional to the unit normal $N$, and by
Eq.
\ref{equation-secon-funda-form-shape} the
proportionally
factor is exactly $\left\langle{}A_N X,Y\right\rangle{}$.
Using this and evaluating
 at an orthonormal pair $X,Y$ we obtain
\begin{equation}
\label{equation-gauss-secti-curva-shape}
K(X,Y)=\tilde{K}(X,Y)+\left\langle{}AX,
X\right\rangle{}\left\langle{}AY,Y\right\rangle{}-\left\langle{}AX,Y\right\rangle{}^2
\end{equation}
In this case we are looking at the
determinant of the matrix
\begin{equation}
\label{equation-matrix-of-shape-operator}
\begin{pmatrix}
\left\langle{}AX,X\right\rangle{}&\left\langle{}AX,Y\right\rangle{}\\
\left\langle{}AX,Y\right\rangle{}&\left\langle{}AY,Y\right\rangle{}
\end{pmatrix}
\end{equation}
which is the matrix representation of the
shape operator $A$ to the plane
spanned by $X,Y$. (That's just a particular
case of the simple fact that the
representation of the matrix $A_i^j$ for
$Ae_i=A_i^je_j$ is given by
$\left\langle{}Ae_i,e_j\right\rangle{}=A^k_i\left\langle{}e_k,
e_j\right\rangle{}=A^j_i$.)

\begin{example}[Sectional curvature of
sphere]
\label{example-sectional-curvature-of-sphere}
Using Eq.
\ref{equation-gauss-secti-curva-shape} we can
prove that the sectional
curvature of $S^n$ is 1 once we make sure
that the shape operator of the sphere
is the identity. And this follows from the
fact that the connection on
$\mathbb{R}^n$ is just directional
derivative, see Example
\ref{example-euclidean-connection}.

More explicitly, the normal component of
$\tilde{\nabla}_X N$ will vanish (see
discussion in Section
\ref{section-Riccati-equation}), so that
$\tilde{\nabla}_XN=-A_N X$. But the
connection in Euclidean space is just
differentiating the vector field, that is,
$\tilde{\nabla}_X N=N X$. And
finally, who's $N$? It's $-\text{id}$ if you
chose the inner normal. So
 $\tilde{\nabla}_X N=-X$.
\end{example}

\begin{exercise}
\label{exercise-Gauss-sectional}
Use Eq.
\ref{equation-gauss-secti-curva-shape}
to show that no
hypersurface $M \hookrightarrow
\mathbb{R}^{n+1}$ can have constant negative
curvature.
\end{exercise}

\begin{proof}
After substituting $\tilde{K}(X,Y)=0$, we see
that Eq.
\ref{equation-gauss-secti-curva-shape} is
determinant of the second fundamental form.
 But… how to argue that this determinant
 can't be always negative?
\end{proof}

\begin{example}[Sectional curvature of
hyperbolic space]
\label{example-hyperbolic-space}
\cite{doc}, Chapter VIII, Exercise 3. The
unit normal of hyperbolic space in the
hyperboloid model is also the identity. This
allows for a similar computation,
but Eq.
\ref{equation-gauss-secti-curva-shape}
doesn't hold
exactly because the
metric is not Riemannian. In fact, the normal
vector has norm $-1$, giving that
hyperbolic space has constant curvature $-1$.
\end{example}

\begin{definition}
\label{definition-total-geode-subma}
A submanifold $M\subset\tilde{M}$ is {\it
totally geodesic} if its second
fundamental form vanishes.
\end{definition}

\begin{lemma}
\label{lemma-totally-geodesic-submanifold}
A submanifold is totally geodesic if and only
if every geodesic of $M$ is a
geodesic of $\tilde{M}$.
\end{lemma}

\begin{proof}
The direct implication is clear: if $\alpha
\equiv 0$ and we take any curve
contained in $M$, it will have zero
acceleration in $M$ and $\tilde{M}$
simultaneously. For the converse we check
that all the eigenvalues
of $\alpha$ are zero. Consider any tangent
vector $v$ to $M$ and let $\gamma_v$
be the geodesic in $M$ with initial velocity
$v$. By hypothesis this geodesic
will also be a geodesic of $\tilde{M}$, so
that $\alpha(v,v)=0$. Diagonalizing
$\alpha$ we see it must vanish identically.
\end{proof}

Lemma \ref{lemma-geode-subma-chara}
says what happens in
general: a geodesic of the submanifold
doesn't need to be a geodesic of the
ambient manifold, but it will always have
acceleration normal to the
submanifold.

\begin{exercise}
\label{exercise-fixed-point-set-isome-tota}
\cite{pet}, Proposition 5.6.5. Let $f:M\to M$
be an isometry. Show that the
connected component of the fixed points of
$f$ is a totally geodesic
submanifold.
\end{exercise}

\begin{proof}
The trick is to show that the eigenvectors of
$df$ are in bijection with the
fixed point set of $f$. For this consider an
eigenvector, a geodesic realising
the vector, map it with $f$ and notice it
must coincide with the original
geodesic because $f$ is an isometry and the
point is a fixed point. This shows
that the fixed point set is a submanifold
since it is parametrized as a linear
subspace under geodesic coordinate charts. It
also shows that the fixes point
set is totally geodesic.
\end{proof}

\begin{definition}
\label{definition-minimal-immersion}
An immersion $f:M\to\tilde{M}$ is called {\it
minimal} if the trace of its shape
operator vanishes at all points of $M$ for
any choice of normal vector field.
\end{definition}

\begin{definition}
\label{definition-minim-free-bound-immer}
A minimal surface $f:\Sigma\to
\mathbb{B}^3\subset \mathbb{R}^3$
is of {\it free boundary} if it is orthogonal
to the boundary of $S^2$.
\end{definition}

\begin{exercise}[Mean curvature and Hessian]
\label{exercise-mean-curvature}
Lista 4, Exercício 8. Let $(M,g)$ be a
Riemannian manifold. Suppose that there
is a function $f \in C^\infty(M)$ such that
$0 \in \mathbb{R}$ is a regular
value of $f$ and let $\Sigma:=f^{-1}(0)$.
\begin{enumerate}
\item Let $N=\frac{\nabla f}{|\nabla f| \in
\mathfrak{X}^\perp(\Sigma)}$. Mostre
que
$$
\left\langle{}\alpha_{\Sigma}(X,Y),N\right\rangle{}=-
\frac{\text{Hess}(f)(X,Y)}{|\nabla
f|},
$$
para todos $X,Y \in \mathfrak{X}(\Sigma)$.

\item Show that the mean curvature of
$\Sigma$ is given by
$$
H_{\Sigma}=-\frac{1}{n}\mathsf{div}
\left(\frac{\nabla
f}{|\nabla f|}\right)
$$
\end{enumerate}
{\it Note.} The notation in item (1)
identifies fields using the differential of
the inclusion map, as is customary.
\end{exercise}

\begin{proof}
\begin{enumerate}
\item One one hand,
 $$
\left\langle{}\alpha(X,Y),\nabla
f\right\rangle{}=-\left\langle{}Y,A_{\nabla f}X\right\rangle{}
$$
by Eq. \ref{equation-secon-funda-form-shape}.
On the other hand,
$$
\text{Hess}(f)(X,Y)=\left\langle{}\tilde{\nabla}_X
\nabla f,Y\right\rangle{}=
\left\langle{}-A_{\nabla f}X+\nabla_X^\perp \nabla
f,Y\right\rangle{},
$$
then multiply by $\frac{1}{|\nabla f|}$.

\item To compute the divergence as the trace
of the gradient we need an
orthonormal frame of $M$. So let $E_i$ be the
completion of
$\frac{\nabla f}{|\nabla f|}$ to an
orthonormal frame of $M$ at a point of
$\Sigma$. Then:
\begin{equation}
\label{equation-diver-norma-gradi}
\begin{aligned}
\text{div}\left(\frac{\nabla f}{|\nabla
f|}\right)
&=\sum_{i=1}^{n-1}
\left\langle{}\tilde{\nabla}_{E_i}\frac{\nabla
f}{|\nabla f|},E_i\right\rangle{}
+\underbrace{\left\langle{}\tilde{\nabla}
_{\frac{\nabla
f}{|\nabla f|}}
\frac{\nabla f}{|\nabla f|},\frac{\nabla
f}{|\nabla f|}\right\rangle{}}
_{=\frac{1}{2}\frac{\nabla f}{|\nabla f|}
\left\langle{}\frac{\nabla f}{|\nabla
f|},\frac{\nabla f}{|\nabla f|}\right\rangle{}=0}
\\
&=\sum_{i=1}^{n-1} \left\langle{}-A_{\frac{\nabla
f}{|\nabla f|}}E_i,E_i\right\rangle{}\\
&=-\text{tr}\left(X \mapsto A_{\frac{\nabla
f}{|\nabla f|}}X\right)\\
&=-nH_{\Sigma}(p)
\end{aligned}
\end{equation}
\end{enumerate}
\end{proof}

Notice that since the Laplacian of a function
$f$ is defined as divergence of
gradient, that is, as the trace of $Z\mapsto
\nabla_Z\nabla f$,
we have for the ambient manifold Laplacian
\begin{equation}
\label{equation-immersion-Laplacian}
\Delta_{\tilde{g}}
f=\text{div}_{\tilde{g}}(\nabla f)
=\text{tr}(Z \mapsto \tilde{\nabla}_Z\nabla
f)
=\sum_{i=1}^{n-1}\left\langle{}\tilde{\nabla}_{E_i}
\nabla
f,E_i\right\rangle{}
+\left\langle{}\tilde{\nabla}_{\nabla f}\nabla
f,\nabla f\right\rangle{}
\end{equation}
where we must add an extra term since a frame
for the ambient manifold
also needs the normal vector. Then we have
$$
\Delta_{\tilde{g}}f-\text{Hess}_{\tilde{g}}
(\nabla
f,\nabla f)
=\sum_{i=1}^{n-1}\left\langle{}\tilde{\nabla}_{E_i}
\nabla
f,E_i\right\rangle{}
$$
If we multiply both sides by
$\frac{1}{|\nabla f|}$, on the right hand
side we
arrive at the mean curvature $H$ by Eq.
\ref{equation-diver-norma-gradi} (after
applying Leibniz rule),
 and on the left hand side…
\begin{equation}
\label{equation-mean-curva-level-sets}
H=\frac{1}{|\nabla
f|}\left(\Delta_{\tilde{g}}
f-\text{Hess}_{\tilde{g}}f
\left(\frac{\nabla f}{|\nabla
f|},\frac{\nabla f}{|\nabla f|}\right)\right)
\end{equation}

{\bf Caveat of Eq.
\ref{equation-mean-curva-level-sets}.} This
is
Lucas Ambrozio's formula from his talk on
Spheres with Minimal Equators (see
Seminars, Section
\ref{seminars-semin-secti-spher-minim}).
Now it's almost perfect except for
putting the $\frac{1}{|\nabla f|}$ inside the
Hessian…

\begin{example}[Equatorial spheres]
\label{example-equatorial-spheres}
Show that equatorial spheres are minimal
hypersurfaces of $S^n$. (See
Seminars, Definition
\ref{seminars-definition-equator}.)
\end{example}

\begin{proof}
To use Eq.
\ref{equation-mean-curva-level-sets} we just
need to compute
the submanifold Laplacian $\Delta_gf$ where
$f(x):=\left\langle{}x,v\right\rangle{}$ for a fixed
$v$ which is just the normal vector defining
the equator, i.e. the equator is
just  $\Sigma:=f^{-1}(0)$, the normal plane
to $v$ intersected with $S^n$. First
we need to compute the gradient of $f$, which
by Eq.
\ref{equation-gradient-in-coordinates} is
just
$$
\nabla
f=g^{ij}\partial_ifE_j=\sum_{i,j}g^{ij}v^iE_j
$$
So that's a problem because in practice we
don't know the coefficients of the
metric, nor do we ever want to compute them.
What this really says is: you don't
know what's the gradient. The next step would
be to compute the Laplacian by
differentiating the gradient
according to Eq.
\ref{equation-immersion-Laplacian}:
$$
\begin{aligned}
\Delta_gf&=\sum_{i=1}^{n-1}\left\langle{}\nabla_{E_i}
\nabla
f,E_i\right\rangle{}
\end{aligned}
$$
Which just takes me back to Eq.
\ref{equation-diver-norma-gradi}. So, need
to compute second fundamental form.

The only easy way I see to do this is use the
fact that the geodesics of the
sphere are well-known by Example
\ref{example-geodesics-of-Sn}. We just need
to
note that the equator is in fact a sphere and
that any curve in the equator is a
geodesic of the equator if and only if it is
a geodesic of the whole $S^n$ (cf.
Lemma
\ref{lemma-totally-geodesic-submanifold}).
But this is obvious by
definition of equator: it's the unitary
vectors in a hyperplane passing through
the origin, and it is equipped with the
metric induced from $S^n$, which in turn
is induced from $\mathbb{R}^{n+1}$, so that
geodesics of the equator are of the
form Eq. \ref{equation-geodesics-of-Sn} by
the same reasons as why the geodesics
of $S^n$ of that form.

The other way to do this, following
\cite{doc}, Chapter VI, Exercises 8 (see also
Exercise \ref{exercise-Clifford-torus}) is
computing the shape operator of the
submanifold with respect to
$\mathbb{R}^{n+1}$ first, which is easy
because
Euclidean space is flat, and then computing
the shape operator with respect to
$S^n$. To do this we just need to compute
inner products of the form
$\left\langle{}A_{N_k}E_i,E_j\right\rangle{}$ since these are
exactly the entries of $A_{N_k}$
in the basis $E_i$ (cf. Eq.
\ref{equation-matrix-of-shape-operator}). But
I will
do this in the next example (maybe some day
do it for $S^n$ too).
\end{proof}

\begin{exercise}[Clifford torus]
\label{exercise-Clifford-torus}
\cite{doc}, Chapter VI, Exercises 2 and 8.
Show that the map
$\mathbf{x}:\mathbb{R}^2\to\mathbb{R}^4$
given by
$$
\mathbf{x}(\theta,\varphi)
=\frac{1}{\sqrt{2}}(\cos\theta,\sin\theta,
\cos\varphi,\sin\varphi),
\qquad(\theta,\varphi)\in\mathbb{R}^2
$$
is an immersion of $\mathbb{R}^2$ in the unit
sphere $S^3\subset\mathbb{R}^4$,
whose image is a torus $T^2$ with zero
sectional curvature in the induced
metric. Finally show that it is a minimal
submanifold of $S^3$. This is called
the {\it Clifford torus}.
\end{exercise}

\begin{proof}
This map is an immersion since the
derivatives of sine and cosine cannot vanish
simultaneously. Its image is a torus since
the map can be seen as a map
from $S^1\times S^1$ to $\mathbb{R}^4$.

It is natural to think that this manifold is
the flat torus from Example
\ref{example-flat-torus}. And it is, because
it is diffeomorphic to a product.
We can use the immersion $\mathbf{x}$ on a
small domain as a local
parametrization of $T^2$.  When
differentiating the parametrization we notice
that the tangent vectors are decomposed as a
pair of vectors tangent to one
circle or to the other. Then the metric
induced from $\mathbb{R}^4$ and then by
$S^3$ is exactly the product metric from
Definition
\ref{definition-product-metric}. By Lemma
\ref{lemma-curvatures-of-product-metric}, the
curvature tensor on $S^1\times
S^1$ vanishes identically since the curvature
tensor of any 1-manifold vanishes
by antisymmetry in the first or two last
entries. This means that this torus is
flat. This says that its first fundamental
form has determinant $-1$ when seen
as a submanifold of $S^3$.

To see it is a minimal submanifold of $S^3$
it suffices to show it is totally
geodesic Lemma
\ref{lemma-totally-geodesic-submanifold}. In
other words, we may
equivalently show that $\alpha\equiv0$ or
that the geodesics of this torus are
geodesics of $\mathbb{R}^n$. However, this
time it is not so immediate what are
the geodesics of $T^2$.

The other way to do this, following
\cite{doc}, Chapter VI, Exercises 8 (see also
Exercise \ref{exercise-Clifford-torus}) is
computing the shape operator of the
submanifold with respect to $\mathbb{R}^4$
first, which is easy because
Euclidean space is flat, and then computing
the shape operator with respect to
$S^3$. To do this we just need to compute
inner products of the form
$\left\langle{}A_{N_k}E_i,E_j\right\rangle{}$ since these are
exactly the entries of $A_{N_k}$
in the basis $E_i$ (cf. Eq.
\ref{equation-matrix-of-shape-operator}).

Strategy: we need the following ingredients:
a point $p$ in the submanifold, a
basis $E_i$ of the tangent space of the
submanifold and a basis $N_k$ of the
normal space. For every basic normal vector
$N_k$ we compute its shape operator
as a matrix. Then we pass to $S^n$. We should
only consider the normal vectors
$N_k$ that are tangent to $S^n$---this way
they are normal to the subsubmanifold
$\Sigma$ as a submanifold of $S^n$. Then the
trace of the shape operator with
respect to those normals should be the mean
curvature.

Differentiating $\mathbf{x}$ as a local
parametrization we see that
$$
e_1=\partial_\theta=(-\sin\theta,\cos\theta,
0,0),\qquad
e_2=\partial_\varphi=(0,0,-\sin\varphi,
\cos\varphi)
$$
are an orthonormal basis of $T_pT^2$, while
$$
n_1=\frac{1}{\sqrt{2}}(\cos\theta,\sin\theta,
\cos\varphi,\sin\varphi),
\qquad
n_2=\frac{1}{\sqrt{2}}(-\cos\theta,-
\sin\theta,\cos\varphi,\sin\varphi)
$$
are an orthonormal basis of $T_p^\perp T^2$.
Only $n_2$ is tangent to $S^3$ at
$p$ since $n_1$ is in fact the position
vector $p$, which is normal to $S^3$.

Now observe that by Eqs.
\ref{equation-secon-funda-form-shape} and
\ref{equation-second-fundamental-form},
$$
\left\langle{}A_{n_2}e_i,e_j\right\rangle{}=\left\langle{}n_k,
\alpha(e_i,e_j)\right\rangle{}
=\left\langle{}n_k,\nabla^{\mathbb{R}^{n+1}}_{e_i}e_j
-\nabla^{T^2}_{e_i}e_j\right\rangle{}
$$
Now
$$
\nabla^{\mathbb{R}^{n+1}}_{\partial_\theta}
\partial_\theta=
-\cos\theta\partial_\theta-
\sin\theta\partial_\theta
$$
so that
$$
\left\langle{}n_2,-\cos\theta\partial_\theta-
\sin\theta\partial_\theta\right\rangle{}
=\frac{1}{\sqrt{2}}
$$
An analogous computation for
$\partial_\varphi$ shows that
$$
A_{n_2}=\frac{1}{\sqrt{2}}\begin{pmatrix}
1&0\\
0&-1
\end{pmatrix}
$$
\end{proof}

Here are some comments after reading
\cite{Jost}, Section 5.4. Let $f:M\to N$ be
an isometric immersion. This allows to use
the vectors
$$
e_\alpha:=f_*\left(\frac{\partial }{\partial
x^\alpha}\right)=
\frac{\partial f^i}{\partial
x^\alpha}\frac{\partial }{\partial f^i},
$$
(where $f^i$ are local coordinates of $N$
near $f(x)$ using that $f$ is an
isometry) as a basis of
$T_{f(x)}N$. Then, by Eq.
\ref{equation-second-fundamental-form},
checking that
$f(M)$ is then a minimal submanifold, i.e.
that the trace of its second
fundamental form vanishes, is to say that
\begin{equation}
\label{equation-isometric-immersion-basis}
(\nabla^{N}_{e_\alpha}e_\alpha)^\perp=0
\end{equation}
where we are implicitly summing over
$\alpha$.

Choosing normal coordinates $\partial_\alpha$
at $M$, using Eq.
\ref{equation-norma-coord-vanis-chris} to
obtain
that the tangent part to the submanifold
vanishes at the point, we get
\begin{equation}
\label{equation-minim-subma-norma-coord}
\begin{aligned}
(\nabla^N_{e_\alpha}e_\alpha)^\perp&=
\nabla^N_{e_\alpha}e_\alpha\\
&=\nabla^N
_{\frac{\partial f^i}{\partial
x^\alpha}\frac{\partial }{\partial f^i}}
\frac{\partial f^j}{\partial
x^\alpha}\frac{\partial }{\partial f^j}\\
&=\frac{\partial^2 f^j}{(\partial
x^\alpha)^2}\frac{\partial }{\partial f^j}
+\frac{\partial f^i}{\partial
x^\alpha}\frac{\partial f^j}{\partial
x^\alpha}
\Gamma^j_{ik}\frac{\partial }{\partial
f^j}\quad \text{for }j=1,\ldots,n.
\end{aligned}
\end{equation}
And it is claimed that in arbitrary
coordinates this is just
\begin{equation}
\label{equation-harmonic-map}
\Delta_Mf^j+\gamma^{\alpha
\beta}(x)\Gamma^j_{ik}(f(x))
\frac{\partial f^i}{\partial
x^\alpha}\frac{\partial f^k}{\partial
x^\beta}=0
\quad \text{for }j=1,\ldots,n.
\end{equation}
(It would be a nice exercise to perform the
passage from normal to arbitrary
coordinates.)
\begin{definition}
\label{definition-harmonic-map}
A map of manifolds satisfying Eq.
\ref{equation-harmonic-map} is
called {\it harmonic}.
\end{definition}
Notice that this is not asking that the
coordinate
functions of $f$ must be harmonic. Also
notice that this says that images of
harmonic isometric immersions are minimal by
definition.

\begin{lemma}
\label{lemma-fibers-of-harmonic-submersion}
Let $f:M\to N$ be a smooth (isometric,
right?) submersion of
Riemannian manifolds. $f$ is
harmonic if and only if its fibers are
minimal submanifolds of $M$.
\end{lemma}

\begin{proof}
Here's my idea: do the same trick as in Eqs.
\ref{equation-isometric-immersion-basis} and
\ref{equation-minim-subma-norma-coord} but
now using
$f^{-1}$ instead of $f$, which we may do
since $f$ is a submersion. Since it is
also an isometry, we get the same result,
again with $f^{-1}$ instead of $f$.
Then the claim becomes: $f$ is harmonic if
and only $f$ its local inverse is
harmonic. I think this might follow from
taking the Laplacian of the composition
$f f^{-1}=\text{Id}$, which must vanish
because the identity is harmonic. A
Leibniz rule should appear (modulo some other
metric term? see Eq.
\ref{equation-Laplacian-of-product}).
\end{proof}

\medskip\noindent

Now we prove that the images of minimal
isometric immersions are critical points
of the volume functional.

\begin{exercise}
\label{exercise-minim-isome-immer-criti}
Consider an isometric immersion $f_0:(M,g)
\to(\tilde{M},\tilde{g})$ and a
vector field $\xi \in f^*T\tilde{M}$.

Define a variation by:
$$
\begin{aligned}
f: (-\varepsilon,\varepsilon)\times M
&\longrightarrow \tilde{M} \\
f(s,p) &=\text{exp}_{p}(s \xi)
\end{aligned}
$$

Then $f_s:=f(s,\cdot):M \longrightarrow
\tilde{M}$ is an immersion (isometric at
$s=0$ by hypothesis). Now we define the
volume functional, which simply
computes the volume of $M_s$:
$$
S(s):= \int_{M}\operatorname{Vol}_{f_s
^*\tilde{g}}
=\int_M f_s^*\operatorname{Vol}_{\tilde{g}}.
$$
Compute $S'(0)$ and show that critical points
of $S$ are immersed submanifolds with
vanishing mean curvature for every normal
vector.
\end{exercise}

\begin{proof}
How to express volume in any coordinate
system?
\[\sqrt{\det (f^*
_sg)_{ij}}dx^1\wedge\ldots\wedge dx^n\]
Let's differentiate:
$$
\begin{aligned}
\frac{d}{ds}S(s)&=\frac{d}{ds}
\int_M\operatorname{Vol}_{f^*_s\tilde{g}}
=\frac{d}{ds}\int_M \sqrt{\det (f^*
_sg)_{ij}}dx^1\wedge\ldots\wedge dx^n\\
&=\int_M \frac{d}{ds}\sqrt{\det (f^*
_sg)_{ij}}dx^1\wedge\ldots\wedge dx^n
\end{aligned}
$$
We must differentiate the square root of the
determinant of a matrix. By
Exercise
\ref{exercise-derivative-of-determinant},
$$
\frac{d}{dt}\det A(t)=
\det A(t) \cdot
\operatorname{tr}\left(A(t)^{-1}\cdot
\frac{d}{dt}A(t)\right)
$$
and using that we see that
$$
\frac{d}{dt}\sqrt{\det A(t)}
=\frac{1}{2}\sqrt{\det A(t)}
\cdot\operatorname{tr}\left(A(t)^{-1}\cdot
\frac{d}{dt}A(t)\right)
$$
So we put $A(s)=(f^*_s\tilde{g})_{ij}$. The
good news is that square root of the
determinant part is exactly the local
coordinate function of
$\operatorname{Vol}_{f_0^*\operatorname{Vol}
_{\tilde{g}}}=\operatorname{Vol}_g$.
Then we only have to integrate that other
function, which hopefully is related
to the mean curvature.

Let's compute $$
\begin{aligned}
\frac{d}{ds}\Big|_{s=0}(f_s^*\tilde{g})_{ij}
&=\frac{d}{ds}\Big|_{s=0}f_s^*\tilde{g}
\left(\partial_i,\partial_j\right)
=\frac{d}{ds}\Big|_{s=0}\tilde{g}
(f^s_*\partial_i,f^s_*\partial_j)\\
&=\tilde{g}\left(\nabla^f_{\partial_s}
f^s_*\partial_i,f^s_*\partial_j\right)
+\tilde{g}
\left(f^s_*\partial_i,\nabla^f_{\partial_s}
f_*^s\partial_j\right)\Big|_{s=0}\\
&=\tilde{g}\left(\nabla^f_{\partial_s}
f^0_*\partial_i,f_*^0\partial_j\right)
+\tilde{g}\left(f^0_*\partial_i,
\nabla^f_{\partial_s}f_*^0\partial_j\right)\\
&=\tilde{g}\left(\nabla^f_{\partial_i}
f^0_*\partial_s,f_*^0\partial_j\right)
+\tilde{g}\left(f^0_*\partial_i,
\nabla^f_{\partial_j}f_*^0\partial_s\right)
\qquad  \text{symmetry lemma} \\
&=\tilde{g}(\nabla^f_{\partial_i}\xi,
f_*^0\partial_j)
+\tilde{g}(f^0_*\partial_i,
\nabla^f_{\partial_j}\xi)\end{aligned}
$$
Where
$\nabla^f$ is the pullback connection. The
point is that we arrived at the
variational field $\xi$. If we assume (for
now) that $\xi$ is normal to $M$ then
upon differentiation we get: $$
\nabla^f_{\partial_i}\xi=-A_{\xi}\partial_i+N
$$
where the normal component $N$ vanishes since
it is orthogonal to the basic
tangent vectors of $M$. We conclude that
$$
\begin{aligned}
\frac{d}{ds}\Big|_{s=0}(f_s^*\tilde{g})_{ij}
&=\tilde{g}\left(-A_{\xi}\partial_i,
\partial_j\right)
+\tilde{g}\left(\partial_i,-A_{\xi}
\partial_j\right)\\
&=\tilde{g}(\xi,\alpha(\partial_i,\partial_j)
)
+\tilde{g}(\alpha(\partial_j,\partial_i),\xi)
\\
&=2\tilde{g}(\xi,\alpha(\partial_i,
\partial_j))\qquad\qquad
\text{$\alpha$ is symmetric}\\
&=-2\tilde{g}(A_{\xi}\partial_i,\partial_j)
\end{aligned}
$$

Now we have to multiply by the inverse
matrix, which is just the inverse of the
metric in $M$ because we are at $s=0$. Then
we take trace and obtain the mean
curvature, defined as the trace of the shape
operator:
$$
\begin{aligned}
\operatorname{tr}\Big(g^{ij}g(A
\partial_j,\partial_k)\Big)
=g^{ij}g(A \partial_j,\partial_i)
=g^{ij}g(A^\ell_j\partial_\ell,\partial_i)
=g^{ij}A^\ell_jg_{\ell i}
=\delta^{i}_jA^j_i
=A^j_j
\end{aligned}
$$
We conclude that
$$
\frac{d}{ds}\Big|_{s=0}S(s)=-\int_M H
\operatorname{Vol}_{g}
$$
\end{proof}

\medskip\noindent

Here are two other important equations in the
theory of isometric immersions.
The {\it Codazzi equation} is
\begin{equation}
\label{equation-Codazzi}
(R_{\nabla^f}(X,Y)Z)^\perp
=(\nabla^\perp_X\alpha)(Y,Z)-(\nabla^\perp_Y,
\alpha)(X,Z)
\end{equation}

Codazzi equation may help us prove… something
cool about bundles I don't quite
remember.

And the {\it Ricci equation} is
\begin{equation}
\label{equation-Ricci}
\left\langle{}R_{\nabla^f}(X,Y)\xi,\eta\right\rangle{}
=\left\langle{}R_{\nabla^\perp}(X,Y)\xi,\eta\right\rangle{}
-\left\langle{}[A_\xi,A_\eta]X,Y\right\rangle{}
\end{equation}

Ricci equation may help prove that the normal
bundle of a high codimension
submanifold is trivial.

\begin{theorem}[Fundamental theorem of
submanifolds]
\label{theorem-funda-theor-subma}
See \cite{pet}, Exercise 3.4.21. Let $M$ be a
simply connected Riemannian
manifold and
$(E^p,\widetilde{\left\langle{}\cdot,\cdot\right\rangle{}},
\tilde{\nabla})$
be a
bundle over $M$ with $\tilde{\nabla}$
compatible with
$\widetilde{\left\langle{}\cdot,\cdot\right\rangle{}}$.
Let $\tilde{\alpha}:\mathfrak{X}(M)\times
\mathfrak{X}(M)\to E$ be a symmetric
tensor. Suppose there is a constant $c$ such
that
$(E,\tilde{\nabla},\tilde{\alpha})$ satisfies
Gauss equation for constant
curvature
\ref{equation-Gauss-constant-curvature},
Codazzi equation
\ref{equation-Codazzi} and Ricci equation
\ref{equation-Ricci}.
Then there exists a unique isometric
immersion $f:M \to \mathbb{Q}_c^{n+p}$ and a
bundle isomorphism $\varphi:E \to
T^\perp_pM$ such that
$\varphi\circ\tilde{\alpha}=\alpha_f$.
\end{theorem}

\medskip\noindent
Now I put some exercises.

\begin{exercise}[Singularities of a Killing
vector field]
\label{exercise-singu-killi-vecto-field}
\begin{reference}
\cite[Exercise 12, Chapter VI]{doc}
\end{reference}
Let $X$ be a Killing vector field on a
Riemannian manifold $M$.
Let $N=\{p \in M: X(p)=0\}$. Show that
\begin{enumerate}
\item If $p\in N$, $V \subset M$ is a normal
neighbourhood of $p$,
and $q \in N \cap V$, then the radial
geodesic segment
joining $p$ to $q$ is contained in $N$.
Conclude that $\gamma \cap V\subset N$.

\item If $p \in N$, there exists a
neighbourhood $V \subset M$
of $p$ such that $V \cap N$ is a submanifold
of $M$.

\item The codimension of any connected
component of $N$ in $M$ is even.
Use that a sphere admitting a
nowhere-vanishing vector field
has odd dimension.
\end{enumerate}
\end{exercise}

\begin{proof}
\begin{enumerate}
\item Consider the flow $\varphi_t$ of $X$
for small $t$.
Since it is an isometry, it must map
geodesics to geodesics.
Since both $p$ and $q$ are in $N$, the two
geodesics
$\gamma$ and $\varphi_t(\gamma)$ intersect at
$p$ and at $q$.
But the exponential is bijective in $V$, so
this cannot happen
unless $X$ vanishes throughout the segment of
$\gamma$ joining $p$ and $q$.
This means that $X$ vanishes throughout
$\gamma$ since
we can take any point in $\gamma$ not between
$p$ and $q$,
apply the flow, which would give another
geodesic but
whose starting segment coincides with
$\gamma$,
and by uniqueness of ODEs such a geodesic
must be $\gamma$ itself.

\item Por indução usando o inciso anterior.

\item Let $p \in N$ and $q$ in a normal
neighbourhood of $p$ in $M$.
Since $N$ is a submanifold we have the
decomposition
$T_qM\simeq T_q N\oplus T_q^\perp N$.
Consider the submanifold
$\tilde{N}:=\text{exp}_q T_q^\perp N$,
and a geodesic sphere contained in this
submanifold.
$X$ cannot vanish along this sphere.
Since $X$ is Killing, it is orthogonal to
such sphere
by $\frac{d}{dt}\left\langle{}X,\gamma'\right\rangle{}=0$
and $\left\langle{}X(q),\gamma'(q)\right\rangle{}=0$.
Thus we have constructed a nowhere-vanishing
vector field
in a hypersphere of $\tilde{N}$,
which shows that $\dim \tilde{N}=\text{codim
}N$ is even.
\end{enumerate}
\end{proof}

\section{Harmonic maps}
\label{section-harmonic-maps}

{\bf Goal.} Understand harmonic maps $T^2 \to
S^3$. Then gauge theory and
algebraic geometry (spectral data).

Let $M$ be a compact Riemann surface, $G$ be
a compact Lie group with a
bi-invariant metric and $f:M\to G$ a smooth
map.

\begin{example}[To keep in mind]
\label{example-torus}
$M=T^2$ and $G=\text{SU}(2)\cong S^3$.
\end{example}

\begin{definition}
\label{definition-connection}
A {\it connection} is a map
$$
\nabla^A:\Omega^{0}(M;G)\to \Omega^{1}(M;E)
$$
where
$\Omega^{p}(M;E)=\Gamma(M,\Lambda^{p}
T^*M\otimes
E)$.
\end{definition}

Consider the pullback bundle of
$(TG,\nabla^{LC})$ over $M$, which we denote
by
$(E,\nabla^A)$.

Locally, we can express
$$
\nabla^A_{\text{loc}}=d+A
$$
where $A$ is the {\it connection matrix}, a
mtrix of 1-forms.

Extend  $\nabla^A$ to an operator
$$
d_A:\Omega^{p}(M;E)\to \Omega^{p+1}(M;E)
$$
by Leibniz rule
$$
d_A(\alpha\wedge\beta=d\alpha\wedge\beta+(-1)
^{|\alpha|}\alpha\wedge
d_A\beta
$$
And we have a {\it curvature} given by
$$
F_A=d^2_A\in\Omega^{2}(M,;\text{End}(E))
$$

\begin{lemma}
\label{lemma-derivative-of-differential}
$d_A(df)=0$
\end{lemma}

\begin{proof}
If $\alpha\in\Omega^{1}(M;E)$ and
$X,Y\in\mathfrak{X}(M)$,
$$
d_A(\alpha)(X,Y)=\nabla_X^A(\alpha(Y))-
\nabla_Y^A(\alpha(X))-\alpha([X,Y])
$$
So for $\alpha=df$… (use that LC is
torsion-free!)
\end{proof}

\begin{definition}
\label{definition-Hodge-star}
The {\it Hodge star operator} is
$$
d^*_A=-*d_A*
$$
where
$$
\begin{aligned}
*: \Omega^{p}(M;E) &\longrightarrow
\Omega^{2-p}(M;E) \\
\alpha\otimes s &\longmapsto *\alpha\otimes s
\end{aligned}
$$
since our manifold is dimension 2.
\end{definition}

\begin{definition}
\label{definition-energy-functional}
$$
E(f)=\frac{1}{2}\int_M|df|^2d\text{Vol}
$$
\end{definition}

\begin{definition}
\label{definition-harmonic-map}
A {\it harmonic map} of smooth manifolds
$f:M\to N$ is a critical point of the
energy functional.
\end{definition}

\begin{lemma}
\label{lemma-criti-point-energ-funct}
The critical points of $E$ satisfy
$$
\tau(f):=\text{tr}_g(\nabla df)=0
$$
where $\nabla df$ is the {\it induced
connection} on $T^*M\otimes E$, i.e.
$$
\nabla(df)(X,Y)=\nabla_X^A(df(Y))-
df(\nabla_X^gY)
$$
\end{lemma}

\begin{lemma}
\label{lemma-harmonic}
$f$ is harmonic if and only if $d_A^*(df)=0$.
\end{lemma}

\begin{proof}
Choose coordinates such that $\nabla^g_XY$
vanishes at the point. Then
$$
\tau(f)=\nabla^A_{e_1}(df(e_1))+
\nabla^A_{e_1}(df(e_2))
$$
Let's compute $d_A^*(df)=-*d_A*(df)$. First
we find that $*df=e^2s_1-e^1s_2$ for
a dual frame $e^i$ and now no longer writing
the tensor product. Then we get
that
$$
\begin{aligned}
d_A(*df)&=de^2\otimes s_1-e^2\nabla^A(s_1)-
\end{aligned}
$$

\end{proof}

\begin{remark}
\label{remark-harmonic-functions}
In the case of a function
$f:\Omega\subset\mathbb{C}\to \mathbb{R}$ we
obtain
the usual Laplacian (Exercise!).
\end{remark}

\section{Riccati equation}
\label{section-Riccati-equation}

Take any point and any tangent vector at that
point. Take the geodesic and take
the geodesic sphere of raduis say $r$ with
center in the point. It's
a hypersurface. Normal vectors are all
multiples of the unit normal vector. But
who's that, the unit normal? Ah, it's the
speed of the geodesic.

So? The normal derivative of any vector field
tangent to the sphere with respect
to the unit normal $\gamma'(r):=\nu$ is…
$$
0=X\left\langle{}\nu,\nu\right\rangle{}=2\left\langle{}\nabla^\perp_X
\nu,\nu\right\rangle{}
$$
that is, normal derivative of $\nu$ with
respect to $X$ is perpendicular to
$\nu$. But normal derivative is normal so
it's actually a multiple of $\nu$,
which makes $\nu$ have norm zero or else
$\nabla^\perp_X \nu=0$.

Now consider everyday Jacobi variation
\ref{proposition-everyday-jacobi-field}
 $f(s,t):=\text{exp}_p(t(v+sw))$ with Jacobi
 field
$J(t)=d_{tv}\text{exp}_p(tw)$. Differentiate
with using the pullback connection
along $f$:
$$
\begin{aligned}
\nabla_{\frac{d}{dt}}^\gamma
J(t)&=\nabla_{\partial_t}^ff_s=
\nabla_{\partial_t}^f
\partial_sf\\
&=\nabla_{\partial_s}^f \partial_t
f=\nabla_{J}^f \gamma' = -A_{\gamma'}f
\end{aligned}
$$
using that by our discussion above the normal
derivative component vanishes. We
conclude that
\begin{equation}
\label{equation-jacob-field-deriv-shape}
J'=AJ
\end{equation}
Finally we differentiate that to obtain
$$
J''=A'J+AJ'=A'J+A^2J
$$
that is,
\begin{equation}
\label{equation-Riccati}
A'+A^2+R_{\gamma'}=0
\end{equation}
since we can let $J(r)$ be any vector
whatsoever and make this construction
work.
\begin{exercise}
\label{exercise-limit-of-shape-operator}
Show that $\lim_{s\to0}sA(s)=\text{Id}$.
\end{exercise}

\begin{proof}[Solution]
By Rafael. First we present a naive
computation. We will use that
$A(t)(J(t))=J'(t)$,  $J(0)=0$ and $J'(0)\neq
0$. Let $\varepsilon>0$ and let
$s>0$ be such that
$$
\frac{\|J(s)-J(0)-sJ'(0)\|}{s}
<\frac{\varepsilon}{2}
$$
and
$$
\|J'(s)-J'(0)\|<\frac{\varepsilon}{2}.
$$
Then
$$
\begin{aligned}
\left\|A(s)(J(s))-\frac{J(s)}{s}\right\|&
=\left\|J'(s)-\left(\frac{J(s)-J(0)}{s}
\right)\right\|\\
&\leq
\|J'(s)-J'(0)\|+\frac{\|J(s)-J(0)-sJ'(0)\|}
{s}\\
&\leq\varepsilon/2+\varepsilon/2=\varepsilon
\end{aligned}
$$
This shows that $A=\frac{1}{s}J$ for small
$s$, but we have differentiated as if
we were in $\mathbb{R}^n$. A similar
computation but picking $s$ such that
$$
\frac{|\left\langle{}J(s),e_i(s)\right\rangle{}-\left\langle{}J(0),
e_i(0)\right\rangle{}
-s\left\langle{}J'(0),e_i(0)\right\rangle{}|}{s}<\varepsilon/
2
$$
and
$$
|\left\langle{}J'(0),e_i(0)\right\rangle{}-\left\langle{}J'(s),e_i(s)
\right\rangle{}|<\varepsilon/2
$$
for a parallel frame $e_i$ along $\gamma$
allows to conclude as desired.
\end{proof}

\section{Hopf-Rinow Theorem}
\label{section-Hopf-Rinow-theorem}

\begin{definition}
\label{definition-non-extendable}
A Riemannian manifold $M$ is {\it extendable}
if there is another Riemannian manifold
$\tilde{M}$
such that $M$ is isometric to a proper open
subset of $\tilde{M}$.
\end{definition}

\begin{theorem}[Hopf-Rinow]
\label{theorem-Hopf-Rinow}
The following are equivalent for any $p \in
M$:
\begin{enumerate}
\item The domain of the exponential map
$\text{exp}_p$ is $T_pM$.
\item Closed and bounded subsets are compact.
\item $(M,d)$ is a complete metric space.
\item $M$ is geodeiscally complete (i.e.
every geodesic is maximally defined in
all of $\mathbb{R}$).
\item Existence of exhaustions.
\end{enumerate}
Moreover, each of the former conditions imply
that there exists a minimizing
geodesic joining any two points.
\end{theorem}

\begin{exercise}
\label{exercise-divergent-curves}
\begin{reference}
\cite[Exercise 3.5, Chapter VII]{doc}
\end{reference}
A {\it divergent curve} in a Riemannian
manifold $M$
is a smooth map $\gamma:[0,\infty) \to M$
such that for any compact $K \subset M$
there exists $t_0\in(0,\infty)$
such that $\gamma(t) \not\in K$ for $t>t_0$.
The length of $\gamma$ is defined by
$$
\ell(\gamma)=\lim_{t \to \infty}
\int_0^t|\gamma'(t)|dt.
$$

Show that $M$ is complete if and only if
the length of any divergent curve is
infinite.
\end{exercise}

\begin{proof}
Suppose $M$ is complete and $\gamma$ is a
divergent curve
with finite length. Choose a ball $B$ of
radius $\ell(\gamma)$
containing $\gamma$. The closure of $B$ is
compact,
but $\gamma$ is contained in it.

For the converse …
\end{proof}

\begin{exercise}
\label{exercise-ray}
A geodesic $\gamma:[0,\infty) \to M$
in a Riemannian manifold $M$ is a {\it ray
starting at $\gamma(0)$}
if it is minimizing between $\gamma(0)$ and
$\gamma(s)$
for all $s \in (0,\infty)$.
Let $M$ be complete, non-compact and $p \in
M$.
Construct a ray starting at $p$.
\end{exercise}

\begin{proof}
Since $m$ is not compact and complete,
we can pick a sequence $p_n \in m\setminus
\bar{b}_n(p)$
and minimizing geodesics $\gamma_n$ from $p$
to $p_n$.
we would like to consider a curve by taking
limit
and show it's minimizing.
however pointwise convergence would not
ensure smoothness of the limit curve.
instead, we consider the normalized initial
velocity vectors
of the curves $\gamma_n$ and take the limit
inside
the compact unit sphere in the tangent space
and apply the exomponential map.
This defines a geodesic $\gamma$ which
is minimizing for any $s \in (0,\infty)$.
Indeed, supposing that our geodesics
are normalized, we have
$$
\begin{aligned}
s &= \ell(\gamma(s))\\
&< s + d(\gamma_n(s),\gamma(s))\\
&< \ell(\gamma_n(s)) +
d(\gamma_n(s),\gamma(s)),
\end{aligned}
$$

\noindent
To conclude consider a curve $\alpha$
connecting $p$ and $\gamma(s)$,
we extend it to a curve $\alpha_n$ connecting
$p$ and $\gamma_n(s)$.
Since $\gamma_n$ is minimizing,
$$
\ell(\gamma_n(s))\leq \ell(\alpha_n)
=\ell(\alpha)+d(\gamma_n(s),\gamma(s)).
$$
The conclusion follows since we showed that
the
length of $\gamma$ up to $\gamma(s)$
is smaller than the length of $\alpha$ up to
$\gamma(s)$
plus a term that converges to zero as $n\to
\infty$.
\end{proof}

\begin{exercise}
\label{exercise-contraction-and-completeness}
Let $M$ and $\overline{M}$ be Riemannian
manifolds
and $f:M \to \overline{M}$ a diffeomorphism.
Suppose $\overline{M}$ is complete and there
is a constant $c>0$
such that
$$
|v|\geq c|d_pf(v)|
$$
for all $p \in M$ and $v \in T_pM$.
Show that $M$ is complete.
\end{exercise}

\begin{proof}
Let $p_n$ be a Cauchy sequence in $M$.
The given inequality is used to prove
that $f(p_n)$ is also Cauchy in
$\overline{M}$,
and thus converges to a point $f(p)$.
Convergence of the preimage follows from
continuity of $f$.
\end{proof}

\begin{exercise}
\label{exercise-local-isometry-hopf-rinow}
\begin{reference}
\cite[Exercise 8, Chapter VIII]{doc}
\end{reference}
Let $M$ be a complete Riemannian manifold,
$\overline{M}$ connected Riemannian manifold
and $f:M \to \overline{M}$ a local isometry.
Suppose that any two points of $\overline{M}$
can be connected be a unique geodesic
of $\overline{M}$.
Show that $f$ is injective and surjective
(and thus a global isometry).
\end{exercise}

\begin{proof}
For injectivety consider the fiber of
any point of $\overline{M}$. Since $M$
is complete, we can connect any two distinct
points in the fiber by a geodesic of $M$.
Since $f$ is a local isometry, it maps
this geodesic to a closed geodesic of
$\overline{M}$, which is impossible:
two nearby points along a closed geodesic
have to geodesics joining them, the short
way and the long way along the geodesic.

For surjectivity suppose there is a point
outside the image of $f$. Connect such a
point with a point inside the image of $f$
with the geodesic of $\overline{M}$
that connects them. At the point in the image
of $f$,  $f^{-1}$ takes this geodesic segment
to a geodesic segment of $M$,
which can be extended
throughout $\mathbb{R}$ since $M$ is
complete. But then the whole geodesic in
$\overline{M}$ would be contained in the
image of the geodesic in $M$ since both
are determined by their initial velocities
and point.
\end{proof}

\begin{exercise}
\label{exercise-bounded-flow-speed}
\begin{reference}
\cite[Exercise 11, Chapter VII]{doc}
\end{reference}
Let $M$ be a complete Riemannian manifold
and let $X$ be a smooth vector field on $M$.
Suppose that $|X(p)|<c$ for all $p \in M$.
Show that the integral curves of $X$
are defined for all $t$.
\end{exercise}

\begin{proof}
The basic idea is:
why if the manifold is complete
every geodesic,
characterized by having constant
velocity,
is maximally defined in $\mathbb{R}$?
The same trick should work in this case.

Suppose that an integral curve $\varphi$
is defined maximally at $(a,b)$.
Choose a sequence converging to $b$,
which is a Cauchy sequence in $\mathbb{R}$.
This gives a Cauchy sequence in $M$
along $\varphi$ by continuity.
By completeness of $M$, $\varphi(t_n) \to p$
for some $p \in M$.
If $X(p) \neq 0$ we can integrate locally
to find a flow $\psi$ at $p$,
which must coincide with $\varphi$ in
the intersection of both curves,
thus extending $\varphi$.
If $X(p)=0$…
\end{proof}

\begin{exercise}
\label{exercise-homogeneous-manifold}
A Riemannian manifold $M$ is
{\it homogeneous} if for any $p,q \in M$
there is an isometry of $M$ that takes
$p$ to $q$. Show that any homogeneous
manifold is complete.
\end{exercise}

\begin{proof}
Let $\gamma$ be any geodesic. Consider
a geodesic ball $B_\varepsilon(\gamma(0))$
of radius $\varepsilon$
about $\gamma(0)$. Let $t \neq 0$
such that $\gamma(t)$ is well defined.
Take $\gamma(0)$ to $\gamma(t)$
with an isometry. Such an isometry
also takes $B_\varepsilon(\gamma(0))$
to $B_\varepsilon(\gamma(t))$.
This shows we can extend $\gamma(t)$
to $\gamma(t+\varepsilon/2)$
for any $t$ in the domain of definition
of $\gamma$, i.e. $\gamma$ is a complete
geodesic.
\end{proof}

\section{Hadamard's theorem}
\label{section-Hadamard}

\begin{theorem}[Hadamard]
\label{theorem-Hadamard}
If $M$ is a complete simply connected
manifold with $K\leq 0$, then the
exponential map is a diffeomorphism.
\end{theorem}

\begin{proof}

\end{proof}

\section{Manifolds of constant sectional
curvature}
\label{section-manif-const-secti-curva}

The following theorem and proof where
explained by Dimitri Korshunov (IMPA) at
Geometric Structures Seminar on July 3, 2025.

{\bf Abstract.} Crystallographic group is a
discrete co-compact subgroup of the
group of isometries of a Euclidean space. A
theorem of Bieberbach says that it
contains a subgroup of translations of finite
index. Or, in other words, any
compact flat Riemannian manifold is a finite
quotient of a flat torus. We will
prove this theorem following a streamlined
geometric argument of Vinberg.

\begin{theorem}[Bieberbach]
\label{theorem-Bieberbach}
In particular, any flat manifold is a
quotient of a torus.
\end{theorem}

\begin{proof}[Sketch of proof]
First we show that it suffices to find only
one pure translation, that is, an
inductive argument shows that if we have a
pure translation, then we can find
others. Perhaps this is the hardest part.

Then we obtain some bounds in operator norm
for the pure translations, this part
looks straightforward.  One of these bounds
is that $\|[A,B]-\text{Id}\|\leq
2\|\text{Id}-A\|\|\text{Id}-B\|$. The other
of the bounds is roughly as follows.
For $g=Ax+a$ and $h=Bx+b$, we have that
$[g,h]=[A,B]x+c$ for some $c$ whose norm
we shall bound by $r$ later.

We use the first bound to show that if two
elements in the group are close to
the identity, then their commutator will be
even closer. This shows that there
is a finite set of small generators
$g_1,\ldots,g_k$. Each of them is of the
form $g_i=A_ix+a_i$.

To arrive at a contradiction suppose that
there are no pure translations to
obtain a contradiction.

Let $r(g)$ be the norm of the translational
part of a transformation $g$. Let
$r$ be the maximum of the translational parts
of the $g_i$.
Let $S:=\{g:r(g)<r\}$. It's a finite set. Let
$g_0$ be minimal in $S$. After
some discussion we show that the linear part
of $g_0$ commutes with the linear
part of $g_i$ for all $i$. In fact, this
implies that $g_0$ itself is in the
center of the group.

Finally we obtain a contradiction with
compactness… simple argument.
\end{proof}

\section{First Variation Formula}
\label{section-first-variation}

\noindent
The notion of variation is beautiful because
it lets us look at a curve by placing it
inside a smooth family of nearby curves. The
important object is not only the family
itself, but the vector field along the
original curve which records, at each time
$t$, the infinitesimal direction in which the
curve is being moved. This is the
variational field.

\begin{definition}
\label{definition-variation}
A {\it variation} of a curve
$\gamma:[0,a]\to M$ is a smooth map
$f:(-\varepsilon,\varepsilon)\times[0,a]\to
M$ such that $f(0,t)=\gamma(t)$. We write
$\gamma_s(t):=f(s,t)$ for the curves in the
variation, and
$f_s=f_*\partial_s$ and $f_t=f_*\partial_t$.
The {\it variational field} of $f$ along
$\gamma$ is
$$
V(t)=f_s(0,t).
$$
\end{definition}

\noindent
The energy of a curve is similar
to the length, but instead of integrating
the norm of the speed,
we integrate the norm squared.
For a given variation, the energy
function computes the energy of each of
the curves in the family.

\begin{definition}
\label{definition-energy}
The {\it energy function} of a variation
$f:(-\varepsilon,\varepsilon)\times[0,a]\to
M$ is
$$
E(s)=\int_0^a\left\langle{}f_t(s,t),f_t(s,t)\right\rangle{}\,dt.
$$
\end{definition}

\noindent
Recall Schwarz inequality from functional
analysis:
$$
\left(\int_0^afg\right)^2\leq
\int_0^af^2\int_0^ag^2
$$
Setting $f\equiv1$ and $g=f_t$ we obtain
\begin{equation}
\label{equation-energy-bounds-length}
\ell(\gamma)\leq aE(\gamma)
\end{equation}

\begin{theorem}[First Variation Formula]
\label{theorem-first-variation-formula}
If  $f:(-\varepsilon,\varepsilon)\times[0,a]
\to M$ is a variation of
the geodesic $\gamma:[0,a]\to M$, then
\begin{equation}
\label{equation-first-variation-formula}
\frac{1}{2}E'(0)=-\int_0^a\left\langle{}V,
\gamma''\right\rangle{}+\left\langle{}V,\gamma'\right\rangle{}|_{0}^a
+\sum_{i=1}^{k-1}\left\langle{}V,\gamma'(t_{i}^-)-
\gamma'(t_i^+)\right\rangle{}
\end{equation}
\end{theorem}

\begin{proof}
If $\gamma$ is smooth,
$$
\begin{aligned}
E'(s)&=
2\int_0^a\left\langle{}f_{ts},f_t\right\rangle{}
=2\int_0^a\left\langle{}f_s,f_{tt}\right\rangle{}-\left\langle{}f_s,
f_t\right\rangle{}|_{0}^a
\end{aligned}
$$
using that
$$
\left\langle{}f_s,f_t\right\rangle{}_t=\left\langle{}f_{st},
f_t\right\rangle{}+\left\langle{}f_s,f_{tt}\right\rangle{}
$$
If $\gamma$ is not smooth, the integral of
$\left\langle{}f_s,f_t\right\rangle{}_t$ must be
done part by part, obtaining the sum in Eq.
\ref{equation-first-variation-formula}.
\end{proof}

\begin{exercise}
\label{exercise-first-var-submfd-dist}
\begin{reference}
\cite[XIX]{doc}
\end{reference}
Let $M$ be a complete Riemannian manifold,
and let $N$ be a closed submanifold of $M$.
Let $q \in M\setminus N$. Show that
there exists a point $p_0 \in N$
such that the distance $d(q,N)$ is
realised at $p_0$, and that
any minimizing geodesic joining $p_0$ to $q$
is orthogonal to $N$ at $p_0$.
\end{exercise}

\begin{proof}
The existence of $p_0$ is a consequence of
completeness of $M$ and closedness of $N$.
By definition of $d(q,N)$ as an infimum, for
each $n \in \mathbb{Z}_+$ we can choose a
point $p_n \in N$ such that
$$
d(q,p_n)<d(q,N)+\frac{1}{n}.
$$
Thus the sequence $(p_n)$ is contained in a
closed metric ball centered at $q$. Since $M$
is complete, Hopf-Rinow implies that this
closed ball is compact. Hence, after passing
to a subsequence, $p_n$ converges to some
$p_0 \in M$. Since $N$ is closed and
$p_n \in N$, we have $p_0 \in N$. Finally,
by continuity of the distance function,
$$
d(q,p_0)=\lim_n d(q,p_n)=d(q,N).
$$

Let $\gamma:[0,a] \to M$ be a minimizing
geodesic from $q$ to $p_0$, parametrized by
arc length. Take $v\in T_{p_0}N$ and choose a
curve $c:(-\varepsilon,\varepsilon)\to N$
with $c(0)=p_0$ and $c'(0)=v$. Extend $v$ to
a vector field $V$ along $\gamma$ such that
$V(0)=0$ and $V(a)=v$, and choose a variation
$f$ of $\gamma$ with variational field $V$,
$f(s,0)=q$, and $f(s,a)=c(s)$.

Since $p_0$ realizes the distance from $q$ to
$N$, every curve $f(s,\cdot)$ joins $q$ to a
point of $N$, so
$$
\ell(f(s,\cdot))\geq d(q,N)=a.
$$
By Schwarz inequality,
$E(s)\geq \ell(f(s,\cdot))^2/a\geq a=E(0)$.
Thus $E'(0)=0$. The first variation formula
gives
$$
0=\left\langle{}V,\gamma'\right\rangle{}|_0^a
=\left\langle{}v,\gamma'(a)\right\rangle{}.
$$
Since $v\in T_{p_0}N$ was arbitrary, $\gamma$
is orthogonal to $N$ at $p_0$.
\end{proof}

\begin{exercise}
\label{exercise-endpoint-variation}
Let $\gamma:[0,a]\to M$ be a smooth curve and
let $c:(-\varepsilon,\varepsilon)\to M$ be a
smooth curve with $c(0)=\gamma(a)$. Show that
there exists a variation
$f:(-\delta,\delta)\times[0,a]\to M$ of
$\gamma$ such that
$$
f(s,0)=\gamma(0),\qquad f(s,a)=c(s).
$$
Moreover, show that its variational field
$V(t)=f_s(0,t)$ may be chosen so that
$V(0)=0$ and $V(a)=c'(0)$.
\end{exercise}

\section{Second Variation Formula}
\label{section-second-variation}

\begin{theorem}[Second Variation Formula]
\label{theorem-second-variation-formula}
Let $f:(-\varepsilon,\varepsilon)\times[0,a]
\to M$ be a variation of a smooth
geodesic $\gamma$ with perhaps piecewise
smooth variation field $V$. Then
\begin{equation}
\label{equation-second-variation-formula}
E''(0)=\int_0^a |V'|^2-\int_0^a \left\langle{}R_{\gamma'}V,V\right\rangle{}+
\left\langle{}\gamma'(t),\nabla_{\partial_s}V(t)
\right\rangle{}|_{0}^a+
\sum_{i}\left\langle{}V(t_i),V'(t^-_i)-V'(t^+_i)
\right\rangle{}
\end{equation}
\end{theorem}

\begin{proof}
$$
\begin{aligned}
E''(s)&
=2\left(\int_0^a\left\langle{}f_s,f_{tt}\right\rangle{}-
\left\langle{}f_s,f_t\right\rangle{}|_{0}^a\right)_s
&=\int_0^a\left\langle{}f_{ss},f_{tt}\right\rangle{}+
\int_0^a\left\langle{}f_s,f_{t
ts}\right\rangle{}
-\left\langle{}f_{ss},f_t\right\rangle{}|_{0}^a-\left\langle{}f_s,
f_{ts}\right\rangle{}|_{0}^a
\end{aligned}
$$
Now substitute
$$
R_{\gamma'}V=R(V,\gamma')\gamma'=f_{t
ts}-f_{tst}
$$
and
$$
\left\langle{}f_s,f_{
ts}\right\rangle{}|_{0}^a=\int_0^a\left\langle{}f_s,f_{
ts}\right\rangle{}_t
=\int_0^a\left\langle{}f_{st},f_{ts}\right\rangle{}+
\int_0^a\left\langle{}f_s,f_{tst}\right\rangle{}
$$
If the variational field is not
differentiable we must be careful when
applying
the fundamental theorem of calculus.
\end{proof}

\begin{definition}
\label{definition-index-form}
The {\it index form} is
\begin{equation}
\label{equation-index-form}
I_a(V,W):=\int_0^a
\left\langle{}V',W'\right\rangle{}-\int_0^a\left\langle{}R_{\gamma'}V,
W\right\rangle{}
\end{equation}
for any $V,W \in \mathfrak{X}_\gamma$.
\end{definition}

\begin{remark}
\label{remark-index-form-alternative-formula}
Using that the connection is metric,
$$
\begin{aligned}
\left\langle{}V,W'\right\rangle{}'&=\left\langle{}V',W'\right\rangle{}+
\left\langle{}V,W^{\prime\prime}\right\rangle{}\\
\implies
\left\langle{}V',W'\right\rangle{}&=\left\langle{}V,W'\right\rangle{}'-
\left\langle{}V,W^{\prime\prime}\right\rangle{}
\end{aligned}
$$
and substituting in the definition we obtain
that
\begin{equation}
\label{equation-index-form-alter-formu}
\begin{aligned}
I_a(V,W)&=\int_0^a\left\langle{}V,W'\right\rangle{}'-\int_0^a
\left\langle{}V,W^{\prime\prime}\right\rangle{}-\int_0^a
\left\langle{}R_{\gamma'}V,W\right\rangle{}\\
&=\left\langle{}V,W'\right\rangle{}|_{0}^a-\int_0^a
\left\langle{}V,W^{\prime\prime}\right\rangle{}-\int_0^a\left\langle{}R_{\gamma'}
W,V\right\rangle{}\\
&=\left\langle{}V,W'\right\rangle{}|_{0}^a-\int_0^a
\left\langle{}V,W^{\prime\prime}+R_{\gamma'}W\right\rangle{}
\end{aligned}
\end{equation}
\end{remark}

\section{Bonnet-Myers Theorem}
\label{section-Bonnet-Myers-theorem}

\begin{theorem}[Bonnet-Myers]
\label{theorem-Bonnet-Myers}
If $M$ is a complete manifold with
$\text{Ric} \geq \frac{1}{r^2}$, then
$\text{diam}M\leq \pi r$. In particular, $M$
is compact.
\end{theorem}

\begin{proof}[Sketch of proof]
Let $p,q \in M$ be points with $d(p,q)>\pi
r$. We will construct a Jacobi
field vanishing at the endpoints, so that
these points are conjugate and thus
any geodesic $\gamma$ joining them cannot be
minimizing.

Let $w_i$ be vectors in $T_pM$ which along
with $\gamma'(0)$ form an orthonormal
basis. Define $W_i$ to be the parallel
transport of $w_i$ along $\gamma$ and
$J_i:=\sin(\pi t/\ell)W_i$.

Computing the second variation formula, which
reduces to the index of $J_i$,
which we may write as in Eq.
\ref{equation-index-form-alter-formu} to
obtain
$$
\begin{aligned}
\frac{1}{2}E''(0)&=-\int_0^\ell\left\langle{}J,J''-
R_{\gamma'}J\right\rangle{}
\approx\int_0^\ell \sin(\pi
t)(1+K(W_i,\gamma'(t))
\end{aligned}
$$
Summing over all indices $i$ we arrive at the
Ricci curvature, which must be
greater or equal than $1/r^2$ by hypothesis.
In comparing with the length of the
curve $\ell>\pi r$, we find that the
integrand must be
negative, which forces one of the indices to
be negative.
\end{proof}

\begin{lemma}
\label{lemma-bonne-myers-impli-compa-unive}
\cite{doc}, Chapter IX, Corollary 3.2. Let
$M$ be a complete Riemannian manifold
with $\text{Ric}\geq \delta>0$. Then the
universal covering of $M$ is compact
and $M$ has finite fundamental group.
\end{lemma}

\begin{proof}

The universal cover with the pullback metric
satisfies the same curvature
hypothesis as $M$, making it compact by
Bonnet-Myers Theorem
\ref{theorem-Bonnet-Myers}. To confirm that
the deck transformation group is
finite notice that if we had an infinite
orbit, we could find a limit point
whose projection to the base would not admit
a regular neighbourhood.
\end{proof}


\begin{exercise}[Exercício 3]
\label{exercise-posit-scala-but-no-posit}
Encontre um exemplo de variedade suave que
admite alguma métrica Riemanniana com
curvatura escalar positiva, mas não admite
uma métrica Riemanniana com Ricci
positivo.
\end{exercise}

\begin{proof}
Considere $S^2\times S^1$. A métrica produto
tem curvatura escalar positiva,
pois é a soma das curvaturas escalares em
cada fator, onde a curvatura escalar
de $S^2$ é constate $1$ quanto que a
curvatura escalar de $S^1$ é constante
zero. Suponha que $S^2\times S^1$ admite uma
métrica com $\text{Ric}>0$. Como
$S^2\times S^1$ é compacta, o fibrado
unitário dela é compacto. Isso significa
que existe uma constante $\delta$ tal que
$\text{Ric}\geq \delta>0$. O
recobrimento universal de $S^2 \times S^1$
com a métrica de
recobrimento satisfaz as mesmas hipóteses de
curvatura, de modo que podemos
aplicar o teorema de Bonnet-Myers para
concluir que é compacto e que o grupo de
transformações de coberta, ou seja, o grupo
fundamental de $S^2\times S^1$ é
finito.
\end{proof}

\section{Weinstein's theorem}
\label{section-Weinstein-theorem}

\begin{theorem}[Weinstein]
\label{theorem-Weinstein}
Let $M^n$ be oriented, compact, $K>0$ and $f
\in \text{Iso}(M)$. Suppose that
if $n$ is even, $f$ preserves orientation,
and if $n$ is odd, $f$ reverts
orientation. Then $f$ has a fixed point.
\end{theorem}

\begin{proof}
       A fórmula da segunda variação é a
       fonte de inspiração, olhe:
$$
E''(0)=\int |\dot V|^2-\int\left\langle{}R_ {\dot
\gamma}V,V\right\rangle{}+\left\langle{}\dot\gamma,
\nabla_{\partial_s}f_s\right\rangle{}|_{0}^a
$$
e pense: o que preciso para tirar algo bacano
dessa formulinha?

De alguma maneira nos damos conta de que isso
tem a ver com os pontos fixos de
uma isometria $f$. Suponha que $f$ não tem
pontos fixos e busquemos uma
contradição. Defina
$$
\begin{aligned}
g: M &\longrightarrow \mathbb{R} \\
g(q) &=d(q,f(q))
\end{aligned}
$$

então como $M$ é compacta e essa função é
contínua, ela tem um ponto mínimo que
chamamos de $p$. Temos que:
$$
0<g(p)=\operatorname{min}g\leq g(q)\qquad
\forall q.
$$

Vamos usar a segunda fórmula da variação para
achar um ponto que avaliado em $g$
fica menor do que $g(p)$. Seja $\gamma$ a
geodésica minimizante entre $p$ e
$f(p)$ parametrizada por comprimento de arco
(que existe porque $M$ compacta
implica completa).

Agora pegue um vetor $v \in T_pM$ unitário e
ortogonal a $\dot \gamma$ (para que
o termo com $R$ na fórmula da seg. var. fique
$K$), transporte ele paralelamente
ao longo de $\gamma$ para obter $V$, e pegue
a variação por geodésicas
\[h(s,t)=\operatorname{exp}_{\gamma(t)}(sV(t)
)=\gamma_{V(t)}(s).\]

Então a 2nda fórmula fica beleza: como o
campo é paralelo, o termo $|\dot V|^2$
morre, o termo da curvatura fica $K$ e o
termo com $\nabla_{\partial_s}f_s$
morre também (porque as curvas verticais são
geodésicas).

Então como estamos em $K>0$, acaba que
$E''(0)<0$, então existe uma vizinhança
de $0$  onde as curvas $f(s,t)$ tem menor
energia que $f(0,t)$ (para toda $s$
nessa  vizinhança). Agora lembre que a
$\ell^2 \leq dE$, e que no caso de
geodésicas minimizantes se da a igualdade.
Então fica que
$$
\frac{1}{d}\ell^2(c^s)\leq
E(s)<E(0)=\frac{1}{d}\ell(\gamma)^2=\frac{1}
{d}d(p,f(p))^2
=\frac{1}{d}(\operatorname{min}g)^2
$$
lado esquerdo. Defina $\beta(s)$ como a curva
``vertical" no tempo 0, ie.
$\beta(s)=f(s,0)$. Queremos mostrar que
$$
\frac{1}{d}d(\beta(s),f(\beta(s))^2\leq
\frac{1}{d}\ell^2(c^s)
$$
Já que por definição, do lado esquerdo temos
$\frac{1}{d}g(\beta(s))^2$, dando
uma contradição com a equação anterior. Para
confirmar essa desigualdade e
concluir a prova só devemos mostrar que a
curva $c^s$ liga $\beta(s)$ e
$f(\beta(s))$, pois a distância entre esses
dois pontos é menor do que o
comprimento de qualquer curva que liga esses
pontos.

Ou seja, queremos que
$$
c^s(0)=\beta(s),\qquad
c^s(d)=f(\beta(s))\qquad \forall s
$$
Ou seja, queremos que
$$
f \circ \beta=\gamma_{V(d)}
$$
Note que $\beta(s)=\gamma_{V(0)}(s)$. Então
derive em $s=0$: \[f \circ
\beta=f\circ\gamma_{V(0)}=\gamma_{V(d)}\iff
f_{*,p}V(0)=V(d)=PV(0)\] onde $P$ é
o transporte paralelo ao longo de $\gamma$ de
$p=\gamma(0)$ a $\gamma(d)$.  Isso
implica que queremos que
\[P^{-1}f_{*,p}V(0)=V(0)\]

Com um desenho muito convincente (ver
Manfredo p. 225, esse agumentinho tá
tranquilo) fica que como $f$ é uma isometria
\[f_{*,p}\gamma'(0)=\gamma'(d)\]
Decorre daí que
$P^{-1}f_{*,p}\gamma'(0)=P^{-1}\gamma'(d)=
\gamma'(0)$.
Então
podemos fixar nossa atenção em
$$
A:=P^{-1}f_{*,p}|_{\gamma'(p)^\perp}\in
\mathsf{O}(n-1)
$$
Nosso objetivo é analisar o que acontece
quando $A$ tem um ponto fixo, $V(0)$,
pois essa é uma condição que queremos que
seja verdade para que $f\circ \beta =
\gamma_{V(d)}$.

Agora segue o argumento dos eigenvalores que
diz assim. Defina
$A:=P^{-1}f_{*,p}$. Suponha que $A$ tem um
ponto fixo. Os autovalores vem dados
assim:
$$
\operatorname{Spec}A=\{\underbrace{\lambda_1,
\overline{\lambda_1},\lambda_2,
\overline{\lambda_2}}_{\text{par}
},-1,\ldots,-1,1,\ldots,1\}
$$
Ou seja, os eigenvalores que não são 1 ou
$-1$ não afetam o determinante. Daí
concluímos que
\begin{itemize}
\item Se $n-1$ é ímpar, $f$ preserva
orientação. Por que? Deixa os complexos do
lado que não importam. Então pensa: se só
tenho $-1$s, então tenho uma
quantidade ímpar de $-1$s porque a dimensão é
ímpar. Respira. Nesse caso $f$
inverte orientação. Mas: $f$ deve ter pelo
menos um $1$! Então a quantidade de
$-1$s  é par. Então não: $f$ preserva
orientação.
\item se $n-1$ é par, $f$ reverte orientação.
\end{itemize}
E essas são as condições que precisamos para
que $A$ tenha um ponto fixo e tudo
funcione. Então bota isso aí na hipótese do
teorema.
\end{proof}

\section{Synge's theorem}
\label{section-Synge-theorem}

\begin{theorem}[Synge]
\label{theorem-Synge}
Let $M^n$ be compact and $K>0$. Then
\begin{enumerate}
\item If $n$ is even,
$$
\pi_1(M)=\begin{cases}
	1\qquad & \text{if $M$ is orientable}\\
	\mathbb{Z}_2\qquad & \text{if $M$ is not
	orientable}
\end{cases}
$$
\item If $n$ is odd, then $M^n$ is
orientable.
\end{enumerate}
\end{theorem}

\begin{proof}
First suppose $n$ is even and $M$ is
orientable. Since $M$ is compact we must
have that $K\geq \delta>0$. Now consider the
universal cover $\tilde{M}$. Then
it also satisfies the curvature bound (with
the pullback metric),
making it compact as well by Bonnet-Myers
\ref{theorem-Bonnet-Myers}.
Now pick a deck transformation, which must be
orientation-preserving if we choose the
orientation on $\tilde{M}$ (simply
connected is always orientable) making the
projection orientation-preserving.
Then $f$ must have a fixed point by
Weinstein's theorem \ref{theorem-Weinstein},
which applies since we have shown that
$\tilde{M}$ is compact.

However, deck transformations cannot have
fixed points unless they are the
identity: this is a consequence of the Unique
Lifting Property
 \ref{proposition-unique-lifting-property},
 which is a
slightly more general statement than the
uniqueness of curve lifts that we
used for Klinenberg's lemma exercise
\ref{exercise-l7-1}.
More exactly: given a covering space and two
maps from any space into the base,
these coincide if they agree at one point
only.

If $M$ is not orientable we can show that the
two-sheeted orientable cover of
$M$ is its universal covering by the same
arguments as above. Then the group of
deck transformations is the group of two
elements $\mathbb{Z}_2$ since the
covering is two-sheeted.

Now suppose $n$ is odd and consider the
two-sheeted orientable cover of $M$.
Then by Weinstein's theorem we can say that
no deck transformation can be
orientation reversing. But since $M$ is
non-orientable, the single deck
transformation generating the group of deck
transformations, which is
$\mathbb{Z}_2$, is not orientation
preserving: if it was, then the quotient of
the two-sheeted orientable covering by the
action of the deck transformations
would induce an orientation on $M$.
\end{proof}

\begin{exercise}
\label{exercise-posit-ricci-but-no-posit}
Encontre um exemplo de variedade suave que
admite uma métrica Riemanniana com
Ricci positivo, mas não admite uma métrica
Riemanniana com curvatura seccional
positiva.
\end{exercise}

\begin{proof}
Considere $S^3\times S^1$. A métrica produto
tem $\text{Ric}>0$, já que ele é a
soma dos tensores de Ricci em cada fator de
acordo ao Exercício
\ref{exercise-curvatures-in-product} (a
curvatura de Ricci de $S^3$ é 1,
enquanto que a curvatura escalar de $S^1$ é
zero). Agora suponha que existe uma
 métrica com $K>0$. Então pelo teorema de
 Synge \ref{theorem-Synge}, como
$S^3\times S^1$ é compacta, de dimensão par e
orientável, o grupo fundamental
dela deveria ser trivial, mas não é o caso.
\end{proof}

\section{Index lemma}
\label{section-index-lemma}

\begin{lemma}[Index]
\label{lemma-index}
\cite{doc} Chapter X, Section 2, Lemma 2.2.
Let $\gamma:[0,a]\to M$ be a
geodesic without conjugate points (this means
$\gamma(t)$ is not conjugate to
$\gamma(0)$ on $(0,a]$). Let
$J\in\mathfrak{X}_\gamma^J$ with
$\left\langle{}J,\gamma'\right\rangle{}=0$ and let $V$ be a
piecewise differentiable vector
field along $\gamma$ with
$\left\langle{}V,\gamma'\right\rangle{}=0$. Suppose
$J(0)=0=V(0)$ and
that there is  $t_0\in(0,a]$ such that
$J(t_0)=V(t_0)$. Then
$$
I_{t_0}(J,J)\leq I_{t_0}(V,V)
$$
and equality holds if and only if $V=J$ on
$[0,t_0]$.
\end{lemma}

\section{Rauch's comparison theorem}
\label{section-Rauch}

We motivate Rauch's comparison theorem with a
result from ordinary
differential equations.

\begin{theorem}[Sturm]
\label{theorem-Sturm}
Suppose $f,\tilde{f},K,\tilde{K}:[0,a]\to
\mathbb{R}$ are smooth functions
satisfying the following conditions.
$\tilde{f} > 0$ in $(0,a]$,
$f(0)=\tilde{f}(0)=0$,
$f'(0)=\tilde{f}'(0)$ and
$$
f''+Kf=0,\qquad
\tilde{f}''+\tilde{K}\tilde{f}=0.
$$
If $K\leq \tilde{K}$, then $f/\tilde{f}$ is
non decreasing. Moreover, if there
is any point $t \in [0,a]$ where
$f(t)=\tilde{f}(t)$, then $f \equiv
\tilde{f}$
and $K \equiv \tilde{K}$.
\end{theorem}

\begin{proof}
Since we want to show $f/\tilde{f}$ is non
decreasing we are interested in
computing the sign of the derivative of this
quotient. We thus calculate
$f'\tilde{f}-f \tilde{f}'$. We only have
information about second derivatives,
so we differentiate to obtain
$$
(f'\tilde{f}-f
\tilde{f}')'=(\tilde{K}-K)f\tilde{f}
$$
We may integrate to find that
\begin{equation}
\label{equation-Sturm}
f'\tilde{f}-f \tilde{f}'=\int_0^t
(f'\tilde{f}-f \tilde{f}')'
=\int_0^t (\tilde{K}-K)f\tilde{f}
\end{equation}
If we show that $f\geq 0$ we get our result.
So suppose that $t_0$ is the
smallest number $t_0 \in (0,a]$ where $f$ is
zero. Then $f'(t_0)<0$ since before
that $f$ is positive, which in turn is
because $f'(0)=\tilde{f}'(0)$ and
$\tilde{f}$ is strictly positive. But this
contradicts Equation
\ref{equation-Sturm} since the integrand was
positive up to $t_0$.

Now suppose that $f(t_0)=\tilde{f}(t_0)$ for
some $t_0 \in (0,a]$. Then Equation
\ref{equation-Sturm} vanishes at $t_0$, and
we conclude that $K \equiv
\tilde{K}$ in $[0,t_0]$. Somehow this implies
that $f$ also equals $\tilde{f}$…

Here's ChatGPT in action: define
$h(t):=f(t)-\tilde{f}(t)$. Notice that on
$(0,t_0]$ it satisfies $h''+Kh=0$ with
$h(0)=h'(0)=0$. But so does constant zero
function, hence  $h\equiv0$. To extend to all
of $[0,a]$ use uniqueness of
ODEs again.
\end{proof}

\begin{theorem}[Rauch]
\label{theorem-Rauch}
Let $\gamma$ be a geodesic of $M$ and
$\tilde{\gamma}$ a geodesic of
$\tilde{M}$. Suppose that $J \in
\mathfrak{X}^J_\gamma$ and $\tilde{J}\in
\mathfrak{X}^J_{\tilde{\gamma}}$ satisfy the
following conditions. $J(0)=0$,
$\tilde{J}(0)=0$, $|J'(0)|=|\tilde{J}'(0)|$
and
$\left\langle{}J,\gamma'\right\rangle{}=\left\langle{}\tilde{J},
\tilde{\gamma}'\right\rangle{}$.
Also,
$$
J''+K J=0,\qquad
\tilde{J}''+\tilde{K}\tilde{J}=0
$$
If $\tilde{\gamma}$ has no conjugate points
and $\tilde{K}\geq K$, then
$|J|/|\tilde{J}|$ is nondecreasing. Moreover,
if $|J|$ and $|\tilde{J}|$
coincide in one point, then $|J|\equiv
|\tilde{J}|$ and $K \equiv \tilde{K}$.
\end{theorem}

\begin{proof}
Let $f:=|J|^2$ and
$\tilde{f}:=|\tilde{J}|^2$. I think the main
difference with
Sturm is that the condition on Jacobi fields
is not a condition on the
derivatives of the functions $f$ and
$\tilde{f}$ we just defined. Note that
$$
f'=2\left\langle{}J',J\right\rangle{},\qquad
J''=2(\left\langle{}J',J'\right\rangle{}+\left\langle{}J'',J\right\rangle{}).
$$
While we may substitute $-KJ$ on the second
bracket, we still have the first
one. Perhaps this is the reason to define
$U(t):=J(t)/|J(r)|$ and
$\tilde{U}(t):=\tilde{J}(t)/|\tilde{J}(r)|$.
But wait, what's that $r$? There's
a trick in this proof.

Notice that $U$ and $\tilde{U}$ are Jacobi
fields and $|U(r)|=|\tilde{U}(r)|$.
\end{proof}

\begin{proposition}
\label{proposition-contraction}
Sejam $p \in M$, $\tilde{p} \in \tilde{M}$ e
$i:T_pM \to T_{\tilde{p}}\tilde{M}$
uma isometria. Então $f:=\text{exp}^{-1}_p
\circ i \circ
\text{exp}_{\tilde{p}}$ é uma contração
métrica, i.e. $|f_*w|\leq |w|$
para todo $w \in T_pM$.

Mas ainda, se $\text{exp}_p^{-1}$ está
definida numa bola totalmente
convexa, $f$ é uma contração métrica.
\end{proposition}

\begin{proof}
Write an abitrary $w\in T_p M$ as a Jacobi
variational field of the variation
$f(s,t)=\text{exp}_p(t(v+sw)$. The resulting
field on $\tilde{M}$ is a
Jacobi field with respect to $iw$ using that
$i$ is an isometry. The conditions
of Rauch theorem are satisfied and the
desired inequality is obtained.

For the metric result we just integrate along
a curve realising the distance and
compute.
\end{proof}

\section{Moore's theorem}
\label{section-Moore}

\begin{theorem}[Moore]
\label{theorem-Moore}
Let $M^n \subset \tilde{M}^{n+p}$ be a
compact submanifold of a Hadamard
manifold $\tilde{M}$. If $K\leq \tilde{K}$,
then $p\geq n$.
\end{theorem}

\section{Focal points}
\label{section-focal points}

\begin{definition}
\label{definition-focal-point}
Let $M$ be a submanifold of $\tilde{M}$ and
$p$ a point not in $M$. We say $p$
is a {\it focal point} of $M$ if there is a
geodesic $\gamma$ orthogonal to $M$
and a Jacobi field
$J\in\mathfrak{X}^J_\gamma$ such that
$$
J'(0)\in T_{\gamma(0)}M,\qquad
J'(0)+A_{\gamma'(0)}J(0)\in
T_{\gamma(0)}^\perp M
$$
\end{definition}

\begin{proposition}
\label{proposition-focal-point-singu-norma}
The focal points of $M$ are the singularities
of $\text{exp}|_{T^\perp M}.$
\end{proposition}

\section{Morse index theorem}
\label{section-morse-index}

\begin{definition}[Index of index form]
\label{definition-index-index-form}
$$
i(I_a):=\text{sup}\{\dim L \subset
\mathfrak{X}_\gamma:I_a |_{L \times L}<0\}
$$
\end{definition}

For the proof of Morse Index Theorem
\ref{theorem-Morse-index} we will use the
following objects
\begin{equation}
\label{equation-index-theorem-spaces}
\begin{aligned}
\mathcal{V}_a&=\{V \in
\mathfrak{X}_\gamma:\text{d.p.p., }V(0)=0,
V(a)=0\}\\
\mathcal{V}_t^+&:=\{V \in
\mathcal{V}_t:V(t_i)=0\}\\
\mathcal{V}_t^-& :=\{V \in
\mathcal{V}_t:V|_{[t_j,t_{j+1}]}\in
\mathfrak{X}_{\gamma|_{[t_j,t_{j+1}]}}^J\}
\end{aligned}
\end{equation}

\begin{proposition}
\label{proposition-kernel-of-index-form}
\cite{doc}, Chapter X, Proposition 2.3. The
kernel of the index form i
$$
\Ker I_t:=\{V \in
\mathcal{V}_t:I_t(V,W)=0\forall
W\in\mathcal{V}_t\}=
V_t \cap \mathfrak{X}_\gamma^J
$$
that is, it is composed of smooth Jacobi
fields along $\gamma$.
\end{proposition}

\begin{proof}[Sketch of proof]
The point here is to show that if $V \in
\Ker$ then it is (1) a Jacobi
 field in each interval where $V$ is smooth,
 and (2) smooth in the whole
 interval.

We use the fact that $V$ is in the kernel to
show that (1) the norm of
$V-R_{\gamma'}V$ is zero, and (2) the left
and right derivatives of $V$
coincide. The first part involves multiplying
by a positive function $f$ that
vanishes at the singular points; this will
make the ``singular part" of the
index form, i.e. the sum, vanish, so that we
get only with the integral part,
which is the norm of the Jacobi equation,
which will vanish since we are in the
kernel. The second part I'm not sure---why do
we need to show that $V$ is
differentiable, or $C^2$ (since it's the
solution of an ODE).
\end{proof}

\begin{lemma}
\label{lemma-index-form-degen-iff-conju}
\cite{doc}, Chapter X, Corollary 2.4. The
index form is degenerate if and only
if the initial and final points of $\gamma$
are conjugate. In this case, the
nullity, i.e. dimension of $\Ker I_a$
coincides with the nullity of $\gamma(a)$
as a conjugate point of $\gamma(p)$, that is,
the dimension of the space of
Jacobi fields vanishing at the endpoints.
\end{lemma}

\begin{proof}
Just by Proposition
\ref{proposition-kernel-of-index-form} and by
the fact that
$\mathcal{V}_a$ is defined as the smooth
vector fields vanishing in the
endpoints.
\end{proof}

\begin{proposition}
\label{proposition-V-is-direct-sum}
\cite{doc}, Chapter XI, Proposition 2.5.
$\mathcal{V}_a=\mathcal{V}^+ \oplus
\mathcal{V}^-$.
\end{proposition}

\begin{proof}
Assume that the partition
$0<t_1<\ldots<t_k=a$ of $[0,a]$ is such that
every
segment $\gamma|_{[t_{i-1},t_i]}$ is
contained in a totally geodesic
neighbourhood, so that no such segment can
have conjugate points.

Let $V \in \mathcal{V}_a$ and pick a field $J
\in \mathcal{V}^-$ defined as the
piecewise Jacobi field with pairwise boundary
conditions $J(t_j)=V(t_j)$. The
assumtion that $\gamma|_{[t_{i-1},t_i]}$ has
no conjugate points ensures the
uniqueness of such $J$. Then
$V-J\in\mathcal{V}^+$ and we obtain the
result.
\end{proof}

\begin{theorem}[Morse Index]
\label{theorem-morse-index}
O índice $i(I_t)$ é finito e igual ao número
de pontos conjugados em $[0,t)$ ao
longo de $\gamma$.
\end{theorem}

\begin{proof}
\begin{enumerate}
\item Define
$$
\begin{aligned}
\mathcal{V}_a&=\{V \in
\mathfrak{X}_\gamma:\text{d.p.p. }V(0)=0,
V(a)=0\}\\
\mathcal{V}_t^+&:=\{V \in
\mathcal{V}_t:V(t_i)=0\}\\
\mathcal{V}_t^-& :=\{V \in
\mathcal{V}_t:V|_{[t_j,t_{j+1}]}\in
\mathfrak{X}_{\gamma|_{[t_j,t_{j+1}]}}^J\}
\end{aligned}
$$
\item Note que
$$
\mathcal{V}_t=\mathcal{V}_t^+ \oplus
\mathcal{V}_t^-
$$
\item Note que
$$
\mathcal{V}_t^-=\bigoplus_{j=1}^{i-1}
T_{\gamma(t_j)}M
$$
porque para qualquer escolha de vetores
 $v_1 \in T_{\gamma(t_1)}M,\ldots,v_{i-1}\in
 T_{\gamma(t_{i-1})}M$ existe
 um único campo $J \in \mathcal{V}^-$ tal que
 $J(t_j)=v_j$, por ser a solução
 da EDO com condições de contorno.
\item Usamos Proposition
\ref{proposition-kernel-of-index-form}.
\item …
\end{enumerate}
\end{proof}



\begin{theorem}[Jacobi's theorem]
\label{theorem-Jacobi-theorem}
Se $\gamma$ é minimizante, então não pode ter
pontos conjugados. More generally,
if $\gamma:I \to M$ is a geodesic and $
\gamma(a)$ is conjugate to $\gamma(0)$
through $\gamma$ then for every $\delta>0$,
$I_{a+\delta}\not \geq 0$.
\end{theorem}

\begin{proof}
Se $\gamma$ é minimizante, a segunda fórmula
da variação é positiva para todo
campo em $\mathcal{V}=\{V\in
\mathfrak{X}_\gamma:V \text{ d.p.p. },
V(a)=0,V(b)=0\}$. Ou seja, a forma do índice
é uma forma bilinear positiva, e
portanto tem signatura zero, e pelo Teorema
do Índice de Morse
\ref{theorem-Morse-index} o número de pontos
conjugados é zero.
\end{proof}

It is possible to prove this theorem without
using Morse Index Theorem
\ref{theorem-Morse-index} by constructing a
variation where the index form is
negative, see Exercise \ref{exercise-l7-4}.

\section{Cut locus}
\label{section-cut-locus}

\begin{definition}
\label{definition-cut-point}
O {\it cut point} de  $p$ ao longo de
$\gamma$ é $\gamma(\rho(\gamma))$
onde
$$
\rho(\gamma)=\text{sup}\{t:\gamma|_{[0,t]}
\text{
é minimizante}\}
$$

\end{definition}

\begin{proposition}
\label{proposition-cut-point-chara}
\cite{doc} Chapter XIII, Proposition 2.2. Let
$\gamma$ be a minimizing geodesic joining $p$
and $q$. Then $q$ is the cut
point of $p$ if and only if either of the
following conditions hold:
\begin{enumerate}
\item $q$ is the first conjugate point of $p$
along $\gamma$.
\item There is another, distinct, minimizing
geodesic $\sigma$ joining $p$ and
 $q$.
\end{enumerate}
\end{proposition}

\begin{proof}[Sketch of proof]
($\implies$) Since $q:=\gamma(r)$ is the cut
point of $p$, we know that for every point
$\gamma(r+\varepsilon)$ along $\gamma$ after
$q$ there must be some other geodesic
$\gamma_\varepsilon$ which realizes the
distance between $p$ and
$\gamma(r+\varepsilon)$

Taking the limit as $\varepsilon \to 0$ we
obtain another geodesic
(because we are taking limit inside $S^n
\subset T_pM$, so we get convergence;
and because ``things will behave well under
the limit")
which puts us in the second condition of the
theorem,
unless the limit geodesic is the same as
$\gamma$ setwise.

Now there's a critical observation:
the two geodesics $\gamma$ and
$\gamma_\varepsilon$ meet at the point
 $\gamma(r+\varepsilon)$, so you can write
 $\gamma_\varepsilon(r+\delta(\varepsilon))=
 \gamma(r+\varepsilon)$
 for some function $\delta(\varepsilon)$
 (that could be negative, doesn't
 matter).
 But since $\gamma_\varepsilon$ is the
 minimizing
 geodesic from $\gamma(0)$ to this point,
 then
 $|\delta(\varepsilon)|<\varepsilon$.

Right, the point is that defining
$$
u_\varepsilon:=(r+\varepsilon)v,\qquad
u'_\varepsilon:=r+\delta(\varepsilon))
v_\varepsilon
$$
we see that $\text{exp}_p$ is not injective,
since these two vectors
hit that point.

This being true for every $\varepsilon>0$,
there is no neighbourhood of $vr$
where $\text{exp}_p$ is injective---much less
a diffeomorphism.
So it cannot be a local diffeomorphism, and
we know that critical points of
$\text{exp}_p$ are conjugate points of $p$.
(Remember: if
$\text{exp}_p$ has a critical point, its
differential has kernel,
and the map $d_{tv}\text{exp}_pw$, for $w$ in
the kernel, gives
a Jacobi field. See Proposition
\ref{proposition-critical-points-exp})
\end{proof}

\begin{exercise}
\label{exercise-diametros}
Calcule o diâmetro de $S^2$, $\mathbb{T}^2$,
$\mathbb{R}P^{2}$.
\end{exercise}

\begin{proof}[Solution]
O diâmetro de $S^2$ pode ser calculado via o
teorema de Bonnet Myers
\ref{theorem-bonnet-myers}: nenhuma geodésica
é minimizante depois de atingir
comprimento $\pi r$, e temos uma geodésica
que atinge esse comprimento: qualquer
uma!

O diâmetro de $\mathbb{T}^2$ é $1$. Isso é
por simples geometria
euclidiana: é o diâmetro do cubo! Por
definição, a métrica de  $\mathbb{T}^2$ é
a induzida pela projeção quociente.

O diâmetro de $\mathbb{R}P^{2}$ é $\pi/2$.
Qualquer geodésica que liga dois
pontos a distância $\pi/2$ é minimizante,
pois estamos na métrica esférica e
podemos pegar cartas esféricas onde dois
pontos a essa distância estão contidos.
Por outro lado, se tivéssemos dois pontos a
distância maior, a geodésica
esférica que percorre o ponto antípoda ao
inicial, chega no ponto final mais
rapidamente que inicial; portanto geodésicas
de comprimento maior que $\pi/2$
não são minimizantes.
\end{proof}

\begin{lemma}
\label{lemma-cut-point-reflexive}
If $q$ is the cut point of $p$ through
$\gamma$, then $p$ is the cut point of
$p$ through $\gamma$.
\end{lemma}

\begin{lemma}
\label{lemma-not-cut-locus-minimizing}

\end{lemma}

\begin{proposition}
The function
$$
\begin{aligned}
	\rho: T^1_p &\longrightarrow [0,+\infty] \\
	v &\longmapsto \text{cut point along
	}\gamma_v
\end{aligned}
$$
is continuous.
\end{proposition}

\begin{lemma}
\label{lemma-cut-locus-closed}
Because it's the image of the unit sphere in
$T_pM$, which is closed.
\end{lemma}

\begin{lemma}
\label{lemma-cut-locus-topology}

\end{lemma}

\begin{lemma}
\label{lemma-outsi-cut-locus-exist-minim}
\cite{doc} Chapter XIII, Corollary 2.8. If $q
\in M\setminus C_m(p)$, there exists a
 unique minimizing geodesic joining
$p$ and $q$.
\end{lemma}

\begin{proof}
By Hopf-Rinow \ref{theorem-Hopf-Rinow}, we
know that
there exists a minimizing geodesic joining
any two points. If there was more
than one such geodesic, by
Proposition \ref{proposition-cut-point-chara}
we get that $q$ is the cut point of  $q$
along either of these geodesics,
a contradiction.
\end{proof}

The former Lemma
\ref{lemma-outsi-cut-locus-exist-minim}
shows that the exponential map $\text{exp}_p$
 is injective outside the cut locus of $p$.
This motivates the following definition.

\begin{definition}
\label{definition-injectivity-radius}
The {\it injectivity radius} of a manifold
$M$ is
$$
i(M)=\text{inf}\{d(p,C_m(p)):p \in M\}
$$
\end{definition}

\section{Bishop-Gromov Theorem}
\label{section-bishop-gromov}

\begin{theorem}[Bishop-Gromov]
\label{theorem-Bishop-Gromov}
Let $p \in M$ and $0\leq t \leq
i(p)=d(p,C_m(p)$. Denote $B_t(p)$ the ball of
raduis $t$ centered in $p$ and $B_{t,k}$ the
ball of radius $t$ in a space of
constant curvature $k$.

If $\text{Ric}\geq k$, then
$\text{Vol}(B_t(p))/\text{Vol}(B_{t,k})$ is a
non
increasing function.

Moreover, if there are numbers  $0<s<r\leq
i(p)$ such that
$$
\frac{\text{Vol}(B_s(p))}{\text{Vol}(B_{s,k}}
=
\frac{\text{Vol}(B_r(p))}{\text{Vol}(B_{r,k})
}
$$
then $B_r(p)=B_{r,k}$ isometrically.
\end{theorem}

\begin{proof}
Use the exponential map a parametrization of
the geodesic sphere with radius
$r$ and centre in $p\in M$. (This
parametrization works at least locally, but
why
have we said in lecture that ``it's a global
chart''?) Then
$$
\text{Vol}(S^n_r(p))=\int_{S^n_r(p)}
\text{Vol}_{S^n_r(p)}=
\int_{S^n_r(0)}\text{exp}_p^*\text{Vol}
_{S^n_r}=
\int_{S^n_r(0)}|\det d_{rv}
\text{exp}_p|\text{Vol}_{S^n_r(0)}
$$
This determinant can be given by a basis of
tangent vectors to $S_r(0)$, each of
which yields a Jacobi field via Proposition
\ref{proposition-everyday-jacobi-field}.
These Jacobi fields are the columns
of the matrix $\mathbb{J}$, which is the
Jacobian matrix of the exponential we
are interested in computing.

Now we want to differentiate this with
respect to $r$ because we aim to apply
Sturm's theorem \ref{theorem-Sturm}. By
Exercise
\ref{exercise-derivative-of-determinant}, we
know that the derivative of the
determinant of a one-parameter family of
invertible matrices $\mathbb{J}(r)$ is
given by $\det
\mathbb{J}(r)\text{tr}(\mathbb{J}^{-1}(r)
\mathbb{J}'(r))$.

By the discussion in Section
\ref{section-Riccati-equation}, we know that
$AJ=J'$ where $A$ is the shape operator with
respect to the unit normal
$\gamma'(r)$ and $J$ is a Jacobi field
defined by
\ref{proposition-everyday-jacobi-field}.
This gives $A=J'J^{-1}$.

Thediagonal entries of the matrix
$\mathbb{J}$ are vectors of this kind, so
that
 its trace, which depends only on the
 diagonal by Exercise
\ref{exercise-trace-indep-coord}, is
precisely the trace of
the shape operator $A$, which we have defined
as $H/(n-1)$ where
$H$ is called the mean curvature.

This data translates to the following
equation
\begin{equation}
\label{equation-Bishop-Gromov-proof1}
\det\mathbb{J}'=\det\mathbb{J}H
\end{equation}
Now recall that any self-adjoint operator can
be expressed uniquely as a
multiple of the identity plus a symmetric
traceless matrix (ref?). This allows
to write
\begin{equation}
\label{equation-Bishop-Gromov-proof2}
A=H\text{Id}+A_0
\end{equation}
with $A_0$ symmetric and traceless.
Differentiating Eq.
\ref{equation-Bishop-Gromov-proof2} and using
that the shape operator satisfies
Riccati equation \ref{equation-Riccati} we
shall obtain that
$$
H'+H^2+\mathcal{R}=0
$$
where
$\mathcal{R}:=\text{Ric}(\gamma')+|A_0|^2$.
Then
\begin{equation}
\label{equation-Sturm-type1}
\det\mathbb{J}''+\mathcal{R}\det\mathbb{J}=0,
\det\mathbb{J}(0)=,\det\mathbb{J}'(0)=1
\end{equation}
\bigskip
Notice that in the case of constant curvature
$k$, Eq.
\ref{equation-Bishop-Gromov-proof1} becomes
\begin{equation}
\label{equation-Bishop-Gromov-proof3}
\det\overline{\mathbb{J}}'=
\det\overline{\mathbb{J}}k
\end{equation}
which in turn yields
\begin{equation}
\label{equation-Sturm-type2}
\det\overline{\mathbb{J}}''+
k\det\overline{\mathbb{J}}=0
\end{equation}
Now we want to compare the volume of the
sphere in $M$ with the volume of a
sphere of the same raduis in the space of
constant curvature
$\mathbb{Q}_k^{n+1}$:
$$
\begin{aligned}
\frac{\text{Vol}(S_r^{n}(p))}{\text{Vol}
(S^{n}_{r,k})}&=
\frac{\int_{S^n(0)}\det
\mathbb{J}}{\int_{S^n(0)}
\det\overline{\mathbb{J}}}
=\frac{1}{\text{Vol}(S^n(0))}
\int_{S^n(0)}\frac{\det\mathbb{J}}
{\det\overline{\mathbb{J}}}
\end{aligned}
$$
By Sturm Theorem \ref{theorem-Sturm}, which
we may apply since both
$\det\mathbb{J}$ and
$\det\overline{\mathbb{J}}$ satisfy Eqs.
 \ref{equation-Sturm-type1} and
 \ref{equation-Sturm-type2}, and moreover
they both vanish at $r=0$ and their
derivatives are $1$ at $r=0$ by Exercise
\ref{exercise-Jacobi-tensor-derivative}, and
of course, since we are supposing
that $\text{Ric}M\geq k$, we conclude that
the integrand is a non {\bf
decreasing} function.
\begin{remark}[Pregunta]
\label{remark-pregunta2}
Según Sturm, el cociente de las funciones en
cuestión, en este caso,
$\frac{\det\mathbb{J}}
{\det\overline{\mathbb{J}}}$
debería ser una función no
decreciente, pero el resultado que buscamos
es que sea no {\bf creciente}.
\end{remark}
\bigskip
Now we turn to computing the ratios of the
volumes of balls of radius $r$.
$$
\begin{aligned}
\frac{\text{Vol}(B_r(p))}{\text{Vol}(B_{r,k})
}&=
\frac{\int_{S^n(0)}\int_0^r\det\mathbb{J}(t)
dt\text{Vol}S^n(0)}
{\text{Vol}(S^n(0))\int_0^r
\det\overline{\mathbb{J}}(t)dt}\\
&=\frac{1}{\text{Vol}(S^n(0))}
\frac{\int_0^r}{}
\end{aligned}
$$
\begin{remark}[Preguntas]
\label{remark-pregunta3}
I have two questions:
\begin{enumerate}
\item We have said in lecture that we may
apply Fubini's theorem by Gauss'
lemma. Why? By Gauss' Lemma the normal
vectors are orthogonal to the spheres,
the constructions related to Riccati equation
are valid (cf. Section
\ref{section-Riccati-equation}). By why does
this allow to use Fubini?
\item We can pull out the volume in the
denominator because the volumes of
spheres in the space of constant curvature
somehow does not depend on the radius
 $r$. But how exactly? I think  $\det
 \overline{\mathbb{J}}(r)$ does depend
on $r$ (in the space of constant curvature),
see Eq.
\ref{equation-Bishop-Gromov-proof3}.
\end{enumerate}
\end{remark}
Then the equation continues to
$$
\begin{aligned}
&=\frac{1}{\text{Vol}(S^n(0))}
\int_{S^n(0)}\frac{\int_0^r \det
\mathbb{J}(t)dt \text{Vol}S^n(0)}
{\int_0^r\det\overline{J}dt}\\
&=\frac{1}{\text{Vol}(S^n(0))}
\int_{S^n(0)}\frac{\int_0^r \frac{\det
\mathbb{J}(t)}
{\det\overline{\mathbb{J}}(t)}
\det\overline{\mathbb{J}}(t)dt
\text{Vol}S^n(0)}
{\int_0^r\det\overline{\mathbb{J}}dt}\\
&=\frac{1}{\text{Vol}(S^n(0))}
\int_{S^n(0)}\frac{\int_0^r \frac{\det
\mathbb{J}(t)}
{\det\overline{\mathbb{J}}(t)}d\mu}{\mu[0,r]}
\text{Vol}S^n(0)
\end{aligned}
$$
where we have introduced a measure
$\mu:=\det \overline{\mathbb{J}}dt$.

By our discussion before, the integrand is a
non decreasing function, almost as
required.

Finally, notice that by the rigidity part on
Sturm's Theorem
\ref{theorem-Sturm}, the condition
$$
\frac{\text{Vol}(B_s(p))}{\text{Vol}(B_{s,k}}
=
\frac{\text{Vol}(B_r(p))}{\text{Vol}(B_{r,k})
}
$$
implies that
$\det\mathbb{J}=\det\overline{\mathbb{J}}$
and $k=\mathcal{R}$.
This makes $A_0=0$ from Eq.
\ref{equation-Bishop-Gromov-proof2} so that
$A$ is
proportional to the identity. By some mixture
of Eqs.
\ref{equation-jacob-field-const-curva}
and
\ref{equation-jacob-field-deriv-shape}, this
means that the
Jacobi fields along $\gamma$ are of the form
Eq.
\ref{equation-jacob-field-const-curva-2}. In
turn, this makes the map
of Proposition \ref{proposition-contraction}
an isometry. {\bf Who is $i$ then?}
\end{proof}

\section{Cheng's Theorem}
\label{section-Cheng-theorem}

Cheng's theorem says what happens in case of
equality in the conditions of
Bonnet-Myers Theorem
\ref{theorem-Bonnet-Myers}.
\begin{theorem}[Cheng]
\label{theorem-Cheng}
Let $M$ be a complete Riemannian manifold
with $\text{Ric}\geq 1$. If
$\text{diam}M=\pi$ then $M=S^n$
isometrically.
\end{theorem}

\begin{proof}
By Bonnet-Myers Theorem
\ref{theorem-Bonnet-Myers}, $M$ is a compact
manifold
with diameter $\pi$. This means that $M$ is
the closure of any ball of radius
$\pi$. Then, by Bishop-Gromov Theorem
\ref{theorem-Bishop-Gromov}, it's enough
to show the rigidity condition holds for
radius $\pi$. That is, we must show
that there is $s<\pi$ such that
$$
\frac{\text{Vol}(B_s(p))}{\text{Vol}(B_{s,1}}
=
\frac{\text{Vol}(B_\pi(p)}{\text{Vol}B_{\pi,
1}}
$$
Consider two points $p_1$ and $p_2$ at
distance $\pi$. Then the balls of radius
$\pi/2$ with center at these points cannot
intersect. This says that
$$
\text{Vol}M\geq B_{\pi/2}(p_1)+B_{\pi/2}(p_2)
$$
But since the volume of $M$ is exactly the
volume of $B_\pi(p)$ for any $p$, we
see that
$$
\frac{\text{Vol}(B_\pi(p))}{2}\leq
B_{\pi/2}(p_i),\qquad i=1,2
$$
\end{proof}

\section{Toponogov's theorem}
\label{section-Toponogov}

\begin{definition}
\label{definition-geodesic-triangle}
\cite{doc}, p. 285. A {\it geodesic triangle}
is the set formed by three
minimizing geodesics $\gamma_i$ so that
$\gamma_{i+1}(0)=\gamma_i(\ell(\gamma_i))$
cyclically. (Just to note that for
some authors geodesic triangles are
minimizing by definition. Looks like for
\cite{meyers} it's not the case.)
\end{definition}

First notice that the following result can be
proved using Rauch's theorem
\ref{theorem-Rauch}:

\begin{theorem}[Toponogov, local version]
\label{theorem-Toponogov-local}
Let $o,p_1,p_2 \in B$ where $B$ is a totally
convex ball of $M$. Let $\gamma_1$
and $\gamma_2$ be the normalized geodesics
joining $o$ to $p_1$ and to $p_2$.

Suppose $\tilde{o}, \tilde{p}_1,\tilde{p}_2
\in \tilde{B}$ is a triple on
another manifold $\tilde{M}$ such that the
distances from $\tilde{o}$ to
$\tilde{p}_1$ and from $\tilde{o}$ to
$\tilde{p}_2$ coincide with those in $M$.

Then for any $t \in [0,\ell(\gamma_1)]$ and
$s \in [0,\ell(\sigma)]$,
$$
\tilde{d}(\tilde{\gamma}_1(t),\tilde{\gamma}
_2(s))\leq
d(\gamma_1(t),\gamma_2(s))
$$
\end{theorem}

\begin{proof}
This is an application of Proposition
\ref{proposition-contraction} for $i$ as
given by the condition of mapping
$\gamma'_1(0)$ to $\gamma_2'(0)$ and the
angle
between them to the corresponding vectors and
angle on $\tilde{M}$.
\end{proof}

The global version of that Theorem
\ref{theorem-Toponogov-local} is the
following. Notice it coincides with our
intuition from comparing plane triangles
with spherical ones; spherical are fatter.

\begin{theorem}[Toponogov, hinge version]
\label{theorem-Toponogov-hinge}
$M$ complete, $K \geq k$. $\gamma_1,\gamma_2$
normalized geodesics with
$\gamma_1(0)=\gamma_2(0)$. Suppose $\gamma_1$
is minimizing.

If $k>0$ suppose additionally that
$\ell(\gamma_2) \leq \pi/\sqrt{k}$ (because
that's where conjugate points appear).

Suppose that
$\tilde{\gamma}_1,\tilde{\gamma}_2$ is a
corresponding hinge in
$\mathbb{Q}_k^2$, that is,
$\ell(\gamma_i)=\ell(\tilde{\gamma}_i)$ and
$\angle(\gamma_1'(0),\gamma_2'(0)=
\angle\tilde{\gamma}_1'(0),\tilde{\gamma}
_2'(0)$.
Then
\begin{equation}
\label{equation-Toponogov-hinge1}
d(\gamma_1(\ell_1),\gamma_2(\ell_2)) \leq
\tilde{d}(\tilde{\gamma}_1(\ell_1),
\tilde{\gamma}_2(\ell_2))
\end{equation}
\end{theorem}

This is not intuitive for me: the distance
between the remaining side should be
fatter in the manifold that is more curved.

The following version is equivalent to the
hinge version and is intuitive for
me:

\begin{theorem}[Toponogov]
\label{theorem-Toponogov-hinge2}
Let $M$ be complete with $K\geq k$. If
$\gamma_i$ is a minimizing geodesic
triangle in $M$, then there is a unique
minimizing geodesic triangle
$\tilde{\gamma}_i$ in $\mathbb{Q}_k^2$ with
$\ell(\tilde{\gamma}_i)=\ell(\gamma_i)$ for
$i=0,1,2$, and such that
\begin{equation}
\label{equation-Toponogov-hinge2}
\tilde{d}(\tilde{\gamma}_1(t),\tilde{\gamma}
(s))\leq
d(\gamma_1(t),\gamma_2(s)
\end{equation}
for all $t\in[0,\ell(\gamma_1)]$ and
$s\in[0,\ell(\gamma_2)]$.
\end{theorem}

Another equivalent version:
\begin{theorem}[Toponogov]
\label{theorem-Toponogov-hinge3}
Let $M$ be complete with $K\geq k$. If
$\gamma_i$ is a minimizing geodesic
triangle in $M$, then there is a unique
minimizing geodesic triangle
$\tilde{\gamma}_i$ in $\mathbb{Q}_k^2$ with
$\ell(\tilde{\gamma}_i)=\ell(\gamma_i)$ for
$i=0,1,2$ and such that
\begin{equation}
\label{equation-Toponogov-hinge3}
\tilde{d}(\tilde{o},\tilde{\gamma}(t))\leq
d(o,\gamma_0(t))
\qquad \forall t\in[0,\ell(\gamma_0)]
\end{equation}
\end{theorem}
All these versions are equivalent and we will
only prove the following one:
\begin{theorem}[Toponogov, metric version]
\label{theorem-Toponogov-metric}
Let $M$ be a complete manifold and $K \geq
k$. Let $o,p_1,p_2 \in M$, distinct
points. Let $\gamma_1$ and $\gamma_2$ be
normalized geodesics joining
$o$ to $p_1$ and to $p_2$. Suppose $\gamma_1$
is minimizing (in the prove we
show that $\gamma_2$ is minimizing too).
Let $\gamma_0$ be a nonconstant geodesic
joining $p_1$ to $p_2$ and suppose
$$
\ell(\gamma_0) \leq \ell(\gamma_1)
+\ell(\gamma_2)
$$
(which is a weaker condition than asking that
$\gamma_0$ is minimizing as well).

If $k>0$ suppose additionally that
$\ell(\gamma_0) \leq \pi/\sqrt{k}$.

Then there exists a geodesic triangle
$\{\tilde{\gamma}_i:i=1,2,0\}$ in
$\mathbb{Q}_k^2$ such that
$\ell(\tilde{\gamma}_i) = \ell(\gamma_i)$ and
\begin{equation}
\label{equation-Toponogov-metric}
\tilde{d}(\tilde{o},\tilde{\gamma}_0(t))
\leq  d(o,\gamma_0(t))
\qquad \forall t \in[0,\ell(\gamma_0)]
\end{equation}
\end{theorem}

\begin{proof}[Outline of proof]
This is the one we proved in class. Define
distance functions
$\rho:=d(\sigma,\cdot)$ and
$\tilde{\rho}:=\tilde{d}(\tilde{\sigma},
\cdot)$.
\end{proof}

\begin{lemma}
\label{lemma-equality-Toponogov}
If equality holds in Eq.
\ref{equation-Toponogov-hinge1}, then
$\gamma_1$ and
$\gamma_2$ are the boundary of a totally
geodesic surface $\Delta \subset M$
with $K_\Delta \equiv k$.
\end{lemma}

\begin{proof}
Using the map from Cartan's lemma. Use
Rauch's theorem \ref{theorem-Rauch} to
compare $M$ with $\mathbb{Q}_k$, and again to
compare $\Delta$ with $M$. As you
might expect, use Gauss' equation
\ref{equation-Gauss} to show that the second
fundamental form of $\Delta$ vanishes. The
essence of the proof is that
$\alpha(\gamma',\gamma')=0$ for the geodesics
emanating from $p$, which is more
or less straightforward, and then showing
that also $\alpha(\sigma',\sigma')=0$
for geodesics emanating from another vertex
of the triangle, which requires to
``look at things from another point of
view''.
\end{proof}

\begin{theorem}[Toponogov, angle version]
\label{theorem-Toponogov-angle}
\cite{Cheeger-Ebin} Theorem 2.2.(A). Let $M$
be a complete manifold with
$K_M\geq H$.

Let $(\gamma_1,\gamma_2,\gamma_3)$ determine
a geodesic triangle in $M$. Suppose
$\gamma_1,$ $\gamma_3$ are minimal and if
$H>0$, suppose $\ell(\gamma_2)\leq
\frac{\pi}{\sqrt{H}}$. Then $M^H$, the simply
connected 2-dimensional space of
constant curvature $H$, there exists a
geodesic triangle
$(\tilde{\gamma}_1,\tilde{\gamma}_2,
\tilde{\gamma}_3)$
such that
$\ell(\gamma_i)=\ell(\tilde{\gamma}_i)$ and
$\tilde{\alpha}_1 \leq\alpha_1$,
$\tilde{\alpha}_3 \leq \tilde{\alpha}_3$.

Except in case $H>0$ and
$\ell(\gamma_i)=\frac{\pi}{\sqrt{H}}$ for
some $i$, the
triangle in $M^H$ is uniquely determined.
\end{theorem}

\section{Gromov's theorem}
\label{section-Gromov-theorem}

\begin{theorem}[Gromov]
\label{theorem-Gromov}
$M^n$ complete, $K \geq 0$, then $\pi$ admits
a set of $3^n$ generators. In
particular, it is finitely generated.
\end{theorem}

\begin{proof}
Consider the universal cover $\tilde{M}$ of
$M$ with the pullback metric. Fix a
point $x \in \tilde{M}$ and consider the set
$$
\{f \in \Gamma:d(x,f(x))<r\},\qquad r>0
$$
By \ref{proposition-deck-trans-act-prope}
every point must have a disjoint
neighbourhood …
\end{proof}

\section{Cartan's Theorem}
\label{section-Cartan-theorem}

\begin{definition}
\label{definition-free-homotopy-set}
The {\it free homotopy set} is the set of
homotopy classes of loops based on any
point of $M$.
\end{definition}

\begin{definition}
\label{definition-geodesic-loop}
A {\it geodesic loop} is a geodesic whose
starting and ending points are the
same. It may not be differentiable at such
point.
\end{definition}

\begin{definition}
\label{definition-closed-geodesic}
A {\it closed geodesic} a smooth geodesic
loop.
\end{definition}

\begin{theorem}[Cartan]
\label{theorem-Cartan}
Let $M^n$ be compact, and $\omega \in
\hat{\pi}_1M$ nontrivial. Then there is a
closed geodesic $\gamma$ such that
$[\gamma]=\omega$.
\end{theorem}

\begin{proof}
Consider a sequence of loops whose lengths
converge to the number
$$
\ell:=\text{inf}\{\ell(c):c \in \omega\}>0
$$
Consider a sequence $c_n$ of curves in
$\omega$ such that their lengths converge
to $\ell$. Then apply Arzelà-Ascoli Theorem
\ref{theorem-Arzela-Ascoli}
 to ensure they converge to some curve $c$
 and show this is a geodesic
in $\omega$.

First suppose that each of the $c_n$ is a
piecewise geodesic parametrized by
arclength (why can we do this?). We may then
pick the supremum $L$ of all the
lengths since we are assuming that their
lengths decrease. Then
$$
d(c_n(t),c_n(s)) \leq \ell
\left(c_n|_{[t,s]}\right) \leq  N |t-s|
$$
This says that $\{c_n\}$ is equicontinuous
(the same $N$ works for all $n$) as a
family of maps from
\end{proof}

\section{Preissman's Theorem}
\label{section-Preissman-theorem}

\begin{definition}
\label{definition-translation}
An isometry $f \in \text{Iso}(M)$ is called a
{\it translation} if there exists
a geodesic of $M$ such that
$f(\gamma)=\gamma$. That is, $f \circ \gamma$
is a
reparametrization of $\gamma$.
\end{definition}

\begin{theorem}[Preissman]
\label{theorem-Preissman}
Let $M^n$ be compact with $K<0$. If $1 \neq
H \subset \pi_1(M)$ is abelian,
then $H=\mathbb{Z}$.
\end{theorem}

\begin{theorem}[Preissman]
\label{theorem-m-compa-k-negat-then-pi1}
If $M$ is compact and $K<0$, then
$\pi_{1}(M)$ is not abelian.
\end{theorem}

\section{Byers' Theorem}
\label{section-Byers-theorem}

\begin{lemma}
\label{lemma-deck-fixes-geode-then-not}
Let $M^n$ be a complete manifold with $K<0$.
Suppose that every element of
$\text{Deck}(\pi)$, where $\pi$ is the
universal covering, fixes the same
 geodesic $\tilde{\gamma}$. Then $M$ is not
 compact.
\end{lemma}

\begin{theorem}[Byers]
\label{theorem-Byers}
Let $M$ be compact, $K<0$ and $H$ soluble
subgroup of $\pi_1(M)$, $H\neq \{e\}$.
Then $H$ is cyclic infinite. Moreover,
$\pi_1(M)$ does not have any cyclic
subgroup of finite index.
\end{theorem}

\section{Cheeger-Gromoll Siplitting Theorem}
\label{section-chegg-gromo-split-theor}

\begin{theorem}[The Strong Maximum Principle]
\label{theorem-strong-maximum-principle}
\cite{pet} Theorem 7.1.7. If $f:(M,g) \to
\mathbb{R}$ is continuous and
subharmonic, then $f$ is constant in a
neighbourhood of every local maximum. In
particular, if $f$ has a global maximum, then
$f$ is constant.
\end{theorem}

\begin{proof}
If $\Delta f>0$ then we can find a support
function. Notice that Prof. Florit
shows the theorem for functions satisfying
$\Delta f \leq 0$.
\end{proof}

The following lemma follows from the lemma
above and gives a bound on the
Laplacian of the distance function.

\begin{lemma}[Calabi]
\label{lemma-Calabi}
If $\text{Ric} \geq 0$, for
$\rho:=d(p,\cdot)$, it holds that
$\Delta\rho \leq (n-1)/\rho$ on $M\setminus
C_m(p)\cup\{p\}$.
\end{lemma}


\begin{theorem}[Cheeger-Gromoll]
\label{theorem-splitting}
Let $M$ be complete with  $\text{Ric}\geq0$.
If $M$ has a line, then $M$ is
isometric to $N\times\mathbb{R}$ for some
Riemannian manifold $N$.
\end{theorem}

\section{Exercises: basic constructions}
\label{section-exercises-constructions}

\begin{exercise}
\label{exercise-trace-indep-coord}
Show that the trace is independent of
coordinates.
\end{exercise}

\begin{proof}[Solution]
Let $A$ be an endomorphism of a vector space.
By definition
$$
\text{tr}A=A_i^i
$$
Now if $P$ is any invertible matrix,
$$
\text{tr}PAP^{-1}=P^j_iA_j^k\overline{P}_k^i=
P^j_i\overline{P}_k^iA_j^k
=\delta_k^jA_j^k=A_j^j
$$
where $\overline{P}$ denotes the indices of
the inverse of $P$.
\end{proof}

The following two lemmas are preparatory
results for Exercise
\ref{exercise-derivative-of-determinant}.

\begin{lemma}
\label{lemma-for-derivative-of-determinant1}
Let
$\gamma_1,...,\gamma_m:(a,b)\to\mathbb{R}^n$
be smooth curves and define
$f:(a,b)\to\mathbb{R}$ by
$f(t)=\det(\gamma_1(t),\ldots,\gamma_n(t))$.
Then
$$
f'(t_0)
=\sum_{k=1}^{m}\det(\gamma_1(t_0),\ldots,
\gamma'_k(t_0),\ldots,\gamma_m(t_0)).
$$
\end{lemma}

\begin{proof}
This follows from multilinearity of the
determinant using the definition of
derivative of a map $g:\mathbb{R}^n \to
\mathbb{R}^m$ as given by
$$
g(x+h)-g(x)=L(x)h+\varepsilon(x;h)
$$
where $\varepsilon$ is $o(h)$, meaning that
$\lim_{h \to 0}
\frac{\varepsilon(x;h)}{h}=0$.

More exactly, since $\gamma_k$ are smooth we
may write
$$
\gamma(t_0+h)=\gamma_k(t_0)+h\gamma'_k(t_0)+
\varepsilon_k(h)
$$
We substitute into the determinant…
\end{proof}

\begin{lemma}
\label{lemma-for-derivative-of-determinant2}
Let $M \in \text{Mat}_{n\times
n}(\mathbb{R})$ be any matrix. Let
$f:\mathbb{T}\to \mathbb{R}$ be defined by
$f(t):=\det(\text{Id}+tM)$. Then
$f'(0)=\text{tr}(M)$.
\end{lemma}

\begin{proof}
This follows easily from Lemma
\ref{lemma-for-derivative-of-determinant1}.
\end{proof}

\begin{exercise}
\label{exercise-derivative-of-determinant}
Let $A(t)$ be a one parameter family of
matrices. Compute $\det A(t)'$.
\end{exercise}

\begin{proof}[Naive solution]
The determinant of a matrix is given by the
product of its eigenvalues
$\lambda_i$. Taking logarithms we see that
$$
\text{log}\det
A(t)=\text{log}\prod_{i}\lambda_i=\sum_{i}
\text{log}\lambda_i
$$
and differentiating we obtain
$$
\frac{\det A(t)'}{\det A(t)}=\sum_i
\frac{\lambda_i'}{\lambda_i}
=\text{tr} A'(t)A(t)^{-1}.
$$
with the caveat that we must confirm whether
the functions $\lambda_i$ are
differentiable. (Still pending…)
\end{proof}

\begin{proof}[Solution]
Consider the determinant as a function
$\det:\text{GL}(n,\mathbb{R})\to\mathbb{R}$.
Is total derivative at
$A \in \text{GL}(n,\mathbb{R})$ is a linear
map
$\text{GL}(n,\mathbb{R})\to \mathbb{R}$ which
when evaluated at a matrix
$H\in\text{GL}(n\mathbb{R})$ is defined by
$$
\det(A)'H=\lim_{t\to0}\frac{\det(A+tH)-
\det(A)}{t}
$$
Now define $f(t):=\det(A+tH)$ so that the
latter equation reads $f'(0)$. Then
observe that
$$
f(t)=\det(A+tH)=\det(A(\text{Id}+tA^{-1}H)=
\det(A)\det(\text{Id}+tA^{-1}H).
$$
By Lemma
\ref{lemma-for-derivative-of-determinant2} we
conclude that
$f(t)=\det(A)\text{tr}(A^{-1}H)$. Letting
$H=A(t)'$ we obtain the result since
by the chain rule
$\det(A(t))'=\det(A(t))'A(t)'$.
\end{proof}

\begin{exercise}
\label{exercise-levi-civit-conne-produ}
Let $P$ and $Q$ be Riemannian manidols. Sow
that the Levi-Civita connection of
the product metric is given by
$\nabla_{Y_1+Y_2}(X_1+X_2)=\nabla^P_{Y_1}X_1+
\nabla^Q_{Y_2}X_2$.
\end{exercise}

\begin{exercise}
\label{exercise-total-geode-subma-produ}
Let $P$ and $Q$ be Riemannian manifolds. Show
that for any $p \in P$ and $q \in
Q$, $P\times \{q\}$ and $\{q\}\times P$ are
totally geodesic
submanifolds of $P \times Q$.
\end{exercise}

\begin{proof}[Solution]
Let $v \in T_{(p,q)}(P\times\{q_0\})$ and
consider the geodesic
 $\text{exp}_p^{P\times Q}(tv)$. Then
its acceleration vanishes since the
Levi-Civita connection of the product metric
is given by Exercise
\ref{exercise-levi-civit-conne-produ}, and
the
component along $Q$ vanishes.
\end{proof}

\section{Exercises: global differential
geometry}
\label{section-exerc-globa-diffe-geome}

\begin{exercise}
\label{exercise-Klingenberg-lemma}
\cite{doc}, Chapter X, Exercise 1. Seja $M$
uma variedade Riemanniana
completa com curvatura seccional $K \leq
K_0$ onde $K_0$ é uma constante
positiva. Sejam $p,q \in M$ e seja $\gamma_0$
e $\gamma_1$ duas geodésicas
distinas unindo $p$ a $q$ com
$\ell(\gamma_0) \leq  \ell(\gamma_1)$.
Admita que $\gamma_0$ é homotópica a
$\gamma_1$, isto é, existe uma família
contínua de curvas $\alpha_t$, $t \in [0,1]$
tal que $\alpha_0 = \gamma_0$
e $\alpha_1=\gamma_1$. Prove que existe $t_0
\in [0,1]$ tal que
$$
\ell(\gamma_0)+\ell(\alpha_{t_0}) \geq
\frac{2\pi}{\sqrt{K_0} }
$$
\end{exercise}

\begin{exercise}
\label{exercise-Hadamard-via-Kingenberg}
Use Klingenberg lemma from Exercise
\ref{exercise-Klingenberg-lemma} to prove
Hadamard's theorem.
\end{exercise}

\begin{proof}
Suponha que $M$ é uma variedade completa,
simplesmente conexa e com curvatura
positiva. Por completitude, sabemos que para
qualquer $p \in M$ o domínio de
$\text{exp}_p$ é todo $T_pM$. Como $K\leq 0$,
sabemos que $M$
não pode ter pontos conjugados, de modo que
$\text{exp}_p$ é um
difeomorfismo local. Também por completitude
sabemos que
$\text{exp}_p$ é sobrejetiva: qualquer ponto
$q$ está ligado a $p$
mediante uma geodésica, e essa geodésica
coincide com uma geodésica partindo de
$p$ dada como a imagem de uma reta em $T_pM$
sob $\text{exp}_p$.

Portanto, o desafio é provar que
$\text{exp}_p$ também é injetiva.
Suponha que não é o caso, i.e. considere dois
pontos $v_1$ e $v_2$ em
$T_pM$ tais que
$\text{exp}_pv_1=\text{exp}_pv_2:=q$. As
imagens das retas geradas por $v_1$ e $v_2$
sob $\text{exp}_p$ são
duas geodésicas $\gamma_1$ e $\gamma_2$ em
$M$ ligando $p$ e $q$.

Como $M$ é simplesmente conexa, existe uma
homotopia entre $\gamma_1$ e
$\gamma_2$. Podemos aplicar o lema de
Klingenberg para obter uma curva
$\alpha_{K_0}$ tal que $$
\ell(\gamma_1)+\ell(\alpha_{K_0}) \geq
2\pi/\sqrt{K_0} $$ para qualquer $K_0>0$.
Isso mostra que o comprimento das
curvas na homotopia não está limitado, o que
não é possível já que a homotopia é
uma função contínua definida em um compacto,
pelo qual a sua imagem debe ser
compacto e portanto limitado. {\bf Errado!
Por exemplo, um círculo percorrido
miles de vezes tem comprimento
arbitrariamente grande e está contido num
compacto.}
\end{proof}

\begin{exercise}
\label{exercise-geode-sem-ponto-conju-e}
$\gamma$ geodésica sem pontos conjugados em
$[0,a]$. Então $\gamma$ é localmente
minimizante.
\end{exercise}

\begin{proof}
Temos dois casos:
\begin{enumerate}
\item Se $E''(0)<0$ usamos o teorema do
índice. Isto é, não é possível que
 $E''(0)<0$ porque isso implica que o índice
 da forma do índice é diferente
 de zero, ou seja, existem pontos conjugados.
\item Se $E''(0)=0$,
\begin{enumerate}
\item[(0)] Como $\gamma$ não tem pontos
conjugados, então $I$ não tem
nulidade em todo  $\mathcal{V}$.
\item $I_{\mathcal{V}^-}$ não tem índice nem
nulidade. Segue do ponto anterior +
algum passo na prova.
\item $I_{\mathcal{V}^-}$ não tem índice.
\item 2+3 dão que $I_{\mathcal{V}^-}>0$.
\item Logo $I>0$.
\end{enumerate}
\end{enumerate}
\end{proof}

\begin{exercise}
\label{exercise-two-submanifolds}
\cite{doc} Chapter IX, Exercise 5. Sejam
$N_1$ e $N_2$ duas subvariedades
 fechadas e disjuntas de uma variedade
 Riemanniana compacta.
\begin{enumerate}
\item Mostre que a distância entre $N_1$ e
$N_2$ é realizada por uma geodésica
	$\gamma$ perpendicular a ambas $N_1$ e
	$N_2$.
\item Mostre que, para qualquer variação
ortogonal $h(t,s)$ de $\gamma$, com
$h(0,s) \in N_1$ e $h(\ell,s) \in N_2$,
tem-se para a fórmula da segunda
variação a seguinte expressão
$$
\frac{1}{2}E''(0)=I_{\ell}(V,V)+
\left\langle{}V(\ell),A_{\gamma'(\ell)}^{(2)}V(\ell)
\right\rangle{}
-\left\langle{}V(0),A^{(1)}_{\gamma'(0)}(V(0))\right\rangle{}
$$
\end{enumerate}
\end{exercise}

\begin{proof}
\begin{enumerate}
\item Since both submanifolds are compact,
there exists a minimizing geodesic
$\gamma$ joining them. Then we can apply the
first variation formula. Consider a
variation that fixes $\gamma$ everywhere but
in a small neighbourhood of one of
the contact points. We realise by Eq.
\ref{equation-first-variation-formula}
that it must be orthogonal. A similar
procedure proves the other endpoint also
arrives orthogonally.

\item By the previous item, we
know that $\gamma'$ is orthogonal to the
curves $h(0,s)$ and $h(\ell,s)$. Then
we can differentiate using the formula for
normal sections at the endpoints of
$\gamma$:
$$
\tilde{\nabla}_{V}\gamma'=A^{(i)}_{\gamma'}V+
\nabla^{\perp,i}_V\gamma'
$$
When computing
$V\left\langle{}\gamma',\gamma'\right\rangle{}$ we obtain
$\left\langle{}\nabla^{\perp,i}_V\gamma',
\gamma'\right\rangle{}=0$
since the shape operator is
tangent to the submanifolds. If $M_1$ and
$M_2$ are hypersurfaces we conclude
that the normal vector
$\nabla^{\perp,i}_V\gamma'$ vanishes since it
must be a
multiple of $\gamma'$. Otherwise I'm not sure
why this term would vanish.

Now, other than the Index form, in the Second
Variation Formula
\ref{equation-second-variation-formula} we
have
$\left\langle{}\nabla_{\partial_s}V,\gamma'\right\rangle{}$.
Notice that
$$
\left\langle{}f_s,f_t\right\rangle{}_s=\left\langle{}f_{ss},
f_t\right\rangle{}+\left\langle{}f_s,f_{st}\right\rangle{}
$$
The left-hand-side will vanish because the
variation is orthogonal, giving the
equality we want (modulo a sign) since
$f_{st}$ is precisely
$\tilde{\nabla}_V\gamma'$ by the Symmetry
Lemma.
\end{enumerate}
\end{proof}

\begin{exercise}
\label{exercise-inter-minim-hyper}
Let $(M^n,g)$ be a complete Riemannian
manifold with $\text{Ric}_M>0$.
Suppose $M_1$, $M_2$ are minimal (compact?)
complete hypersurfaces. Show that
$M_1 \cap M_2 \neq \emptyset$.
\end{exercise}

\begin{proof}
Let $\gamma$ be a minimizing geodesic joining
$p \in M_1$ to $q\in M_2$ and
minimizing the distance between $M_1$ and
$M_2$. Let $e_1,\ldots,e_{n-1}\in
T_pM$ be such that along with $\gamma'(0)$
form an orthonormal basis of $T_pM$.
Let $E_i$ be the parallel transport of $e_i$
along $\gamma$. Notice that
$E_i(\ell)$ is tangent to $M_2$ since it is
orthogonal to $\gamma'(\ell)$ and
$M_2$ is a hypersurface, but it is not
immediate that the variational curve will
be contained in $M_2$ if we define the
variation using the exponential map. This
is essential for the second variational
formula to vanish since $\gamma$ is
minimizing among curves joining $M_1$ and
$M_2$.

We need to produce the variation not from the
vector field but from the
variational curves. Take rectifying
coordinates for $M_1$ at $p$. Define the
curve $c^i$ as the $i$-th coordinate curve on
the whole neighbourhood of $p$
along $\gamma$.

Now take the velocity of $c^i$ for some
$t>0$. Parallel transport this vector
along $\gamma$ until it reaches a rectifying
neighbourhood of $q$. Along the
path of $\gamma$ we can use the exponential
to define a variation. Once in the
rectifying neighbourhood of $q$, we must
define by hand the variational curves
to ensure that the curve at $q$ is contained
in $M_2$.

Suppose that $\gamma(t_0)$ is in the
rectifying neighbourhood of $q$. Consider a
smooth function $\varphi:[t_0,\ell]\to
\mathbb{R}$ that is zero at $t_0$ and $1$
at $\ell$. We can smoothly deform the
exponential curves to a curve contained in
$M_2$ at $t=\ell$ by defining
$$
c^s_i(t):=(1-\varphi)\text{exp}_{\gamma(t)}
sE_i(t)+\varphi
sE_i(t)
$$
Notice that the variational vector is
orthogonal to $M_2$ at $q$, so that we may
apply Exercise
\ref{exercise-two-submanifolds}. Thus we can
express the second
variation formula in such a way that when
summing over all indices we obtain the
mean curvatures of each submanifold, which
vanish, thus obtaining a
contradiction with the hypothesis that
$\text{Ric}>0$.
\end{proof}

\begin{exercise}
\label{exercise-even-dimen-posit-k-compa}
Prove que $M^{2n},K_M>0$ compacta, então ela
é simplesmente conexa.
\end{exercise}

\begin{exercise}
\label{exercise-odd-dimen-posit-k-orien}
$M^{2n+1}$, $K_M>0$ compacta $\implies$
orientável.
\end{exercise}

\begin{proof}
I want to show that $M^{2n+1}$ is
diffeomorphic to its double orientable cover
 $\tilde{M}$. Since $M$ is odd-dimensional,
 so is $\tilde{M}$.
\end{proof}

O exercício anterior usa Exercício 6 da lista
6. Ou, equivalentemente,

\begin{exercise}
\label{exercise-inver-orien-short}
Se $M^{2n+1}$ tem $K_M >0$ e $\gamma:S^1 \to
M$ inverta orientação, então é
possível encurtar $\gamma$ na sua classe de
homotopia.
\end{exercise}

\begin{exercise}
\label{exercise-produ-compa-does-not-admit}
Let $N,M$ be compact. Then $N \times M$ does
not admit a metric with $K<0$.
\end{exercise}

\begin{proof}
If both $N$ and $M$ are simply connected then
they cannot be compact by
Hadamard's theorem. Assume first that both
are not simply connected. Then each
has a nontrivial homotopy class, and such a
pair generates an abelian subgroup
of $\pi_1(M \times N)$ that is not isomorphic
to $\mathbb{Z}$, a contradiction
with Preissman Theorem
\ref{theorem-Preissman}.

Now suppose that $N$ is simply connected and
$M$ is not. The universal cover of
$M\times N$, if it admits a metric with
negative sectional curvature,
must be a Hadamard manifold, so that
$\tilde{M} \times N \cong \mathbb{R}^n$.
Then the induced map on cohomology is
injective map ({\bf why?})
$\pi^* :H^{k}(M)\to H^{k}(\tilde{M} \times
N)=H^{k}(\mathbb{R}^n)=0$. This means
by Poincaré duality that $0=H^{n}(M)\cong
H_0(M)\neq 0$, a contradiction.
\end{proof}

\begin{exercise}
\label{exercise-SnxR-no-positive-Ric}
Show that $S^n \times \mathbb{R}$ does not
admite a complete metric with
$\text{Ric}>0$.
\end{exercise}

\begin{proof}
If it did, then by Preissman theorem
\ref{theorem-Preissman} then $S^n\times
\mathbb{R}\cong N \times \mathbb{R}$, but the
upshot is that that the latter
isomorphism is an isometry, and the
$\mathbb{R}$ factor is flat, so we can't
have strictly positive Ricci.
\end{proof}

\begin{exercise}
\label{exercise-compl-witho-compa-set-isom}
Let $(M^n,g)$ be complete and such that there
exists a compact set
 $K\subset M$ such that $M\setminus K \cong
 \mathbb{R}^n \setminus \tilde{K}$
for some compact set $\tilde{K} \subset
\mathbb{R}^n$. If $\text{Ric}M \geq 0$,
then $M^n \cong \mathbb{R}^n$.
\end{exercise}

\begin{proof}
Since $M\setminus K$ is isometric to
$\mathbb{R}^n\setminus \tilde{K}$, there
must be a line (oops! this is not obvious).
But suppose for now that there {\it
is} a line on $M$. Then by the Splitting
theorem \ref{theorem-splitting} there
is an isomorphism $M \cong N \times
\mathbb{R}$. ({\it My original idea:}
 If we manage to show that $N$ also has a
 line we would have that
 $M \cong N_2 \times \mathbb{R}^2$.  if we
 can further do this process and find
 that all $N_i$ have lines,
then we end up with $M \cong N_k \times
\mathbb{R}^k$ for some manifold $N_k$
without lines.)

Let's fix some notation. Let
$f:M\setminus K \to \mathbb{R}^n
\setminus\tilde{K}$ be the given isometry and
$\varphi:N^{n-1}\times\mathbb{R} \to M$ the
isometry given by the Splitting
theorem. Consider the restriction $f \circ
\varphi |_{ N^{n-1}\times \{ a\}}$
for some $a \in \mathbb{R}$. Observe que
$N^{n-1}\times \{ a\}$ é totalmente
geodésica (cf. Exercise
\ref{exercise-total-geode-subma-produ})
, e portanto também a sua imagem sob  $f
\circ \varphi$. Toda subvariedade
totalmente geodésica de $\mathbb{R}^n$ é um
plano, e acabamos.

However, notice that it is not immediate from
the
isomorphism $f$ that $M$ has a line. Fix a
line $L$ in $\mathbb{R}^n$. Then its
inverse image under $f$ is minimizing in $M
\setminus K$, but there could be a
curve that minimizes distance between some
two points of $f^{-1}(L)$ that is not
contained in $f^{-1}(L)$.

Suppose that $f^{-1}(L)$ is not minimizing in
$M$. Then there are two points $j$
and $-j$ in $f^{-1}(L)$ and a minimizing
geodesic $\sigma$ joining them and
intersecting $K$. Then the path that joins
any point before $j$ and any point
after $j$ on $\sigma$, before $\sigma$
intersects $K$, is a minimizing segment
contained in $M\setminus K$ and should thus
be contained in $f^{-1}(M)$. This
shows that $f^{-1}(K)$ accumulates near $K$,
which is not possible.
\end{proof}

\begin{exercise}
\label{exercise-S1xS1xR-admite-curvatura--1}
Mostre que $S^1 \times S^1 \times \mathbb{R}$
{(\bf não?)} admite uma métrica de
 curvatura $K\equiv-1$. {\it Dica.} Teorema
 de Preissman
\ref{theorem-Preissman}.
\end{exercise}

\begin{proof}
Porque $\mathbb{Z} \times \mathbb{Z}$ é
abeliano mas não é isomorfo a
$\mathbb{Z}$.
\end{proof}

\begin{exercise}
\label{exercise-rpnxr-admit-curva-posit}
$\mathbb{R}P^{n}\times \mathbb{R}P^{m}$
admite uma métrica com
curvatura positiva?
\end{exercise}

\begin{proof}
Se $m,n$ são os dois ímpares, são
orientáveis, então o produto é orientável e
de
dimensão par, e pelo Teorema de Synge
\ref{theorem-Synge} deve ser
simplesmente conexa, absurdo.

Se as duas tem dimensão par, considere o
recobrimento duplo orientável, que deve
ser compacto (pode usar Bonnet-Myers
\ref{theorem-Bonnet-Myers}), com curvatura
positiva e dimensão par, e portanto é
simplesmente conexo. Ou seja, trata-se do
recobrimento universal. Porém, ele é um
recobrimento duplo e portanto a fibra
tem dois elementos, e desse jeito o
$\text{Deck}\cong
\pi_{1}(\mathbb{R}P^{n}\times\mathbb{R}P^{m})
$
deve ser $\mathbb{Z}_2$,
mas não é.

E se uma é de dimensão par e a outra ímpar… é
igualzinho que no caso anterior.
\end{proof}

\begin{exercise}
\label{exercise-Sn-Tn-CPn-existem-curvaturas}
Considere $S^n,\mathbb{T}^n$ e
$\mathbb{C}P^{n}$. Existem métricas com
$K>0$, $K
\leq 0$, $K<0$, $K \geq 0$, $\text{Ric}>0$,
$\text{Ric} \geq 0$? Faça uma
tabela.
\end{exercise}

\begin{exercise}
\label{exercise-Hadamard-manifold-sphere}
Let $M^n$ be a Hadamard manifold and
$S^{n-1}_r \subset M$. Mostre que
 a curvatura média satisfaz $H_S\geq
 \frac{1}{r}$. {\it Dica.} A curvatura média
da esfera em $\mathbb{R}^n$ é $\frac{1}{r}$.
\end{exercise}

\begin{proof}
Considere uma geodésica que liga $p$, o
centro da esfera, e o bordo dela.
Considere um referencial de campos de Jacobi
$J_i$. Lembre que cada um dele
satisfaz $A J=J'$. Então
$$
H=\sum_{i}\left\langle{}A
\frac{J_i}{|J_i|},J_i\right\rangle{}=
$$
\end{proof}

\begin{exercise}
\label{exercise-product-of-compact-and-Rn}
$M^n$ compact, $\text{Ric} \geq 0$, then
$\tilde{M}=N \times \mathbb{R}^n$ for
some compact manifold $N$, where  $\tilde{M}$
is
the universal cover of $M$.
\end{exercise}

\begin{proof}
If $\tilde{M}$ is compact we are done.
So suppose it is not.
Then there exists a
ray. We want to construct
a line to use Splitting Theorem
\ref{theorem-splitting}.
Consider the fundamental domain of the cover
$\tilde{M}$,
which exists since
$M$ is compact: we may take a finite covering
of normal neighbourhoods of $M$,
so that the preimage under $\pi$ is a compact
set containing at least one
representative of each class.

Let $p_i$ be a sequence along the ray
$\gamma$
so that $p_i$ goes to infinity.
For each of them there exists a
deck transformation $F_i$ that puts $p_i$
inside
$K$. Consider the ray $F_i \circ \gamma$.
\end{proof}

\begin{exercise}
\label{exercise-rays-semispaces}
Lista 8, Exercício 12. Let $M$ be a complete
Riemannian manifold not compact
with $K \geq 0$. Let $\sigma$ be a ray of
$M$. Define for $t>0$ the {\it
semispace}
$$
H_t^{\sigma}:=M\setminus
\bigcup_{s>0}\Big(\sigma(t+s),s).
$$
Show that $H_t^{\sigma}$ is totally convex,
that is, is a geodesic segment
 $\alpha:[0,1]\to M$ has its endpoints in
 $H_t^{\sigma}$, i.e.
$\alpha(0),\alpha(1) \in H_t^{\sigma}$, then
that segment is completely
contained in $H_t^{\sigma}:\alpha([0,1])
\subset H_t^{\sigma}$. {\it Hint.}
Suppose that it's false and use Toponogov.
\end{exercise}

\begin{proof}
Suppose that there is a geodesic
$\gamma:[0,1]\to M$ with $\gamma(0),\gamma(1)
\in M$, but it intersects $H_t^{\sigma}$.
Consider the triangle joining $p$,
$\sigma(t+s)$ and some other point.

Choose the other point so that you obtain a
contradiction using the structure of
$H_t^{\sigma}$ of embedded balls---if a point
is in any ball, it must be in all
the next ones.
\end{proof}

\section{Lista 3}
\label{section-lista-3}

\begin{exercise}[Curvas minimizantes]
\label{exercise-curvas-minimizantes}
\begin{enumerate}
\item Seja $\gamma$ uma curva suave por
partes parametrizada por comprimento de
 arco, conectando $p$ a $q$, pontos de uma
 variedade Riemanniana $(M,g)$.
 Mostre que se $d(p,q)=\ell(\gamma)$, então
 $\gamma$ é uma geodésica.
 Dizemos que $\gamma$ realiza a distância
 entre $p$ e $q$.
\item Suponha que $\gamma,\sigma:[0,2]\to M$
são geodésicas distintas
 e satisfazem: $\gamma(0)=\sigma(0):=p$,
 $\gamma(1)=\sigma(1):=q$,
$\gamma$ e $\sigma$ realizam a distância
entre $p$ e $q$.
 Mostre que $\gamma$ não realiza a distância
 entre $p$ e $\gamma(1+s)$
para nenhum $s>0$.
\end{enumerate}
\end{exercise}

\begin{proof}
\begin{enumerate}
\item Primeiro suponha que $\gamma$ é suave.
Podemos ver que trata-se de uma
geodésica usando seu campo de velocidades
como campo variacional na primeira
fórmula da variação:
$$
E'(0)=\int_0^a |\gamma'|=0 \implies
\gamma'\equiv 0
$$
Uma geodésica quebrada não pode ser
minimizante porque em uma vizinhança
totalmente normal (cf.
\ref{proposition-total-norma-neigh})
do ponto singular podemos achar uma geodésica
radial (suave)
 que acurta a distância entre qualquer ponto
 antes do ponto singular
 e qualquer ponto depois do ponto singular.
\item Se as duas geodésicas se intersectam
não tangencialmente, podemos
construir um caminho mais curto entre $p$ e
$\gamma(1+s)$ ao longo de $\sigma$,
a menos de suavizar a quina perto de $q$. Se
as curvas se intersectam
 tangencialmente devem ser a mesma.
\end{enumerate}
\end{proof}

\section{Lista 5}
\label{section-lista-5}

\begin{exercise}
\label{exercise-l5-5}
\cite{doc} Capítulo XII, Exercício 6. Uma
geodésica $\gamma:[0,\infty)\to M$ em
uma variedade Riemannana $M$ é um {\it raio
partindo de $\gamma(0)$} se ela é
minimizante entre $\gamma(0)$ e $\gamma(s)$
para todo  $s \in (0,\infty)$.
Admita que $M$ é completa, não compacta, e
seja $p \in M$. Mostre que $M$ contém
um raio partindo de $p$.
\end{exercise}

\begin{proof}
Como $M$ é compacta e completa, ela não pode
ser limitada. Então
existe uma sequência de pontos $p_i$  tal que
$\lim_{i \to \infty}
d(p_i,p)=\infty$. Para cada ponto podemos
pegar uma geodésica minimizante
$\gamma_n$ ligando $p$ e $p_n$. Podemos
associar a cada geodésica um vetor
unitário $v_i \in S^n \subset T_pM$ tal que
$\gamma_i(t)=\text{exp}_p(tv_i)$.

Como $S^n \subset T_pM$ é compacta existe um
vetor $v$ limite da sequência
$v_i$. A geodésica
$\gamma(t):=\text{exp}_p(tv)$ é um raio, pois
por
continuidade de $\text{exp}_p$ e da distância
Riemanniana,
$$
d(\gamma(0),\gamma(t))=d\left(p,\text{exp}
_p(tv)\right)
=d\left(p,\text{exp}_p\left(t\lim_{i \to
\infty} v_i\right)\right)
=\lim_{i \to \infty} d(p,\gamma_i (t))
$$
Pegando $i$ suficientemente grande, teremos
que $\gamma_i$ é minimizante entre
$p$ e $\gamma_i(t)$. Isso significa que
$d(p,\gamma_i(t))$ está dada como o
comprimento de $\gamma_i$ até esse ponto, e
pegando o limite concluímos que o
comprimento de $\gamma$ é a distância entre
$\gamma(0)$ e $\gamma(t)$.
\end{proof}

\begin{definition}
\label{definition-line-riemannian-manifold}
Let $(M,g)$ be a Riemannian manifold. A
geodesic $\gamma:\mathbb{R} \to M$ is
called a {\it line} if
$d(\gamma(t),\gamma(s))=|t-s|$ for all $t,s
\in
\mathbb{R}$.
\end{definition}

\begin{exercise}
\label{exercise-l5-6}
Mostre que toda métrica completa em $S^n
\times \mathbb{R}$ admite uma linha.
\end{exercise}

\begin{proof}
Considere o conjunto $S^n\times
\mathbb{R}\setminus(S^n \times\{0\})$.
Trata-se de um aberto disconexo, e portanto
não fechado, nem compacto nem
limitado.  Dentro desse conjunto podemos
pegar duas sequências de pontos $p_i$ e
$q_i$, onde a segunda coordenada de $p_i$ tem
signo positivo, e a segunda
coordenada de $q_i$ tem  signo negativo para
toda $i$. Considere a geodésica
$\gamma_i$ que liga $p_i$ com $q_i$, que deve
passar pela esfera $S^n\times
\{0\}$. Então obtemos uma sequência $v_i$ de
vetores no fibrado tangente
unitário de $S^n\times\{0\}$ dadas como as
velocidades das geodésicas
$\gamma_i$. Como tal fibrado é compacto,
achamos um limite $v$ dessa sequência
que por um argumento análogo ao exercício
anterior, converge a um vetor cuja
geodésica associada é uma linha. ($\gamma$ é
minimizante, e como está
parametrizada por comprimento de arco obtemos
a condição dada na Definição
\ref{definition-line-riemannian-manifold}.
\end{proof}

\section{Lista 6}
\label{section-lista-6}

\begin{exercise}
\label{exercise-posit-scala-but-no-posit}
Encontre um exemplo de variedade suave que
admite alguma métrica Riemanniana com
curvatura escalar positiva, mas não admite
uma métrica Riemanniana com Ricci
positivo.
\end{exercise}

\begin{proof}
Se o problema for para $\text{Ric}\geq
\delta>0$, poderíamos usar o teorema de
Bonnet-Myers \ref{theorem-Bonnet-Myers} como
segue. Um produto de uma esfera com
uma reta tem curvatura escalar positiva, mas
não admite uma métrica com
$\text{Ric}\geq\delta>0$ porque não é
compacto.

{\bf Ideia errada:} Considere $S^1\times
S^1$. É claro que na métrica produto
temos curvatura escalar positiva (erro! As
variedades de dimensão 1 tem tensor
de curvatura nulo pela antisimetria nas
primeiras (ou últimas) duas entradas!)
pelo Exercício
\ref{exercise-curvatures-in-product}. Como é
uma superfície, a
curvatura de Ricci coincide com a curvatura
seccional. (Ricci é o promédio das
curvaturas seccionais.) A curvatura seccional
é a curvatura Gaussiana. Se
$S^1\times S^1$ tivesse uma métrica com
curvatura Gaussiana positiva, pelo
teorema de Gauss-Bonnet teria característica
de Euler positiva, mas a sua
característica de Euler é zero. No nosso
curso, aplicamos o teorema de Synge
\ref{theorem-Synge} para uma variedade
orientável de dimensão par com curvatura
escalar positiva, que deve ter grupo
fundamental trivial, uma contradição com o
fato de que $\pi_{1}(S^1\times
S^1)=\mathbb{Z}\times\mathbb{Z}$.

Nem tudo está perdido: considere
$\mathbb{R}P^{2}$. Na métrica induzida pela
projeção quociente desde $S^2$ temos
curvatura escalar constante 1, e é uma
superfície, de modo que a curvatura seccional
coincide com a curvatura de Ricci,
aplicamos o teorema de Synge
\ref{theorem-Synge}, mas agora
$\mathbb{R}P^{2}$
não é orientável, então seu grupo deveria ser
$\mathbb{Z}_2$… pois é.

Depois reparei que esse argumento está
totalmente errado: no caso das
superfícies nem só a curvatura escalar
coincide com a Ricci… também com a
seccional!

 $S^3\times S^1$. Então sim: pelo teorema de
Synge \ref{theorem-Synge}, uma variedade
orientável de dimensão par com
curvatura {\it escalar} positiva deveria ter
grupo fundamental trivial, que não
é o caso. Mas já não é verdade que a
curvatura de Ricci coincide com a curvatura
 seccional; poderíamos ter um plano com
 curvatura seccional negativa e ainda ter
Ricci positivo. Temos que descartar a
possibilidade de ter um plano com
curvatura seccional não positiva.
\end{proof}

\begin{exercise}
\label{exercise-posit-ricci-but-no-posit}
Encontre um exemplo de variedade suave que
admite uma métrica Riemanniana com
Ricci positivo, mas não admite uma métrica
Riemanniana com curvatura seccional
positiva.
\end{exercise}

\begin{proof}
Considere $S^3\times S^1$. A métrica produto
tem $\text{Ric}>0$, já que ele é a
soma dos tensores de Ricci em cada fator de
acordo ao Exercício
\ref{exercise-curvatures-in-product} (a
curvatura de Ricci de $S^3$ é 1,
enquanto que a curvatura escalar de $S^1$ é
zero). Agora suponha que existe uma
 métrica com $K>0$. Então pelo teorema de
 Synge \ref{theorem-Synge}, como
$S^3\times S^1$ é compacta, de dimensão par e
orientável, o grupo fundamental
dela deveria ser trivial, mas não é o caso.
\end{proof}

\begin{exercise}
\label{exercise-two-total-geode-subma}
Let $(M,g)$ be a complete, connected
Riemannian manifold with positive
curvature. Let $M_1$ and $M_2$ be two totally
geodesic submanifolds such that
$\dim M_1+\dim M_2\geq n$. Show that $M_1\cap
M_2\neq \emptyset$.
\end{exercise}

\begin{proof}[Solução]
Suponha que a interseção não é vazia. Então
existe uma geodésica minimizante não
trivial que minimiza a distância entre $A$ e
$B$. Sabemos que essa geodésica
intersecta tanto $A$ quanto $B$
ortogonalmente. Queremos ver que $E''(0)<0$,
ou
seja que $\gamma$ não pode ser minimizante,
uma contradição.

Suponha que conseguimos construir uma
variação de $\gamma$ tal que todos os
pontos iniciais $f(s,0)$ estão em $A$, e o
mesmo acontece com os pontos finais,
i.e. $f(s,\ell) \in B$. Então as derivadas
das curvas $f(s,0)$ e $f(s,\ell)$ em
$s=0$ são perpendiculares a  $\gamma'(0)$ e
$\gamma'(\ell)$. Ou seja, temos que
$\left\langle{}\gamma'(0),\nabla_{\partial_s}
V\right\rangle{}|_{0}^\ell=0$.

Suponha ainda que essa variação tem campo
variacional paralelo. Então acabou
porque a segunda fórmula da variação diz que
\[E''(0)=-\int_0^\ell
\left\langle{}R_{\gamma'}V,V\right\rangle{}<0\]

Vamos construir essa variação. O passo
inicial é fácil: pegamos qualquer vetor
$v \in T_a A$ e transportamos paralelamente
ao longo de $\gamma$ para obter o
campo paralelo $V \in \mathfrak{X}_\gamma$.
Como $A$ é totalmente geodésica, a
variação
$f(s,t):=\text{exp}_{\gamma(0)}(sV(0))$ fica
dentro de $A$.

Para concluir basta ver que $ V(\ell) \in
T_bB$, já que $B$ também é totalmente
geodésica. Isso segue da última hipótese.
Podemos realizar esse processo para
cada vetor básico de  $T_a A$, obtendo $\dim
A$ vetores linearmente
independentes em $T_bB$ (já que o transporte
paralelo é uma isometria). Se
nenhum deles ficasse em $T_bB$, poderíamos
construir um espaço de dimensão $\dim
A$ estritamente contido em $T_b M\setminus
T_b B$, absurdo pois isso implicaria
que $\dim A< \text{codim} B$, porém, $\dim A+
\dim B \geq n$.

\textbf{Edit.} Depois de reler a prova do
exercício seguinte, parece que não era
necessário usar que $\gamma$ intersecta
ortogonalmente $A$ e $B$, pois a
variação que usei foi por geodésicas, então
$\nabla_{\partial_s}V=0$.
\end{proof}

\begin{exercise}
\label{exercise-free-homotopy-variation}
Seja $M^{2n}$ uma variedade Riemanniana de
dimensão par, completa, orientável e
com curvatura seccional $K>0$. Seja $\gamma$
uma geodésica fechada em $M$ de
comprimento $\ell(\gamma)$. Mostre que
existem curvas livremente homotópicas a
$\gamma$ em $M$, arbitrariamente próximas de
$\gamma$, que possuem comprimento
menor que $\ell(\gamma)$.
\end{exercise}

\begin{proof}
Considere qualquer vector unitário ortogonal
a $\dot \gamma$, defina $V\in
\mathfrak{X}_{\gamma}''$ como sendo o
transporte paralelo de $v$ ao longo de
$\gamma$. Note que $\dot V=0$. Defina a
variação
$f(s,t)=\text{exp}_{\gamma(t)}sV(t)$. Note
que
$\nabla_{\partial_s}f_s=0$.

A segunda fórmula da variação nos diz que
\[E''(0)=-\int_0^a \left\langle{}R_{\dot
\gamma}V,V\right\rangle{}=-\int_0^a K<0\] já que
$K>0$. Como $\gamma$ é uma geodésica,
temos que para $s$ pequeno
\[\frac{1}{a}\ell^2(c^s)\leq E(s) < E(0)=
\frac{1}{a}\ell^2(\gamma)\] Note que essa
variação pode não preservar o
fechamento das curvas. Para resolver isso
olhemos ao seguinte desenho:

Fica claro que se $V(0)=V(a)$ as curvas $c^s$
são fechadas para toda $s$. Então
a condição que precisamos é que
\[P^\gamma_{0,a}V(0)=V(a)\] onde $P$ é o
transporte paralelo. Dito de outra forma,
basta ver que $P$ tem um ponto
fixo---além de $\gamma'(0)=\gamma'(a)$, que
já é um ponto fixo. Esse argumento é
uma imitação da prova do teorema de
Weinstein: como $P$ preserva orientação, o
determinante dele deve ser positivo, ou seja,
o produto dos seus autovalores
deve ser positivo. Restringindo-nos ao
subespaço $\gamma'(0)^\perp$, que tem
dimensão ímpar, concluímos que para que o
produto dos autovalores de $P$ seja
positivo, deve ter pelo menos um que seja
$1$.
\end{proof}

\section{Lista 7}
\label{section-lista-7}

\begin{exercise}
\label{exercise-minim-impli-no-conju-point}
Seja $\gamma:[0,a]\to M$ uma geodésica em uma
variedade Riemanniana $M$.

 Prove que se $\gamma$ é minimizante, então
 $\gamma$ não possui pontos
 conjugados em  $(0,a)$. Encontre um exemplo
 de geodésica $\gamma:[0,a] \to M$
sem pontos conjugados que não é minimizante.
\end{exercise}

\section{Lista 8}
\label{section-lista-8}

\begin{exercise}
\label{exercise-l8-1}
Prop. 2.12 do capítulo XIII, \cite{doc}. Seja
$p \in M$.
Suponha exista um ponto $q \in C_m(p)$ que
realiza a distância de $p$ a
$C_m(p)$. Então:
\begin{enumerate}
\item ou existe uma geodésica
minimizante $\gamma$ de $p$ a $q$ ao longo da
qual $q$ é
conjugado a $p$,
\item ou existem exatamente duas geodésicas
minimizantes $\gamma$ e $\sigma$ de $p$ a
$q$; além disto,
$\gamma'(\ell)=-\sigma'(\ell)$,
$\ell=d(p,q)$.
\end{enumerate}
\end{exercise}

\begin{proof}
No final da secção 1 do Capítulo XIII está
especificado que as variedades que se
consideram no capítulo são completas.
Portanto podemos supor que existe
uma geodésica minimizante $\gamma$ ligando
$p$ e $q$.

Então podemos aplicar a Proposição
\ref{proposition-cut-point-chara}:
 um ponto $q$ é o cut point de $p$ ao longo
 de uma
geodésica minimizante $\gamma$ se e somente
se
alguma das seguintes condições é verdadeira:
(a) $q$ é o primeiro ponto conjugado a $p$ ao
longo de $\gamma$, ou
(b) existem duas geodésicas minimizantes
ligando $p$ e $q$.

Se (a) é verdadeira, terminamos. Se (b) é
verdadeira, temos que existe outra
geodésica minimizante $\sigma$ ligando $p$ e
$q$.

{\bf Intento pessoal (não funcionou):}
considere a variação por geodésicas
$$
f(s,t):=\text{exp}_{\gamma(t)}
(s\text{exp}_{\gamma(t)}^{-1}\sigma(t))
$$
cujo campo de Jacobi
$$
J(t)=d_{t\text{exp}_{\gamma(t)}^{-1}\sigma(t)
}
\text{exp}_p(\text{exp}_{\gamma(t)}^{-1}
\sigma(t))
$$
se anula em $t=0,1$. Pensei que esse campo de
Jacobi estava bem definido pelo
Corolário 2.8, Cap. XIII (Lema
\ref{lemma-outsi-cut-locus-exist-minim}), que
assegura
 que $\text{exp}_p$ é injetiva fora do cut
 locus de $p$. Porém, precisaria
 que a exponencial ao longo de $\gamma$,
 $\text{exp}_{\gamma(t)}$, for injetiva fora
 do cut locus de $p$.
Essa condição seria garantida se
$d(p,q)=i(M)$, ou seja, se a exponencial
 $\text{exp}_{p'}$ for injetiva na bola de
 raio $d(p,q)$ para qualquer
$p' \in M$. Em conclusão: minha variação não
está bem definida.
(Se estivesse, com
isso terminaria o exercício, pois obtemos que
o único caso em que $p$ não é
conjugado a $q$ é se
$\gamma'(\ell)=-\sigma'(\ell)$, i.e. se o
campo de Jacobi
associado à variação é nulo.)

\bigskip

A variação certa é produzida no \cite{doc}
supondo que
$\gamma'(\ell)\neq -\sigma'(\ell)$; isso
implica que existe um vetor
 $V \in T_q M$ tal que
$\left\langle{}\gamma'(\ell),V\right\rangle{}< 0$ e
$\left\langle{}\sigma'(\ell),V\right\rangle{}< 0$.

Uma maneria simples de conferir a existência
desse vetor $V$ é olhando para o
plano gerado por $\gamma'(\ell)$ e
$\sigma'(\ell)$ (usando que são linearmente
independentes). O conjunto de vetores $V$
nesse plano satisfazendo
$\left\langle{}\gamma'(\ell),V\right\rangle{}<0$ é um
semiespaço, e o mesmo acontece com
 os vetores satisfazendo
 $\left\langle{}\sigma'(\ell),V\right\rangle{}<0$. Esses dois
 semiespaços
devem ter interseção não vazia precisamente
porque
$\gamma'(\ell)\neq -\sigma'(\ell)$.

Agora usamos o fato de que $p$ e $q$ não são
conjugados para
achar uma vizinhança do vetor $\ell
\gamma'(0)$ onde $\text{exp}_p$ é um
difeomorfismo e assim levantar alguma
curva $r$ que realize o vetor $V$ ao espaço
tangente,
 digamos $v:(-\varepsilon,\varepsilon) \to
 T_q M$.

A variação então está dada por
$$
f(s,t):=\text{exp}_p\left(\frac{t}{\ell}v(s)
\right)
$$
Podemos aplicar a primeira fórmula da
variação, na qual: o termo com a
integral é zero porque $\gamma$ é uma
geodésica; o termo da sumatoria é zero
porque trata-se de uma variação suave; e o
primeiro termo dos extremos é zero
porque a variação fixa o ponto de partida.
Obtemos que
$$
E'(0)=\left\langle{}V,\gamma'(0)\right\rangle{}<0
$$
Note que Manfredo escreve a equação anterior
usando o funcional de distância.
Isso também é válido: de fato a primeira
fórmula da variação pode ser
deduzida de maneira análoga (feito em sala
para variações próprias)
 para o funcional de distância no caso de
 geodésicas parametrizadas por
comprimento de arco (cf. Teo. 6.3
\cite{ler}).

O fato da derivada do funcional de
comprimento ser negativa nos diz que para
valores perto de $s=0$ podemos achar curvas
com menor comprimento. Mais
precisamente, existe um $s>0$ tal que
$\ell(f(s,\cdot))<\ell(\gamma)$.

É claro que o mesmo procedimento genera uma
variação $\tilde{f}$ de $\sigma$, e
podemos supor que a mesma $s>0$ faz
$\ell(\tilde{f}(s,\cdot))<\ell(\sigma)$.

Concluímos seguindo o argumento do Professor
Manfredo: se
$\ell(\gamma_s)=\ell(\sigma_s)$, então temos
duas geodésicas com o mesmo
comprimento ligando $p$ e
$\gamma_s(\ell)=\sigma_s(\ell)$; isso
significa
que existe um ponto $\tilde{t} \in (0,\ell]$
tal que $\gamma_s(\tilde{t})$ é
 o cut point de $p$. Porém, isso contradiz o
 fato de que $q$ realiza a
distância entre $p$ e o seu cut locus.

Finalmente, se
$\ell(\gamma_s)<\ell(\sigma_s)$, segue que o
cut point de $p$
 ao longo de $\sigma_s$ está a distância
 menor do que $q$, que de novo contradiz
a nossa hipótese.
\end{proof}

\begin{exercise}
\label{exercise-l8-2}
Proposição 2.13, Cap. XIII \cite{doc}. Se a
curvatura seccional $K$ de uma
variedade Riemanniana completa $M$ satisfaz
$$
0< K_{\text{min}}\leq K\leq K_{\text{max}},
$$
Então
\begin{enumerate}
\item $i(M) \geq \pi/\sqrt{K_{\text{max}}}$,
ou
\item existe uma geodésica fechada $\gamma$
em $M$, cujo comprimento é menor do
que o de qualquer outra geodésica fechada em
$M$, tal que
$$
i(M)=\frac{1}{2}\ell(\gamma)
$$
\end{enumerate}
\end{exercise}

\begin{proof}
Suponha que (1) não é verdadeiro.

Note que pelo Teorema de Bonnet-Myers
\ref{theorem-Bonnet-Myers}, $M$ é compacta e
portanto $C_m(p)$ é compacto para
todo $p$. Isso nos permite achar dois pontos
$p$ e $q$ cuja distância é o raio
de injetividade  $i(M)$.

Então podemos usar o exercício anterior para
$p$ e $q$.  Primeiro vamos ver o
que acontece no caso da segunda possibilidade
daquele exercício, i.e. que
existam exatamente duas geodésicas ligando
$p$ e $q$, digamos $\gamma$ é
$\sigma$, tais que
$\gamma'(\ell)=-\sigma'(\ell)$ onde
$\ell=d(p,q)$.  Considere
$\gamma * \overline{\sigma}$, onde $*$ é a
concatenação de curvas e
$\overline{\sigma}$ é a curva $\sigma$
percorrida em sentido contrário. É claro
que $\gamma * \overline{\sigma}$ é uma
geodésica fechada de comprimento $2
d(p,q)$. Note que por definição
$d(p,q)=\ell=i(M)$, como queríamos.

Essa geodésica fechada tem a distância mínima
entre as geodésicas fechadas, já
que se tivéssemos alguma outra com distância
menor, podemos pegar um ponto
qualquer nela e o seu cut point ficaria a
distância exatamente a metade do
comprimento do laço (pois existem duas
geodésicas que chegam nele: cada
metade do laço partindo em direções opostas
do ponto inicial escolhido),
contradizendo o fato de que $i(M)=d(p,q)$.

Portanto, basta descartar o primeiro caso do
exercício anterior. Para chegar a
 uma contradição, suponha que $p $ é
 conjugado a $q$, i.e. que existe um campo
 de Jacobi $J$ que se anula em $p$ e em $q$.
 Podemos usar o teorema de Rauch
\ref{theorem-Rauch} para comparar esse campo
com $\tilde{J}$,
um campo em $S^n_{K_{\text{max}}}$, a esfera
de curvatura constante
 $K_{\text{max}}$. Como $K \leq
 K_{\text{max}}$,
concluímos que $\tilde{J}$ deve se anular em
$\ell=d(p,q)=i(M)$. Absurdo, pois
estamos supondo que (1) não é verdadeiro,
i.e. que
 $i(M) < \pi/\sqrt{K_{\text{max}}}$.
\end{proof}

\begin{exercise}
\label{exercise-l8-3}
\cite{doc}, Capítulo XIII, Proposição 3.4. Se
a curvatura seccional $K$ de uma
variedade Riemanniana $M^n$, compacta,
orientável e de dimensão par, satisfaz
$0<K \leq 1$, então $i(M) \geq \pi$.
\end{exercise}

\begin{proof}
Para obter uma contradição, suponha que
existe um ponto $p \in M$ tal que
 $d(p, C_m(p))<\pi$. Como $M$ é compacta,
 sabemos que $C_m(p)$ é compacto e
portanto existe uma geodésica $\gamma$
ligando $p$ com $q \in C_m(p)$ e
realizando a distância $d(p,C_m(p))$.

Pelo teorema de Rauch, sabemos que qualquer
geodésica não pode ter pontos
conjugados antes de alcançar comprimento
$\pi$. Explicitamente, se $J$ é um
campo de Jacobi ao longo de alguma geodésica
$\gamma$ tal que $J(0)=0$ e
$J(\ell(\gamma))=0$, comparando com um campo
$\tilde{J}$ em $S^n$ tal que
 $\tilde{J}(0)=0$, $|\tilde{J}'(0)|=|J'(0)|$
 e
$\left\langle{}J,\gamma'\right\rangle{}=\left\langle{}\tilde{J},
\tilde{\gamma}\right\rangle{}$,
concluímos que
 $|\tilde{J}|\leq |J|$ ao longo de $\gamma$.
 Como as geodésicas de $S^n$ não
 tem pontos conjugados antes de atingir
 comprimento $\pi$, concluímos que
$J$ não pode se anular antes desse ponto.

Isso mostra que, como estamos supondo que
$d(p,C_m(p))<\pi$, devemos estar no
caso (2) do Exercício \ref{exercise-l8-1},
i.e. existem exatamente duas
geodésicas $\gamma$ e $\sigma$ ligando $p$ e
$q$, e
$\gamma'(\ell)=-\sigma'(\ell)$. Repetindo o
procedimento do Exercício anterior
 \ref{exercise-l8-3}, podemos usar essas duas
 geodésicas para achar uma
geodésica fechada de comprimento $2i(M)$,
cujo comprimento é menor do que o
comprimento de qualquer outra geodésica
fechada.

\begin{remark}
\label{remark-curva-posit-impli-maior-que}
Na prova do teorema de Synge
\ref{theorem-Synge}, o Professor
Manfredo afirma que ``Como $M$ é compacta e
tem curvatura positiva, $K\geq
\delta>0$"; que não me parece imediato, pois
a curvatura seccional não é uma
função definida em $M$. Porém, não preciso
mostrar isso para garantir a
existência do loop geodésico (via o exercício
anterior); apenas é necessário
que $M$ seja compacta para garantir a
existência do ponto $q$.
\end{remark}

\medskip

Para concluir seguimos a prova de \cite{doc}.
Vamos a usar o exercício 6 da
lista 6, (cf.
\ref{exercise-free-homotopy-variation}), onde
mostrei que, nestas condições, i.e.
$M$ de dimensão par, completa, orientável e
com curvatura seccional positiva,
existem curvas livremente homotópicas a
qualquer geodésica fechada $\gamma$, que
possuem comprimento menor que $\ell(\gamma)$.
Lembre que essas curvas foram
construídas mediante a função exponencial
aplicada a um campo paralelo ao longo
de $\gamma$, de forma que são curvas suaves;
vou usar esse fato depois.

A ideia é chegar numa contradição mostrando
que existe uma terceira geodésica
ligando $p$ e $q$. Essa geodésica se obtém
como o limite de uma sequência de
geodésicas associada à variação dada pelo
Exercício 6 da Lista 6.

Seja $c_s$ uma das curvas da variação com
comprimento estritamente
menor do que $\ell(\gamma)$. Pegue o ponto
inicial $c_s(0):=\tilde{p}_s$ e o
ponto $\tilde{q}_s$ em $c_s$ tal que
$d(\tilde{p}_s,\tilde{q}_s)$ é máxima ao
longo de $c_s$. Existe uma geodésica
$\tilde{\gamma}_s$ que liga
$\tilde{p}_s$ e $\tilde{q}_s$.

A seguir mostrarei que tomando limite quando
$s \to 0$, obteremos uma terceira
geodésica minimizante $\tilde{\gamma}_0$ que
liga $p$ e $q$, que não é possível
de novo pelo Exercício \ref{exercise-l8-1}.

De fato, podemos dar explicitamente
$\tilde{\gamma}_0$ como sendo $exp_q(tw)$
onde $w$ é o vetor limite das velocidades
iniciais de cada $\tilde{\gamma}_s$,
supondo que elas são de tamanho 1 para
assegurar a existência do limite por
compacidade do fibrado tangente unitário.
Note que a geodésica obtida minimiza a
distância entre $p$ e $q$: a geodésica liga
$p$ com  $q$ porque, por um lado, o
limite dos pontos iniciais é $p$ por
definição, e, por outro lado, porque os
pontos finais são os pontos que maximizam a
distância dentro de cada geodésica.
A geodésica limite é minimizante por
continuidade da função distância.

Como cada curva $c_s$ é diferenciável, temos
um vetor tangente a ela em cada
ponto $\tilde{q}_s$. Como $\tilde{\gamma}$ é
minimizante, quando aplicamos a
fórmula da primeira variação,vemos que
$\tilde{\gamma}'_s$ é ortogonal a
$c_s'(\tilde{q}_s)$. Como o a métrica é
contínua, concluímos que
$\tilde{\gamma}_0'(q) \perp \gamma'(q)$.
\end{proof}

\begin{exercise}
\label{exercise-l8-11}
Seja $M$ uma variedade Riemanniana completa,
de curvatura seccional não
negativa. Sejam $\gamma,\sigma:[0,\infty) \to
M$ geodésicas tais que
$\gamma(0)=\sigma(0)$. Se $\gamma$ é um raio
e
$\angle(\gamma'(0),\sigma'(0))<\frac{1}{2}
\pi$,
então $\lim_{t \to \infty}
\rho(\sigma(0),\sigma(t))=\infty$ (i.e.
$\sigma$ vai para o infinito).
\end{exercise}

\begin{proof}[Meu intento de prova]
Seja $t \in \mathbb{R}$. Considere uma
geodésica minimizante $\tau_t$ ligando
$\gamma(t)$ com $\sigma(t)$. Defina o angulo
$\beta_t:=\angle(-\gamma'(t),\tau'_t(0))$ e o
número $s_t$ como sendo o tempo em que
$\tau_t$ chega em $\sigma$, i.e.
$\tau_t(s_t)=\sigma(t)$. Ver Figura

Como $0 \leq  K$, podemos comparar a
``hinge'',
i.e. a informação do angulo e vetores
incidentes no ponto $\gamma(t)$,
$(-\gamma'(t),\tau_t'(0),\beta)$ com uma
hinge no espaço
euclidiano. Obtemos que
$$
\rho_{\mathbb{R}^n}(\tilde{\gamma}(0),
\tilde{\tau}_t(s_t))\leq
\rho(\gamma(0),\tau_t(s_t))=\rho(\sigma(0),
\sigma(t))
$$
Podemos calcular a distância
$\rho_{\mathbb{R}^n}(\tilde{\gamma}(0),
\tilde{\tau}_t(s_t))$
usando a lei de
cosenos euclidiana desde que conheçamos
$\beta_t$ e $\ell(\tau_t)$:
$$
\rho_{\mathbb{R}^n}(\tilde{\gamma}(0),
\tilde{\tau}_t(s_t))^2=
t^2+\ell(\tau_t)^2-2t\ell(\tau_t)\cos \beta_t
$$
Para calcular $\ell(\tau_t)$ podemos usar o
Teorema de Toponogov de novo, dessa
vez comparando a hinge $(\gamma'(0),
\sigma'(0),\alpha)$ com uma correspondente
no espaço euclidiano:
Obtemos que
$$
\rho_{\mathbb{R}^n}(\tilde{\gamma}(t),
\tilde{\sigma}(t))\leq
\ell(\tau_t)
$$
Como $\alpha$ é constante conforme $t$ muda,
podemos fixar $\tilde{\gamma}$ e
$\tilde{\sigma}$ para fazer a nossa
comparação. Conforme  $t$ avança, nos
movemos a velocidade constante ao longo de
$\tilde{\gamma}$ en direção ao
infinito, e embora não sabemos se a distância
$\rho(\sigma(0),\sigma(t))$ cresse
 ou diminue, a distância do nosso interesse
$\rho_{\mathbb{R}^n}(\tilde{\gamma}(t),
\tilde{\sigma}(t))$
vai sempre crescendo, e concluimos que
diverge ao infinito.

Portanto o problema acaba se conseguimos
achar uma cota inferior positiva para
$\beta_t$. Parece que isso é equivalente a
provar que a distância entre
$\gamma$ e $\sigma$ está acotada
inferiormente por uma constante positiva…

A lei de cosenos do lado Euclideano nos diz
que
$$
\rho_{\mathbb{R}^n}(\tilde{\gamma}(t),
\tilde{\sigma}(t))^2=
t^2+|\tilde{\sigma}(t)|^2-2t
|\tilde{\sigma}(t)|\cos \alpha
\leq t^2+|\tilde{\sigma}(t)|^2-2t
|\tilde{\sigma}(t)|
$$
Supondo que
$\tilde{\gamma}(0)=\tilde{\sigma}(0)=0 \in
\mathbb{R}^n$.
Parece que para medir essa distância
deveríamos conhecer
$|\tilde{\sigma}(t)|=\rho(\sigma(0),\sigma(t)
)$...
 o que parece levar o problema de volta ao
problema inicial.
\end{proof}

\begin{proof}[Prova de \cite{Cheeger-Ebin}]
Por desigualdade triangular, para quaisquer
$t,s \in \mathbb{R}$,
$$
\rho(\sigma(t),\sigma(0)) \geq
\rho(\gamma(s),\sigma(t))-\rho(\gamma(s),
\gamma(0)).
$$
Como o lado esquerdo não depende de $s$,
vemos que basta tomar o limite quando
$s\to \infty$ e mostrar que esse limite está
inferiormente limitado por um
número que depende linearmente de $t$, i.e.
$t \cos\alpha$.

Considere uma geodésica minimizante
$\tau_{s,t}$ ligando $\gamma(s)$ e
$\sigma(t)$. Podemos usar o teorema de
Toponogov \ref{theorem-Toponogov-angle}
na versão de ângulos (que está enunciada em
\cite{Cheeger-Ebin}) já que duas das
geodésicas no triângulo
$\gamma(0),\gamma(s),\sigma(t)$ são
minimizantes por
hipótese. Como a cota inferior da curvatura é
0 não temos problemas com o
comprimento de $\sigma$.

Obtemos que existe um triângulo Euclidiano
com lados do mesmo comprimento tal
que os ângulos correspondentes são menores o
iguais que os ângulos com que
começamos.

Denote $A:=\rho(\gamma(0),\gamma(s))=s$,
$B:=\ell(\sigma(t))=t$ e
 $C:=\rho(\gamma(s),\sigma(t))=\ell(\tau_{s,
 t})$.
Nosso objetivo é limitar inferiormente $C-A$
conforme  $s \to \infty$.
 Aplicamos a lei de cosenos no lado
 Euclidiano para obter
$$
\begin{aligned}
C^2&=A^2+B^2-2AB\cos\tilde{\alpha}\\
\implies C^2-A^2&=B^2-2AB\cos\tilde{\alpha}\\
\implies
C-A&=\frac{B^2-2AB\cos\tilde{\alpha}}{C+A}
\end{aligned}
$$
Onde só $A=s$ e $C$ dependem de $s$. O
comportamento assintótico desse quociente
quando $s \to \infty$ é claro
geometricamente: como $A$ e $C$ são dois
lados de
um triângulo, eles crescem
proporcionadamente. Formalmente, como $C \leq
A+B$,
obtemos $C+A \leq 2A+B$. Daí
$$
\begin{aligned}
\frac{1}{C+A}&\geq \frac{1}{2A+B}\\
\implies \frac{B^2-2AB\cos
\tilde{\alpha}}{C+A}&\geq
 \frac{B^2-2AB \cos \tilde{\alpha}}{2A+B}
=\frac{A}{A}\frac{B^2/A-2B\cos
\tilde{\alpha}}{2+B/A}
\end{aligned}
$$
que (a menos de analizar o que acontece se o
numerador é negativo…)
 converge a $B\cos \tilde{\alpha}$ conforme
 $s \to \infty$.
 Concluímos que $C-A$ deve
convergir a
$-t\cos\tilde{\alpha}\leq-t\cos\alpha$.
Tirando o signo obtemos a
desigualdade desejada.
\end{proof}

\clearpage

% --- End differential-geometry.tex ---

% --- Begin algebraic-geometry.tex ---
\chapter{Algebraic Geometry}
\label{chapter-algebraic-geometry}





\section{Algebraic sets}
\label{section-algebraic-sets}

\noindent
The {\it affine space} $\mathbb{A}^n$ is the
Cartesian
product of the field $n$ times.
An {\it algebraic set} is the zero set of an
ideal
in $k[x_1,…,x_n]$.

$V(I)$ and lots of properties…

\begin{lemma}
\label{lemma-irreducible-variety-prime-ideal}
An algebraic set is irreducible if and only
if
its defining ideal is prime.
\end{lemma}

\begin{proof}
First notice that in
the case that $Y$ is an irreducible
subvariety,
this ideal is prime: if $fg \in I_Y(U)$,
then $V(f) \cup  V(g) \supset Y\cap U$
so that $(V(f) \cup  V(g))\cap Y = Y\cap U$,
and then it must be true that $V(f)\cap
Y\subseteq V(g)\cap Y$
or that $V(f)\cap Y \supseteq V(g)\cap Y \cap
U$
since $Y$ is irreducible, and thus, say,
$V(f)\supseteq Y \cap U$,
that is, $I(V(f))\subseteq Y \cap U$
so that $f \in I(Y \cap U)$.
\end{proof}

\begin{theorem}[Weak Nullstellensatz]
\label{theorem-weak-nullstellensatz}
(Misha's formulation.) Let $A \subset
\mathbb{C}^n$
be an affine variety and $\mathcal{O}_A$ the
ring of polynomial
functions on $A$.
Then every maximal ideal in $A$
is an ideal of a point $a \in A$.
\end{theorem}

\noindent
The idea behind the radical idea is that
we can take roots of elements:

\begin{definition}
\label{definition-radical-ideal}
$I \subset R$ is {\it radical}
if for every $x \in R$ we have $x^n \in I
\implies x \in I$.
\end{definition}

\begin{definition}
\label{definition-nilpotent}
An element $a \in R$ is called {\it
nilpotent}
if $a^n = 0$.
\end{definition}

\begin{exercise}
\label{exercise-no-nilpotents-iff-radical}
$A/I$ has no nilpotents if and only if $I$ is
a radical ideal.
\end{exercise}

\noindent
We may understand the Nullstellensatz
as saying that algebraic subsets of an affine
variety
are the radical ideals of its coordinate
algebra.

\begin{theorem}[Nullstellensatz]
\label{theorem-nullstellensatz}
$$
I(V(I))=\sqrt{I}
$$
where $I$ is an ideal of $k[x_1,…,x_n]$.
\end{theorem}

\begin{proof}
Prove the Weak Nullstellensatz first
and then use Rabinowitsch trick
(which uses localization).
\end{proof}

\noindent
(Prove this:)
$k[x_1,…,x_n]/I$ has no nilpotents if and
only if
$I = \sqrt{I}$.

\medskip\noindent
Now we discuss how to decompose a variety in
irreducible components.

\begin{definition}
\label{definition-irreducible-variety}
An algebraic set is {\it irreducible}
if it cannot be decomposed as a union of two
properly
contained algebraic sets.
\end{definition}

\begin{lemma}
\label{lemma-irred-iff-coord-ring-reduc}
An affine variety is irreducibly if and only
if
its coordinate ring has no zerodivisors.
\end{lemma}

\begin{proof}
A zero divisor gives an expression $fg=0$.
\end{proof}

\begin{exercise}
\label{exercise-surje-map-irred-varie}
If $\varphi:A \to B$ is a surjective morphism
of varieties
and $B$ is irreducible, so is $A$.
\end{exercise}


To prove the decomposition in
irreducible ideals we must discuss
Noetherianity.

\begin{definition}
\label{definition-noetherian-ring}
A ring $R$ is {\it noetherian}
if any ascending chain of ideal stabilizes.
\end{definition}

\begin{theorem}[Hilbert's basis theorem]
\label{theorem-hilberts-basis}
If $R$ is a Noetherian ring, so is $R[x]$.
Equivalently, every ideal in $R[x]$ is
finitely generated.
\end{theorem}

\begin{lemma}
\label{lemma-noeth-iff-ideal-finit-gener}
A ring $R$ is Noetherian
if and only if
all its ideals are finitely generated.
\end{lemma}

\begin{proof}
Try!
\end{proof}

\begin{theorem}[Hilbert's basis theorem,
Misha's version]
\label{theorem-hilberts-basis-misha}
Any finitely generated ring is Noetherian.
\end{theorem}

\noindent
The idea is to use this theorem
to show that any algebraic set may be
decomposed
as a finite union of irreducible varieties.
The idea of the proof is to reduce the
original algebraic set
to a union of two proper subvarieties, and
repeat the process until we arrive at
irreducible varieties.

\begin{definition}
\label{definition-noeth-topol-space}
A topological space is {\it Noetherian}
if any descendending sequence of closed
subsets stabilizes.
\end{definition}

\begin{definition}
\label{definition-zariski-topology}
{\it Zariski topology} is the topology
whose closed sets are algebraic sets.
\end{definition}

\begin{exercise}[Misha]
\label{exercise-varie-zaris-compa}
Any variety is Zariski compact.
\end{exercise}

\noindent
To prove that any variety has a decomposition
as a finite union of irreducible components,
we first show pointwise that every point
is contained in a finite union of irreducible
components.
With existence at hand we just prove
in a similar way the existence of the
decomposition
for all the variety.

\medskip\noindent
I don't know where to put this yet:

\begin{definition}
\label{definition-degree}
\begin{reference}
\href{https://en.wikipedia.org/wiki/%
Degree_of_an_algebraic_variety}{Wikipedia}.
\end{reference}
The {\it degree} of an affine or projective
variety $X$ of dimension $n$ is the number of
points
of intersection of $X$ and $n$ hyperplanes in
general position.
\end{definition}

\begin{exercise}
\label{exercise-hypersurfaces-of-pn}
\begin{reference}
\cite[Exercise I.2.8]{hart}
\end{reference}
A projective variety $Y \subseteq
\mathbb{P}^n$
has dimension $n-1$ if and only if
it is the zero set of a single irreducible
homogeneous polynomial
$f$ of positive degree.
\end{exercise}

\begin{proof}
$Y$ has dimension $n-1$ if and only if there
is a chain
of irreducible subsets $Y \supset Y_2\supset…
\supset Y_{n-1}$
if and only if there is a chain of
prime ideals $I(Y) \subset I_2
\subset…\subset I_{n-1}$.
But $S:=k[x_0,…,x_n]$ has Krull dimension
$n+1$,
so that $I(Y)$ has height 1
(not considering the irrelevant ideal) and
thus principal
since $S$ is a UFD.
\end{proof}

\begin{proof}
($\implies$) Suppose $Y$ ir irreducible
and $\dim Y=n-1$. Pick any nonzero
homogeneous $g \in I(Y)$.
Factor $g$ into homogeneous irreducible
facts and choose one such facto $f$.
Then $f$ is homogeneous, irreducible and
vanishes on $Y$,
so $Y \subseteq V(f)$.
Both $Y$ and $V(f)$ are irreducible
projective
varieties and have the same dimension $n-1$,
hence $Y=V(I(f))$.

($\impliedby$) Suppose $Y=V(f)$ for some
nonconstant homogeneous polynomial $f$.
Then it has dimension $n-1$.
\end{proof}

\section{Regular functions, rational
functions and morphisms}
\label{section-regular-functions}

\noindent
We start with the intuitive definition, yet
different from Harshorne's,
which works well for affine varieties.
A {\it regular} or {\it polynomial map}
is $\varphi:V \subset \mathbb{A}^n \to W
\subset \mathbb{A}^m$
if there exist $P_1,…,P_m \in k[x_1,…,x_n]$
such that
$\varphi(a_1,…,a_n)=(P_1(a_1,…,a_n),…,
P_m(a_1,…,a_n))$
and $\varphi(V) \subset W$.
If $\varphi$ is polynomial then there is an
induced
map in coordinate rings
$\varphi^*:k[W] \to k[V]$ given by pullback.
This map is well defined since
$\varphi^*(I(W))\subset I(V)$.

A polynomial map $\varphi:V \to W$ is an {\it
isomorphism}
if it has a left and right inverse that is
also a polynomial function.
$V \cong W \iff k[V] \cong k[W]$ for affine
varieties.

\medskip\noindent
Here is Hartshorne's definition:

\begin{definition}
\label{definition-regular-function}
Let $Y$ be a quasi-affine variety
(i.e. an open set of an affine variety).
A function $Y \to k$ is called {\it regular}
at $P \in Y$
if there is an open set $U \ni P$ and
$g,h \in k[x_1,…,x_n]$ such that
$f=g/h$ on $U$. $f$ is {\it regular} if it is
regular
at all points of $Y$.
\end{definition}

\noindent
And the definition for projective varieties:

\begin{definition}
\label{definition-regul-funct-proje}
Let $Y \subset \mathbb{P}^N$ be a
(quasi)projective variety.
A function $f:Y \to k$ is regular at $P \in
Y$
if there is an open set $U \ni P$ such that
$f=g/h$ for two homogeneous polynomials of
the same degree,
with $h$ nonzero in $U$.
\end{definition}

\noindent
Note that although homogeneous polynomials
are not functions on $\mathbb{P}^n$
(because we can rescale points, but the
functions will give different outputs),
the quotient of two homogeneous polynomials
is (because the rescaling factor
cancels out).

\begin{definition}
\label{definition-morphism}
A {\it morphism} of varieties $\varphi:V \to
W$
is a function that pulls back polynomial
functions
$f:W \to k$ to polynomial functions
$\varphi^*f:V \to k$.
\end{definition}

\begin{exercise}
\label{exercise-morphisms-def}
$\varphi$ is a morphism of affine varieties
if and only if
it is a polynomial function.
\end{exercise}

\begin{proof}
If $\varphi$ is polynomial then it is a
morphism:
if $f:W \to k$ is polynomial,
then
$\varphi^*f(a_1,…,a_n)=f(P_1(a_1,…,a_n),…,
P_m(a_1,…,a_n))$
is also polynomial.
Conversely, if $\varphi$ is a morphism
we see that its coordinate functions are
polynomials
by computing $y_i\varphi^*$ where $y_i:W \to
k$
is $i$-th coordinate function (which is
clearly a polynomial map).
\end{proof}

\medskip\noindent
Now we introduce the notion of local ring
(stalk).
Here's Olivier's (intuitive) definition:
Let $V$ be an affine variety (:= irreducible
algebraic set)
and $p \in V$.
$$
\mathcal{O}_p(V)
:= \{f \in k(V):\exists g,h \in k[V]
\text{ s.t. }f=g/h \text{ and } h(p) \neq
0\}\subset k(V).
$$

Compare with Hartshorne's:

\begin{definition}
\label{definition-local-ring}
Let $Y$ be a variety (affine or projective)
and $P \in Y$.
The {\it local ring} of $P$ on $Y$ is
the ring of germs of regular functions on $Y$
near $P$.
\end{definition}

\begin{exercise}
\label{exercise-local-ring-is-local}
The local ring is local with maximal ideal
given by functions vanishing at $P$.
\end{exercise}


\begin{exercise}
\label{exercise-stalk-affine-varieties}
Show that if $V$ is an affine variety,
$\mathcal{O}_p(V) =
k[V]_{\mathfrak{m}_p}$
where $\mathfrak{m}_p$ is the set of regular
functions
that vanish in $p$.
\end{exercise}

\begin{proof}
Consider the obvious map: take a quotient of
classes
$[g]/[h]$ in the localization of the
coordinate ring of $V$
and map it to the rational function $g/h$.
This is well defined because if $g-g' \in
I(V)$,
then they obviously coincide in any open set
of $V$
because they coincide in all of $V$, hence
forming the same germ.
The key fact is that $\mathcal{O}_p(V)$ is
the germs of rational
functions of $V$.
It is injective because if a germ $g/h$ is
zero, then there's
an open of $V$ where $g$ vanishes, implying
that $g=0$
since it is a polynomial (or holomorphic)
map.
It is obviously surjective.
\end{proof}

\begin{definition}
\label{definition-degree-zero-localization}
Let $S=\bigoplus_{d\geq 0} S_d$ be a
graded ring and let $f \in S$ be
homogeneous.
The localization $S_f$ is again graded by
\[
\deg(a/f^m)=\deg(a)-m\deg(f)
\]
for homogeneous $a \in S$.
We write $(S_f)_0$ for its degree-zero
part.
Thus $(S_f)_0$ consists of the fractions
$a/f^m$ such that
\[
\deg(a)=m\deg(f).
\]
More generally, if $\mathfrak{p}$ is a
homogeneous prime of $S$, we write
$S_{(\mathfrak{p})}$ for the degree-zero
part of the localization of $S$ at the
homogeneous elements outside
$\mathfrak{p}$.
\end{definition}

\begin{proposition}
\label{proposition-standard-open-proj}
Let $Y=\operatorname{Proj} S$ and let
$f \in S$ be homogeneous.
Then the open set
\[
D_+(f)=\{\mathfrak{p}\in Y:f\notin
\mathfrak{p}\}
\]
is affine, with coordinate algebra
\[
\mathcal{O}_Y(D_+(f))=(S_f)_0.
\]
In other words, restricting to the chart
where $f$ does not vanish corresponds to
localizing at $f$ and taking degree zero.
\end{proposition}

\begin{theorem}
\label{theorem-functions-projective-case}
\begin{reference}
\cite[Theorem 3.4]{har}
\end{reference}
Let $Y \subseteq \mathbb{P}^n$ be a
projective variety
with homogeneous coordinate ring $S(Y)$.
Then:
\begin{enumerate}
\item $\mathcal{O}(Y)=k$,
\item
$\mathcal{O}_P=S(Y)_{(\mathfrak{m}_P)}$,
where $\mathfrak{m}_P$ is the ideal generated
by the homogeneous
$f \in S(Y)$ not vanishing in $P$,
\item $K(Y)\cong S(Y)_{((0))}$.
\end{enumerate}
\end{theorem}

\noindent
Now let's discuss the equivalence of
categories between
affine varieties and finitely generated
reduced algebras.

\begin{proposition}
\label{proposition-coordinate-algebra}
The coordinate algebra of an affine variety
$V$,
i.e. $k[x_1,…,x_n]/I(V)$ is isomorphic to
the ring of regular functions on $V$, i.e.
$\mathcal{O}(V)$.
\end{proposition}
The following proposition is also at
\href{https://stacks.math.columbia.edu/tag/%
01I1}{01I1}

\begin{proposition}
\label{proposition-irreducible-prime}
Let $V,W \subset \mathbb{A}^n$ are
irreducible algebraic set.
There is a bijection
$$
\Hom_{k-alg}(k[W],k[V]))
\xrightarrow{bij}\underbrace{\Mor(V,W)}_{
\text{polynomial maps}}.
$$
\end{proposition}

\begin{proof}
We map
$$
\begin{aligned}
\tilde{\varphi}:
\underbrace{k[W]}_{=k[y_1,…,y_m]/I(W)}
&\longrightarrow
\underbrace{k[V]}_{=k[x_1,…,x_n]/I(V)} \\
 y_1&\longmapsto f_1\\
&\vdots \\
y_m & \longmapsto f_m
\end{aligned}
$$

\noindent
to
$$
\begin{aligned}
\varphi: V &\longrightarrow W \\
(x_1,…,x_n) &\longmapsto (f_1(x_1,…,x_n), …,
f_m(x_1,…,x_n)).
\end{aligned}
$$

\noindent
We check that this is well-defined, i.e.
that it doesn't depend on the choice of
representatives,
namely the $f_i$. But choosing other $f_i$
which differ
by an element of the ideal would not make a
difference
when evaluating at $V$ because the ideal is
precisely $I(V)$.

The inverse map we already described above.
\end{proof}

\noindent
This shows that we have an equivalence of
categories:

\begin{lemma}
\label{lemma-algeb-varie-categ-equiv}
There is an equivalence of categories between
finitely generated reduced
$k$-algebras
and irreducible algebraic sets.
\end{lemma}


\medskip\noindent
Now let's discuss rational functions.
Again let's start with a more reasonable
explanation.
Since $k[V]$ for an irreducible affine
algebraic set is a domain,
there exists a fraction field, which is the
set of {\it rational functions}.
Functions here look like $f/g$ for $f,g \in
k[V]$,
which are not defined in $V(g)$.
Since this can lead to different ways
of defining the same function
--- the correct definition shall be that
we take the greatest set where the function
can be defined.

\begin{definition}
\label{definition-rational-ring}
The {\it ring of rational
functions} is the set of regular functions
defined
on open sets with the identifications that
$(U,f)\sim(V,g)$
if they coincide in an intersection. (Like
the stalk,
but there's no point.)
\end{definition}

Compare this to Voisin's way of describing
a meromorphic function:
it's a function that can be written locally
as a quotient of two holomorphic functions.



\section{Dimension: Krull dimension,
transcendence degree}
\label{section-dimension}

\begin{definition}
\label{definition-irreducible-topological}
A nonempty subset $Y$ of a topological space
$X$
is {\it irreducible} if it cannot be
expressed
as a union of two proper subsets,
each of which is closed in $Y$.
\end{definition}

\begin{definition}
\label{definition-topol-space-dimen}
If $X$ is a topological space,
the {\it dimension} of $X$ is the supremum of
all integers $n$
such that there exists a chain
$Z_0 \subset Z_1\subset … \subset Z_n$
of distinct irreducible closed subsets of
$X$.
The {\it dimension of an algebraic variety}
is its dimension as a topological space.
\end{definition}

\begin{definition}
\label{definition-krull-dimension}
In a ring $A$, the {\it height} of a prime
ideal $\mathfrak{p}$
is the supremum of all integers $n$
such that there exists a chain
$\mathfrak{p}_0
\subset\mathfrak{p}
_1\subset…\subset\mathfrak{p}_n=\mathfrak{p}$
of distinct prime ideals.
The {\it Krull dimension} of $A$ is the
supremum of the heights
of all prime ideals.
\end{definition}

\begin{proposition}
\label{proposition-dimension}
If $Y$ is an affine algebraic set,
the dimension of $Y$ is equal to the
dimension
of its coordinate ring.
\end{proposition}

\begin{theorem}
\label{theorem-dimen-trans-degre}
Let $k$ be a field and let $B$ be an integral
domain
which is a finitely generated $k$-algebra.
Then:
\begin{enumerate}
\item the dimension of $B$ is equal to the
transcendence degree
of the quotient field $K(B)$ of $B$ over $k$;
\item For any prime ideal $\mathfrak{p}$ in
$B$, we have
$$
\text{height }\mathfrak{p}+ \dim B /
\mathfrak{p}=\dim B.
$$

\end{enumerate}
\end{theorem}

\begin{exercise}
\label{exercise-affine-twisted-cubic}
Let $Y\subseteq \mathbb{A}^3$ be the set
$Y= \{(t,t^2,t^3:t \in k\}$.
Show that $Y$ is an affine variety of
dimension 1.
Find generators for the ideal $I(Y)$.
Show that $A(Y)$ is isomorphic to a
polynomial ring in one variable over $k$.
We say that $Y$ is given by the {\it
parametric representation}
$x=t$, $y=t^2$, $z=t^3$.
\end{exercise}

\begin{proof}
It's immediate that $x^2-y$ and $x^3-z$ must
be in
the ideal defining $Y$. Let
$J=(x^2-y,x^3-z)$.
Consider the ring map
$$
\begin{aligned}
\varphi:k[x,y,z] &\longrightarrow k[t] \\
x &\longmapsto t\\
y &\longmapsto t^2\\
z &\longmapsto t^3
\end{aligned}
$$

\noindent
Since $\varphi$ is given literally by the
parametrization of $Y$,
it follows from definition that the ideal of
$Y$
is the kernel of $\varphi$.

It's clear the in the kernel of $\varphi$
lie the polynomials $x^2-y$ and $x^3-z$.
But how to make sure that $J$ is the complete
kernel?

The procedure by AI is the following.
First notice that there is an isomorphism
$k[x,y,z]/J \cong k[x]$
given by $f(x,y,z)+J \mapsto f(x,x^2,x^3)$.
This shows that $k[x,y,z]/J$ is a domain, and
thus $J$ is prime.
Further, the projection $k[x,y,z] \to
k[x,y,z]/J$,
which obviously has kernel $J$,
is exactly the map I introduced before from
$k[x,y,z]$ to $k[t]$ when we see $k[x,y,z]/J$
as $k[t]$.
Since the kernel of such map is the ideal of
$Y$ by definition,
we are done.
\end{proof}

\section{Zariski tangent space, smooth
points}
\label{section-multi-smoot-point-bezou}

\noindent
Let's discuss Zariski tangent space. The
simplest notion of nonsingularity is what we
expect:

\begin{definition}
\label{definition-nonsingular-affine}
Let $Y \subseteq \mathbb{A}^n$ be ann affine
variety,
and let $f_1,…,f_t \in R=k[x_1,…,x_n]$ be a
set of generators
for the ideal of $Y$.
$Y$ is {\it nonsingular at a point $P \in Y$}
if the rank of the matrix
$\left(\frac{\partial f_i}{\partial
x_j}(P)\right)$
is $n-r$, where $r$ is the dimension of $Y$.
$Y$ is {\it nonsingular} if it is nonsingular
at every point.
\end{definition}

\noindent
It can be checked that this definition is
independent of the choice of generators. But
we look for a criterion that is more
algebraic.

\begin{definition}
\label{definition-residue-field}
Let $R$ be a ring and $\mathfrak{m}$ a
maximal ideal of $R$.
The {\it residue field} is  $R/\mathfrak{m}$.
\end{definition}

\begin{definition}
\label{definition-regular-local-ring}
Let $R$ be a noetherian local ring with
maximal ideal $\mathfrak{m}$
and residue field $k:=R/\mathfrak{m}$.
$R$ is a {\it regular local ring} if
$\dim_k \mathfrak{m}/\mathfrak{m}^2=\dim R$.
\end{definition}

\noindent
Upshot. $\mathfrak{m}$ is the ideal of
functions vanishing at $P$.
This means that the constant term of $h \in
\mathfrak{m}$ is zero.
Differentiate $h$.
Now the constant term is what originally was
the linear part.
If $h'$ has no constant term, that is, if $h$
vanishes at $P$
up to order 1,
then $h$ can be seen as a product $h=fg$ for
$f,g \in \mathfrak{m}$.
Indeed, the smallest nonzero term of $h$ is
of degree 2,
so we can make it look like a product of two
things of degree 1.
Conversely, products $h=fg$ vanish up to
order 1 since
the first-order coefficient is given by the
derivative
(Taylor expansion)
and $dh=df\cdot g+fdg =0$ at $P$.

Then we can interpret
$\mathfrak{m}/\mathfrak{m}^2$ as follows.
Among all the functions which vanish at $P$,
i.e. inside $\mathfrak{m}$,
we identify any two functions if their
difference
is in $\mathfrak{m}^2$,
that is, if their difference vanishes up to
order 1 at $P$,
that is, if their linear parts coincide at
$P$.

Thus, $\mathfrak{m}/\mathfrak{m}^2$ is the
space
of linear parts of polynomials vanishing at
$P$.

\begin{theorem}
\label{theorem-nonsi-jacob-iff-regul-local}
Let $Y \subseteq \mathbb{A}^n$ be an affine
variety.
Let $P \in Y$ be a point.
Then $Y$ is nonsingular at $P$
if and only if
the local ring $\mathcal{O}_{P,Y}$ is a
regular local ring.
\end{theorem}

\begin{proof}
Still hard to grasp completely,
but an important insight is that the map
$$
\begin{aligned}
\nabla: k[x_1,…,x_n]  &\longrightarrow k^n \\
 f &\longmapsto \nabla f=\left(\frac{\partial
 f}{\partial x_1},
…,\frac{\partial f}{\partial x_n}\right)
\end{aligned}
$$

\noindent
has kernel $\mathfrak{m}^2$ and is surjective
when restricted to $\mathfrak{m}$
(because $\nabla(x_i-p_i)=(0,…,0,1,0,…0)$, so
we construct a base of $k^n$)
so we have the isomorphism
$$
\mathfrak{m}/\mathfrak{m}^2=k^n,
$$
which we understand as the space of linear
directions at $P$.
\end{proof}

\begin{definition}
\label{definition-nonsingular}
Let $Y$ be any variety,
$Y$ is {\it nonsingular} at a point $P \in Y$
if the local ring $\mathcal{O}_{P,Y}$ is a
regular local ring,
that is,
if $\dim_k\mathfrak{m}/\mathfrak{m}^2=\dim
\mathcal{O}_{P,Y}$.
\end{definition}

\medskip\noindent
Now we discuss singularity theory for plane
curves.

Upshot: the multiplicity of a point
is how many tangent lines there are at the
point.

Smooth points of a plane curve are those with
multiplicity 1.
The multiplicity of a plane curve $F$ at the
origin
is the degree of the lowest nonzero
homogeneous term $F_m$
of $F=F_m+F_{m+1}+…+F_n$.
The tangent lines at the origin are the
linear factors of $F_m$, which decomposes in
linear components
by the following homogenization trick:

\begin{exercise}
\label{exercise-factor-in-linear-terms}
Let $f_m$ be a homogeneous polynomial of
degree $m$.
Show that $f_m$ is a product of $m$ linear
factors
$f_m=\prod (b_ix+c_iy)$.
\end{exercise}

\begin{proof}
Upshot: fix $y=1$, factor as product of
linear coefficients
using fundamental theorem of algebra, and
then rehomogenize.
The last is because a 1-variable polynomial
$F$ is homogeneous of degree $m$ if and only
if
$F(\lambda x)=\lambda^mF(x)$. So if you have
$F(x)=f(x,1)$
you rehomogenize by putting
$y^mF(x/y)=y^mf(x/y,1)$.
Which is strange because $y^mF(x/y)$ is in
$k[x,y,1/y]$,
but that doesn't matter because when you plug
it in you get
 $$
y^mf(x/y,1)=y^mc\prod(\frac{x}{y}-\alpha_i)=
c\prod(x-\alpha_iy),
$$
(where $c$ is a constant coming from making
the product having no
coefficients on $x$ after using fundamental
theorem of algebra)
which is already a product of linear
polynomials.
We conclude by showing that in fact
$y^mf(x/y,1)=f(x,y)$,
but that's just literally
$$
y^mf(x/y,1)=y^m\sum_{i}a_i \frac{x^i}{y^i}
=\sum_ia_i x^iy^{m-i}.
$$
\end{proof}

\noindent
Thus tangent lines can have multiplicity.
A {\it node} is an ordinary double point,
which means that there are two tangent lines
and they are not the same.
A {\it cusp} is a double point with a double
tangent line.

\begin{exercise}[Intersection multiplicity]
\label{exercise-intersection-multiplicity}
If $Y,Z \subseteq \mathbb{A}^2$ are two
distinct curves,
given by equations $f$ and $g$, and if $P \in
Y \cap Z$,
we define the {\it intersection multiplicity}
$(Y \cdot Z)_P$
of $Y$ and $Z$ at $P$ to be the length of the
$\mathcal{O}_P$-module $\mathcal{O}_P/(f,g)$.
\begin{itemize}
\item Show that $(Y\cdot Z)_P$ is finite,
and $(Y\cdot Z)_P \geq  \mu_P(Y) \mu_P(Z)$.
\end{itemize}
\end{exercise}


\medskip\noindent
The upshot about Bézout theorem is that the
``size'' of the intersection
of two curves is the product of their
degrees.
In nice cases this is just the number of
intersection points.
In even nicer cases one of the curves is a
line,
in which case the product of degrees is
just the degree of the curve,
so that we compute the degree of a curve by
{\it the number of intersection points with a
hyperplane}.

But the notion of ``size'' of the
intersection
has to be made precise, which is done via the
notion of multiplicity.
The idea is that intersecting with a line
$\ell=ax +by+c$ where, say,  $b \neq 0$,
can be done by substituting $y$
for $-\alpha x-b$.
This makes the ``resultant''
$f_\ell(x):=f(x,\alpha x+\beta)$
be a polynomial in one variable,
which can be factored as a product of linear
terms, each of the roots of which
are the intersection points,
and are now equipped with a multiplicity.

\begin{definition}
\label{definition-multiplicity-line}
The {\it multiplicity} or {\it intersection
index}
of $f\in k[x,y]\setminus k$ and
the line $\ell=ax+by+c$ (where $a \neq 0$ or
$b\neq 0$) is
$$
(\ell,f)_p =
\begin{cases}
0\qquad &\text{if }p \not\in V(f,\ell) \\
\infty \qquad &\text{if }p \in V(\ell) \text{
and }V(\ell) \subset V(f)\\
m_i \qquad & \text{if }p=(x_i,\alpha
x_i+\beta)
\end{cases}
$$
where $f_\ell(x):=f(x,\alpha x+\beta)$
factors as $c\prod(x-x_i)^{m_i}$.

If $V(\ell)\not \subset V(f)$, the integer
$$
\text{deg}(f)-\sum_{i=1}^rm_i
$$
is the {\it multiplicity of $\ell,f$ at
infinity}.
\end{definition}

\begin{exercise}
\label{exercise-asymptotic-direction}
If $f=\sum_{i=0}^df_i \in k[x,y]$ for $f_i$
homogeneous of degree $i$,
an {\it asymptotic direction} of $f$ is
$ax+by$
which divides $f_d$.

Show that  $m_\infty \neq 0$ if and only if
$y-\alpha x$
is an asymptotic direction of $f$.
\end{exercise}

\begin{proof}
This uses a theorem using the resultant…
\end{proof}

\begin{theorem}[Bézout]
\label{theorem-bezout}
Let $k$ be an algebraically closed field.
Let  $C,P \subset \mathbb{P}_k^2$ be two
curves without
irreducible components in common
and suppose that $C \cap D= \{p_1,…,p_r\}$
with intersection multiplicities  $m_1,…,m_r$
in $p_1,…,p_r$.
Then
$$
\text{deg}f\text{deg}g=\sum_{i=1}^r m_i.
$$
\end{theorem}


\begin{proposition}
\label{proposition-multiplicity-point}
Let $f \in k[x,y]\setminus k$ and $p \in
V(f)$.
There exists an integer $m_p(f)\geq 1$ called
the {\it multiplicity of $p$ in $f$}
such that for every line $\ell$ passing
through  $p$
 $$
(\ell,f)_p \geq m
$$
and $(\ell,f)_p>m$ for at most $m$ lines and
at least one line.
\end{proposition}

\begin{definition}
\label{definition-tangent-line}
Let $p=(x_0,y_0) \in V(f)$.
Write $f(x+x_0,y+y_0)=f_m(x,y)+\text{lower
degree terms} $,
i.e. translate $f$ so that the origin is in
its zero set
and write as homogeneous polynomials.
Then put
$f_m(x,y)=\prod_{i=1}^r(a_ix+b_iy)^{m_i}$.
The lines $\ell_i=a_ix+b_iy$ are called
{\it tangent lines} to $f$ and each has a
multiplicity $m_i$.
\end{definition}

\begin{definition}
\label{definition-smoot-ordin-point-curve}
A point $p \in V(f)$ is {\it smooth} if
$m_p(f)=1$
and  {\it ordinary} if all its tangent lines
have multiplicity 1.
\end{definition}

\noindent
Thus, ordinary points are singular,
but these are mild singularities.
Essentially they are points where $m$
straight lines intersect transversally.

Definition of intersection index of two plane
curves
without common components.
Definition of intersection index of two plane
curves
with common components.
Properties of intersection index (7 of them).
These 7 properties characterize uniquely the
intersection index.

The multiplicity of a root $\alpha$ of a
power series
$f\in k[[x]]$,
i.e. the greatest integer such that
$(x-\alpha)^r|f$,
equals $\dim_k k[[x]]/(x^r)=r$.
(\href{https://www.youtube.com/live/%
EspNBiSU8F4?si=UEcdxrysAS47ZA9H&t=5230}{
Here}'s why.)
This leads to the alternative definition
of intersection index as
$\dim_k
k[[x,y]]/(f(x-x_0,y-y_0),g(x-x_0,y-y_0))$.
This satisfies the 7 properties of the
intersection index.


\noindent
\begin{proposition}[Jacobian criterion]
\label{proposition-jacobian-criterion}
The singular locus of $V(f)$ is
$V(f,\partial_x,\partial_y)$.
\end{proposition}



\begin{theorem}[Plücker formulas]
\label{theorem-plucker-formulas} Let $C
\subset \mathbb{P}^2$ be an irreducible curve
of degree $d$ whose only singularities are
$\delta$ nodes and $\chi$ ordinary cusps
(cusps such that the tangent line intersects
with multiplicity 3 at the singular point).
Let \begin{itemize} \item $i$ be the number
of inflection lines (lines $\ell$ such that
$(\ell,C)_p \geq 3$, and, in this case, do
not intersect another point with multiplicity
$\geq 2$), and

\item $d^*$ the number or lines tangent to
$C$ passing through a point $p \not\in C$
such that $p$ is not in any tangent lines at
the singular points, tangent inflection lines
or bitangent lines (lines that are tangent to
two different points). \end{itemize}

\noindent Then $$
\begin{aligned}
d(d-1)&=d^*+2\delta+3\chi\\
3d(d-2)&=i+6\delta+8\chi. \end{aligned}
$$
\end{theorem}

\section{Rational maps and blow up}
\label{section-rational-maps-and-blow-up}

\noindent
The category of varieties and dominant
rational maps is equivalent to the category
of finitely generated field extensions of $k$
with the arrows reversed.

\begin{exercise}
\label{exercise-cuspidal-curve-blow-up}
Let $Y$ be the cuspidal curve $y^2=x^3$ in
$\mathbb{A}^2$.
Blow up the point $O=(0,0)$, let $E$ be the
exceptional curve,
and let $\tilde{Y}$ be the strict transform
of $Y$.
Show that $E$ meets $\tilde{Y}$ in one point,
and that $\tilde{Y} \cong \mathbb{A}^1$.
\end{exercise}


\section{Normalization}
\label{section-normalization}

\noindent
All bounded meromorphic functions are
holomorphic.

\begin{definition}
\label{definition-integral-closure}
Let $R$ be a domain.
An element of the fraction field
$\text{Frac}(R)$
is {\it integral} if it is the root of a
monic polyomial
with coefficients in $R$.

$R$ is {\it integrally closed} if it
coincides with its integral closure,
that is, the set of
elements of $\text{Frac}(R)$
which are integrally closed over $R$.
\end{definition}

\begin{definition}
\label{definition-normal-variety}
A variety is called {\it normal at $P$}
is $\mathcal{O}_P$ is integrally closed.
It is called {\it normal} if every point
is normal.
\end{definition}

\begin{exercise}
\label{exercise-normal-varieties}
If $Y$ is an affine variety,
$\mathcal{O}_P$ is intergrally closed
for all $P \in Y$ if and only if
$\mathcal{O}_Y$ is integrally closed.
\end{exercise}

\begin{proof}
For the converse
let $f/g \in \text{Frac}\mathcal{O}_P$
be integral.
Then there exists $P \in \mathcal{O}_P[t]$
monic killing $f/g$.
We want to show that $f/g \in \mathcal{O}_P$
for which it's enough to put $f/g$ in
$\mathcal{O}_Y$ using the closure-integrality
of $Y$.
Suppose
$$
P=\sum \frac{f_i}{g_i}t^i.
$$
Since this polynomial is not in
$\mathcal{O}_Y[t]$, we cannot use it directly
to prove integrality of $f/g$;
in fact $f/g$ might not be integral at all.
So we get rid of the fractions
considering the least common multiple
of the $g_i$.
But if we just multiply by the lcm
we get a polynomial that is not monic.
We must consider
$$
P\cdot lcm^{\text{deg}P}
\left(\frac{f}{g}\right)
=\sum lcm^{\text{deg}P}\frac{f_i}{g_i}
\left(\frac{f}{g}\right)^i
=\frac{f_i}{g_i}
\left(lcm \frac{f}{g}\right)^{\text{deg}P}
+\frac{f_i}{g_i}lcm
\left(lcm\frac{f}{g}\right)^{\text{deg}P-1}
+…
+\frac{f_0}{g_0}lcm^{\text{deg}P},
$$
which is monic over $\mathcal{O}_Y$
and kills $lcm \frac{f}{g}$,
which implies that $\frac{f}{g}
\in \mathcal{O}_Y$
and thus $g$ is never zero
(or simply a unit of the coordinate algebra),
and thus $f/g$ is in $\mathcal{O}_P$.

The direct implication is actually easier.
Let $f/g \in \text{Frac}\mathcal{O}_Y$
be integral, i.e. there is a polynomial
$P \in \mathcal{O}_Y[t]$ killing $f/g$.
Pick any point $p \in Y$.
Then $P$ can be taken to be a polynomial
in $\mathcal{O}_p[t]$, which in fact
kills the fraction
$\overline{f/g} \in \text{Frac}\mathcal{O}_p$
in the local ring of $p$.
Since $\mathcal{O}_p$ is integrally closed,
$f/g \in \mathcal{O}_P$.
Thus, thinking of $f/g$ as a meromorphic
function, we see that $g(p) \neq  0$ for all $p$,
meaning the fraction $f/g$ is in fact
regular. Alternatively, thinking of $f/g$
as an element of the localization by the
maximal ideal $\mathfrak{m}_P$,
saying that $f/g \in \mathcal{O}_P$
means that $g \not\in \mathfrak{m}_P$
for all $P$, which is to say that $g$
is in no maximal ideal of $k[Y]$,
which means that $g$ is a unit.
\end{proof}

\begin{exercise}
\label{exercise-normal-varieties-examples}
\begin{enumerate}
\item Show that every conic in $\mathbb{P}^2$
is normal.
\item Show that the quadric surfaces
$Q_1,Q_2$ in $\mathbb{P}^3$ given by
$Q_1:xy-zw$ and $Q_2:xy=z^2$ are normal.
\item Show that the cuspidal cubic
$y^2=x^3$ in $\mathbb{A}^2$
is not normal.
\end{enumerate}
\end{exercise}

\begin{proof}
\begin{enumerate}
\item
\item
\item The element $x^2/y$ has a minimal
polynomial in the local ring at the origin,
namely $P(t)=t-x$, but it clearly
does not lie in the local ring since $y$
vanishes at $(0,0)$.
\end{enumerate}
\end{proof}

\begin{exercise}
\label{exercise-normalization}
Let $Y$ be an affine variety.
Show that there is a normal affine
variety $\tilde{Y}$, and a morphism
$\pi:\tilde{Y}\to Y$,
with the property that whenever $Z$
is a (affine) normal variety and 
$\varphi:Z \to Y$ is a dominant morphism,
then there is a unique morphism
$\theta:Z \to \tilde{Y}$ such that
$\varphi=\pi \circ \theta$.
\end{exercise}

\begin{proof}
Let $\tilde{Y}=\Spec \tilde{A}$
where $\tilde{A}$ is the integral closure
of $A=\mathcal{O}_Y$.
Then by Exercise
\ref{exercise-normal-varieties}
it is normal.
Now if the dominant map $\varphi:Z \to Y$
induces an injective algebra morphism
$\varphi^\#:A \to B$
and $B$ is integrally closed,
we can put $\theta^\#:\tilde{A} \to B$
by the inclusion.
Since any element $x \in \tilde{A}$
is the root of a unique monic polynomial
$P \in A[t] \subset B[t]$,
and since $\theta^\#(x)$ must be
the root of a monic polynomial in $B[t]$
because $B$ is integrally closed,
we conclude that $\theta^\#(x)$
is also a root of $P$.
… but what if the roots are different?

A correct solution is as follows.
The inclusion $A \xymatrix{\ar@{^{(}->}[r]&}
B$ induces an inclusion on fraction fields
$\text{Frac}A \xymatrix{\ar@{^{(}->}[r]&}
\text{Frac} B$.
The key observation is that by definition
an element in the integral closure 
$\tilde{A}$ of $A$ is an element in
the fraction field $\text{Frac} A$.
This says that any integral element
can already be seen as an element in
$\text{Frac}B$.
But since $B$ is integrally closed and the
minimal polynomial of any such element
has coefficients in $A \subset B$,
we are done: an integral element of $A$
is mapped to an element of $B$ making
the following diagram commute
$$
\xymatrix{
A\ar[dr]_{\pi^\sharp}
\ar@{^{(}->}[r]^{\varphi^\sharp}
&  B\\
&  \tilde{A}\ar@{-->}[u]^{\theta^\sharp}
}
$$
For uniqueness, suppose that
$\theta':\tilde{A} \to B$ is another
map making the above diagram commute.
Then for any $x=a/b$ integral over $A$ 
we have
$\theta'(b)\theta'(x)=\theta'(a)$.
But this coincides on $A$
with $\theta^\sharp$
we introduced before, giving
$\theta^\sharp(b)(\theta'(x)-\theta^\sharp)
=0$. Since there are no zerodivisors on $B$,
we obtain $\theta\equiv\theta^\sharp$.
\end{proof}

\begin{exercise}
\label{exercise-projectively-normal-vars}
A projective variety $Y \subseteq
\mathbb{P}^n$ is {\it projectively normal} 
(with respect to the given embedding)
if its homogeneous coordinate ring
$S(Y)$ is integrally closed.
\begin{enumerate}
\item If $Y$ is projectively normal
then $Y$ is normal.
\item There are normal varieties 
(i.e. stalks are integrally closed)
in projective space which are not 
projectively normal.
For example, let $Y$ be the twisted
quartic curve in $\mathbb{P}^3$
given parametrically by
$(x,y,z,w)=(t^4,t^3u,tu^3,u^4)$.
\item Show that the twisted quartic curve
$Y$ above is isomorphic to $\mathbb{P}^1$,
which is projectively normal. Thus projective
normality depends on the embedding.
\end{enumerate}
\end{exercise}

\begin{proof}
\begin{enumerate}
\item Take an affine chart $D_+(f)$.
By Proposition
\ref{proposition-standard-open-proj}
the coordinate algebra of $Y$
at this open set is given by
$(S(Y)_f)_0$, the degree-0 elements
in the localization $S(Y)_f$.
By Commutative Algebra Proposition
\ref{prop-localization-integrally-closed},
localization of an integrally closed
ring remains integrally closed,
thus $S(Y)_f$ is integrally closed.
Let $A=(S(Y)_f)_0$.
If $x \in \text{Frac}(A)$ is integral
over $A$, then it is integral over
$S(Y)_f$.
Thus $x \in S(Y)_f$.
But $x$ is a quotient of degree-0
elements, hence has degree 0.
Therefore $x \in (S(Y)_f)_0=A$.

\item By the next item, the twisted quartic
is isomorphic to $\mathbb{P}^1$, hence is
normal.
It remains to see that it is not
projectively normal in this embedding.
Its homogeneous coordinate ring is
\[
R=k[t^4,t^3u,tu^3,u^4]\subset k[t,u].
\]
The element $t^2u^2$ lies in
$\text{Frac}(R)$, since
\[
t^2u^2=\frac{(t^3u)(u^4)}{tu^3}.
\]
Moreover it is integral over $R$, since
\[
(t^2u^2)^2=(t^3u)(tu^3)\in R.
\]
But $t^2u^2\notin R$: the exponent
$(2,2)$ is not in the semigroup generated
by
\[
(4,0),(3,1),(1,3),(0,4).
\]
Thus $R$ is not integrally closed.

\item Let
\[
\varphi:\mathbb{P}^1\to Y
\]
be the map
\[
[t:u]\longmapsto [t^4:t^3u:tu^3:u^4].
\]
We construct its inverse.
On the open subset of $Y$ where
$(x,y)\neq (0,0)$, define
\[
[x:y:z:w]\longmapsto [x:y].
\]
On the open subset where
$(z,w)\neq (0,0)$, define
\[
[x:y:z:w]\longmapsto [z:w].
\]
These opens cover $Y$.
On the image of $\varphi$ we have
\[
[x:y]=[t^4:t^3u]=[t:u]
\]
and
\[
[z:w]=[tu^3:u^4]=[t:u].
\]
Thus the two local formulae agree on
overlaps and define a morphism
$Y\to \mathbb{P}^1$ inverse to
$\varphi$.
Hence $Y\cong \mathbb{P}^1$.
Finally, the standard embedding of
$\mathbb{P}^1$ has homogeneous
coordinate ring $k[t,u]$,
which is integrally closed.
Thus $\mathbb{P}^1$ is projectively
normal.
\end{enumerate}
\end{proof}



\section{Divisors, invertible sheaves}
\label{section-divisors-line-bundles}

\begin{definition}
\label{definition-weil-divisor}
\begin{reference}
\href{https://stacks.math.columbia.edu/tag/%
0BE2}{0BE2},
\cite[p. 130]{hart}
\end{reference}
A {\it Weil divisor} on the variety $X$ is a
formal sum of irreducible
subvarieties over the integers.

More formally, a {\it prime divisor} is an
integral closed
subscheme $Z \subset X$ of codimension 1.
A {\it Weil divisor} is a formal sum $D=\sum
n_Z Z$ where
the sum is over prime divisors of $X$
and the collection
$\{Z|n_Z \neq  0\}$ is locally finite.
If all coefficients are non-negative we say
the divisor is {\it effective}.
\end{definition}

\noindent
A reasonable notion of order of vanishing
should then be introduced. In the case of
Riemann surfaces
this can be done using the order of
vanishings
of the denominator and numerator functions in
$f=g/h$ where $f \in K(X)^*$.
A more general construction is a bit more
involved.

\begin{definition}
\label{definition-principal-weil-divisor}
\begin{reference}
\href{https://stacks.math.columbia.edu/tag/%
0be0}{0be0}
\end{reference}
Let $X$ be a Noetherian integral scheme.
Let $f \in R(X)^*$.
The {\it principal Weil divisor associated to
$f$}
is the Weil divisor
$$
\text{div}(f)=\text{div}_X(f)
=\sum \text{ord}_Z(f)[Z]
$$
where the sum is over prime divisors and
$\text{ord}_Z(f)$ is as in Definition ?.
This makes sense by Lemma ? [a lemma about
finiteness of the sum].
\end{definition}

\begin{definition}
\label{definition-weil-divisor-class-group}
\begin{reference}
\href{https://stacks.math.columbia.edu/tag/%
0BE4}{0BE4}
\end{reference}
Let $X$ be a locally Noetherian scheme.
The {\it Weil divisor class group} of $X$,
denoted
$\text{Cl}(X)$, is the quotient
of the group of Weil divisors by the subgroup
of principal
Weil divisors.
\end{definition}

\noindent
The next section,
\href{https://stacks.math.columbia.edu/tag/%
02SE}{02SE},
shows there is an injective homomorphism
$\text{Pic}(X) \to \text{Cl}(X)$
which is also surjective if $X$ is normal
and the local rings of $X$ are UFDs
(which happens for smooth varieties).

\medskip\noindent
Earlier, in
\href{https://stacks.math.columbia.edu/tag/%
01WQ}{01WQ},
we have

\begin{definition}
\label{definition-cartier-divisor}
Let $S$ be a scheme.
\begin{itemize}
\item A {\it localy principal closed
subscheme} of $S$
is a closed subscheme whose sheaf of ideals
is locally generated by a single element.

\item An {\it effective Cartier divisor} on
$S$
is a closed subscheme $D \subset S$ whose
ideal sheaf
$\mathcal{I}_D \subset \mathcal{O}_S$ is an
invertible
$\mathcal{O}_S$-module.
\end{itemize}
\end{definition}

\noindent
Followed by Lemma
\href{https://stacks.math.columbia.edu/tag/%
01WS}{01WS},
which essentially says that a Cartier
divisor is cut locally by a single function.
After this we would show the well-known
construction of sheaf associated to an
effective Cartier divisor
and eventually construct all group operations
and subtleties
to arrive at Lemma
\href{https://stacks.math.columbia.edu/tag/%
01X0}{01X0},
which says that effective Cartier divisors on
a scheme
are the same as invertible sheaves with fixed
regular global section.

\medskip\noindent
Here's two basic facts I'm not sure how to
prove:
\begin{enumerate}
\item
$$
\mathcal{O}_{\mathbb{P}^n}(1)\cong
\mathcal{O}_{\mathbb{P}^n}(H)
$$
where $\mathcal{O}_{\mathbb{P}^n}(1)$ is
defined
as dual of the tautological bundle,
and $H$ is any hyperplane (zeroes of a
homogeneous degree-1 polynomial
in $n+1$ variables) seen as an effective
Cartier divisor,
so that $\mathcal{O}_{\mathbb{P}^n}(H)$ is
its associated line bundle.

\item
$$
H^0(\mathbb{P}^n,\mathcal{O}_{\mathbb{P}^n}
(1))
\cong
k[x_0,…,x_n]_1,
$$
that is, the space of sections of
$\mathcal{O}_{\mathbb{P}^n}(1)$
is isomorphic to the space of homogeneous
polynomials
of degree 1 in $n+1$ variables
(each of which is a hyperplane in
$\mathbb{P}^n$).
I think this might be
\href{https://www.youtube.com/live/%
bXhgeRZMu8M?si=5BPdnECUREsCUwpV&t=4601}{
here}.
\end{enumerate}

\begin{exercise}
\label{exercise-degree-d-surface}
Show that a hypersurface in $\mathbb{P}^3$
given by the zeroes of a homogeneous
polynomial
of degree $d$ has degree $d$,
i.e., if $S \subset \mathbb{P}^3$, $S=V(f)$
with
$f$ homogeneous of degree $d$,
then
 $$
\mathcal{O}_S(1).\mathcal{O}_S(1)=d.
$$
\end{exercise}

\begin{proof}
The exercise is kind of by definition no?
The intersection form counts the size of the
intersection when
the two classes are just two hypersurfaces
intersecting transversally (as is the case of
two
hyperplane sections in general position,
both of which have $\mathcal{O}_S(1)$ as
associated line bundle). The intersection
has $d$ points because its a line in
$\mathbb{P}^3$
intersecting a degree-$d$ polynomial.
And that's it.
\end{proof}

\begin{exercise}
\label{exercise-intersection-lines}
Let $S \subset \mathbb{P}^3$ be a surface
of degree $d$ and $L \subset S$ a line.
Compute $L^2$ in $S$. Check that when $d=3$,
$L^2=-1$!
\end{exercise}

\begin{proof}
We apply the genus formula:
$$
(L.L+K_S)=2p_a-2
$$
By $L$ being a line we know that $p_a=0$
and that the intersection number of $L$
with any divisor of degree $d'$ will be $d'$.
Indeed, $L.d'H=d'(L.H)=d'$
because $L.H$ is the degree of the line
(by definition of degree, namely
the intersection number of $L$ with as many
hyperplanes as its dimension).
Thus
$$
\begin{aligned}
-2=(L.L+K_S)&=L^2+(L.K_S)=L^2+(L.\mathcal{O}
_S(d-4))
=L^2+d-4\\
\iff L^2=-d+2.
\end{aligned}
$$

\noindent
where the canonical bundle of $S$
is computed using adjunction formula
$K_S=(K_{\mathbb{P}^3}+N)|_S
=(\mathcal{O}_{\mathbb{P}^3}(-4)+\mathcal{O}
_{\mathbb{P}^3}(d))$.
\end{proof}

\begin{exercise}[Bézout]
\label{exercise-bezout}
Show that if $C_d$ and $C_{d'}$ are two plane
curves
(i.e. contained in $\mathbb{P}^2$)
of degrees $d$ and $d'$ then
$C_d.C_{d'}=dd'$.
\end{exercise}

\begin{proof}
Just notice that the line bundles associated
to
$C_d$ and $C_{d'}$ are $dH$ and $d'H$,
then use linearity of the bilinear of form.
To check why the divisor associated to $C_d$
is indeed $dH$ note that the intersection
number
of $C_d$ and a line $H$ is $d$ because $C_d$
is given by the roots of a polynomial of
degree $d$,
thus $C_d.H=d$, and since the Picard group of
$\mathbb{P}^2=\mathbb{Z}H$
we must have $C_d \sim dH$.
\end{proof}
\section{Adjunction formulas}
\label{section-adjunction formulas}

There are several statements called
adjunction formula
in different texts. All of them concern
``subvarieties'',
that is, closed embedded subschemes.

\begin{exercise}[Genus formula for a curve on
a surface]
\label{exercise-genus-formu-curve-surfa}
Let $C \to X$ be a closed embedded subscheme
of dimension $1$ (as a topological space,
i.e. pure dimension)
inside a smooth surface $X$.
Then
$2p_a-2=(\mathcal{O}_X(C),\mathcal{O}_X(C))$.
\end{exercise}

\begin{proof}
Consider the ideal sheaf exact sequence
$$
\xymatrix{
0\ar[r]&\mathcal{O}_X(-C)\ar[r]&\mathcal{O}
_X\ar[r]&\mathcal{O}_C\ar[r]&0
}
$$
This sequence splits since there is an
obvious inverse morphism
to the inclusion $\mathcal{O}_X(-C)\to
\mathcal{O}_X$, namely
mapping a function $f$ to
Then $\mathcal{O}_X\cong \mathcal{O}_X(-C)
\oplus \mathcal{O}_C$.
 $\chi(\mathcal{O}_X)=\chi$
\end{proof}


\section{Linear systems}
\label{section-linear-systems}

\noindent
A basic fact is that
$$
H^0(\mathcal{O}(d)) \cong
\{\text{homogeneous polynomials of degree
$d$}\}.
$$
And more generally,
$$
|D|=\mathbb{P}H^0(\mathcal{O}(D)).
$$

\medskip\noindent
Following \cite[Section 5.2]{Fulton}, let
$P_1,…,P_n$ be points in $\mathbb{P}^2$.
Define
$$
V(d,r_1P_1,…,r_nP_n):=
\{\text{curves $F$ of degree $d$}:
m_{P_i}F\geq r_i\forall i\}.
$$
This space has an interesting property:
when we blow up the $n$ points, we obtain
$$
H^0(X,\mathcal{O}(dH-\sum r_iE_i))
=
\{f \in
H^0(\mathbb{P}^2,\mathcal{O}(d)):m_{P_i}(f)
\geq  r_i\}.
$$

Recall (proof?) that the canonical divisor of
the blow up is
$K_{Bl}=\pi^* K_{\mathbb{P}^2}+E_1+…+E_n=-3H
+ \sum E_i$.
And the projectivization of the sections of
this canonical divisor
are exactly $V(-3,P_1,…,P_n)$, which makes
no sense because there are no curves of
degree $-3$,
that is the canonical divisor of the blow up
has no sections.

\begin{exercise}
\label{exercise-hodge-numbers-blow-up-9pts}
Let $C_1,C_2$ be two cubic curves in
$\mathbb{P}^2$.
Compute the Hodge numbers of the blowup of
the 9 points
of intersection.
\end{exercise}

\begin{proof}
The solution is:
$$
\begin{aligned}\xymatrix@C=.5em@R=.5em{
&  &   1\\
&  0&  &  0\\
0&&  10&  &  0\\
&  0&  &  0\\
&  &  1
}\end{aligned}
$$
Which follows from:
\begin{itemize}
\item $h^{0,0}=1$ because $X$ is connected.
\item $h^{1,0}=0$ because $X$ is simply
connected.
\item $h^{2,0}=0$ because the canonical
bundle of $X$ has no sections.
\item $h^{1,1}=10$ because in the exponential
exact sequence
we get $\text{Pic} \cong H^2(X,\mathbb{Z})$
(by vanishing irregularity and genus)
and we know
$\text{Pic}=\mathbb{Z}H \oplus
\mathbb{Z}E_1\oplus…\oplus \mathbb{Z}E_9$.
\end{itemize}
\end{proof}

\section{Ample line bundles}
\label{section-ample-line-bundles}

First is this lemma that comes from
modules.tex. I think these sets $X_s$ are
the base points of the bundle. Because look:
image of $s$ just means consider
the section $s$ of the line bundle as a germ
near $x$. Now a line bundle is a
locally free rank-1 $\mathcal{O}_X$-module,
so its sections, like $s$, may be
multiplied by germs of functions in the
maximal ring $\mathfrak{m}_x$, i.e. the
functions that vanish at $x$. So $X_s$ is the
vanishing locus of the section
$s$. If $s(x)\neq 0$, obviously $s
\not\in\mathfrak{m}_x\mathcal{L}_x$, so $x\in
X_s$. Conversely, I would like to
show that if $s(x)=0$ then
$s\in\mathfrak{m}_x\mathcal{L}_x$ but I'm not
sure
how. It's like: a vector field with a zero
can be multiplied by a function that
vanishes at the point, sure, but what's this
function?

\begin{lemma}
\label{lemma-s-open}
From modules.tex.
\begin{slogan}
A (local) trivialisation of a linebundle
is the same as a (local) nonvanishing
section.
\end{slogan}
Let $X$ be a ringed space. Assume that each
stalk $\mathcal{O}_{X, x}$
is a local ring with maximal ideal $\mathfrak
m_x$.
Let $\mathcal{L}$ be an invertible
$\mathcal{O}_X$-module.
For any section $s \in \Gamma(X,
\mathcal{L})$ the set
$$
X_s = \{x \in X \mid \text{image }s \not\in
\mathfrak m_x\mathcal{L}_x\}
$$
is open in $X$. The map $s :
\mathcal{O}_{X_s} \to \mathcal{L}|_{X_s}$
is an isomorphism, and there exists a section
$s'$
of $\mathcal{L}^{\otimes -1}$ over $X_s$ such
that $s' (s|_{X_s}) = 1$.
\end{lemma}

\begin{proof}
Suppose $x \in X_s$.
We have an isomorphism
$$
\mathcal{L}_x \otimes_{\mathcal{O}_{X, x}}
(\mathcal{L}^{\otimes -1})_x
\longrightarrow
\mathcal{O}_{X, x}
$$
by Lemma
\ref{lemma-constructions-invertible}.
Both $\mathcal{L}_x$ and
$(\mathcal{L}^{\otimes -1})_x$
are free $\mathcal{O}_{X, x}$-modules of rank
$1$. We conclude
from Algebra, Nakayama's Lemma
\ref{algebra-lemma-NAK} that
$s_x$ is a basis for $\mathcal{L}_x$. Hence
there exists
a basis element $t_x \in
(\mathcal{L}^{\otimes -1})_x$
such that $s_x \otimes t_x$ maps to $1$.
Choose an open neighbourhood $U$ of
$x$ such that $t_x$ comes from a section $t$
of $\mathcal{L}^{\otimes -1}$ over $U$ and
such that
$s \otimes t$ maps to $1 \in
\mathcal{O}_X(U)$.
Clearly, for every $x' \in U$ we see that $s$
generates
the module $\mathcal{L}_{x'}$. Hence $U
\subset X_s$.
This proves that $X_s$ is open. Moreover, the
section
$t$ constructed over $U$ above is unique, and
hence
these glue to give the section $s'$ of the
lemma.
\end{proof}

Recall from Modules, Lemma
\ref{modules-lemma-s-open}
that given an invertible sheaf $\mathcal{L}$
on a locally ringed
space $X$, and given a global section $s$ of
$\mathcal{L}$
the set $X_s = \{x \in X \mid s \not \in
\mathfrak m_x\mathcal{L}_x\}$
is open. A general remark is that
$X_s \cap X_{s'} = X_{ss'}$, where $ss'$
denote
the section $s \otimes s' \in \Gamma(X,
\mathcal{L} \otimes \mathcal{L}')$.

\begin{definition}
\label{definition-ample}
\begin{reference}
\cite[II Definition 4.5.3]{EGA}
\end{reference}
Let $X$ be a scheme.
Let $\mathcal{L}$ be an invertible
$\mathcal{O}_X$-module.
We say $\mathcal{L}$ is {\it ample} if
\begin{enumerate}
\item $X$ is quasi-compact, and
\item for every $x \in X$ there exists an $n
\geq 1$
and $s \in \Gamma(X, \mathcal{L}^{\otimes
n})$ such
that $x \in X_s$ and $X_s$ is affine.
\end{enumerate}
\end{definition}

\begin{exercise}
\label{exercise-ample-bundle-on-K3}
Let $L$ be an ample bundle on a K3 surface
$M$. Prove that
$\mathcal{L}^{\otimes 2}$ is globally
generated (that is, for each $x\in M$
there exsits a section $h \in
H^{0}(L^{\otimes 2})$ which does not vanish
in
$x$).
\end{exercise}

\begin{proof}
This just asks that the $n$ in Definition
\ref{definition-ample} is $2$ for all
$x\in X$. Because, again, $x\in X_s$ means
that $s(x)\neq 0$ because if it was,
then we could somehow write $s$ as a product
of a vanishing function on
$\mathfrak{m}_x$ and a local frame of
$\Gamma(X,\mathcal{L})$. But I guess for
the exercise do this: a line bundle is {\it
ample} if there is $n$ such that the
canonical embedding (cf Lemma
\ref{lemma-map-into-proj}) is an embedding,
i.e.
that $\mathcal{L}^{\otimes n}$ is {\it very
ample}. (Interestingly, the notion
very ampleness is defined in morphisms.tex.)
\end{proof}

Now we pass to the part where ampleness gives
you an {\bf open immersion} to
some projective space. Because, it's only
very ampleness that gives an
embedding, right? (Actually I think here in
stacks project there are no
embeddings but closed immersions.)

\begin{definition}
\label{definition-gamma-star}
From modules.tex. Let $(X, \mathcal{O}_X)$ be
a ringed space.
Given an invertible sheaf $\mathcal{L}$ on
$X$ we define
the {\it associated graded ring} to be
$$
\Gamma_*(X, \mathcal{L})
=
\bigoplus\nolimits_{n \geq 0} \Gamma(X,
\mathcal{L}^{\otimes n})
$$
Given a sheaf of $\mathcal{O}_X$-modules
$\mathcal{F}$ we set
$$
\Gamma_*(X, \mathcal{L}, \mathcal{F})
=
\bigoplus\nolimits_{n \in \mathbf{Z}}
\Gamma(X,
\mathcal{F} \otimes_{\mathcal{O}_X}
\mathcal{L}^{\otimes n})
$$
which we think of as a graded $\Gamma_*(X,
\mathcal{L})$-module.
\end{definition}

\begin{lemma}
\label{lemma-map-into-proj}
Let $X$ be a scheme.
Let $\mathcal{L}$ be an invertible
$\mathcal{O}_X$-module.
Set $S = \Gamma_*(X, \mathcal{L})$ as a
graded ring.
If every point of $X$ is contained in one of
the
open subschemes $X_s$, for some $s \in S_{+}$
homogeneous, then
there is a canonical morphism of schemes
$$
f : X \longrightarrow Y = \text{Proj}(S),
$$
to the homogeneous spectrum of $S$ (see
Constructions, Section
\ref{constructions-section-proj}).
This morphism has the following properties
\begin{enumerate}
\item $f^{-1}(D_{+}(s)) = X_s$ for any $s \in
S_{+}$ homogeneous,
\item there are $\mathcal{O}_X$-module maps
$f^*\mathcal{O}_Y(n) \to \mathcal{L}^{\otimes
n}$
compatible with multiplication maps, see
Constructions, Equation
(\ref{constructions-equation-multiply}),
\item the composition
$S_n \to \Gamma(Y, \mathcal{O}_Y(n)) \to
\Gamma(X, \mathcal{L}^{\otimes n})$
is the identity map, and
\item for every $x \in X$ there is an integer
$d \geq 1$
and an open neighbourhood $U \subset X$ of
$x$
such that $f^*\mathcal{O}_Y(dn)|_U \to
\mathcal{L}^{\otimes dn}|_U$
is an isomorphism for all $n \in \mathbf{Z}$.
\end{enumerate}
\end{lemma}

\section{Riemann-Roch theorem for curves}
\label{section-riemann-roch}

\noindent
Let $D=\sum n_p p$ be a divisor on a curve.
How many rational functions $z \in K^*$
satisfy
$\text{div}(z) \geq -D$? That is, rational
functions
which can only have poles of order $\leq n_p$
at the points $p$ with $n_p>0$,
and can only have zeros of order $\geq
|n_p|$
at the points $p$ with $n_p<0$.

For example, fix a point $p$ in the curve.
Then $\mathcal{O}_C(p)$ is the space
of rational functions with a simple pole at
$p$.
$\mathcal{O}_C(2p)$ is the space of rational
functions
with a pole of order at most 2 at $p$, and so
on.
Note that $\mathcal{O}_C(p) \subset
\mathcal{O}_X (2p)
\subset \mathcal{O}(3p)…$.
That is, the more poles we allow, this space
of functions will be larger.

In classical (Riemann's) notation, this space
was denoted
$L(D)$, and its dimension $\ell(D)$.
Modern notation reads $H^0(\mathcal{O}_C(D))$
instead of $L(D)$ and
$h^0(\mathcal{O}_C(D))$
instead of $\ell(D)$.

Riemann-Roch formula for curves is
\begin{equation}
\label{equation-riemann-roch-curves}
h^0(D)-h^0(K-D)=\text{deg}(D)-g+1,
\end{equation}

\noindent
where $K$ is a divisor associated to the
curve canonically.
This is done by choosing a meromorphic
section $s$
of the bundle $\Omega_C^1$
and taking $K:=\text{div}(s)$.
Then one shows that this is independent of
the choice of $s$.

Note that by Serre duality
$h^1(D)=h^0(K-D)$.


\section{Schemes}
\label{section-shemes}

For a definition of presheaf,
see Categories, Definition
\ref{categories-definition-presheaf}.

\begin{definition}
\label{definition-sheaf}
Let $X$ be a topological space.
\begin{enumerate}
\item A {\it sheaf $\mathcal{F}$ of sets on
$X$} is a presheaf
of sets which satisfies the following
additional property: Given
any open covering $U = \bigcup_{i \in I} U_i$
and any collection
of sections $s_i \in \mathcal{F}(U_i)$, $i
\in I$ such that
$\forall i, j\in I$
$$
s_i|_{U_i \cap U_j} = s_j|_{U_i \cap U_j}
$$
there exists a unique section $s \in
\mathcal{F}(U)$ such that
$s_i = s|_{U_i}$ for all $i \in I$.
\item A {\it morphism of sheaves of sets} is
simply a
morphism of presheaves of sets.
\item The category of sheaves of sets on $X$
is denoted
$\Sh(X)$.
\end{enumerate}
\end{definition}


\medskip\noindent
Let $X$ be a topological space. Let $x \in X$
be a point.
Let $\mathcal{F}$ be a presheaf of sets on
$X$.
The {\it stalk of $\mathcal{F}$ at $x$} is
the set
$$
\mathcal{F}_x
=
\colim_{x\in U} \mathcal{F}(U)
$$
where the colimit is over the set of open
neighbourhoods
$U$ of $x$ in $X$. The set of open
neighbourhoods is
partially ordered by (reverse) inclusion:
We say $U \geq U' \Leftrightarrow U \subset
U'$.
The transition maps in the system are
given by the restriction maps of
$\mathcal{F}$.
See Categories, Section
\ref{categories-section-posets-limits}
for notation and terminology regarding
(co)limits over systems.
Note that the colimit is a directed colimit.
Thus it is easy to describe $\mathcal{F}_x$.
Namely,
$$
\mathcal{F}_x
=
\{
(U, s)
\mid
x\in U, s\in \mathcal{F}(U)
\}/\sim
$$
with equivalence relation given by $(U, s)
\sim (U', s')$ if and only if
there exists an open $U'' \subset U \cap U'$
with $x \in U''$ and
$s|_{U''} = s'|_{U''}$. Given a pair $(U, s)$
we sometimes denote
$s_x$ the element of $\mathcal{F}_x$
corresponding to the equivalence
class of $(U, x)$. We sometimes use the
phrase
``image of $s$ in $\mathcal{F}_x$'' to denote
$s_x$.
For example, given two pairs $(U, s)$ and
$(U', s')$ we sometimes
say ``$s$ is equal to $s'$ in
$\mathcal{F}_x$'' to indicate
that $s_x = s'_x$. Other authors use the
terminology
``germ of $s$ at $x$''.

\medskip\noindent
Now let's discuss closed immersions.

\noindent
Closed immersions of ringes spaces are in
\href{https://stacks.math.columbia.edu/tag/%
01C1}{01C1}.

\begin{definition}
\label{definition-closed-immersion}
A map of ringed spaces
$i:(Z,\mathcal{O}_Z) \to (X,\mathcal{O}_X)$
is a {\it closed immersion}
if $i$ is injective,
its image is closed (I think these two are
``closed immersion of topological spaces'')
and the induced map $\mathcal{O}_X \to
i_*\mathcal{O}_Z$
is surjective; denote its kernel by
$\mathcal{I}$.
To make sure that if $(X,\mathcal{O}_X)$ is a
scheme
then so is $(Z,\mathcal{O}_Z)$ with add the
extra condition
that the $\mathcal{O}_X$-module $\mathcal{I}$
is
locally generated by sections.
\end{definition}

\medskip\noindent
Now let's discuss restricting sheaves to
subschemes.

Upshot about restricting sheaves to
subschemes:
the sheaves on $Z$ correspond essentially
(in the categorical sense)
to sheaves on $X$ with support on $Z$.

More explicitly, from
\href{https://stacks.math.columbia.edu/tag/%
01QY}{01QY}:
the pushforward of the inclusion (a closed
immersion of schemes)
$i:Z \to X$, with kernel on induced map
$\mathcal{O}_Z \to \mathcal{O}_X$
denoted by $I$, is an exact, fully faithful
functor
$$
i_*:\QCoh(\mathcal{O}_Z) \to
\QCoh(\mathcal{O}_X)
$$
with essential image the quasi-coherent
$\mathcal{O}_X$ modules
$\mathcal{G}$ such that
$\mathcal{I}\mathcal{G}=0$.

It makes sense to think of this operation as
$\mathcal{F}|_Z=i^*\mathcal{F}=\mathcal{F}
\otimes\mathcal{O}_Z$
for $\mathcal{F}\in\QCoh(X)$.
Indeed: tensoring by the ideal sheaf would
kill the information on $Z$,
and tensoring with the structure sheaf does
the opposite
--- it's the sheaves on $X$ that are no fun
outside $Z$,
so basically sheaves on $Z$.

\medskip\noindent
Here's a more detailed discussion:

In understanding the cohomology of sheaves
when
we restrict them to submanifolds (see Complex
Geometry
Exercise
\ref{complex-geometry-exercise-%
deformations-curve-in-k3},
I found Stacks Project
\href{https://stacks.math.columbia.edu/tag/%
01AW}{01AW}
section on closed immersions and abelian
sheaves.
First we have a few propositions stating that
sheaves on the subscheme $Z \subset X$ can be
pushed
to sheaves on $X$.

And then there's the converse:
what I would like to call the restriction
functor $\mathcal{H}$
but rather is called the {\it sub-module of
sections with support in $Z$}.
In Remark
\href{https://stacks.math.columbia.edu/tag/%
01AY}{01AY}
(and then in Remark
\href{https://stacks.math.columbia.edu/tag/%
0G6N}{0G6N}
for ringed spaces; these seem to be analogue
constructions…)
we see that for an abelian sheaf
$\mathcal{F}$ on $X$
we can define an abelian sheaf
$\mathcal{H}_Z(\mathcal{F})$
which although is a sheaf on $X$
we can just think it's a sheaf on $Z$ -
that's beacause
$$
\mathcal{H}_Z(\mathcal{F})(U)
=\{s \in \mathcal{F}(U))| \text{the support
of $s$ is contained in }Z\cap U\}.
$$
We can just think that it's the functions
that vanish outside the subvariety.
Which doesn't really make sense:
a function vanishing outside the submanifold
will most likely
vanish also in the submanifold! Just by
continuity,
the submanifold is in the closure of the open
set!
I guess what I mean to say is that I don't
mind if
the functions do not vanish outside $Z$…
I just want to consider them as functions on
$Z$…

\medskip\noindent
Here's another attempt in understanding how
to restrict sheaves to submanifolds.
First and foremost, lacking a reference this,
I understand the symbol $\mathcal{F}|_Z$ for
a sheaf
on $X$ and a subscheme  $Z \subset X$ to mean
the
pullback of $\mathcal{F}$ by the inclusion.
Then look at
\href{https://stacks.math.columbia.edu/tag/%
01QY}{01QY}
(more basically, as above,
\href{https://stacks.math.columbia.edu/tag/%
01AW}{01AW}…
in fact there appear to be several version of
this statement)
to recall that sheaves of the subscheme $Z$
should correspond
to sheaves on $X$ with support on $Z$.
{\bf So essentially $\mathcal{F}|_Z$ should
be a version of $\mathcal{F}$
but with support on $Z$.}
The key point here is that this is non other
than
$\mathcal{F} \otimes \mathcal{O}_Z$.
Why? Because when multiply with the ideal
sheaf of $Z$ we obtain
$(\mathcal{F} \otimes
\mathcal{O}_Z)\mathcal{I}=0$
since sections of $\mathcal{I}$
vanish in
$\mathcal{O}_Z=\mathcal{O}_X/\mathcal{I}$.
That is, the support of $\mathcal{F} \otimes
\mathcal{O}_Z$ indeed
contained in $Z$: if the sections become zero
when
multiplied by sections of $\mathcal{I}$
(sections of $\mathcal{I}$ are only nonzero
outside $Z$)
it means that all the action was inside $Z$.
that outside $Z$

I think the upshot here is that structure
sheaves
act as some sort of ``contour'' for the
geometric object:
they vanish {\it outside} the geometric
object,
they tell us what is the shape of the object
by describing its surroundings, rather than
its interior.
This is exactly how algebraic/analytic sets
are defined,
and this why, when we want to restrict a
sheaf
to a subscheme, we tensor with the structure
sheaf:
tensor by the ideal sheaf would kill the
information on $Z$,
and tensoring with the structure sheaf does
the opposite.














\medskip\noindent
Now let's dicuss Proj of a graded ring.

\noindent
(Stacks Project
\href{https://stacks.math.columbia.edu/tag/%
01M3}{01M3},
Proj of a graded ring.)
This section may be taken as the preamble to
the question: what is projective space
anyway?
Yes, projective space is Proj of a graded
ring $S$.
And yes, projective space is a scheme.
So in particular it's locally affine.

(Stacks Project
\href{https://stacks.math.columbia.edu/tag/%
01MX}{01MX},
Functoriality of Proj.)
This is reminiscent of the olden days.
The point is that Proj is covariant,
i.e. for rings  $S' \subset S$ we have
$\text{Proj}S \subset \text{Proj}S$
(under the right conditions,
which is the point of this section),
meaning that, indeed,
{\bf Projs are projective schemes}.






\medskip\noindent
Noe let's discuss invertible sheaves on Proj
(twists!).

\noindent
As for the construction of the twisted
sheaves:
for a projective space $\text{Proj} S$,
we just compute the twisted $S$ module
$S(m)$ by the formula $S(m)_d=S_{m+d}$
and turn that into a sheaf
$\widetilde{S(m)}$.

Here's the construction of the twisted
sheaves $\mathcal{O}_X(n)$.
The point is that there is a good way
to pass from an $S$-module $M$ to an
$\mathcal{O}_X$-module, called
$\widetilde{M}$.
So basically you just define the twists
by $M(d)_n=M_{n+d}$ and apply that
construction.

Pick an element of degree $d$.
Then think of the module generated by this
element.
The least degree piece of such a module
is the degree $d$ piece of the original
module,
that is we have $M(-d)$.











\medskip\noindent
Lost note on logarithmic derivations:

\noindent
Just to record a notion of logarithmic
derivation
from Vladimiro Benedetti, Danielle Faenzi and
Simone.

Let $F$ be a homogneous polynomial of degree
$d$.
The {\it logarithmic derivations} of $F$ are
$$
\text{Der}_U(-\text{log}F)
=\{\partial \in \text{Der}_U
:\partial F \in \left\langle{}F\right\rangle{}_U\}
$$
where $U=\mathbb{C}[x_1,…,x_e]$.










\medskip\noindent
Lost notes on gonality.

\noindent
From a talk by Juliana Coelho at Bandoleiros
2025:

A curve $C$ is {\it $k$-gonal} if there is a
map $C \to \mathbb{P}^1$ of degree $k$.

A curve is {\it stable} if it is nodal and
$\text{Aut}(C)$ is finite.
Equivalently, the rational components of $C$
meet the
rest of the curve in at least three points.

We may study the moduli of $k$-stable curves,
or the space of $k$-gonal structures.

\section{Reduced schemes}
\label{section-reduced}

\noindent
So far a reduced scheme,
for me,
is just a scheme where there are
no ``repetitions'':
$(x^2)$ is not reduced because
it gives the same information as $(x)$.
Algebraically this means that the
coordinate ring is reduced,
which is defined by asking
that there are no nilpotent elements,
i.e. elements such that some power vanishes.
(There's no definition of reduced ring in
Stacks Project
since it's taken as part of the basic
algebraic knowledge).

The definition is on stalks but the first
lemma
shows it's equivalent to define it on any
open set.
This is in virtue of some ``injectivity''
lemma
from the ring at $U$ to the product of all
stalks parametrized in $U$;
essentially that if a section vanishes when
projected to all
the stalks then it's zero.

\begin{definition}
\label{definition-reduced}
Let $X$ be a scheme. We say $X$ is {\it
reduced} if every local ring
$\mathcal{O}_{X, x}$ is reduced.
\end{definition}


\begin{proof}
Assume that $X$ is reduced.
Let $f \in \mathcal{O}_X(U)$ be a section
such that $f^n = 0$.
Then the image of $f$ in $\mathcal{O}_{U, u}$
is zero for
all $u \in U$. Hence $f$ is zero, see
Sheaves, Lemma
\ref{sheaves-lemma-sheaf-subset-stalks}.
Conversely, assume that $\mathcal{O}_X(U)$ is
reduced
for all opens $U$. Pick any nonzero element
$f \in \mathcal{O}_{X, x}$.
Any representative $(U, f \in
\mathcal{O}(U))$  of $f$ is nonzero and
hence not nilpotent. Hence $f$ is not
nilpotent in $\mathcal{O}_{X, x}$.
\end{proof}

\section{Flat morphisms}
\label{section-flat-morphisms}

\noindent
The essential technical property
for for defining flatness
is the preservation of exact sequences.
Right-exactness is true in general;
it follows from currying in category
of commutative rings.
But the functor $-\otimes_R N$ does not
preserve injectivity of maps---
that's the point of flatness.

\begin{lemma}[Internal Hom for R-modules]
\label{lemma-hom-from-tensor-product}
For any three $R$-modules $M, N, P$,
$$
\Hom_R(M \otimes_R N, P) \cong \Hom_R(M,
\Hom_R(N, P))
$$
\end{lemma}

\begin{proof}
An $R$-linear map $\hat{f}\in \Hom_R(M
\otimes_R N, P)$ corresponds to an
$R$-bilinear map $f : M \times N \to P$. For
each $x\in M$ the mapping $y\mapsto f(x, y)$
is $R$-linear by the universal
property. Thus $f$ corresponds to a
map $\phi_f : M \to \Hom_R(N, P)$. This map
is $R$-linear since
$$
\phi_f(ax + y)(z) =
f(ax + y, z) = af(x, z)+f(y, z) =
(a\phi_f(x)+\phi_f(y))(z),
$$
for all $a \in R$, $x \in M$, $y \in M$ and
$z \in N$. Conversely, any
$f \in \Hom_R(M, \Hom_R(N, P))$ defines an
$R$-bilinear
map $M \times N \to P$, namely $(x, y)\mapsto
f(x)(y)$.
So this is a natural one-to-one
correspondence between the
two modules
$\Hom_R(M \otimes_R N, P)$ and $\Hom_R(M,
\Hom_R(N, P))$.
\end{proof}

\begin{lemma}[Tensor product is right exact]
\label{lemma-tensor-product-exact}
Let
$$
\begin{aligned}
M_1\xrightarrow{f} M_2\xrightarrow{g} M_3 \to
0
\end{aligned}
$$
be an exact sequence of $R$-modules and
homomorphisms, and let $N$ be any
$R$-module. Then the sequence
\begin{equation}
\label{equation-2ndex}
M_1\otimes N\xrightarrow{f \otimes 1}
M_2\otimes N \xrightarrow{g \otimes 1}
M_3\otimes N \to 0
\end{equation}
is exact. In other words, the functor $-
\otimes_R N$ is
{\it right exact}, in the sense that
tensoring
each term in the original right exact
sequence preserves the exactness.
\end{lemma}

\begin{proof}
For every $R$-module $P$
we apply the functor $\Hom(-, \Hom(N, P))$ to
the first exact
sequence. We obtain
$$
0 \to
\Hom(M_3, \Hom(N, P)) \to
\Hom(M_2, \Hom(N, P)) \to
\Hom(M_1, \Hom(N, P))
$$
which is exact by Lemma \ref{lemma-hom-exact}
(1).
By Lemma \ref{lemma-hom-from-tensor-product}
this becomes the sequence
$$
0 \to \Hom(M_3 \otimes N, P) \to
\Hom(M_2 \otimes N, P) \to \Hom(M_1 \otimes
N, P)
$$
which is therefore also exact. Then using
Lemma \ref{lemma-hom-exact} (1) again, we
arrive at the desired exact sequence.
\end{proof}

\begin{remark}
\label{remark-tensor-product-not-exact}
However, tensor product does NOT preserve
exact sequences in general.
In other words, if $M_1 \to M_2 \to M_3$ is
exact, then it is not necessarily true that
$M_1 \otimes N \to M_2 \otimes N \to M_3
\otimes N$
is exact for arbitrary $R$-module $N$.
\end{remark}

\begin{example}
\label{example-tensor-product-not-exact}
Consider the injective map $2 : \mathbf{Z}\to
\mathbf{Z}$
viewed as a map of $\mathbf{Z}$-modules.
Let $N = \mathbf{Z}/2$. Then the induced map
$\mathbf{Z} \otimes \mathbf{Z}/2 \to
\mathbf{Z} \otimes \mathbf{Z}/2$
is NOT injective. This is because for
$x \otimes y\in \mathbf{Z} \otimes
\mathbf{Z}/2$,
$$
(2 \otimes 1)(x \otimes y) = 2x \otimes y = x
\otimes 2y = x \otimes 0 = 0
$$
Therefore the induced map is the zero map
while $\mathbf{Z} \otimes N\neq 0$.
\end{example}

\begin{definition}
\label{definition-flat-module}
For $R$-modules $N$, if the
functor $-\otimes_R N$ is exact, i.e.
tensoring
with $N$ preserves all exact
sequences, then $N$ is said to be {\it flat}
$R$-module.
We will discuss this later in Section
\ref{section-flat}.
\end{definition}

\medskip\noindent
Epiphany:
in the category of commutative rings
pushout is tensor product.
So think of $\Spec $ as a functor
from $\mathit{CRing}^{\text{op}}$ to $\Sch$,
then pushout goes to pullback,
and what's an example of a pullback?
Fibre!
So, the coordinate ring of a fiber
is essentially given by the residue field
at the point that parametrizes it!
(tensored with the coordinate ring
of the deformation space).

There is a lot of information on Stacks
Project about flatness.
It looks like the heart of the concept is
captured in the commutative-algebraic notion
of preserving
exact sequences:

\begin{definition}
\label{definition-flat}
Let $R$ be a ring.
\begin{enumerate}
\item An $R$-module $M$ is called {\it flat}
if whenever
$N_1 \to N_2 \to N_3$ is an exact sequence of
$R$-modules
the sequence $M \otimes_R N_1 \to M \otimes_R
N_2 \to M \otimes_R N_3$
is exact as well.
\item An $R$-module $M$ is called {\it
faithfully flat} if the
complex of $R$-modules
$N_1 \to N_2 \to N_3$ is exact if and only if
the sequence $M \otimes_R N_1 \to M \otimes_R
N_2 \to M \otimes_R N_3$
is exact.
\item A ring map $R \to S$ is called {\it
flat} if
$S$ is flat as an $R$-module.
\item A ring map $R \to S$ is called {\it
faithfully flat} if
$S$ is faithfully flat as an $R$-module.
\end{enumerate}
\end{definition}

\medskip\noindent
Recall that a module $M$ over a ring $R$ is
{\it flat} if the functor
$-\otimes_R M : \text{Mod}_R \to
\text{Mod}_R$ is exact. A ring map
$R \to A$ is said to be {\it flat} if $A$ is
flat as an $R$-module.

\section{Hilbert polynomial}
\label{section-Hilbert-polynomial}

\noindent
The following lemma will be made obsolete by
the more general
Lemma
\ref{lemma-numerical-polynomial-from-euler}.

\begin{lemma}
\label{lemma-hilbert-polynomial}
Let $k$ be a field. Let $n \geq 0$. Let
$\mathcal{F}$ be a coherent sheaf
on $\mathbf{P}^n_k$. The function
$$
d \longmapsto \chi(\mathbf{P}^n_k,
\mathcal{F}(d))
$$
is a polynomial.
\end{lemma}

\begin{proof}
We prove this by induction on $n$. If $n =
0$, then
$\mathbf{P}^n_k = \Spec(k)$ and
$\mathcal{F}(d) = \mathcal{F}$.
Hence in this case the function is constant,
i.e., a polynomial
of degree $0$. Assume $n > 0$. By
Lemma
\ref{lemma-euler-chara-exten-base-field}
we may assume $k$ is infinite. Apply
Lemma \ref{lemma-exact-sequence-induction}.
Applying Lemma
\ref{lemma-euler-characteristic-additive}
to the twisted sequences
$0 \to \mathcal{F}(d - 1) \to \mathcal{F}(d)
\to i_*\mathcal{G}(d) \to 0$
we obtain
$$
\chi(\mathbf{P}^n_k, \mathcal{F}(d)) -
\chi(\mathbf{P}^n_k, \mathcal{F}(d - 1)) =
\chi(H, \mathcal{G}(d))
$$
See Remark
\ref{remark-exact-seque-induc-cohom}.
Since $H \cong \mathbf{P}^{n - 1}_k$
by induction the right hand side is a
polynomial.
The lemma is finished by noting that any
function
$f : \mathbf{Z} \to \mathbf{Z}$ with the
property that the map
$d \mapsto f(d) - f(d - 1)$ is a polynomial,
is itself a polynomial.
We omit the proof of this fact (hint: compare
with
Algebra, Lemma
\ref{algebra-lemma-numerical-polynomial}).
\end{proof}

\begin{definition}
\label{definition-hilbert-polynomial}
Let $k$ be a field. Let $n \geq 0$. Let
$\mathcal{F}$ be a coherent sheaf
on $\mathbf{P}^n_k$. The function
$d \mapsto \chi(\mathbf{P}^n_k,
\mathcal{F}(d))$ is called the
{\it Hilbert polynomial} of $\mathcal{F}$.
\end{definition}

\noindent
The Hilbert polynomial has coefficients in
$\mathbf{Q}$ and not
in general in $\mathbf{Z}$. For example the
Hilbert polynomial
of $\mathcal{O}_{\mathbf{P}^n_k}$ is
$$
d \longmapsto {d + n \choose n} =
\frac{d^n}{n!} + \ldots
$$
This follows from the following lemma and the
fact that
$$
H^0(\mathbf{P}^n_k,
\mathcal{O}_{\mathbf{P}^n_k}(d)) = k[T_0,
\ldots, T_n]_d
$$
(degree $d$ part) whose dimension over $k$ is
${d + n \choose n}$.

\begin{lemma}
\label{lemma-hilbert-polynomial-H0}
Let $k$ be a field. Let $n \geq 0$. Let
$\mathcal{F}$ be a coherent sheaf
on $\mathbf{P}^n_k$ with Hilbert polynomial
$P \in \mathbf{Q}[t]$.
Then
$$
P(d) = \dim_k H^0(\mathbf{P}^n_k,
\mathcal{F}(d))
$$
for all $d \gg 0$.
\end{lemma}

\begin{proof}
This follows from the vanishing of cohomology
of high enough twists
of $\mathcal{F}$. See
Cohomology of Schemes,
Lemma
\ref{coherent-lemma-coherent-projective}.
\end{proof}


\medskip\noindent
For completeness I include earlier notes
from \cite{har} on the matter.

The fact that $M$ is finitely generated is
what makes the following two
definitions make sense.

\begin{definition}
\label{definition-Hilbert-function}
The {\it Hilbert function} of a finitely
generated graded $S=k[x_0,\ldots,x_r]$
-module $M$ is
$$
H_M(d)=\dim_kM_d
$$
\end{definition}

\begin{definition}
\label{definition-sysygy}
Define $F_0$ to be the free $S$-module on the
generators of $M$. Elements in the
 kernel $M_1$ of the inclusion are called
 {\it sysygies}. By Hilbert's basis
theorem, $M_1$ is also finitely generated, so
we may choose a set of generators
and repeat this process.
\end{definition}

\begin{theorem}[Hilbert Syzygy Theorem]
\label{theorem-Hilbert-syzygy}
\begin{reference}
\cite[Theorem 1.1]{sys}
\end{reference}
Any finitely generated $S$-module $M$ has a
finite graded free resolution
$$
\xymatrix{
0\ar[r]&F_m\ar[r]^{\varphi_m}&\ar[r]&F_{m-q}
\ar[r]&\cdots\ar[r]&
F_1\ar[r]^{\varphi_1}&F_0
}
$$
Moreover, we may take $m\leq r+1$, the number
of variables in $S$.
\end{theorem}

\begin{lemma}
\label{lemma-Hilbert-function}
Suppose that $S=k[x_0,\ldots,x_r]$ is a
polynomial ring. If the graded
$S$-module $M$ has finite free resolution
$$
\xymatrix{
0\ar[r]&F_m\ar[r]^{\varphi_m}&F_{m-1}
\ar\cdots[r]&F_1\ar[r]^{\varphi_1}&
\ar[r]&F_0
}
$$
with each $F_i$ a finitely generated free
module,
$F_i=\bigoplus_{j}S(-a_{i,j})$, then
\begin{equation}
\label{equation-Hilbert-function}
H_M(d)=\sum_{i=0}(-1)^i\sum_{j}\binom{r+d-
a_{i,j}}{r}
\end{equation}
\end{lemma}

\begin{lemma}
\label{lemma-hilbe-funct-becom-polyn}
There is a polynomial $P_M(d)$ called the
{\it Hilbert polynomial} such that, if
$M$ has free resolution as above, then
$P_M(d)=H_M(d)$ for
$d\geq\text{max}_{i,j}\{a_{i,j}-r\}$.
\end{lemma}

\begin{proof}
When $d$ satisfies this bound then the
binomial coefficients in Eq.
\ref{equation-Hilbert-function} are
polynomials of degree $r$ in $d$.
\end{proof}

\begin{theorem}[Hilbert-Serre]
\label{theorem-Hilbert-Serre}
\begin{reference}
\cite[I, Theorem 7.5]{hart}
\end{reference}
Let $M$ be a finitely generated graded
$S=k[x_0,\ldots,x_n]$. Then there exists
a unique polynomial $p_M$ such that
$p_M(\ell)=\dim S_\ell$ for large enough
$\ell$.
\end{theorem}

\begin{definition}
\label{definition-Hilbert-polynomial}
\begin{reference}
\cite[I, p. 52]{har}
\end{reference}
The polynomial $P_M$ of Hilbert-Serre Theorem
\cite{Hilbert-Serre} is the {\it
Hilbert polynomial} of the finitely generated
$k[x_0,\ldots,x_n]$-module $M$.
\end{definition}

\begin{definition}
\label{definition-degre-proje-varie}
\begin{reference}
\cite[p. 52]{hart}
\end{reference}
If $Y\subset \mathbb{P}^n$ is an algebraic
set of dimension $r$, we define the
{\it degree of $Y$} to be $r!$ times the
leading coefficient of the Hilbert
polynomial of the homogeneous coordinate ring
$S(Y)$.
\end{definition}

\begin{exercise}
\label{exercise-very-ample-bundl-self-inte}
\begin{reference}
\cite[V, Exercise 1.2]{har}
\end{reference}
Let $H$ be a very ample divisor on the
surface $X$, corresponding to a
projective embedding
$X\subseteq\mathbb{P}^N$. If we write the
Hilbert
polynomial of $X$ as
$P(z)=\frac{1}{2}az^2+bz+c$, show that
$a=H^2$,
$b=\frac{1}{2}H^2+1-\pi$, where $\pi$ is the
genus of a nonsingular curve
representing $H$, and $c=1+p_a$.
\end{exercise}

\section{Nakai-Moishezon Criterion}
\label{section-Nakai-Moishezon-criterion}

\begin{theorem}[Nakai-Moishezon Criterion]
\label{theorem-Nakai-Moishezon-criterion}
\begin{reference}
\cite[V, Theorem 1.10]{hart}
\end{reference}
A divisor $D$ on the surface $X$ is ample if
and only if $D^2>0$ and $D.C>0$ for
all irreducible curves $C$ in $X$.
\end{theorem}

\begin{proof}
The direct implication is easy: since $D$ is
ample,  $mD$ is very ample for some
$m$, so that $m^2D^2$ is the
self-intersection number of $mD$. By exercise
\ref{exercise-very-ample-bundl-self-inte},
 $D^2$ is the leading coefficient of the
 Hilbert polynomial of $X$ as a
subscheme of $\mathbb{P}^n$. This means that
$D^2$ is twice the leading
coefficient of the Hilbert polynomial of a
projective variety for large enough
$m$, so that it must be a positive number
(it's the dimension of one of the
graded components of the coordinate ring of
the surface).
\end{proof}

\section{Hilbert scheme}
\label{section-Hilbert-scheme}

{\bf Upshot \cite[p. 6]{HarrMorr}.}
We wish to parametrize subschemes of a
projective space (or
perhaps a more general scheme?). Since there
are too many such subschemes we
restrict ourselves to schemes with a given
Hilbert polynomial, since the latter
``encodes the most important numerical
invariants of schemes''. The Hilbert
scheme is introduced via a theorem by
Grothendieck
as the object that represents the functor
$\mathbf{Hilb}_{P,r}$
that maps a reduced scheme $B$ to the
set of proper flat families
$$
\xymatrix{
\mathcal{X}\ar@{^{(}->}[r]^i\ar[rd]&
\mathbb{P}^r
\times B \ar[d]^{\pi_B}\\
& B
}
$$
with $\mathcal{X}$ having Hilbert polynomial
$P$.

\begin{theorem}[Grothendieck, '66]
\label{theorem-Grothendieck}
The functor $\mathbf{Hilb}_{P,r}$ is
representable by a projective scheme
$\mathcal{H}_{P,r}$.
\end{theorem}

SEE Hilbert schemes of subschemes.

\section{Deformation theory}
\label{section-deformation-theory}

\noindent
\cite{Sernesi-Overview} has a great way of
motivating deformation
theory via the moduli problem.
Here let's just put the following important
meaning for
``deformation theory'' (as opposed to
studying flat families!):

\begin{quote}
``{\it Deformation theory} is the study of
infinitesimal deformaitons
as a tool to understand the local structure
of the moduli space.
The goal is to be able to describe the
restriction
of the universal family to a small
neighborhood of $m \in \mathcal{M}$,
or, more precisely, its restriction to the
germ of $M$ at $m$.

What is interesting here is that we can study
order and infinitesimal
deformations even though the functor $F$ is
not representable or simply we don’t
yet know it is. This is the most frequent
case. Such an investigation will
reveal the infinitesimal properties at $[X]$
of a yet unknown global structure
on $M$ which will be hopefully understood at
a subsequent stage of the
investigation.  In order words it turns out
to be possible and convenient to
separate the {\it global moduli problem} from
the {\it local moduli problem},
and deformation theory studies the latter,
with the purpose of constructing a
family of deformations of a given object
parametrized by the spectrum of a local
ring, and having properties as close as
possible to a universal property.''
\end{quote}




\noindent
I start by reading Stacks Project.

\medskip\noindent
The first notion is {\it thickening of ringed
spaces},
which I ultimately think of as a closed
subscheme
$(X,\mathcal{O}_X) \to (X',\mathcal{O}_{X'})$
with a nilpotent ideal sheaf,
which in an imprecise way means
that $\mathcal{I}_X^n=0$ for some $n$,
and in a precise way its given on sections.

\medskip\noindent
The following definition is from
[lucas-defos], which in turn comes from
\cite{Sernesi-deformations}

\begin{definition}
\label{definition-deformation}
Let $X$ be an algebraic $\mathbb{C}$-scheme.
\begin{enumerate}
\item A {\it deformation} of $X$ is a
Cartesian diagram $\xi$
$$
\xymatrix{
X\ar[r]\ar[d]&\mathcal{X}\ar[d]^{\pi}\\
\Spec \mathbb{C}\ar[r]^s&S
}
$$
where $\pi$ is a flat surjective morphism of
algebraic $\mathbb{C}$-schemes and
$S$ is connected. (Recall that flatness
accounts for ``continuity''.)
\item A {\it local deformation} of $X$ is a
deformation $\xi$ where
$S=\Spec A$ for $A$ a noetherian local
$\mathbb{C}$ algebra with residue
field $\mathbb{C}$.
\item An {\it infinitesimal deformation of
$X$} is a local deformation with $A$
an artinian local $\mathbb{C}$-algebra with
residue field $\mathbb{C}$. $X$ is
called {\it rigid} if all infinitesimal
deformations are trivial.
\item An {\it inifinitesimal deformation of
order $n$} is an infinitesimal
deformation when $S=\Spec
(\mathbb{C}[t]/(t^{n+1})$.
\end{enumerate}
\end{definition}

{\bf Upshot.} An interpretation of the
so-called {\it dual numbers}
$k[t]/(t^2)$ (see
\cite{Hartshorne-deformation}) as the tangent
space of
something is thinking of Taylor polynomials:
after quotienting by $t^2$ we loose
the tails of the polynomials and are left
with the first derivative information
only - this justifies the use of the ``first
order deformations'' term.

There is the following interpretation in
[continued-fractions] p. 39: the space
of first order deformation classes of $X$ is
$D(\mathbb{C}[t]/(t^2)$. This is
said to ````represent'' the tangent space
$\mathbb{T}^1_X$ of the hypothetical
deformation space of $X$''. (I put double
quotations because the word
``represent'' is quoted in the original
text.) Further, if $X$ is nonsingular
and compact, then
$\mathbb{T}^1_X=H^{1}(X,T_X)$.

Which basically I interpret as: the dimension
of the deformation space of a
smooth compact variety is $H^{1}(X,T_X)$.

\medskip\noindent
After a few months of the previous text in
this section
I think I have understood the construction of
Kodaira-Spencer map for embedded
deformations,
which is \cite[Proposition
3.2.1(ii)]{Sernesi-deformations}.

But I will actually briefly outline a
particular case
of this, which is easier: \cite[Examples
3.2.4(ii)]{Sernesi-deformations}.
This the case when $X \subset Y$ is a Cartier
divisor.
Then $X$ is cut out locally by a set of
functions $f_i$
defined on an affine cover $U_i$ of $Y$.
Then we have a line bundle defined by the
transition functions $f_{ij}=f_i/f_j$.
This is the line bundle $\mathcal{O}_Y(X)$,
which
by the adjunction formula is also the normal
bundle $N_{X/Y}$
of $X$ in $Y$.

Now consider a first order deformation
$$
\xymatrix{
X\ar[r]\ar[d]
&\mathcal{X}\subset Y\times
B_\varepsilon\ar[d]^\pi\\
0\ar[r]
&B_\varepsilon
}
$$

The great thing about first order
deformations
is that the structure sheaf of $\mathcal{X}$
on an affine open $U_i$ of $Y$
looks like the tensor product
$\mathcal{O}_{U_i}\otimes
\mathbb{C}[\varepsilon]/(\varepsilon^2)$
(\href{https://stacks.math.columbia.edu/tag/%
01I4}{01I4}).
By definition of tensor product this just
says
that sections look like $f+\varepsilon g$
for $f,g \in \Gamma(\mathcal{O}_{U_i})$.

If $\mathcal{X}$ is a Cartier divisor on $Y
\times B_\varepsilon$
(it is; but I'm missing a formal proof;
the intuition is that an ``infinitesimal
thickening'' $\mathcal{X}$
of $X$ has the same underlying topological
space
but different scheme structure;
\href{https://stacks.math.columbia.edu/tag/%
01JL}{01JL}
\href{https://stacks.math.columbia.edu/tag/%
01I4}{01I4}),
it is given as the vanishing locus
of some local sections $f_i+\varepsilon g_i$.
(Perhaps the $f_i$ coincide with the
distinguished
section of $X$.)
The line bundle associated to $\mathcal{X}$
is given by the cocycle
$$
F_{ij}=f_{ij}+\varepsilon g_{ij}
$$
for $g_{ij}=\frac{g_i-f_{ij}g_j}{f_j}$.
Indeed, note that
$\frac{1}{f_j+\varepsilon
g_j}=\frac{1}{f_j}-\varepsilon
\frac{g_j}{f^2_j}$
and compute $F_{ij}=(f_i+\varepsilon
g_j)/(f_j+\varepsilon g_j)$.

Then
$$
f_i+\varepsilon g_i=(f_{ij}+\varepsilon
g_{ij})(f_j+\varepsilon g_j)
$$
gives
$$
g_i=f_{ij}g_j+g_{ij}f_j
$$
and since $f_j$ vanishes along $X$,
we are left with $g_i=f_{ij}g_j$
which is how a section of the sheaf $N_{X/Y}$
is constructed by definition of sheaf.


\begin{example}
\label{example-modul-space-nonsi-riema}
\begin{reference}
\cite[Example 3.1]{continued-fractions}
\end{reference}
Fix $g \geq 2$. The dimension of the
deformation space of a nonsingular
projective curve $X$ is $3g-3$. This is ``the
dimension of the moduli space of
curves of genus $g$''.

We can compute this number by Riemann-Roch
formula on the bundle $-K_X$.
Indeed, since $X$ is a curve,
$\Omega_X^1=K_X$ and thus $-K_X=T_X$. We get
$$
\begin{aligned}
h^0(-K)-h^0(K-(-K))&=\text{deg}(-K)-g+1\\
&=-2g+2-g+1\qquad \substack{\text{degree
additive and}\\ \text{deg}K=2g-2}\\
&=-3g+3
\end{aligned}
$$
Now by Serre duality
$h^0(K-(-K))=h^1(-K)=h^1(T_X)$, and it turns
out that a
Riemann surface of genus $g \geq 2$ has no
holomorphic vector fields, so that
$h^0(-K)=h^0(T_X)=0$.
\end{example}

\begin{exercise}
\label{exercise-deformations-of-curves-in-K3}
Let $C$ be a smooth genus $g$ curve which can
be embedded in a K3 surface  $M$,
and $X$ the family of all deformations of $C$
in $M$.
\begin{enumerate}
\item Prove that $\dim X \leq  g$.
\item Let $\mathcal{X}_g$ be the space of all
curves of genus $g$ (smooth?)
which can be possibly embedded to a K3
surface. Prove that each irreducible
component $Z$ of $\mathcal{X}_g$ satisfies
$\dim_\mathbb{C} \leq  g+19$. Deduce
that there exists a compact complex curve
which cannot be embedded in a K3
surface.
\end{enumerate}
\end{exercise}

\begin{exercise}
\label{exercise-}
Let $C$ be a smooth curve embedded in a K3
surface $X$. Show that the dimension
of the deformation space of $C$ is $\leq g$.
\end{exercise}

\begin{proof}
% \begin{enumerate}
% \item
The deformation space of a variety is the
space of isomorphism classes of
deformations as explained above. It turns out
that there is a way to associate
1-cocyles of the tangent sheaf to
deformations, so that in fact the deformation
space $\text{Def}_1$ is isomorphic to
$H^{1}(X,\mathcal{T}_X)$ for any variety
$X$.

For our curve $C$ we thus know that the
dimension of the space of deformations
(deformations not necessarily contained in
$X$) is $h^1(\mathcal{T}_C)$.  The
family of deformations of $C$ that are
contained in $X$ is the Hilbert scheme of
curves with fixed Hilbert polynomial $P(t)$
after quotienting by $\mathbb{C}s$,
where $s$ is the section whose vanishing
locus is  $C$.
This says that the number we are looking for
is $h^{0}(X,\mathcal{O}(C))-1$.

Before showing that in fact
$h^0(X,\mathcal{O}(C))=1+g$,
let's explain why this Hilbert scheme
construction actually works.
Intuitively, any curve lying in the same
component
as $C$ would be deformation-equivalent to $C$
--- at least we know that it's Hilbert
polynomial is the same.
But it's not clear what's the deformation
involving these two
(this question should be addressed!).
Further, this doesn't explain Vladimiro's
suggesion of ``taking quotient by
$\mathbb{C}s$''.
To me it looks like I would just compute the
dimension
of the tangent space of $\text{Hilb}^P(M)$ at
$C$.

Now I will show that
$h^0(X,\mathcal{O}(C))=1+g$ (see
\cite[Lemma 1.2.1, Remark 1.2.2]{huk}).

Consider the ideal sheaf exact sequence
twisted by
$\mathcal{O}(C)$:
$$
\xymatrix{
0\ar[r]&\mathcal{O}_X\ar[r]&\mathcal{O}(C)
\ar[r]
&\mathcal{O}_X(C) \otimes
\mathcal{O}_C=\mathcal{O}(C)|_{C}\ar[r]&0
}
$$
By $X$ being a K3 we know that
$H^{1}(\mathcal{O}_X)=0$, so that we have the
short exact sequence in cohomology
$$
\xymatrix{
0\ar[r]&H^{0}(\mathcal{O}_X)\ar[r]&H^{0}
(\mathcal{O}(C))\ar[r]
&H^{0}(\mathcal{O}(C)|_{C})\ar[r]&0
}
$$
so that
$h^0(\mathcal{O}(C))=1+h^1(\mathcal{O}(C)
|_{C})$.
A version of adjunction formula says
$\omega_C \cong (K_X
\otimes\mathcal{O}(C)|_{C}$, and using that
$K_X=\mathcal{O}_X$ we obtain
$h^0(\mathcal{O}(C)|_{C})=h^0(\omega_C)=g$.

To

\medskip\noindent
For the record I put other thoughts I went
through in solving this exercise.

Recall from \cite[p. 146]{gri} that the
normal bundle
$\mathcal{N}$ of a hypersurface of a smooth
variety $X$ satisfies
$\mathcal{N}^\vee\cong\mathcal{O}_X(-C)|_{C}
$.
Taking duals we get that
$\mathcal{N}\cong\mathcal{O}_X(C)|_{C}$.

By adjunction formula
$2g-2=(\mathcal{O}(C),\mathcal{O}(C))$.
Applying Riemann-Roch
to $\mathcal{O}(C)$ (which is by definition
the dual of the ideal sheaf of $C$,
which is a line bundle on $X$), we obtain
that
$\chi(\mathcal{O}(C))=2+\frac{1}{2}
(\mathcal{O}(C),\mathcal{O}(C))$.
Then
$\chi(\mathcal{O}(C))=g+1$.

Now recall that
$\chi(\mathcal{O}(C))=h^0(\mathcal{O}(C)
-h^1(\mathcal{O}(C))+h^2(\mathcal{O}(C))$. By
Serre duality and $X$ being a K3
surface we see that $h^2(\mathcal{O}(C))\cong
h^{0}(\mathcal{O}(-C))$, which is
the ideal sheaf of $C$. Any section of such a
sheaf would vanish along $C$, and
since $X$ is compact we conclude there cannot
be any such section.

Now we show that also $h^1(\mathcal{O}(C))=0$
to conclude that
$h^0(\mathcal{O}(C))=g+1$.

We thus conclude that $h^0(\mathcal{O}(C))$




Since $C$ is smooth we can use the normal
exact sequence
$$
\xymatrix{
0\ar[r]&\mathcal{T}_C\ar[r]&\mathcal{T}
_X|_{C}\ar[r]&\mathcal{N}\ar[r]&0
}
$$
Taking Euler characteristic we see that
$\chi(\mathcal{T}_C)+\chi(\mathcal{N})=
\chi(\mathcal{T}_X|_{C})$.


This means that we should be done once we
compute $\chi(\mathcal{T}_X|_{C})$.
For this we can use Riemman-Roch formula for
coherent sheaves on a curve, which
tells us that
$$
\text{deg}(\mathcal{T}_X|_{C})=
\chi(\mathcal{T}_X|_{C})
-\text{rk}(\mathcal{T}_X|_{C})
\cdot\chi(\mathcal{O}_C)
$$
But of course we know that since $C$ is a
curve, by Serre duality we get that
$\chi(\mathcal{O}_C)=1-g$. (Indeed:
$h^1(\mathcal{O}_C)=h^0(\Omega^1_C):=p_a(C)
$.)
And now the question is what is
the degree. Apparently this is just the first
Chern class $c_1$ of the
restricted tangent bundle. So what is it? And
what is $h^0(\mathcal{T}_C)$ if
the genus is 0 or 1, and that's it.

I
know that $h^0(\mathcal{T}_X)=0$ and
$h^1(\mathcal{T}_X)=20$ by $X$ being a K3
surface, but I'm not sure what happens when
we restrict to $C$.

If $g \geq 2$ we know that
$h^0(\mathcal{T}_C)=0$, so that
$\boxed{\chi(\mathcal{T}_C)=h^1(\mathcal{T}
_C)}$.
\bigskip
To compute the
latter Euler characteristic, which is given
by definition by
$\chi(\mathcal{T}_X|_{C})=h^0(\mathcal{T}
_X|_{C})-h^1(\mathcal{T}_X|_{C})$,
we
first note that $h^0(\mathcal{T}_X|_{C})=0$,
because this is the dimension of
the global holomorphic vector fields on $X$
restricted to $C$, which is constant along
$C$ since $X$ is
smooth. And then this number is actually zero
by Hodge numbers of a K3 surface.

and taking cohomology long exact sequence we
obtain
$$
\xymatrix{
\cdots\ar[r]&H^{0}(T_X)\ar[r]&H^{0}
(\mathcal{N})\ar[r]&H^{1}(T_C)\ar[r]&\\
H^{1}(T_X)\ar[r]&H^{1}(\mathcal{N})\ar[r]&
\cdots
}
$$
\medskip\noindent
Or is it $h^1(\mathcal{N})$? See
\cite{HarrMorr}. If this was the case, then
we
can use adjunction formula as above to get
that
$2g-2=(\mathcal{N},\mathcal{N})=\text{deg}_C
\mathcal{N}$. Then we may find
$h^0(\mathcal{N})$ via Riemann-Roch:
$$
h^0(\mathcal{N})-h^0(K_C-\mathcal{N})=
\text{deg}\mathcal{N}-g+1
$$
Note that
$h^0(N_C-\mathcal{N})=h^1(-K_C+N+K_C)=
h^1(\mathcal{N})$
via Serre
duality, and by Riemann-Roch on a surface as
above we see that
$$
g+1=\chi(\mathcal{N})=h^0(\mathcal{N})-
h^1(\mathcal{N})
\implies
h^1(\mathcal{N})=-g-1+h^0(\mathcal{N})
$$
so that
$$
\begin{aligned}
h^0(\mathcal{N})-(-g-1+h^0(\mathcal{N}))&=
\text{deg}\mathcal{N}-g+1\\
\implies \text{oops! I lost
$h^0(\mathcal{N})$ in this operation…}
\end{aligned}
$$
Maybe if I just use normal exact sequence and
realise that
$$
h^0(\mathcal{N})=h^0(\mathcal{T}_X|_{C})-
h^0(\mathcal{T}_C)
$$
I know that $h^0(\mathcal{T}_C)=0$ for $g>1$,
so the question is how to compute
the restricted holomorphic vector fields.
% \end{enumerate}
\end{proof}

\medskip\noindent
Look for the ``Perfect obstruction theory''.
This is used by Bruzzo for deformations of
(super)stacks,
work with E. Paiva and ---.

\medskip\noindent
Until I have a better candidate of where to
put this…

\begin{definition}
\label{definition-extension}
\begin{reference}
\cite[p.7]{Sernesi-deformations}
\end{reference}
Let $A \to R$ be a ring homomorphism
(i.e. $R$ is an $A$-algebra).
An {\it $A$-extension of $R$ by $I$} is a
short exact sequence
$$
\xymatrix{
0\ar[r]
&I\ar[r]
&E\ar[r]^\varphi
&R\ar[r]
&0.
}
$$
where $E$ is an $R'$-algebra
and $\varphi$ is a morphism of $A$-algebras
whose kernel is $I$.
We also require that $I^2=(0)$;
this condition implies that $I$ has the
sturcture of $R$-module.
\end{definition}

\begin{example}
\label{example-dual-numbers-extension}
$$
\xymatrix{
0\ar[r]
&(\varepsilon)\ar[r]
&A[\varepsilon]\ar[r]
&A\ar[r]
&0
}
$$
where $A[\varepsilon]$ is the algebra of dual
numbers over $A$.
\end{example}

\section{Cotangent modules}
\label{section-cotangent-modules}

\noindent
My aim in this section is to cover the
minimal
theory needed to apply results in \cite{jan2}
to my deformation problem. The main theorem
I expect to use is the following:

\begin{theorem}
\label{theorem-deformations-threefolds}
\begin{reference}
\cite[Theorem 5.7]{jan2}
\end{reference}
If $\mathcal{K}$ is a 3-dimensional manifold
then
$$
\dim T^1_{\mathbb{P}(K)}
=11d_3+5e_3+3e_4+e_{\geq 5}+c_{\geq
6}+5f_1^{(3)}+2f_1^{(4)}+h^2(\mathcal{K}).
$$
\end{theorem}

\noindent
I need to do some interpretation first.
From \cite[p. 1]{jan3}:

\begin{quotation}
  ``For each $\mathbb{C}$-algebra $A$,
there is a cohomology theory providing
modules $T^i_A$.
However, only $T^1_A$ and $T^2_A$ are
relevant for the deformation
theory of $\Spec A$. The module $T^1_A$
collects the infinitesimal
deformations of $A$ and $T^2_A$ contains the
obstructions
for lifting these deformations to decent
parameter spaces.
Eventually, $\text{Der}_{\mathbb{C}}(A,A)$
will be called $T_A^0$."
\end{quotation}

\cite[p. 3]{jan3}:
\begin{quotation}
``In general, for a finitely generated
$\mathbb{C}$-algebra $A$,
the modules $T_A^i$ allow the following ad
hoc definitions:
Let $P=\mathbb{C}[x_0,…,x_n]$ mapping onto
$A$ so that $A \simeq P/I$
for an ideal $I$. The $T_A^1$ is the cokernel
of the natural map
$\text{Der}_{\mathbb{C}}(P,P) \to
\Hom_P(I,A)$.
Moreovr, […] and we obtain $T^2_A$ as the
cokernel
of the induced map
$\Hom_P(P^m,A)\to\Hom_A(R/R_0,A)$.''
\end{quotation}

\begin{remark}[dani]
\label{remark-dani1}
Let's think more about that definition of
$T^1_A$
as the cokernel of the ``natural map''
$\text{Der}_\mathbb{C}(P,P) \to \Hom_P(I,A)$.

Pick a derivation $\chi \in
\text{Der}_{\mathbb{C}}(P,P)$.
We want a map $I \to A=P/I$.
Pick an element $a \in I$ and map it to
$\chi(a)$.
And that's it. It might be zero or not.
If the $\chi$ we picked maps any $a \in I$ to
some element in $I$,
then this $\chi$ is in the kernel of the
natural map.
So $T^1_A$ is the image of the natural map
modulo the kernel.
\end{remark}

\noindent
Unfortunately this doesn't say much about the
following crucial fact:
the $T^i_A$ modules are
$\mathbb{Z}^{n+1}$-graded.
But here goes an explanation.
We use the notation that $e_p$ is a ``face
not in the complex'',
which of course corresponds to a generating
monomial $x^p$
of the SR ideal $I$.
It has a degree, of course (given by the
``support'' of $p$).
Now here's a rather sloppy step but:
since the cotangent modules are built upon
$I$,
who is given by $e_p$,
we just notice that the degree of $e_p$
along with a (very natural, too) notion of
degree among relations,
induce a grading in all the cotangent modules
$T_A^i$.
Thus, the grading is morally just the degree
of the monomials,
which is in fact the support of $p$, a vector
$\mathbf{c} \in \mathbb{Z}^{n+1}$.

Next in line is the decomposition
$\mathbf{c}=\mathbf{a}-\mathbf{b}$
into positive minus negative part. But how
exactly…?








\section{Fano varieties}
\label{section-Fano-varieties}

\begin{definition}
\label{definition-Fano-variety}
A {\it Fano variety} is a projective variety
with $-K_X$ ample.
\end{definition}

\noindent
The idea of the Fano index is that the
larger the index, the simpler is the variety.
In a talk by Jõao, it is defined as the
larger integer
dividing $-K_X$ in $\text{Pic}(X)$.
\begin{definition}
\label{definition-Fano-index}
$$
r(X):=\text{min}\{r:\frac{c_1(X)}{r}\in
H^{2}(X,\mathbb{Z})\}
$$
\end{definition}

\begin{exercise}
\label{exercise-fano-vanis-highe-cohom}
By Kodaira vanishing theorem
\ref{theorem-Kodaira-vanishing},
you can show that the cohomology $H^{i}(X,L)$
for
a Fano variety $X$ vanishes. You just have to
put $L=\mathcal{O}(k)$ with $k\geq
-r$, where $r$ is the Fano index.
\end{exercise}

\begin{exercise}
\label{exercise-Pic-H2-Fano}
Show that  $\text{Pic}(X)\cong
H^{2}(X,\mathbb{Z})$ holds for Fano
varieties.
\end{exercise}

\begin{remark}
\label{remark-deriv-categ-fano-3-folds}
\begin{reference}
Marcos Jardim, CIMPA 2025 Florianópolis,
Lecture 2.
\end{reference}
If $H^3(X,\mathbb{Z})=0$ of a Fano 3-fold,
then its derived category is
generated by 4 elements.
\end{remark}

\medskip\noindent
Araujo-Druel em 2017 classify
Fano foliations relating the Fano index (of
the foliation)
with its rank.
João found conditions
(integrability of the leaves)
to describe del Pezzo foliations
(a kind of Fano foliation)
with Picard rank 1.
Araujo (2019) classifies the leaves of a
algebraically integrable
del Pezzo foliation (with no restriction on
the (log
canonica) singularities!).


\section{Quivers}
\label{section-quivers}

\begin{definition}
\label{definition-quiver}
A {\it quiver} is a set of vertices $Q_0$, a
set of arrows $Q_1$ equipped with
the maps of source $s$ and target $t$ that to
each arrow they assign the point
that is source or target of the arrow.
\end{definition}

\begin{definition}
\label{definition-representation-of-quiver}
A {\it representation} of a quiver is a set
of finite dimensional vector spaces
equipped with maps between them realising a
given quiver (incomplete…).
\end{definition}

There is a notion of projective
representation, which I missed to write. But
it
is analogous to the injective representation:

\begin{definition}
\label{definition-injec-repre-quive}
Given a quiver $Q$, the {\it injective
representation} of $Q_0$ is given by, for
$i \in Q_0$,
$$
I(i)_j=\begin{cases}
k\qquad &i=j \\
k^{d'}\qquad &j \neq i
\end{cases}
$$
where $d'$ is the number of paths from $j$ to
$i$.
\end{definition}

\section{Stacks}
\label{section-stacks}

My first definition of stack can be extracted
from

\begin{definition}
\label{definition-superstack}
A {\it superstack} is a stack over
the étale site $\text{SSch}$ of superschemes,
i.e. it is a category fibered in groupoids
over the category of superschemes,
the latter equipped with the
étale topology,
satisfying the descent condition.
\end{definition}

Here are some other definitions:

\begin{definition}
\label{definition-algebraic-stacks}
Let  $\mathfrak{X}$ be a stack over
$\text{Sch}_{\text{ét}}$.
An {\it algebraic space} is
such that there exists morphism
$\mathcal{U} \to \mathfrak{X}$
where $\mathcal{U}$ is a scheme, that is
schematic, étale and injective (check this
one).

$\mathfrak{X} \to y$ is {\it representable}
if
there exists a scheme $\mathcal{U}$ and a map
$\mathcal{U} \to y$ such that the
fibered product
$$
\xymatrix{
\mathcal{U} \times_y
\mathfrak{X}\ar[r]\ar[d]\ar@{}[dr]|-
{\lrcorner}&\mathfrak{X}\ar[d]\\
\mathcal{U}\ar[r]&y
}
$$
is an algebraic space.

Finally, a  stack is {\it algebraic} (resp.
{\it Deligne-Mumford})
is there exists a
representable surjective morphism
$\mathcal{U} \to \mathfrak{X}$
that is smooth (resp. étale).

A {\it stable map} over a projective
variety $X$ is an element of the first
Chow group $\beta \in A_1$, where
 $(C,g)$ is an algebraic curve and
$f:C \to X$ with $[f(C)]=\beta$.
\end{definition}

\noindent
The curves that are points under this map
(contractible) are {\bf stable}.

\section{Moduli}
\label{section-moduli}

\noindent
Pseudo-admissible covers may be the way
Mumford orifinally developed stability
notions.

\noindent
This is a two-session minicourse by Ugo
Bruzzo
in LEGALzinho 2025.

Let $X$ be a smooth complex projective
variety,
$X \overset{j}{\hookrightarrow}\mathbb{P}^n$.
$\mathcal{O}_{\mathbb{P}^n}(1)$ is the class
of hyperplanes of $\mathbb{P}^n$.
$\mathcal{L}=j^*\mathcal{O}_{\mathbb{P}^n}(1)
$.
$c_1(\mathcal{L})\in H^2(X,\mathbb{Z})$ is
a {\it polarization}.

Let $\mathcal{E}$ be a coherent
$\mathcal{O}_X$-module.
For example, $\mathcal{E}$ may be a locally
free $\mathcal{O}_X$-module.
Recall that $H^k(X,\mathcal{E})$ are finite
dimensional vector spaces
and they vanish for $k>\dim X=n$.
Recall that the {\it Euler characterstic} of
$\mathcal{E}$
is given by
$\chi(\mathcal{E})=\sum_{k=0}^n(-1)
^k\dim_{\mathbb{C}}H^k(X,\mathcal{E})$.

The {\it reduced Hilbert polynomial} of
$\mathcal{E}$ is
$$
p_{\mathcal{E}}(k)=\frac{1}{\text{rk}
\mathcal{E}}
\chi(\mathcal{E}\otimes_{\mathcal{O}_X}
\mathcal{L}^k).
$$

\begin{definition}
\label{definition-stable-module}
Let $\mathcal{E}$ be a coherent
$\mathcal{O}_X$-module
with no torsion and positive rank.
$\mathcal{E}$ is {\it stable} if for all $f
\subset \mathcal{E}$
($0 \text{rk} f<\text{rk}\mathcal{E}$),
$$
p_f(k)<p_{\mathcal{E}}(k)\qquad \text{for }
k\gg 0
$$
and {\it semi-stable} if we put $\leq$
instead of $<$.

\begin{lemma}
\label{lemma-stable}
If $\mathcal{E}$ is stable then $\mathcal{E}$
is irreducible
($\mathcal{E} \neq  \mathcal{E}_1 \oplus
\mathcal{E}_2$)
and $\text{Aut}(\mathcal{E})\simeq
\mathbb{C}^*$.
\end{lemma}

\end{definition}

The {\it functor of moduli} is
$$
\begin{aligned}
\mathcal{M}: \text{Sch}_\mathbb{C}
&\longrightarrow \mathsf{Set} \\
S
&\longmapsto
\left\{
\substack{\text{stable family}  \\ \text{of
sheaves with} \\
\text{Hilbert polynomial} f\\
\text{parametrized by $S$}}
\right\}\Big/\sim.
\end{aligned}
$$

\medskip\noindent
Riemann-Roch.
Let $C$ be a smooth, projective, connected
curve
of genus $g$.
Let $\mathcal{E}$ be a coherent sheaf over
$H^2(C,\mathbb{Z})$
with generator of $H^2(X,\mathbb{Z})$
$w=[pt] \in H^2(C,\mathbb{Z})$.
$c_1(\mathcal{E})=\text{deg}\mathcal{E}\cdot
w$.

\begin{theorem}
\label{theorem-euler-characteristic-moduli}
$$
\begin{aligned}
\chi(\mathcal{E})
&=\int_C(1+(1-g)w)(\text{rk}\mathcal{E}+
\text{deg}\mathcal{E}w)\\
&=\int_C
\text{deg}\mathcal{E}+w+(1-g)\text{rk}
\mathcal{E}\cdot
w\\
&=\text{deg}(\mathcal{E})+1-g
\text{rk}\mathcal{E}.
\end{aligned}
$$
\end{theorem}

\begin{example}
\label{example-hilbert-polynomial-sheaf}
\end{example}

\begin{theorem}
\label{theorem-moduli-suave}
Se $(r,d)$, $r>0$ are coprime,
the functor of moduli of stable bundles
over $C$ of rank $r$ and degree $d$
is representable, and the moduli space
$\mathcal{M}_C(\mathcal{E},d)$ is smooth.
\end{theorem}

\noindent
In the non coprime case this does not
necessarily hold.

\begin{definition}
\label{definition-}
$\mathcal{M}$ is a {\it grosseiro} moduli
space
if the funtor is not representable.
\end{definition}

\noindent
This means that although there is not
a universal family,
there is still a correspondence
of the closed points of $\mathcal{M}$ and the
equivalence classes of semistable sheaves
of rank $r$ and degree $d$.

\medskip\noindent
Now let's fix $\dim X=2$.
Let $\mathcal{E}$ be a stable
$\mathcal{O}_X$-module.
Let $r=\text{rk}\mathcal{E}$,
and the {\it discriminant}
$$
\Delta(\mathcal{E})=c_2(\mathcal{E})-
\frac{\mathcal{E}-1}{2r}c_1(\mathcal{E})^2
\in H^4(X,\mathbb{Q})\cong\mathbb{Q}.
$$
Here enter $c_1(\mathcal{E}) \in
H^2(X,\mathbb{Z})$
and $c_2(\mathcal{E}) \in H^4(X,\mathbb{Z})$.

\begin{theorem}[Bogomolov]
\label{theorem-bogomolov}
$\Delta(\mathcal{E}) \geq 0$.
\end{theorem}

\noindent
It could be interesting to study the proof of
this theorem.

\medskip\noindent
Miyaoka (1987). Consider a curve $C$ and a
bundle $E$.
We can for any fibre, which is a vector space
$V$,
consider its projectivization
$V\setminus\{0\}/\mathbb{C}^*=\mathbb{P}(V)$.
Then we can bundlize this construction ad
obtain
$\mathbb{P}E \to C$.
Consider $\mathcal{O}_{\mathbb{P}E}(1)$.
We define
$$
\lambda=rc_1(\mathcal{O}_{\mathbb{P}E}(1))
-\pi^*c_1(E) \in H^2(\mathbb{P}E,\mathbb{Z}).
$$

\begin{theorem}
\label{theorem-semistable-nef}
$E$ is semistable if $\lambda$ is
semicorrente efectiva.
\end{theorem}

\noindent
\begin{definition}
\label{definition-nef}
Consider again a projective smooth connected
curve $\Gamma$ and let $f:\Gamma \to
\mathbb{P}E$.
We say $\lambda$ is {\it nef} if
$$
\int_{\Gamma} f^*\lambda\geq 0.
$$
\end{definition}


\noindent
The following theorem was proved by Nakayama
(who is a different author than that of
Nakayama's lemma)
in 1999 but was unkown to Bruzo and Daniel
HR.
\begin{theorem}[Nakayama '99, Bruzzo-HR '04]
\label{theorem-bruzzo-daniel-nakayama}
The following are equivalent:
\begin{enumerate}
\item $\lambda$ is nef.
\item For all $f:\Gamma \to X$,
$f^*E$ is semi-stable.
\item $E$ is semistable with respect to
the polarization, and  $\Delta(E)=0$.
\end{enumerate}
\end{theorem}

The polarization says how the surface is
embedded in projective space,
it is $H=c_1(\mathcal{L})$. The stability
notion is
$$
\mu(E)=\frac{c_1(E)\cdot
H^{n-1}}{\text{rk}E}.
$$
\medskip\noindent
They have proved that for Higgs bundles 1
$\iff$ 2, $3 \implies 1,2$.
And also for $1,2 \implies 3$ but only for
$r=2$ using the Higgs Grassmannian,
but the problem is very complicated for
general rank.

\medskip\noindent
Second lecture. I think this might be a more
general
definition of the Hilbert polynomial:
$$
\sum_{k=0}^n(-1)^k \dim
H^k(X,\mathcal{E})=\chi(\mathcal{E})
=\int_X\text{ch}(\mathcal{E})\text{td}X
$$
And we also have:
$$
\chi(\mathcal{E}\otimes \mathcal{L}^*)
=\int_X (1+(1-g)w)(r+kr
\underbrace{c_1(\mathcal{L})}_{\text{deg}
\mathcal{L}\cdot
w}
+\text{deg}\mathcal{E}\cdot w.
$$
And then
$$
p_{\mathcal{E}}=k \text{deg} \mathcal{L}
+1-g+\mu(\mathcal{E}).
$$



\medskip\noindent
Now we discuss Higgs bundles.
Let $X$ be ac onnected complex projective
variety.

\begin{definition}
\label{definition-higgs-bundle}
A {\it Higgs bundle} on $X$
is a pair $\mathfrak{C}=(E,\phi)$ with
\begin{itemize}
\item $E$ a vector bundle on $X$
\item $\phi:$ a morphism (connection-like…)
\end{itemize}
\end{definition}

For the following see Nitsuze, FGA explained.
There following scheme is also representable
(I think this is the ``universal quotient'')

We have also consider the Grassmanization of
a vector bundle.
We obtained a short exact sequence
and tried to compute the tensor product with
the holomorphic differentials.
The idea was to pullback the grassman bundle
given a map $f:Y \to X$.
The result is not satisfactory, the pulled
back
thing may be strange. This is called the
{\it Higgs Grassmannian}.


Here's some notion also studied by Campana,
Demailly:
numerically flat bnef bundle $E$.
Then we have that
\begin{lemma}
\label{lemma-proof-hugo-bruzzo}
If $c_1=0$ then numerically flat bundle iff
semistable.
\end{lemma}

\medskip\noindent
This is the beginning of a talk by Alan Muniz
at Bandoleros 2025 titled
``Rank 2-bundles on $\mathbb{P}^3$ with odd
determinant''.
(Work in progress with A. Fontes.)

Problem: Classify rank-2 bundles on
$\mathbb{P}^3$.

We consider bundles $E$ with chern classes
$c_1=-1$, $c_2=2n$ and $c_3=0$.
Let $\mathcal{B}(-1,2n)$ be the moduli space
of $\mu$-stable rank 2-bundles.

\begin{itemize}
\item $\mathcal{B}(-1,2)$ is irreducible of
dimension 11
[Hartshorne-Sols '82, Manolache '81].

\item $\mathcal{B}(-1,4)$ has two components
[Barnica-Manolache '85].

\item $c_2\geq 6$, Tikhomirov (2014)
Fontes-Jardim (2023, 2024, 2025).
$\mathcal{B}(-1,6)$ has $\geq 4$ components.
\end{itemize}

\medskip\noindent
Now let's discuss Serre correspondence.

How can we describe these spaces?
Let $E$ be a vector bundle of rank 2
on $\mathbb{P}^3$ and $\sigma \in H^0(E(r))$.
Then $C=(\sigma_0)\subset \mathbb{P}^3$ is
a curve * equipped
with an ``extension class''
$\zeta \in H^0(\omega_C(4-2r-c_1))$.
In fact, this correspondence
is 1-1, and we call it Serre correspondence.
We have that $\text{deg}C=c_2(E(r))$ and
$p_a(C)=1=\frac{c_2(E(r))(c_1(E(r))-4)}{2}$
To pass from the curve to the bundle
and viceversa we use
$$
\xymatrix{
0\ar[r]
&\mathcal{O}(-r)\ar[r]
&E\ar[r]
&\mathcal{I}_C(c_1+r)\ar[r]
&0
}
$$


\medskip\noindent
Upshot about Deligne-Mumford stacks:
the sit below a stack via a map that is
étale.


















\medskip\noindent
Now let's discuss framed moduli.

\noindent
Framed instantons on the blow-up of
$\mathbb{P}^3$ at a point.

Let $X$ be a projective variety and fix a
polarization $H$.
Let $\mathcal{E}$ be a sheaf on $X$.
A {\it framed pair} is a
$(\mathcal{F},\alpha)$ where
$\mathcal{F}$ is a coherent sheaf on $X$
and $\alpha$ is a {\it framing datum}, i.e.
a map $\alpha:\mathcal{F} \to \mathcal{E}$.

Set
$$
\mathcal{E}(\alpha)
=\begin{cases}
	1\qquad &\alpha\neq 0 \\
	0\qquad &\text{otherwise} .
\end{cases}
$$
Define the Hilbert polynomial by
$$
P(k)=\chi(\mathcal{F} \otimes
\mathcal{O}(kH)).
$$
For fied $\delta \in \mathbb{Q}[t]$ rational
positive,
define
$$
P_{(\mathcal{F},\alpha)}=P_{\mathcal{F}}(k)-
\varepsilon(\alpha)\delta.
$$

We can also define induced framings from
subsheaves $\mathcal{F}' \subset
\mathcal{F}$, by asking that the inclusion
commutes
with the framing datum $\alpha$.
Then we put
$$
\varepsilon(\alpha')
=\begin{cases}
	1\qquad &\mathcal{F}' \subseteq Ker \alpha
	\\
	0\qquad &\text{otherwise} .
\end{cases}
$$
We say that $\mathcal{F}$ is {\it
$\delta$-(semi)stable}
if for any subpair $(\mathcal{F}',\alpha')$
with
$\text{rk}\mathcal{F}'=r'<r=\text{rk}
\mathcal{F}$
one has
$$
\frac{1}{r'}P_{(\mathcal{F}',\alpha')}(k)
\leq  \frac{1}{r}P_{(\mathcal{F},\alpha)}(k).
$$
If $\text{deg}(\delta)\leq \dim X$,
there are moduli functors
$\mathcal{M}^{st}(X,\mathcal{E},\delta)
\subseteq
\mathcal{M}^{\text{ss}}(X,\mathcal{E},\delta)
$.

\begin{remark}
\label{remark-degree-greater-than-dimension}
If $\text{deg} \delta>\dim X$
then $\alpha$ is injective or zero, and
$\mathcal{F}$ is semistable.
This gives quot schemes.
\end{remark}

\noindent
Huybrechts-L. In $\dim X\leq 2$
there exists a quasi-projective fine moduli
space
for $\delta \in \mathbb{Q}[t]$ and $p \in
\mathbb{Q}[k]$.

\medskip\noindent
Then there is a notion of good framing by
Bruzz-M,
depending on $D$, an effective divisor on
$X$.
They proved that having good framing is
enough
to obtain a fine quasi-projective moduli
space.

\begin{definition}
\label{definition-instanton}
A coherent sheaf $\mathcal{E}$ on
$X=\tilde{\mathbb{P}^3}$
is an {\it instanton} if it is
$\mu_{\mathcal{L}}$-semistable,
$c_1(\mathcal{E})=0$ and
$H^0(\mathcal{E})=0=H^0(Ec(-3,2))$,
$H^1(\mathcal{E}(-2,1))=0=H^2(\mathcal{E}(-2,
1))$
and
$H^3(\mathcal{E}(0,-1))=H^2(\mathcal{E}(-1,1)
)$.
\end{definition}

\noindent
Some examples are locally free sheaves of
rank 2 (via Harrshorne-Serre
correspondence),
torsion-free rank 2 constructed via
elementary
transformations.

\medskip\noindent
General result: for rank 2 instantons on Fano
threefolds
that have a plane $E$,
$\mathcal{E}|_{E}$ is $\mu$-semistable
if and only if $h^1(\mathcal{E}(-2,1) \otimes
\mathcal{O}(-E))=0$.




\section{Stability}
\label{section-stability}

\noindent
Stability helps us select the
building blocks of an abelian
category, every object in the
category can be reconstructed
from those.
Wall and chamber structure: inside a
compact set $K \subset
\text{Stab}(\mathcal{T})$,
only finitely many walls appear.
If $C$ is a connected component of
the complement of the walls in $K$,
then the stable objects are constant
on $C$.
Think of an ant walking in
$\text{Stab}(\mathcal{T})$ with a
backpack of objects and a detector:
when it crosses a wall and some object
stops being stable, the red light fires.
Stable objects change only when you
cross a wall.

\begin{center}
\setlength{\unitlength}{1mm}
\begin{picture}(98,78)(0,0)
\scriptsize

% axes
\linethickness{0.35pt}
\put(6,7){\vector(1,0){88}}
\put(88,7){\vector(0,1){68}}
\put(92,10){$\beta$}
\put(91,68){$\alpha$}

% Gieseker/PT chamber boundary
\linethickness{0.45pt}
\qbezier(18,7)(18,13)(24,15)
\qbezier(24,15)(31,16)(38,19)
\qbezier(38,19)(49,23)(60,22)

\qbezier(60,22)(57,27)(57,35)
\qbezier(57,35)(50,39)(47,45)
\qbezier(47,45)(41,53)(52,56)
\qbezier(52,56)(58,57)(68,58)
\qbezier(68,58)(75,62)(77,74)

% diagonal wall, blue in the photo
\linethickness{0.9pt}
\qbezier(22,73)(46,40)(70,7)

% semicircular wall, red in the photo
\linethickness{0.75pt}
\qbezier(45,7)(45,23)(60,22)
\qbezier(60,22)(75,23)(75,7)

% marked point and labels
\linethickness{0.4pt}
\put(59.35,21.35){\circle*{1.3}}
\put(61.5,23.5){$(\beta^\nu,\alpha^\nu)$}

\put(20,36){Gieseker chamber}
\put(47,59){PT chamber}
\put(39,42){$\theta^-_\nu$}

\end{picture}
\end{center}

Let $\mathcal{A}$ be an abelian
category with Grothendieck group
$K_0(\mathcal{A})$.

\begin{definition}
\label{definition-stability-function}
An additive homomorphism
$Z: K_0(\mathcal{A}) \to \mathbb{C}$
is a {\it stability function} if for all
$0 \neq E \in \mathcal{A}$,
$$
\text{Im}Z(E) \geq 0\qquad
\text{and}\qquad
\text{Re}Z(E) <0 \text{ if }
\text{Im}Z(E)=0.
$$
The {\it slope} of $E \in \mathcal{A}$
with respect to $Z$ is
$$
M(E)=
\begin{cases}
-\frac{\text{Re}Z(E)}{
\text{Im}Z(E)}\qquad &
\text{if }\text{Im}Z(E) \neq 0 \\
\infty\qquad &
\text{if }\text{Im}Z(E)=0.
\end{cases}
$$
A nonzero object $E \in \mathcal{A}$
is called {\it stable} if for all
proper non trivial subobjects
$F \subset E$, we have
$$
M(F) < M(E)
$$
and {\it semistable} if we put
$\leq $ instead.
\end{definition}

\begin{definition}
\label{definition-stability-condition}
A pair $(\mathcal{A},Z)$ is called
a {\it stability condition} if
\begin{enumerate}
\item $Z$ is a stability function, and
\item every $0 \neq E \in \mathcal{A}$
has a Harder-Narasimhan filtration.
\end{enumerate}
\end{definition}

\begin{lemma}
\label{lemma-harder-narasimhan-existence}
Let $Z:K_0(\mathcal{A})\to \mathbb{C}$
be a stability function.
Assume that
\begin{enumerate}
\item $\mathcal{A}$ is Noetherian, and
\item the image of $\text{Im}Z$
is discrete in $\mathbb{R}$.
\end{enumerate}
Then every $0 \neq E \in \mathcal{A}$
has a Harder-Narasimhan filtration
with respect to $Z$.
\end{lemma}

\noindent
Let $\mathcal{D}$ be a triangulated
category. Fix a finite rank lattice
$\Lambda$, a surjective homomorphism
$v:K_0(\mathcal{D}) \to \Lambda$
and a norm $\|\cdot\|$ in $\Lambda$.

\begin{definition}
\label{definition-bridgeland-stability}
A {\it Bridgeland stability condition}
$\mathcal{D}$ is a pair $(\mathcal{A},Z)$
where $\mathcal{A}$ is the heart of a
bounded $t$-structure on $\mathcal{D}$,
such that
\begin{enumerate}
\item $(\mathcal{A},Z)$ is a stability
condition,
\item the following support property
holds:
$$
\text{inf}\left\{
\frac{|Z(E)|}{
\|v(E)\|}
:0 \neq  E \in \mathcal{A}
\text{semistable}\right\}
>0.
$$
\end{enumerate}
\end{definition}

\begin{example}
\label{example-stability-curve}
Let $X$ be a smooth projective
curve and $\mathcal{D}=D^b(X)$.
Take $\mathcal{A}=\Cohstack(X)$ as the
heart of the standard
$t$-structure.
Let
$$
Z(E)=-\deg(E)+i\,\mathrm{rk}(E).
$$
Then $\sigma$-stability coincides
with slope stability.
\end{example}

\noindent
Let $Stab^{Br}_\Lambda(\mathcal{D})$
be the set of Bridgeland stability
conditions and let
$Stab_\Lambda(\mathcal{D})$ be the set
of stability conditions on $\mathcal{D}$
with respect to $(\Lambda,v)$.

\begin{theorem}[Bridgeland]
\label{theorem-bridgeland}
The natural map
$$
Stab_\Lambda(\mathcal{D})
\to \Hom(\Lambda,\mathbb{C})
$$
is a local homeomorphism.
$Stab_\Lambda(\mathcal{D})$ is a complex
manifold and
its dimension is the rank of $\Lambda$.
In particular,
$$
Stab_\Lambda^{Br}(\mathcal{D})
\subset
Stab(\mathcal{D}).
$$
\end{theorem}

\noindent
Let $(\mathcal{A},Z)$ be a stability
condition and $M$ be the slope with
respect to $Z$.
For each $\beta \in \mathbb{R}$
we define the following pair of
subcategories on $\mathcal{A}$:
$$
\begin{aligned}
\mathcal{T}^\beta(A,Z)&=\mathcal{T}^\beta
=\{\mathcal{E} \in \mathcal{A}
: M(\mathcal{G})>\beta
\text{ whenever }
\mathcal{E} \to\!\!\!\!\!\to
\mathcal{G}\}\\
\mathcal{F}^\beta(\mathcal{A},Z)
&=\mathcal{F}^\beta
=\{\mathcal{E} \in \mathcal{A}
:M(\mathcal{F})\leq \beta
\text{ for every }
0 \neq  \mathcal{F} \subset \mathcal{E}\}.
\end{aligned}
$$

\begin{definition}
\label{definition-tilt}
The {\it tilt}
$\mathcal{A}^*
=\left\langle{}\mathcal{F}[1],\mathcal{T}\right\rangle{}
\subset\mathcal{D}$
of $\mathcal{A}$ with respect to
$(\mathcal{T}^\beta,\mathcal{F}^\beta)$
is the smallest extension closed full
subcategory of $\mathcal{D}$ containing
$\mathcal{T}^\beta$ and
$\mathcal{F}^\beta[1]$.
\end{definition}

\noindent
$\mathcal{A}^*$ is the heart of a bounded
$t$-structure on $\mathcal{D}$.

\begin{definition}
\label{definition-gieseker-stability}
Let $\mathcal{E}$ be a torsion free
sheaf.
We say that $\mathcal{E}$ is
{\it Gieseker semistable} if
for every subsheaf $F \subset \mathcal{E}$
with $0 < \mathrm{rk}(F) <
\mathrm{rk}(\mathcal{E})$,
$$
\frac{P_F(m)}{\mathrm{rk}(F)}
\leq
\frac{P_{\mathcal{E}}(m)}{
\mathrm{rk}(\mathcal{E})}
\qquad \text{for } m \gg 0.
$$
It is {\it Gieseker stable}
if we replace $\leq$ by $<$.
Here $P_{\mathcal{E}}(m)$ is the
Hilbert polynomial of $\mathcal{E}$.
\end{definition}

\begin{definition}
\label{definition-2-gieseker-stability}
Let $\mathcal{E}$ be a torsion free
sheaf.
We say that $\mathcal{E}$ is
{\it 2-Gieseker semistable} if
for every subsheaf $F \subset \mathcal{E}$
with $0 < \mathrm{rk}(F) <
\mathrm{rk}(\mathcal{E})$,
$$
\frac{P_F(m)-\chi(F)}{\mathrm{rk}(F)}
\leq
\frac{P_{\mathcal{E}}(m)-\chi(\mathcal{E})}
{\mathrm{rk}(\mathcal{E})}
\qquad \text{for } m \gg 0.
$$
It is {\it 2-Gieseker stable}
if we replace $\leq$ by $<$.
\end{definition}

\noindent
The Gieseker chamber is the region
where this is the relevant stability.
The PT chamber is the neighboring region
where the PT condition is the relevant
one.

\begin{definition}[PT stability]
\label{definition-pt-stability}
\end{definition}

\begin{theorem}
\label{theorem-pt-stable-triple}
Assume $\gcd(\mathrm{ch}_0,
\mathrm{ch}_1,\mathrm{ch}_2)=1$.
Then every PT stable object comes from
a stable triple $(B,P,\eta)$:
\begin{enumerate}
\item $B$ is a 2-Gieseker
      semistable sheaf with
      $\mathrm{hd}(B)\geq 1$,
\item $P$ is a $0$-dimensional sheaf,
\item $\eta \in
      \Hom_{D^b(X)}(P,B[2])=
      \Ext^2(P,B)$ is nontrivial.
\end{enumerate}
Moreover,
$$
A_\eta := \operatorname{cone}(\eta)[-1].
$$
\end{theorem}

\begin{remark}
Stable triples generalize rank $1$
PT stable points.
\end{remark}

\begin{theorem}[Jardim et al.]
\label{theorem-single-wall-crossing}
There is a single wall crossing
that takes the DT moduli space to
the PT moduli space.
Equivalently, the Gieseker and PT
chambers are separated by one wall.
\end{theorem}

\begin{proof}
\begin{enumerate}
\item Understand the stable
objects at the boundary,
labelled $\Theta_v^-$
in the picture.
\item Study the asymptotic stable
objects, that is,
what you find if you go to
infinity.
\item Identify what the PT
chamber is.
This may also be an asymptotic
analysis.
\item Show that there are no other
walls between the $\Theta$ wall
and infinity on the Gieseker side,
and no other walls in the
direction of the PT side.
\end{enumerate}
\end{proof}

\medskip\noindent
History.
\begin{itemize}
\item Pandharipande and Thomas
(2009) suggested that DT/PT
correspondence could be a wall
crossing phenomenon in
Bridgeland stability space.
\item Bayer (2009) established
that the rank $1$ case is a
wall crossing for polynomial
stability conditions.
\item Toda.
\item \dots
\end{itemize}

\medskip\noindent
Case studies.

\noindent
Let $X$ be a smooth threefold with
Picard rank $1$.
\begin{itemize}
\item Collapsing wall.
For
$v=(r,0,0,-n)$,
$$
\mathcal{G}(v)=
\mathrm{Quot}^n(
\mathcal{O}_X^{\oplus r})/\!/
\mathrm{GL}(r),
$$
while $\mathcal{T}(v)=\emptyset$.
\item Fake wall.
For $v=(v_0,v_1,v_2,v_{3,\max})$,
$\mathcal{G}(v)=\mathcal{T}(v)$.
This is a wall that changes nothing.
\end{itemize}

\noindent
Now set $X=\mathbb{P}^3$.
\begin{itemize}
\item For
$v=(2,-1,-\tfrac{1}{2},-\tfrac{1}{6})$,
$\mathcal{G}(v)$ is birational to a
$\mathbb{P}^1$-fibration over
$\mathbb{P}^3\times\mathbb{P}^3$,
while $\mathcal{T}(v)\cong\mathbb{P}^3$.
The wall crossing induces a
retraction from the $6$-dimensional
space $\mathcal{G}(v)$ to
$\mathbb{P}^3\times\mathbb{P}^3$,
and
$\mathcal{T}(v)\simeq \mathbb{P}^3$.
\end{itemize}

\medskip\noindent
There is a map
$$
\xymatrix{
\mathcal{G}(-v)\ar[dr]^\gamma
&&
\mathcal{T}(v)\ar[dl]^\tau
\\
&
\mathcal{M}_{\bar\beta,\bar\alpha,s}(v)
&
}
$$
Pavel and Tajakka show that
the PT moduli stack and the
Bridgeland moduli stack at the wall
are projective schemes.
So the map above is a map of
schemes.
\noindent
Formula of the map missing.

\medskip\noindent
Some questions:
\begin{itemize}
\item (Existence.) Given a projective
variety $X$, are there stability
conditions on $D^\text{b}(X)$?
Positive answer for surfaces,
Fano 3folds with Picard rank 1,
abelian 3folds, quintic 3folds,
products of curves, $\mathbb{P}^n$,
\dots
\item (Moduli spaces.) Is
$M_\sigma(v)$ a projective scheme?
Cannot use usual git techniques to
study. Yes for K3, $\mathbb{P}^2$,
and $\mathbb{P}^1\times \mathbb{P}^1$.
The parameter space is a stack
(ref. missing, involves Toda).
\item What do stable objects
look like?
\end{itemize}

\medskip\noindent
\medskip\noindent
Following are very sketchy notes picked
up along the way; don't hessitate to skip

\begin{enumerate}
\item Stable objects in an abelian category
are the ``building blocks":
we can reconstruct the whole category from
them.
\item  An abelian subcategory (hart)
$\subset$ a
triangulated
\item Stability defined via stability f
unction on $\mathcal{A}$.
\item (Example.) $\mathcal{A}=\text{Coh}X$ is
heart
of $D^b(X)$ with stability function.
\item Bridgeland Stabl:= the stability
conditions are a complex manifold
of complex codimension $\text{rk}\Lambda$:
$$
\mathcal{Z}:\text{Stab}(\mathcal{T})
\longrightarrow
\text{Hom}(\Lambda,\mathbb{C})
$$
\item There's a chamber structure;
moduli space changes across chambers.
\item Picture: blue + black are walls.
Q. What are $\beta$ and $\alpha$?
\item thm: bridgeland stable
= gieseker stable ?
\item Q. slope stability
= bridgeland stability?
A. Not always.
\item DT/PT correspondance:
only one wall between PT and G chambers
\medskip
\item Polynomial stability function.
This is an asymptotic version of BS.
\item There are some $\rho$'s.
Arrangements of $\rho_i$ are polynomial
stability conditions on a threefold.
\item Pata-Thomas introduced stability
for rank 1 objects.
Bayer compares the---wall. Q.
Same for Bridge S---only one wall?
\medskip
\item Def. A {\it stable triple}: when
$\text{gcd}(\text{ch}_0,\text{ch}_2,\text{ch}
_3)=1$,
every PT stable object
comes from three conditions (missing).
\item What happens when
you cross the blue wall?
Both $\mathcal{G}$ and $\mathcal{T}$ are
projective.
\medskip
\item For $X$ smooth threefold with
$\text{rk}\text{Pic}=1$, $\mathcal{G}(v)$,
$v=(r,0,0,-n)$, $\mathcal{G}(v)$ is a known
sheaf object and
$\mathcal{T}=\emptyset$.
\item $X$ sm 3 rkpic1, Fake wall;
$\mathcal{G}=\mathcal{T}$.
\item (Extra.) red circle is a wall for a
weaker form of stability.
\end{enumerate}

\begin{definition}[Talk at impa]
\label{definition-slope-stability}
A rank-2 sheaf $\mathcal{F}$ is
{\it semistable (stable)} if
$H^{0}(\mathcal{F}(t))=0$ for $-t \geq(>)
\frac{c_1F}{2}$
\end{definition}

\noindent
Compare with

\begin{definition}[moduli-curves.tex]
\label{definition-semistable}
Let $f : X \to S$ be a family of curves.
We say $f$ is a {\it semistable
family of curves} if
\begin{enumerate}
\item $X \to S$ is a prestable family of
curves,
and
\item $X_s$ has genus $\geq 1$ and
does not have a rational tail for all $s \in
S$.
\end{enumerate}
\end{definition}

\section{Instantons}
\label{section-instantons}


\begin{itemize}
\item The twistor diagram.
$$
\xymatrix{
&  \ell_\infty \ar@{.>}[dr]\ar[dl]\\
\mathbb{C}P^{2} \ar@{^{(}->}[rr]& &
\mathbb{C}P^{3}\ar[dd]
_{\substack{\text{twistor}
\\ \text{map}}}\\ \\
\mathbb{C}^2\ar@{^{(}->}[u]
\mathbb{R}^4 \ar[u]\ar[rr]& &S^4
}
$$

\item Instantons are solutions to some
Yang-Mills solution.
\item Expository paper
of Donaldson arXiv:2205.08639 \item ADHM
construction, 1978. The first appearance of
Algebraic Geometry in Mathematical Physics.
See Hitchin-Kobayashi correspondence.
\item Donaldson: ``unashamedly
computational''.
\item Expository paper by Simon.
\item Take bundle $(E,\nabla)$ with
an anti-self dual
(ADS) connection on $S^4$ and pullback to
$\mathbb{C}P^{3}$ via the twistor map

$$
\tau[x:y:z:w]=[x+jy:z+jw] \qquad \text{note:
$\tau^{-1}(p)=\mathbb{C}P^{1}$}
$$
Then:
\begin{itemize}
\item Restriction to fibres are trivial.
\item Invariant under anti holomorphic
involution (check this formula!)
$[x:y:z:w] \mapsto [-y:x:-w:z]$.
\item $\mathcal{E}$ is also an instanton
bundle.
\item Penrose Transform:
 $H^1(\mathcal{E}(-2)) \cong \Ker \Delta=0$
 where
$\Delta$ is a Laplacian.
\item Definition of instanton sheaf
on $\mathbb{P}^3$ via $c_1(E)=0$ and some
vanishing of cohomologies.
\item Passage from
[differential equations? algebraic geometry?]
to linear algebra: via {\it monads}.
The point is that instanton sheaf is
equivalent to  ``$E$ being the cohomology
of a linear monad''; theorem by Horps
in the 60's, and is the main tool used
by ADHM. Indeed, ADHM equations come from
the cohomology sequences of the so-called
monads.
\item Mathamatical inst bund:=
locally free instanton sheaves.
\end{itemize}
\end{itemize}

\noindent
{\bf Properties.}
\begin{itemize}
\item The only instanton of rank 1 on
$\mathbb{P}^3$ is
the structure sheaf.
\item non-trivial rank 2 locally free
instanton sheaf is  $\mu_0$ stable
\item double dual is locally free and
also instanton
\item non-trivial rank 2 instanton sheaf
is Fieseker stable
therefore it makes sense to define
moduli space of instanton sheaves as
an open subset of
$\mathcal{M}(c)=\mathcal{G}(k,0,2,0)$.
\end{itemize}

Then studied the irreducibility (Tikhomirov)
and smoothness (Jardim-Verbitsky, 2014. Uses
``3rd hyperkähler quotient'') of
$\mathcal{I}(c)$, the moduli space of rank 2
locally free instanton sheaves of charge $c.$
But nobody likes this results; want new
proofs.

In contrast, $\mathcal{M}(c)$ of rank
2-instanton sheaves of charge $c$ is not
irreducible in general! $\mathcal{M}(1)$ and
$\mathcal{M}(2)$ are irreducible,
$\mathcal{M}(3)$ has exactly 2 irreducible
components of dimension 21; $\mathcal{M}(4)$
has 4 irredicuble components: the locally
free is irreducible, and the other 3 that
intersect the closure of the locally free,
$\overline{\mathcal{I}(c)}$. 3 components of
dimension 29 and one of dimension 32

\begin{remark}
\label{remark-gieseker}
In general the $\mu$ moduli space is
not projective, but the Gieseker is.
\end{remark}

\noindent
{\bf Is $\mathcal{M}(c)$ connected?} Use
$\mathbb{C}^*$ action. True for $c \leq 4$;
every component intersects
$\overline{\mathcal{I}(c)}$ in this range!

\begin{definition}[Elementary transformation]
\label{definition-elementary-transformation}
$F$ of rank 2 locally free instanton,
$Q$ of rank 0 instanton with 1 dimensional
sheaf $h^p Q(-2)$ for $p = 0,1$.
So in this conditions if we have an
epimorphism
 $$
F \overset{\varphi,\text{ surj}}{\to}\} Q
$$
we get that $\Ker \varphi$ is an instanton.
\end{definition}

So  we might be interested in
$$
E \hookrightarrow E^{*
*}\xrightarrow{\text{surj.}} Q_E
$$
Let's have a look at
$\mathcal{M}(3)=\overline{\mathcal{I}(3)}\cup
\overline{C(0,1,3}$. Consider a generic line
bundle of degree 0 on a planar cubic (which
is encoded in that we have intersection of
something of cimension 1 and something of
simension 3), $L inn \text{Pic}^0(C)$ so we
have an epimorphism
$$
\xymatrix{
0\ar[r]&E\ar[r]&\mathcal{O}_{\mathbb{P}^3}
\oplus
\mathcal{O}_{\mathbb{P}^3}\ar[r]&(i_*
L)\ar[r]&0
}
$$
yielding a bundle $E$ as previously outlined.

So we are studying the components via
pushforwards of sheaves on complete
intersection curves inside $\mathbb{P}^3$

``when we do alementary transformation of
rational (not rationl?) we get
something on the boundary of locally free.''

\medskip\noindent
{\bf Perverse instanton sheaves.} Like an
instnaton sheaf in monad description but with
some restrictions on the cohomologies. This
leads to definition of {\it $0$-rimensional
instanton}, a perverse instanton such that
$\mathcal{H}^0=0$.

\medskip\noindent
{\bf Instantons and quivers.}

\medskip\noindent
{\bf Framed instantons.} Fix a line $j:L\to
\mathbb{P}^3$ and $E$ a perverse
sheaf; a {\it framing} at $L$ is an
isomorphis [missing].

Apply GIT to the ADHM data to construct a
moduli space
$$
\mathcal{P}(r,c)=\mathcal{V}(r,c)^{\text{st}}
/\!/\text{GL}(V_c)
$$
and it will follow that $\mathbb{P}(?,?)$ is
connected.

\medskip\noindent

Considerations on quaternionic spaces lead to
generalization of what has been discussed so
far to higher dimensions. I.e. an {\it
instanton sheaf} on $\mathbb{P}^n$ is… [other
cohomological conditions] sheaf

Why are instantons interesting? They are the
simplest; may provide examples for Bridgeland
stability.

In Kuznetsov (2012) and Faenzi (2014)
introduced {\it rank 2 instanton bundles on
Fano 3-folds}. An {\it instanton bundle} on
$X$ is a $\mu$-stable … and some Chern class
is called the {\it charge}. An {\it instanton
sheaf} (introduced by Marcos-Gaia) is ….

 ``Since we are imposing $\mu$-stability on
the defintion we can consider the moduli
$\mathcal{I}_X(c) \subset
\mathcal{G}_X(2,-r_X,c,0)$''.

There's also monad representations as an
ingredient.

\medskip\noindent
Here are two questions that invite us to join
the instanton fever:

\medskip\noindent
{\bf Task 1.} Construct rank 2-instanton
sheaves that do not deform into locally free
ones, and obtain the new irreducible
components of $\mathcal{G}(2,-r_X,c,0)$.

\medskip\noindent
{\bf Task 2.} Nonlocally instanton sheaves
that can be deformed into non-locally free
ones: the {\it instanton boundary}
$\overline{\mathcal{I}(c)}/\mathcal{I}(X)
$[formula
right?]


\medskip\noindent
{\bf Recipe to construct your own instanton.}
\begin{enumerate}
\item (Make a bunch of instantons.) Find an
appropriate curve to
use Serre correspondence to find some
rank 2 instanton sheaf:
$$
\xymatrix{
0\ar[r]&\mathcal{O}_{\mathbb{P}^3}(-1)\ar[r]&
\underbrace{E}_{\exists
}\ar[r]&\mathcal{I}_{\sqcup
\text{lines}}\ar[r]&0
}
$$
where $E$ is a locally free instanton
of charge = the number of lines $-1$.

\item (Is your family of instantons generic?)
You look at the $\text{Ext}$s.
``Therefore, the family of instantons only
defines a locally closed subset
within a generically smooth irreducible
component of $\mathcal{G}$''.

\item ``Find suitable rank 0 instanton
sheaves
 to perform an elementary
transformations on the examples obtained in
Step 1.'' But they are non-locally
free.

\item Now I have my instantons, I know they
are locally free: but how to prove
that the elementary transformations deform
to locally free ones? Looks like
the challenge is to prove that the
deformation is locally free.
\end{enumerate}

In the papers by the group there are several
particular cases when the
deformations are locally free. But they don't
have a general result that would
work for 3-folds.

\medskip\noindent
{\bf What you need to call a thing an
instanton.}
\begin{itemize}
\item Minimal cohomology possible;
try to kill as much cohomology as you can.
\item Fixing $c_1$ (which may determine other
Chern classes).
\item Some stability condition
like $\mu$-stability or quiver stability.
Here is an example that is not
$\mu$-semistable:
$T\mathbb{P}^3(-1)\oplus\Omega_{\mathbb{P}^3}
(1)$
\item Whenever possible,
look for a monadic representation.
(The monadic
representation comes from ADHM ---
the beginnings of this theory.
And it's still here!)
\end{itemize}

\section{Jardim-Muniz}
\label{section-jardim-muniz}

\noindent
Reflexive sheaves
are vector bundles except for a small
locus.

\begin{definition}
\label{definition-reflexive}
Let $X$ be an integral locally Noetherian
scheme. Let $\mathcal{F}$
be a coherent $\mathcal{O}_X$-module.
The {\it reflexive hull}
of $\mathcal{F}$ is the
$\mathcal{O}_X$-module
$$
\mathcal{F}^{**} = \SheafHom_{\mathcal{O}_X}(
\SheafHom_{\mathcal{O}_X}(\mathcal{F},
\mathcal{O}_X), \mathcal{O}_X)
$$
We say $\mathcal{F}$ is {\it reflexive}
if the natural map
$j : \mathcal{F} \longrightarrow
\mathcal{F}^{**}$
is an isomorphism.
\end{definition}

\begin{lemma}
\label{lemma-reflexive-torsion-free}
Let $X$ be an integral locally Noetherian
scheme. Let $\mathcal{F}$
be a coherent $\mathcal{O}_X$-module.
\begin{enumerate}
\item If $\mathcal{F}$ is reflexive, then
$\mathcal{F}$
is torsion free.
\item The map $j : \mathcal{F}
\longrightarrow \mathcal{F}^{**}$
is injective if and only if $\mathcal{F}$
is torsion free.
\end{enumerate}
\end{lemma}

\begin{remark}[Talk at IMPA, 11 June]
\label{remark-reflexive-talk}
Torsion could also be defined so that the
sheaf can inject onto its dual. In this talk
we discussed the moduli space of
reflexive/torsion-free sheaves, which turned
out to be parametrized by $c_1, c_2$ and
$c_3$. This was denoted by $R(c_1,c_2,c_3)$.
Actually I think it may have been Manolache
that proved the existence of this moduli
space.

Alan Muniz

Nesta palestra discutiremos a classificação
de feixes reflexivos de posto dois e seus
espaços de módulos. Apresentaremos algumas
ferramentas básicas usadas na construção e
determinação de tais feixes. Aplicaremos
estas técnicas para o caso de feixes com
segunda classe de Chern igual a quatro,
obtido recentemente em colaboração com Marcos
Jardim.
\end{remark}




\medskip\noindent
Now let's discuss
distributions on manifolds.

Here's the abstract from a talk by Marcos
Jardim at Geometric Structures:

``I will revise the work done over the past
10 years with various collaborators on
distributions and foliations on 3-folds,
especially on the projective space, with a
focus on properties of the tangent sheaf and
singular scheme."

Here are two key ideas: if the distribution
is codimension 1 we can write:
$$
\xymatrix{
0\ar[r]&F\ar[r]&TX\ar[r]^{\omega}&I_Z \otimes
L\ar[r]&0
}
$$
where $L$ is a line bundle and
$\omega \in H^{0}(\Omega_X \otimes L)$, and
$Z=\{p:\omega(p)=0\}$.

When codimension is 2 then $\mathcal{D}$ is
given by
a holomorphic vector field
$\nu$: $T_p=\left\langle{}\nu(p)\right\rangle{}$.

It can be encoded as an exact sequence
$$
\xymatrix{
	0\ar[r]&L\ar[r]^{\nu}&TX\ar[r]&N\ar[r]&0
}
$$
where $L$ is a line bundle
and $\nu \in H^{0}(TX \otimes L^\vee)$;
  $Z=\{p | \nu(p) = 0\}$.

\begin{remark}
\label{remark-stauration}
Saturation means that $Z \subset X$ is a
union
of curves and points.
\end{remark}

And again, distributions are parametrized by
Chern classes.

Two interesting open questions:
\begin{enumerate}
\item {\bf Conjeture.} if $\mathcal{D}$ is a
codimension 1 foliation of degree
$d$ on $\mathbb{P}^3$, then $c_2(F)\leq
d^2-d+1$ and
bound is attained
by rational foliations of type $(1,d+1)$.
(True for $d \leq 2$.)
\item {\bf Conjecture (with Pepe Seade).}
$\mathcal{D}$ is a codimension 1
foliation on a smooth projective 3-fold,
then $\text{Sing}\mathcal{D}$ is connected.
\end{enumerate}

\begin{theorem}[Jardim-Muniz]
\label{theorem-jardim-muniz}
Conditions on Chern classes used to
understand moduli space $R(c_1,c_2,c_3)$.
$c_2=4$ gives (?). For $c_3\leq 6$, possible
``spectrum" exists…
\end{theorem}









\section{Campana}
\label{section-campana}

A \textit{special} variety should be thought
of
as the farthest notion possible from being of
general type.

\begin{definition}[Campana]
$X$ complex projective.
$X$ is \textit{special} if $\forall i \leq  p
\leq \dim X$,
for all line bundle $L \subset \Omega^P_X$,
$\kappa(L)<P$ (Kodaira dimension).
\end{definition}

\noindent
A conjecture by Campana says that
$X$ is special if and only if there exists
$f:\mathbb{C}\to X$
Zariski dense.
Compare this with Kobayashi hyperbolicity!
There are no lines on …?

The main tool for proving this kind of
``Bloch-Ochai" statements (i.e.
something of the kind: if $h^0(\Omega_X)>
\dim X$
then there are no nonconstant
$f:\mathbb{C}\to X$
Zariski dense)
is by passing through the Albanese variety.

In a talk [fill details]
we study a theorem of this kind but with the
bound
given by $q^+(X)$ defined as the supremum of
the $h^0(\hat{X},\pi^*\Omega_{X}$
for $\hat{X}\to X$ some (étale) cover.
Can we prove a BO stement with the condition
$q^+(X)>\dim X$?

Notion of special comes in: indeed
 $$
q^+(X)>\dim X \implies X \text{ not special
}\implies
\not\exists f:\mathbb{C}\to X \text{Zariski
dense}.
$$

The following theorem is not at all easy to
prove:
\begin{theorem}[Campana]
\label{theorem-campana}
$\hat{X}\xrightarrow{\text{étale} }X$, then
$\hat{X}$ is special iff $X$ is special.
\end{theorem}

\noindent
My question: why does the abelianess of
the Albanese variety is important in this
construction?





% --- End algebraic-geometry.tex ---

% --- Begin complex-geometry.tex ---
\chapter{Complex Geometry}
\label{chapter-complex-geometry}





\section{Complex structures on vector spaces}
\label{section-complex-structure}

\noindent
Let's discuss complex structures at the level
of linear algebra.

\begin{definition}
\label{definition-compl-struc-vecto-space}
An {\it almost complex structure} on a vector
space $V$
is an endomorphism $I \in \text{End}(V)$ such
that $I^2 = -\text{Id}_V$.
\end{definition}

\begin{lemma}
\label{lemma-almos-compl-struc-then-compl}
A vector space endowed with an almost complex
structure
admits a natural structure of a complex
vector space.
\end{lemma}

\begin{proof}
$(a+ib)v:=av + b I(v)$.
\end{proof}

\noindent
Conversely, given a complex vector space,
multiplication by $i$ produces an almost
complex structure.
Thus the two notions become the same.

\medskip\noindent
Now we introduce the complexification of a
real vector space.
Let $V_{\mathbb{C}}:=V
\otimes_{\mathbb{R}}\mathbb{C}$.
This means that elements of $V_{\mathbb{C}}$
are finite sums of elements of the form $v
\otimes (a+iv)$
for $a,b \in \mathbb{R}$ and $i = \sqrt{-1}$.

Notice that although $(V,I)$ is also a
complex vector space,
it is a different one from $V_\mathbb{C}$.
The complexification $V_\mathbb{C}$
is much more interesting since it gives us
the $(p,q)$-decomosition
when we take the exterior algebra.

Further, $V_\mathbb{C}=V \oplus i V$.
We define {\it conjugation} of vectors in
$V_{\mathbb{C}}$ by
$\bar{v}=\overline{v \otimes \lambda} := v
\otimes \bar{\lambda}$

This means that $\overline{V}$ is just $V$
but
with the scalar product $\lambda
v:=\bar{\lambda}v$.
for any $v \in V_{\mathbb{C}}:=V
\otimes_{\mathbb{R}}\mathbb{C}$.
This makes the equation
$\overline{\Lambda^{1,0}(V)}=\Lambda^{0,1}(V)
$,
which appears in the next proposition,
obvious.
(See \cite[p. 25]{huc}.)

\begin{lemma}
\label{lemma-i-eigenspace-decomposition}
We have a decomposition
$$
V_\mathbb{C}=V^{1,0}\oplus V^{0,1}
$$
given by the eigenspaces of $I$.
$\overline{\Lambda^{1,0}}=\Lambda^{0,1}$
holds by definition of conjugation.
\end{lemma}


Now let's discuss the $p,q$ decomposition.
The upshot is the following:
$$
T^2(V \oplus  W)=
T^2(V)\otimes T^0(W)
\oplus T^1(V)\otimes T^1(W)
\oplus T^0(V)\otimes T^2(W).
$$
Explicitly, the decomposition of a general
element on the LHS is
$$
\begin{aligned}
(v_1,w_1)\otimes(v_2,w_2)
&=(v_1,0)\otimes(v_2,w_2)+(0,w_1)\otimes(v_2,
w_2)\\
&=(v_1,0)\otimes(v_2,0)+(v_1,0)\otimes(0,w_2)
\\
&+(0,w_1)\otimes(v_2,0)+(0,w_1)\otimes(0,w_2)
\end{aligned}
$$
which is valid because the tensor product of
a vector space
$T^2X$ is just the pairs $x \otimes y$
subject
to relations that $(x_1+x_2)\otimes y=x_1
\otimes y+x_2\otimes y$
and addition for $y$ and scalar product.
So we just identify elements
$$
\begin{aligned}
(v_1,0)\otimes(v_2,0)
&\leftrightarrow
v_1 \otimes v_2 \in T^2(V) \otimes
T^0(W):=T^{2,0}(V \oplus W)\\
(v_1,0)\otimes(0,w_2)
&\leftrightarrow
v_1 \otimes w_2 \in T^1(V) \otimes T^1(W):=
T^{1,1}(V\oplus W)\\
(0,w_1)\otimes(v_2,0)
&\leftrightarrow
w_1 \otimes v_2 \in T^1(W) \otimes  T^1(V)
\cong T^1(V) \otimes  T^1(W):=T^{1,1}(V\oplus
W)\\
(0,w_1)\otimes(0,w_2)
&\leftrightarrow
w_1\otimes w_2 \in T^0(V) \otimes
T^2(W):=T^{0,2}(V \oplus  W).
\end{aligned}
$$

\noindent
We have proved that any element in $T^2(V
\oplus W)$
can be written as an element in the vector
space
$T^2(V)\otimes T^0(W)
+ T^1(V)\otimes T^1(W)
+ T^0(V)\otimes T^2(W)$.
Since the intersection of the factors is
disjoint,
the sum can be taken direct.

\medskip\noindent
Notice that we have two copies of
$T^{1,1}(V\oplus W)$.
By the binomial theorem and isomorphisms
of the kind $A \otimes B \cong B \otimes A$,
we readily see that
 $$
T^k(V \oplus W)
\cong
\bigoplus_{p+q=k}\binom{k}{p}T^p(V)\otimes
T^q(W).
$$
Of course the binomial coefficient matters
little
when the underlying field is $\mathbb{C}$.

The next step is to take the quotient on the
LHS
by the relation $v \otimes v=0$,
i.e., the ideal generated by $v \otimes v$.
When we compute $(v_1,w_1)\otimes(v_1,w_1)$
as above,
this relation becomes $v_1 \otimes v_1=0$
and $w_1\otimes w_2=0$ on the first and last
components.
If we want the middle component to vanish as
well
we must ask that the isomorphism
$T^1(W) \otimes  T^1(V)\cong T^1(V) \otimes
T^1(W)$
exchanges sign. Then we obtain that
$$
\Lambda^2(V \oplus W)
=\Lambda^{2,0}(V\oplus W)
\oplus\Lambda^{1,1}(V\oplus W)
\oplus\Lambda^{0,2}(V \oplus W)
$$
by defining $\Lambda^{p,q}(V \oplus
W):=\Lambda^p(V)\otimes \Lambda^q(W)$.

To make sure that this generalizes to all
$p,q$
write, e.g. $p+q=3$,
$$
\begin{aligned}
(v_1,w_1)\otimes(v_1,w_1)\otimes(v_1,w_1)
&=(v_1,0)\otimes(v_1,0)\otimes(v_1,0)\\
&+(v_1,0)\otimes(v_1,0)\otimes(0,w_1)\\
&+…
\end{aligned}
$$

\noindent
so that imposing the relation $v \otimes v$
on $T^3(V \oplus W)$
implies imposing it on every factor.
[Further arguments needed here… but looks
like just boring computations…]

\medskip\noindent
The main result of Hodge theory is that
when we take $T^{1,0}(M)\oplus T^{0,1}(M)$,
the eigenspaces of the complex structure,
and pass to complex singular cohomology of a
Kähler manifold we get
\begin{equation}
\label{equation-hodge-theorem}
H^{k}(M,\mathbb{C})
=\bigoplus_{p+q=k}H^{p,q}(M)
\end{equation}

\noindent
where $H^{p,q}(M)$ is the set of differential
forms
represented by a closed form in
$\Lambda^{p,q}(M)$.

\medskip\noindent
It also happens that
$\overline{\Lambda^{p,q}(V)}=\Lambda^{q,p}(V)
$,
which is due to the fact that
$\overline{\Lambda^{1,0}}=\Lambda^{0,1}$.


\begin{exercise}
\label{exercise-1,1-forms-are-i-invariant}
Show that $(1,1)$-forms are invariant under
the action of
the complex structure $I$.
And conversely, a 2-form is of type $(1,1)$
only if it is $I$-invariant.
\end{exercise}

\begin{proof}
Let $\alpha\wedge\beta \in \Lambda^{1,1}(M)$.
Recall that
$\Lambda^{1,1}(M)=\Lambda^{1,0}(M)
\otimes\Lambda^{0,1}(M)$,
i.e. $\alpha \in \Lambda^{1,0}(M)$ and
$\beta \in \Lambda^{0,1}(M)$.
For any vectors $v,w$ tangent to some point
$p \in M$,
$$
\begin{aligned}
I(\alpha\wedge\beta)(v,w)
&=(\alpha\wedge\beta)(Iv,Iw)\\
&=\alpha(Iv)\beta(Iw)-\beta(Iv)\alpha(Iw)\\
&=-\sqrt{-1}\sqrt{-1}(\alpha(v)\beta(w)-
\beta(v)\alpha(w))\\
&=\alpha(v)\beta(w)-\beta(v)\alpha(w)\\
&=(\alpha\wedge\beta)(v,w).
\end{aligned}
$$

\noindent
The converse follows easily by the direct sum
splitting of $\Lambda^2(M)$ and noting that
a $(2,0)$ or $(0,2)$-form $\gamma$
satisfies $I\cdot\gamma(v,w)=-\gamma(v,w)$
in an analogous computation.
\end{proof}

\noindent
In particular this implies that $\omega$,
the fundamental form of the Riemannian metric
$g$
and the complex structure $I$ is a
$(1,1)$-form.
So this shows that Kähler classes are in
$H^{1,1}(M)$.






\noindent
Now we shall continue discussing complex
linear algebra up
to Lefschetz hard theorem
\ref{theorem-lefschetz-decomposition}
and the Hodge-Riemann relations
\ref{theorem-hodge-riemann-relation}.


\begin{definition}
\label{definition-lefschetz-operator}
Let $(V,g,I)$ be an almost complex vector
space
equipped with an hermitian metric $g$.
The {\it Lefschetz operator} is
$$
\begin{aligned}
L: \Lambda^\bullet(V_{\mathbb{C}}^\vee)
&\longrightarrow
\Lambda^\bullet(V_{\mathbb{C}}^\vee) \\
\alpha &\longmapsto \omega\wedge\alpha.
\end{aligned}
$$
\end{definition}

\noindent
Recall that a metric on a vector space
induces a metric on the dual vector space
by defining the dual basis of an orthonormal
basis
to be orthonormal.
And for higher exterior powers all we do is
declare that wedge products of elements of
the dual orthonormal basis
(in ascending order)
form an orthonormal basis; more precisely
if $e_1,…,e_n$ is orthonormal for $V$ then
the elements of the form
$e_{i_1}\wedge…\wedge e_{i_k}$ with
$i_1<…<i_k$
are an orthonormal basis for
$\Lambda^k(V^\vee)$.

\begin{exercise}
\label{exercise-hodge-star-existence}
Let $V$ be an oriented Riemannian vector
space with
orientation form $\text{Vol}$.
There is a unique operator
$*:\Lambda^\bullet(V^\vee) \to
\Lambda^\bullet(V^\vee)$
called the {\it Hodge star} such that
$$
\alpha\wedge*\beta=g(\alpha,\beta)\text{Vol}
$$
for any $\alpha,\beta \in \Lambda^k(V^\vee)$.
\end{exercise}

\begin{proof}
First prove uniqueness,
then prove that mapping a basic form
$e^{i_1}\wedge… \wedge e^{i_k}$
to $e^{j_1}\wedge…\wedge e^{j_{n-k}}$
where $(i_1,…,i_k,j_1,…,j_{n-k})$ is a
permutation of
$(1,…,n)$ will work (up to sign maybe).
\end{proof}

\noindent
Intuition on the Hodge star operator is that
it maps a form to its ``complement'' with
respect to the volume form.
For instance, $*dx=dy \wedge dz$ in
$\mathbb{R}^3$.

\begin{lemma}
\label{lemma-hodge-star-isometry}
Hodge star is an isometry.
\end{lemma}

\begin{proof}
Consider two basic covectors $e^I$ and $e^J$.
Then $*e^I$ and $*e^J$ are determined by $I$
and $J$.
By definition of the metric on
$\Lambda^\bullet(M)$,
$\left\langle{}e^I,e^J\right\rangle{}$ is zero if $I\neq J$
and 1 otherwise.
Similarly, $\left\langle{}*e^I,*e^J\right\rangle{}$ is zero
if $I\neq J$,
while, if $I=J$, then it will clearly be 1.
\end{proof}

\begin{lemma}
\label{lemma-hodge-star-squared}
\begin{reference}
\cite[Lemma 5.5]{voi}
\end{reference}
$$
*^2=(-1)^{k(n-k)}\text{Id}.
$$
\end{lemma}

\begin{proof}
We prove that
$$
\left\langle{}\alpha-*^2\alpha,\beta\right\rangle{}=0\qquad
\forall \beta \in \Lambda^k(M).
$$
$$
\begin{aligned}
\left\langle{}\alpha,\beta\right\rangle{}\text{Vol}
&=\left\langle{}*\alpha,*\beta\right\rangle{}\text{Vol}\\
&=*\beta\wedge *^2\alpha\\
&=(-1)^{\text{deg}(*\alpha)
\text{deg}(\beta)}*^2\alpha\wedge *\beta\\
&=\left\langle{}(-1)^{(n-k)k}*^2\alpha,
*\beta\right\rangle{}\text{Vol}.
\end{aligned}
$$
\end{proof}

\noindent
In the setting of linear algebra alone,
there is no confusion on which inner product
we use to define adjoint operators.

\begin{definition}
\label{definition-dual-lefschetz}
The adjoint operator of $L$ with respect to
the metric
is called {\it dual Lefschetz operator}, i.e.
it is
the unique operator $\Lambda$ such that
$$
\left\langle{}L \alpha,\beta\right\rangle{}=\left\langle{}\alpha,
\Lambda \beta\right\rangle{}.
$$
\end{definition}

\begin{exercise}
\label{exercise-dual-lefsc-hodge-twist}
$\Lambda=*^{-1}\Lambda *$.
\end{exercise}

\begin{proof}
$$
\begin{aligned}
\left\langle{}L \alpha,\beta\right\rangle{}
&=L \alpha\wedge*\beta\\
&=\omega\wedge\alpha\wedge*\beta\\
&=\alpha\wedge\omega\wedge*\beta\\
&=\alpha\wedge L * \beta\\
&=\alpha\wedge * *^{-1}L * \beta\\
&=\left\langle{}\alpha,*^{-1}L *
\beta\right\rangle{}\text{Vol}.
\end{aligned}
$$
\end{proof}

\noindent
Next we define $H$, another operator on
$\Lambda^\bullet$
called the {\it counting operator},
which is just the identity multipliied by the
codegree of a form:
$$
H|_{\Lambda^kV}:=(k-n)\text{id}.
$$
\begin{proposition}
\label{proposition-sl2-relat-exter-algeb}
Let $(V,\left\langle{}\cdot,\cdot\right\rangle{})$ be an
euclidean vector space
endowed with an compatible almost complex
structure $I$.
Then
$$
[H,L]=2L,\qquad [H,\Lambda]=-2\Lambda,\qquad
[L,\Lambda]=H,
$$
that is, there is a representation of
$\mathfrak{sl}_2$
on $\Lambda^\bullet(V)$.
\end{proposition}

\begin{definition}
\label{definition-primitive-form}
Let $(V,\left\langle{}\cdot,\cdot\right\rangle{})$ be an
euclidean
vector space with a compatible almost complex
structure $I$.
{\it Primitive forms} are the elements of
$\Ker \Lambda$.
The set of primitive forms of degree $k$
is denoted $P^k \subset \Lambda^k(V^\vee)$.
\end{definition}

\begin{theorem}[Lefschetz decomposition]
\label{theorem-lefschetz-decomposition}
$$
\Lambda^k(V^\vee)
=
\bigoplus_i L^i(P^{k-2i})
$$
holds and is called the {\it Lefschetz
decomposition}.
Moreover, the decomposition is orthogonal
with respect to
$\left\langle{}\cdot,\cdot\right\rangle{}$.
\end{theorem}

\noindent
The upshot about this theorem
is that it is proved using the fact
that $\Lambda^\bullet(V^\vee)$ is an
$\mathfrak{sl}_2$-representation.

\begin{definition}
\label{definition-hodge-riemann-pairing}
\begin{reference}
\cite[Definition 1.2.35]{huc}
\end{reference}
Let $(V,\left\langle{}\cdot,\cdot\right\rangle{})$
be an euclidean vector space with a
compatible almost complex structure.
The {\it Hodge-Riemann pairing} is

$$
\begin{aligned}
Q: \Lambda^k(V^\vee)\times \Lambda^k(^\vee)
&\longrightarrow
\Lambda^{2n}(V^\vee)\cong \mathbb{R} \\
Q(\alpha,\beta)&
=(-1)^{\frac{k(k-1)}{2}}
\alpha\wedge\beta\wedge\omega^{n-k}.
\end{aligned}
$$
\end{definition}

\begin{theorem}[Hodge-Riemann bilinear
relation]
\label{theorem-hodge-riemann-relation}
$Q$ is positive definite on $P^{p,q}$, more
precisely,
$$
\sqrt{-1}^{p-q}Q(\alpha,\bar{\alpha})
=(n-(p+q))!\left\langle{}\alpha,
\alpha\right\rangle{}_{\mathbb{C}}>0,
$$
for $\alpha \in P^{p,q}\setminus\{0\}$ and
$p+q\leq n$.
\end{theorem}

\medskip\noindent
Now let's discuss the local theory on complex
manifolds,
which is rather just analysis on an open set
of the complex numbers.
We start defining $\partial,\bar\partial$
and aim at proving Dolbeault theorem,
which is proved exactly in the same way as de
Rham's theorem.

\medskip\noindent
Upshot. We recall by sheaf theory
that the existence of an acyclic resolution
of a sheaf
implies that the cohomology of the resolution
equals the cohomology of the sheaf (see next
paragraph).
de Rham's complex is an acylic resolution
of the constant sheaf
$\underline{\mathbb{R}}$;
acyclicity is due to Poincaré lemma - a local
property.
There some way of viewing the complex of
singular cochains
as an acyclic resolution of the constant
sheaf $\underline{\mathbb{R}}$.
Thus both de Rham cohomology and singular
cohomology
are isomorphic to the cohomology of the
constant sheaf $\underline{\mathbb{R}}$.

The derived functors assign a resolution to
every object
in the abelian category if it has enough
injectives.
An acylcic resolution is such that the
derived functors
of all objects in the resolution vanish.
The cohomology of a sheaf is defined
as the resolution given by its exact
functors.

Dolbeault theorem is exactly the same:
Dolbeault's complex is acylcic by
$\bar\partial$-Poincaré lemma,
making it a resolution of the sheaf
$\Omega^p$.

There is an interesting relationship between
the two results,
which we anticipated since $\bar\partial$
appears to act on a ``piece'' of the exterior
algebra,
fixing another piece.

\medskip\noindent
Let us start by defining $\partial$ and
$\bar\partial$.
Take a real-valued smooth function $f \in
C^\infty(V,\mathbb{R})$.
We know that the de Rham differential is
given
in local coordinates $x^i,y^j$ of our $2n$
real-dimensional manifold by
$$
df=\partial_{x^i}fdx^i + \partial_{y^j}fdy^j.
$$
When we complexify the exterior algebra we
obtain
$\Lambda_{\mathbb{C}}^\bullet(V)
= \Lambda^\bullet(V)
\otimes_\mathbb{R}\mathbb{C}
= \Lambda^\bullet(T^*V
\otimes_\mathbb{R}\mathbb{C})$.
Here we can multiply by complex numbers,
so we can define
$$
\partial_z:= \frac{1}{2}(\partial_x  - i
\partial_y ),\qquad
\partial_{\bar{z}}:=\frac{1}{2}(\partial_x +
i \partial_y)
$$
and obtain
$$
dz=dx+idy,\qquad  d\bar{z}=dx - i dy.
$$

We see that
$$
\begin{aligned}
\partial_zfdz+\partial_{\bar{z}} fd\bar{z}
&=\frac{1}{2}(\partial_xf-i\partial_y
f)(dx+idy)\\
&+\frac{1}{2}(\partial_x f+i \partial_y
f)(dx- i dy)\\
&=\partial_x f dx + \partial_y f dy\\
&=df.
\end{aligned}
$$

\noindent
\begin{definition}
\label{definition-dbar}
\begin{reference}
\cite[Definition 1.3.5]{huc}
\end{reference}
Let $X$ be an almost complex manifold.
(Or vector space! The following
objects may be defined over $V_{\mathbb{C}}$;
we put almost complex manifold since the
definition
for almost complex manifolds shall be
identical)

If  $d:\mathcal{A}^k_{X,\mathbb{C}}\to
\mathcal{A}^{k+1}_{X,\mathbb{C}}$
is the $\mathbb{C}$-linear extension of the
exterior differential, then one defines
$$
\partial:=\Pi^{p+1,q} \circ d
:\mathcal{A}^{p,q}_X \to
\mathcal{A}^{p+1,q}_X,
\qquad
\bar{\partial}:=\Pi^{p,q+1} \circ d
:\mathcal{A}^{p,q}_X \to A^{p,q+1}_X
$$
where $\Pi^{p,q}$ is the projection.
\end{definition}

\noindent
From this definition we see that
$$
\partial f = \partial_z fdz,\qquad
\bar\partial f = \partial_{\bar{z}}fd\bar{z}
$$
for $f$ real or complex-valued.

If we take $f$ complex-valued
and put $f=u+iv$,
the condition $\bar\partial f =0$
becomes the Cauchy-Riemann equations.
Thus a complex-valued function is holomorphic
if and only if
$\bar\partial f=0$, which means that its
complexified de Rham differential is a
$(1,0)$-form,
i.e. a multiple of $dz$ only.

This way we see that since real partial
derivatives commute,
$$
\begin{aligned}
\partial_z \partial_{\bar{z}}
&=\left(\frac{1}{2}(\partial_x - i
\partial_y)\right)
\left(\frac{1}{2}(\partial_x+i
\partial_y)\right)\\
&=\frac{1}{4}(\partial_x^2
+i \partial_x\partial_y-i\partial_y\partial_x
-i^2 \partial_y^2)\\
&=\frac{1}{4}(\partial_x^2+\partial_y^2)\\
&=\frac{1}{4}(\partial_x^2
-i \partial_x\partial_y+i\partial_y\partial_x
-i^2 \partial_y^2)\\
&=\left(\frac{1}{2}(\partial_x + i
\partial_y)\right)
\left(\frac{1}{2}(\partial_x-i
\partial_y)\right)\\
&=\partial_{\bar{z}}\partial_z,
\end{aligned}
$$

\noindent
so that
$$
\begin{aligned}
\partial \bar\partial f
&=\partial(\partial_{\bar{z}}f d\bar{z})\\
&=\partial_z \partial_{\bar{z}}f dz \wedge
d\bar{z}.
\end{aligned}
$$

\noindent
And similarly,

$$
\begin{aligned}
\bar\partial \partial f
&=\partial_{\bar{z}}(\partial_z f dz)\\
&=\partial_{\bar{z}}\partial_zf d\bar{z}
\wedge dz\\
&=-\partial_z\partial_{\bar{z}}fdz\wedge
d\bar{z}
\end{aligned}
$$

\noindent
and thus
$$
\partial\bar\partial=-\bar\partial \partial
$$
on functions.

For an arbitrary $(p,q)$-form
$$
\alpha=f_{I,J}dz^I \wedge d\bar{z}^J,\qquad
|I|=p,|J|=q,
$$
we shall also have
$\partial\bar\partial=\bar\partial\partial$.



\medskip\noindent
We have the following local result:
\begin{lemma}[$\bar\partial$-Poincaré lemma
on the open disc]
\label{lemma-poincare-dbar-local-lemma}
\begin{reference}
\cite[Corollary 1.3.9]{huc}
\end{reference}
If $\alpha \in \mathcal{A}^{p,q}(B)$
is $\bar\partial$ closed, then there exists
$\beta \in \mathcal{A}^{p,q-1}$ such that
$\alpha=\bar\partial \beta$.
\end{lemma}

\begin{proof}
Looks like this is a consequence of complex
analysis.
\end{proof}

\begin{proposition}
\label{proposition-leibniz-for-del}
$$
\begin{aligned}
\partial(\alpha\wedge\beta)&=
\partial\alpha\wedge\beta
+(-1)^{p+q}\alpha\wedge\partial\beta\\
\bar\partial(\alpha\wedge\beta)
&=\partial\alpha\wedge\beta+(-1)^{p+q}
\alpha\wedge\bar\partial\beta.
\end{aligned}
$$
Also,
$$
\partial^2=\bar\partial^2=0,\qquad \partial
\bar\partial=-\bar\partial \partial.
$$

\end{proposition}


\section{Complex manifolds}
\label{section-complex-manifolds}

\noindent
Finally we arrive at complex manifolds.

\begin{definition}
\label{definition-complex-manifold}
A {\it complex manifold} $M$ is a smooth
manifold admitting an open cover
$\{U_\alpha\}$ and coordinate maps
$\varphi_\alpha:U_\alpha\to\mathbb{C}^n$ such
that $\varphi_\alpha\circ\varphi_\beta^{-1}$
is holomorphic on
$\varphi_\beta(U_\alpha\cap
U_\beta)\subset\mathbb{C}^n$ for all
$\alpha,\beta$.
\end{definition}

\begin{example}
\label{example-projective-space}
The standard coordinate charts of projective
space
$\mathbb{P}^n$ are
$$
U_i:=\{(z_0:…:z_n): z_i \neq 0\}
$$
along with the maps
$$
\begin{aligned}
\varphi_i: U_i &\longrightarrow \mathbb{C}^n
\\
(z_0:…:z_n) &\longmapsto
\left(\frac{z_0}{z_i},…,\frac{z_{i-1}}{z_i},
\frac{z_{i+1}}{z_i},…,\frac{z_n}{z_i}\right).
\end{aligned}
$$

\noindent
Transition functions are given by … which are
clearly holomorphic.
\end{example}

\noindent
As a warning note that we can take the charts
to be
$U_i=\{z_i=1\}$,
but when considering $U_i \cap U_j$
we must be careful since this is {\it not}
$\{z_i=1\}\cap \{z_j=1\}$.
For example in $\mathbb{P}^1$, this would say
that the intersection
$U_0 \cap U_1$
is the single point $(1:1)$, but what about
$(1:2)$?

\begin{definition}
\label{definition-almost-complex-manifold}
An {\it almost complex manifold} $M$ is a
smooth manifold together with
an endomorphism of the real tangent bundle
$I:TM \to TM$
such that $I^2=-\text{Id}$.
\end{definition}

\begin{proposition}
\label{proposition-complex-structure}
A complex manifold has a natural almost
complex structure.
\end{proposition}

\begin{proof}
Define bundle maps $I_i$ on local
trivializations $U_i$
by rotating vectors $90$ degrees in
counter-clockwise
direction at every copy of $\mathbb{C}$ in
the tangent space $\mathbb{C}^n$.
The upshot is that since $X$ is a complex
manifold
we can interpret the $I_i$ as multiplication
by $\sqrt{-1}$
in the complex vector space structure of the
tangent space.

Then the changes of coordinates satisfy the
condition
needed to glue to a global endomorphism,
namely
$I_j=g_{ij}I_ig_{ij}^{-1}$, trivially because
the $g_{ij}$ are holomorphic and thus
$\mathbb{C}$-linear.

\noindent
{\bf Addendum.} It looks like we can just
directly
define the $I_i$ by multiplication by
$\sqrt{-1}$.
\end{proof}


\noindent
But what about the converse?
What's the condition for an almost complex
manifold to be complex?

\begin{definition}
\label{definition-integ-almos-compl-struc}
An almost complex structure on a smooth
manifold $X$
is called {\it integrable} if it is induced
from a
complex structure on $X$ as in Proposition
\ref{proposition-complex-structure}.
\end{definition}

\begin{theorem}[Newlander-Nirenberg]
\label{theorem-newlander-nirenberg}
An almost complex structure $I$
on a smooth manifold is integrable if and
only if
$[T^{0,1},T^{0,1}]\subset T^{0,1}$.
\end{theorem}

Other equivalent conditions…

The {\it tensor} is given by
$$
N_I[X,Y]=[X,Y]+I[IX,Y]+I[X,IY]-[IX,IY]
$$




\medskip\noindent
Now let's discuss Dolbeault cohomology.


\begin{proposition}
\label{proposition-dolbeault-proposition}
\begin{reference}
\cite[Proposition 2.6.11]{huc}
\end{reference}
The space $H^0(X,\Omega^p_X)$ of holomorphic
$p$-forms on a complex manifold is the
subspace
$\{\alpha \in
\mathcal{A}^{p,0}(X):\bar\partial=0\}$.
\end{proposition}

\begin{proof}
The problem is local.
We use that $f dz_{i_1}\wedge…\wedge
dz_{i+p}$ is holomorphic
if and only if $f$ is holomorphic, which is
equivalent to
$$
\bar\partial(fdz_{i_1}\wedge…\wedge dz_{i_p}
=\sum \frac{\partial f}{\partial
\overline{z}_i}d\overline{z}_i
\wedge dz_{i_1}\wedge…\wedge dz_{i_p}=0.
$$
\end{proof}

\noindent
For a more general statement we must define
the {\it $(p,q)$-Dolbeault cohomology} by
 $$
H^{p,q}(X)_{\bar\partial}
:=\frac{\Ker
(\bar\partial:\mathcal{A}^{p,q}(X) \to
\mathcal{A}^{p,q+1}(X)}
{\text{Im}(\bar\partial:\mathcal{A}^{p,q-1}
(X)
\to \mathcal{A}^{p,q}(X))}.
$$
Then

\begin{theorem}[Dolbeault]
\label{theorem-dolbeault}
\begin{reference}
\cite[Corollary 2.6.21]{huc}, \cite[Theorem
4.2]{voi}
\end{reference}
The Dolbeault cohomology of $X$ computes the
cohomology
of the sheaf $\Omega^p_X$, i.e.
$H^{p,q}_{\bar\partial}(X)=H^q(X,\Omega^p_X)
$.
\end{theorem}



\section{Kähler manifolds, Kähler identities}
\label{section-kahler}


\begin{definition}
\label{definition-hermitian-structure}
A {\it hermitian structure} on a complex
manifold $M$
is a Riemannian metric $g$ that is also an
isometry for
the complex structure $I$, that is
$g(I\cdot,I\cdot)=g(\cdot,\cdot)$.
\end{definition}

\begin{definition}
\label{definition-fundamental-form}
If $(M,I,g)$ is a hermitian manifold,
$\omega(\cdot,\cdot):=g(I\cdot,\cdot)$
is called the {\it fundamental form}.
$(M,I,g)$ is called {\it Kähler} if its
fundamental form
is closed.
\end{definition}

\noindent
As in with vector spaces endowed with almost
complex structures,
we have the operators $L,\Lambda,*$ defined
on the
cotangent spaces of a complex manifold.
As a corollary we have that Lefschetz
decomposition
passes to exterior algebra
of a hermitian manifold:

\begin{lemma}
\label{lemma-lefsc-decom-manif}
Let $(X,g)$ be an hermitian manifold.
Then there exists a direct sum decomposition
of vector bundles
$$
\Lambda^k(X)=\bigoplus_{i\geq
0}L^i(P^{k-2i}X),
$$
there $P^{k-2i}X$ is the kernel of $\Lambda$.
\end{lemma}

\medskip\noindent
Apart from the linear-algebraic operators,
on manifolds we have other operators.

\begin{definition}
\label{definition-adjoi-de-rham-diffe}
The {\it adjoint de Rham differential} is
$$
d^*:=(-1)^{n(k+1)+1}*\circ d\circ *
:\mathcal{A}^k(M) \to \mathcal{A}^{k-1}(M)
$$
where, as always, $\mathcal{A}^k(M)$ denotes
the
sections of $\Lambda^k(M):=\Lambda^k(T^*M)$.
\end{definition}

\begin{definition}
\label{definition-laplace}
The {\it Laplace operator} is
$$
\Delta=d d^*+d^*d.
$$
\end{definition}

\noindent
If $M$ is a complex manifold,
then
$$
d^*=-*\circ d \circ *.
$$

Similarly we define
$$
\partial^*:=-* \circ \partial \circ *,
\qquad
\bar\partial^*:=-* \circ \bar\partial \circ
*,
$$
which yield the Laplacians
$$
\begin{aligned}
\Delta_\partial&=\partial \partial^*
+\partial^*\partial,\\
\Delta_{\bar\partial}&=\bar\partial
\bar\partial^*+\bar\partial^*\bar\partial.
\end{aligned}
$$

\begin{proposition}[Kähler identities]
\label{proposition-kahler-identities}
Let $X$ be a complex manifold endowed with a
Kähler metric $g$. Then
\begin{enumerate}
\item $[\bar\partial,L]=[\partial,L]=0$,
$[\bar\partial^*,\Lambda]=[\partial^*,
\Lambda]=0$.
\item  Etc.
\item
$\Delta_\partial=\Delta_{\bar\partial}=
\frac{1}{2}\Delta$.
\item $\Delta$ commutes with
$*,\partial,\bar\partial,\partial^*,
\bar\partial^*,L$
and $\Lambda$.
\end{enumerate}
\end{proposition}

\begin{definition}
\label{definition-dc}
On a complex manifold $(M,I)$,
$$
d^c:=I^{-1}d I=-I d I.
$$
\end{definition}

Notice (prove) that
$$
d^c=\frac{i}{2}(\bar\partial-\partial)
$$
so that
$$
dd^c=i(\partial \bar\partial).
$$
We may also define
$$
{d^c}^*:=-* \circ d^c \circ *.
$$




\noindent
Now let's introduce Fubini-Study form.
The following discussion of psh functions is
because
to pass from a potential, i.e. a function,
to a $(1,1)$-form we apply $dd^c$ (or
$\partial\bar\partial$).
Now if the potential is psh,
then the Hessian will be positive definite,
which is just the same as saying that
$dd^cf$ is positive definite $(1,1)$-form.

\begin{definition}
\label{definition-positive-11form}
Let $V$ be a real vector space equipped with
a complex structure $I$.
A $(1,1)$-form $\omega$ is {\it positive}
if $\omega(x,Ix)\geq 0$ for any nonzero $x
\in V$.
(Notice the $I$ is in the second entry!)
It is called {\it Hermitian} or {\it
semipositive}
if $\omega(x,Ix)>0$ for all nonzero $x \in
V$.
\end{definition}

\noindent
See \cite[Lemma 3.1.7]{huc}
to find Huybrecht's definition of a closed
{\it positive definite}
$(1,1)$-form, namely if
$\omega=\frac{i}{2}\sum h_{ij}dz_i \wedge
d\bar{z}_j$
such that $(h_{ij}(x))$ is a positive
definite hermitian matrix for any $x \in X$.

The point is that this condition
is the same as saying that the associated
Riemannian metric is in fact positive
definite.
Indeed, $\omega(\cdot,I\cdot)=g(\cdot,\cdot)$
which is positive definite.

Finally observe that numerical positivity of
the
Kähler class,
i.e. positivity with respect to the
intersection form,
follows exactly from this fact: we shall have
that
$\int_Y \omega^p>0$ for any analytic cycle
$Y$.

\begin{definition}
\label{definition-psh}
A function $f$ on a complex manifold is
{\it plurisubharmonic} if $dd^cf$ is
a positive $(1,1)$-form.
\end{definition}

\noindent
Again: psh is equivalent to the Hermitian
matrix
$$
\left(\frac{\partial^2 f}{\partial
z_i\partial \bar{z}_j}\right)
$$
being positive semidefinite.
ChatGPT: it is also equivalent to: its
restriction to
any complex line being subharmonic ($\Delta
f\geq 0$).

\begin{exercise}
\label{exercise-logl}
Let $z_1,…,z_n$ be the standard coordinates
on $\mathbb{C}^n$,
and $l(z_1,…,z_n):=\sum|z_i|^2$.
Prove that
$dd^c\text{log}l=\frac{dd^cl}{l}-\frac{dl
\wedge d^cl}{l^2}$
on $\Omega=\mathbb{C}^n\setminus 0$.
Prove that the Hermitian form $dd^c
\text{log} l$
is degenerate and can be written as
$(-\sqrt{-1})\sum_{i=1}^{n-1}
\alpha_i\xi_i\wedge \bar{\xi}_i$,
where $\xi_i$ is a local frame in
$\Lambda^{1,0}(\Omega)$
and $\alpha_i$ are positive functions.
\end{exercise}

\begin{proof}
For the first equation we just use the fact
(why?) that both $d^c$ and $d$ satisfy the
chain rule.
For the second equation I computed that
\begin{equation}
\label{equation-ddcl1}
dd^cl=\sqrt{-1}\sum dz_i\wedge d\bar{z}_i.
\end{equation}
But I also need to compute
$d^cl=\frac{\sqrt{-1}}{2}(\bar\partial-
\partial)l$
and $dl=(\partial+\bar\partial)l$.
I find that $dl \wedge d^cl=\sqrt{-1}\partial
l \wedge \bar\partial l$.
But I also found that
$\partial l=\sum \bar{z}_idz_i$ and
$\bar\partial l=\sum z_i d\bar{z}_i$.
So I get
\begin{equation}
\label{equation-ddcl2}
dl\wedge d^cl=\sqrt{-1}
\left(\sum \bar{z}_i dz_i\right)
\wedge
\left(\sum z_i d \bar{z}_i\right)
=\sqrt{-1}\sum\bar{z}_iz_jdz_i\wedge
d\bar{z}_j
\end{equation}

\noindent
Now we join Equations \ref{equation-ddcl1}
and \ref{equation-ddcl2}
to find the coefficient functions with
respect to $dz_i \wedge d\bar{z}_i$.
From Equation \ref{equation-ddcl1} we get
coefficient functions $\delta_{ij}$,
and from Equation \ref{equation-ddcl2} we get
$\bar{z}_iz_j$ so
$$
dd^c\text{log}l=\sqrt{-1}\sum
\left(\frac{\delta_{ij}}{l}
-\frac{\bar{z}_iz_j}{l^2}\right)
dz_i \wedge d\bar{z}_j.
$$
I'm not sure how to prove that  the
coefficient functions
are positive… Or that the form is degenerate…
But fortunately this can be checked at
\cite[p. 118]{huc}.
\end{proof}

\noindent
The point of this exercise is to show that
the form
$dd^c\ell$ is almost an Hermitian form ---
but it has a kernel.
I think the point of the next exercises in
this
\href{http://verbit.ru/IMPA/Surfaces-2025/%
assign-surfaces-2025-04.pdf}{
handout}
is to show that the kernel is actually
the ``radial directions''.
And the point of that is to take quotient
by $\mathbb{C}^*$ action, which is literally
passing to the projective space.

A bit more formally:

\begin{definition}
\label{definition-basic-form}
Let $M$ be a manifold and $\omega \in
\Omega^k(M)$.
Suppose that $\pi:M \to B$ is a surjective
submersion with connected fibers.
We say that $\omega$ is {\it basic} (with
respect
to $\pi$) is there exists a form
$\overline{\omega}\in \Omega^k(B)$ such that
$\pi^*  \overline{\omega}$.
\end{definition}

\begin{exercise}
\label{exercise-basic-forms-characterization}
\begin{enumerate}
\item Show that $\omega$ is basic if and only
if
$i_X \omega=0$ and $L_X\omega=0$
for all vector fields $X$ tangent
to the fibers of $\pi$.
In particular, if $\omega$ is closed,
show that it is basic if
$\Ker (\pi_*) \subseteq \Ker(\omega)$
(pointwise in $M$).

\item Suppose that $\omega$ is a closed
2-form on $M$ and $\Ker(\pi_*)=\Ker(\omega)$.
Show that $\omega=\pi^*\overline{\omega}$
and $\overline{\omega}\in\Omega^2(B)$
is symplectic.

\item (Application to reduction.)
\end{enumerate}
\end{exercise}

\begin{proof}
\begin{enumerate}
\item Since $\pi$ is surjective, for any
vector $\bar{X}$ tangent to $B$ there is
a vector $X$ tangent to $M$ such that
$\pi_*X=\bar{X}$.
Define a form
$\overline{\omega}(\bar{X}^1,…,\bar{X}^k):=
\omega(X^1,…,X^k)$.

To show this does not depend on the choice of
lift of $\bar{X}$
just notice that the tangent space of $M$ can
be written
$T_pM \simeq T_pF_b \oplus \pi^*(T_bB)$ where
$\pi(p)=b$.
Since $\omega$ kills vectors tangent to the
fibers,
$\overline{\omega}$ is well-defined.

Notice that there can be different choices of
$p$ such that
$\pi(p)=b$. We must show that the definition
is also
independent of this choice.
Let $X \in T_pM$ and $Y \in T_qM$ be such
that
$\pi(p)=\pi(q)=b \in B$ and $\pi_*X=\pi_*Y$.
We wish to show that
$\omega(…,X,…)=\omega(…,Y,…)$.

This a little more involved so I'll just
sketch it for now.
If there is a vector field $V$ tangent to the
fibers
whose flow $\varphi$ joins $p$ and $q$, we
are done,
since we pull back $\omega$ along this flow,
and by the condition of the Lie derivative,
i.e. $L_V\omega$ for $V$
tangent to the fiber, we see that
$\varphi^*\omega=\omega$.
(Indeed,
$L_V\omega=\frac{d}{dt}\varphi^*\omega$,
so the derivative is zero and at $t=0$ we
have $\omega$.)

Note that the choice of $Y$ within $T_qM$
is superfluous since we already know that
$\omega$
is constant within every tangent space at a
poi

…
\end{enumerate}
\end{proof}




\begin{exercise}
\label{exercise-kahler-form-conformal-class}
Let $X$ be a connected complex manifold of
dimension $n>1$
and let $g$ be a Kähler metric.
Show that $g$ is the only Kähler metric in
its conformal class,
i.e. if $g'=e^fg$ is Kähler then $f$ is
constant.
\end{exercise}

\begin{exercise}
\label{exercise-lefsc-injec-1-forms}
Let $X$ be a connected complex manifold and
$h$ be a Kähler metric.
Let $\omega$ be the associated Kähler form.
Show that if $\dim X \geq  2$, the wedge
product with $\omega$
is injective on 1-forms.

Show that if $\dim X \geq  2$ and $\phi$ is a
differentiable
function with values in $\mathbb{R}^+$ such
that $\phi h$ is also
a Kähler metric, then $\phi$ is constant.
\end{exercise}

\begin{proof}
Take local symplectic coordinates, i.e. such
that
$\omega = \sum_idz^i \wedge d \bar{z}^i$.
Let $\alpha$ be any 1-form and suppose it is
written
in those coordinates as
$\alpha=\sum_ia_idz^i +
\sum_i\bar{a}_id\bar{z}^i$.
Then
$$
\begin{aligned}
\alpha\wedge\omega&=
\left(\sum_ia_idz^i\right)
\wedge
\left(\sum_i dz^i \wedge d\bar{z}^i\right)
+
\left(\sum_i \bar{a}_i dz^i\right)
\wedge
\left(\sum_i dz^i \wedge d\bar{z}^i\right)\\
&=
\sum_{i,j}a_i dz^i \wedge dz^j \wedge
d\bar{z}^j
+
\sum_{i,j}\bar{a}_i d\bar{z}^i \wedge dz^j
\wedge d\bar{z}^j.
\end{aligned}
$$

\noindent
We see that if $\alpha \wedge \omega=0$
then both its $1,2$ and $2,1$ parts are zero,
which foces $a_i$ and $\bar{a}_i$ to vanish,
i.e. $\alpha=0$.
\end{proof}
\begin{exercise}
\label{exercise-holom-vecto-bundl-kahle}
\begin{reference}
\cite[Exercise 2, Chapter 3]{voi}
\end{reference}
Let $E$ be a holomorphic vector bundle of
rank $r$
over a complex manifold $X$.
Show that if $L$ is a holomorphic line bundle
on $X$,
then $\mathbb{P}(E^* \otimes L^*)=
\mathbb{P}(E^*)$
but the line bundle
$\mathcal{O}_{\mathbb{P}(E^*\otimes L^*)}(1)$
is isomorphic to
$\mathcal{O}_{\mathbb{P}(E^*)}(1) \otimes
\pi^*L$,
where $\pi:\mathbb{P}(E^*)\to X$ is the
structural morphism.
\end{exercise}

\begin{exercise}
\label{exercise-kahler-iff-i-parallel}
$(M,I,g)$ is Kähler if and only if
$\nabla^{LC}I=0$.
\end{exercise}





\section{Hodge theory}
\label{section-hodge-theory}

\noindent
On a compact hermitian manifold $(X,g)$ we
may define
$$
(\alpha,\beta):=\int_Xg_{\mathbb{C}}(\alpha,
\beta)*1,
$$
which is an hermitian product on
$\mathcal{A}_{\mathbb{C}}^*(X)$.

\begin{lemma}
\label{lemma-adjoints}
\begin{reference}
\cite[Lemma 3.2.3]{huc}
\end{reference}
Let $X$ be a compact hermitian manifold.
Then $\partial^*$ and $\bar\partial^*$
are adjoints to $\partial$ and $\bar\partial$
with respect to
$(\cdot,\cdot)$.
\end{lemma}

\begin{proof}
We have
$$
\begin{aligned}
(\partial\alpha,\beta)
&=\int_X\left\langle{}\partial\alpha,
\beta\right\rangle{}_\mathbb{C}\text{Vol}
=\int_X \partial\alpha\wedge*\beta,\\
(\alpha,\partial^*\beta)
&=\int_X\left\langle{}\alpha,-
*\partial*\beta\right\rangle{}_{\mathbb{C}}\text{Vol}
=\int_X\alpha\wedge\partial*\beta.
\end{aligned}
$$

\noindent
Now we also have
$\partial(\alpha\wedge\beta)
=\partial\alpha\wedge\beta+
\alpha\wedge\partial\beta$,
and the integral of the first vanishes by
Stokes theorem.
\end{proof}

\begin{lemma}
\label{lemma-harmonic-iff-closed}
Let $X$ be a compact hermitian manifold.
A form is $\bar\partial$-harmonic if and only
if
it is $\bar\partial$-closed and
$\bar\partial^*$ closed.
The same holds for $\partial$,
and in fact, on any compact oriented
Riemannian manifold $(M,g)$,
a similar result holds for the usual
$d$-Laplacian.
\end{lemma}

\begin{proof}
$$
\begin{aligned}
(\Delta \alpha,\alpha)&=((d
d^*+d^*d)\alpha,\alpha)\\
&=(d d^*\alpha,\alpha)+(d^*d\alpha,\alpha)\\
&=(d^*\alpha,d^*\alpha)+(d\alpha,d\alpha)\\
&=\|d^*\alpha\|^2+\|d\alpha\|^2.
\end{aligned}
$$
\end{proof}

\noindent
The following is the first of the symmetries
of the Hodge diamond:

\begin{lemma}
\label{lemma-hodge-commutes-with-laplacian}
\begin{reference}
\cite[Lemma A.0.14]{huc}
\end{reference}
For any Riemannian manifold $(M,g)$,
not necessarily compact, $*\Delta=\Delta*$
and $*:\mathcal{H}^k(M,g) \cong
\mathcal{H}^{m-k}(M,g)$.
\end{lemma}

\noindent
In the hermitian case, we have for a compact
hermitian connected
manifold $(M,g)$, that the pairing
$$
\begin{aligned}
\mathcal{H}^{p,q}_{\bar\partial}
\times\mathcal{H}^{n-p,n-q}_{\bar\partial}
 &\longrightarrow \mathbb{C} \\
(\alpha,\beta) &\longmapsto
\int_X\alpha\wedge\beta
\end{aligned}
$$

\noindent
is non-degenerate and thus
$$
\mathcal{H}^{p,q}_{\bar\partial}(X,g)
\cong
\mathcal{H}^{n-p,n-q}_{\bar\partial}(X,g),
$$
which is called {\it Serre duality},
the second of the symmetries of the Hodge
diamond.
We shall later generalize this Serre duality
to other vector bundles.

There is a third symmetry of the Hodge
diamond,
given by the Lefschetz operator.
Since $[L,\Delta]=0$, one of the Kähler
identities,
we see that $L$ maps harmonic forms to
harmonic forms.
Further, since $L^{n-k}:\Lambda^k(V^\vee) \to
\Lambda^{2n-k}(V^\vee)$
is bijective we have that the induced map
$$
L^{n-k}:\mathcal{H}^{p,k-p}(X,g)
\to \mathcal{H}^{n+p-k,n-p}(X,g)
$$
is injective. It is also bijective since…
thus it is an isomorphism.

\medskip\noindent
Now we are almost ready for Hodge theorem.

\begin{proposition}
\label{proposition-decom-ortho}
\begin{reference}
\cite[Proposition 3.2.2]{huc}
\end{reference}
Let $(X,g)$ be a compact hermitian manifold.
The degree, bidegree and Lefschetz
decomspositions
are orthogonal with respect to
$(\cdot,\cdot)$.
\end{proposition}

\noindent
The bidegree decomposition passes to
$\partial$ and $\bar\partial$-harmonic forms
on any compact hermitian manifold,
but only for Kähler manifolds does it pass to
$d$-harmonic forms.

\begin{proposition}
\label{proposition-decom-harmo-forms}
Let $(X,g)$ be a hermitian manifold, not
necessarily compact.
\begin{enumerate}
\item $\mathcal{H}^{k}_{\bar\partial}(X,g)
=\bigoplus_{p+q=k}\mathcal{H}^{p,q}
_{\bar\partial}(X,g)$
and
$\mathcal{H}^{k}_\partial(X,g)
=\bigoplus_{p+q=k}\mathcal{H}^{p,q}
_\partial(X,g)$.

\item If $(X,g)$ is Kähler then both
decompositions coincide with
$\mathcal{H}^{k}(X,g)
=\bigoplus_{p+q=k}\mathcal{H}^{p,q}(X,g)$.
In particular,
$\mathcal{H}^k(X,g)_{\mathbb{C}}
=\mathcal{H}^k_{\bar\partial}(X,g)
=\mathcal{H}^k_{\partial}(X,g)$.
\end{enumerate}
\end{proposition}

\noindent
We finally arrive at the Hodge decomposition.
The first result concerns only compact
Riemannian manifolds,
and ``requires some hard, but by now
standard, analysis'':

\begin{theorem}[Hodge decomposition]
\label{theorem-hodge-decomposition}
Let $(M,g)$ be a compact Riemannian manifold.
Then
$$
\mathcal{A}^k(M) =
d(\mathcal{A}^{k-1}(M))
\oplus
\mathcal{H}^k(M,g)
\oplus
d^* (\mathcal{A}^{k+1}(M))
$$
which is orthogonal with respect to
$(\cdot,\cdot)$.
Moreover, the space $\mathcal{H}^k(M)$ is
finite-dimensional.
\end{theorem}

\noindent
And now we also have
\begin{theorem}[Hodge decomposition]
\label{theorem-hodge-decomposition-hermitian}
Let $(X,g)$ be a compact hermitian manifold.
Then
$$
\mathcal{A}^{p,q}(X)=
\partial\mathcal{A}^{p-1,q}(X)
\oplus
\mathcal{H}^{p,q}_\partial(X,g)
\oplus
\partial^*\mathcal{A}^{p+1,q}(X)
$$
and
$$
\mathcal{A}^{p,q}(X)=
\bar\partial\mathcal{A}^{p-1,q}(X)
\oplus
\mathcal{H}^{p,q}_{\bar\partial}(X,g)
\oplus
\bar\partial^*\mathcal{A}^{p+1,q}(X).
$$
If $(X,g)$ is Kähler, then
$\mathcal{H}^{p,q}_\partial(X,g)=\mathcal{H}
^{p,q}_{\bar\partial}(X,g)$.
\end{theorem}

\noindent
Upshot: if $X$ is Kähler we can pass from
usual
harmonic $(p,q)$-forms to
$\bar\partial$-harmonic $(p,q)$-forms,
which takes us to $H^{p,q}(X)$.
To pass from $H^k(X)$ to $\mathcal{H}^k(X)$
and from
$\mathcal{H}^{p,q}_{\partial}(X)$ to
$H^{p,q}(X)$
we have the following two lemmas
which in fact are corollaries from the
corresponding decompositions
(Theorems \ref{theorem-hodge-decomposition}
and
\ref{theorem-hodge-decomposition-hermitian},
respectively).

\begin{lemma}
\label{lemma-de-rham-to-harmonic}
\begin{reference}
\cite[Corollary A.0.17]{huc}
\end{reference}
Let $(M,g)$ be a compact oriented Riemannian
manifold.
The natural map $\mathcal{H}^k(M) \to
H^k(M,\mathbb{R})$
that associates to any harmonic form its
cohomology class is
an isomorphism.
\end{lemma}

\begin{lemma}
\label{lemma-harmo-pq-dolbe-pq}
Let $(X,g)$ be a compact hermitian manifold.
Then the canonical projection
$\mathcal{H}_{\bar\partial}^{p,q}(X,g)\to
H^{p,q}(X)$
is an isomorphism.
\end{lemma}

\medskip\noindent
Before we pass to Hodge theorem we state one
final key result

\begin{lemma}[ddbar Lemma]
\label{lemma-ddbar}
Let $X$ be a compact Kähler manifold.
If $\alpha$ is a $d$-closed $(p,q)$-form,
it is equivalent that $\alpha$ is
\begin{enumerate}
\item $d$-exact,
\item $\partial$-exact,
\item $\bar\partial$-exact,
\item $\partial \bar\partial$-exact.
\end{enumerate}
\end{lemma}

\noindent
There is another version in terms of
$dd^c:=
I^{-1}dI=\frac{1}{2}(\partial-\bar\partial)$.

\begin{lemma}[ddc Lemma]
\label{lemma-ddc}
\begin{reference}
\cite[Lemma 3.A.22]{huc}
\end{reference}
Let $X$ be a compact Kähler manifold.
Any $d$-closed and $d^c$-exact form form is
$dd^c$-exact.
\end{lemma}



\noindent
That's all the ingredients for a
(very rough) idea on why the Hodge theorem
holds.
Upshot:
$$
H^k(X,\mathbb{C})
= \mathcal{H}^k(X)
= \bigoplus_{p+q=k}\mathcal{H}^{p,q}(X)
=
\bigoplus_{p+q=k}\mathcal{H}^{p,q}
_{\bar\partial}(X)
=\bigoplus_{p+q=k}H^{p,q}(X).
$$
\begin{enumerate}
\item Start with singular complex cohomology,
\item pass to harmonic forms
via Lemma \ref{lemma-de-rham-to-harmonic}
(this is where heavy real analysis enters;
what Huybrechts calls Hodge decomposition),
\item decompose harmonic forms in
$(p,q)$-types
via Proposition
\ref{proposition-decom-harmo-forms}
(using Kählerness),
\item use that $\mathcal{H}^{p,q}\cong
\mathcal{H}^{p,q}_{\bar\partial}$
by the same Proposition
\ref{proposition-decom-harmo-forms}
(using Kählerness),
\item go back to Dolbeault cohomology
via Lemma \ref{lemma-harmo-pq-dolbe-pq}
(again Hodge decomposition, now for
$\bar\partial$).
\end{enumerate}

\begin{theorem}[Hodge]
\label{theorem-hodge-theorem}
Let $(X,g)$ be a compact Kähler manifold.
Then there exists a decomposition
$$
H^k(X,\mathbb{C})
= \bigoplus_{p+q=k}H^{p,q}(X)
$$
which does not depend on the Kähler
structure.
Moreover, with respect to complex conjugation
on
$H^\bullet(X,\mathbb{C})=H^\bullet(X,
\mathbb{R})\otimes\mathbb{C}$
we have
$\overline{H^{p,q}(X)}=H^{q,p}(X)$,
and, by Serre duality on harmonic forms,
$H^{p,q}(X) \cong H^{n-p,n-q}(X)^\vee$.
\end{theorem}

\begin{exercise}
\label{exercise-pic-of-pn-is-z}
\begin{reference}
\cite[Exercise 3.2.11]{huc}
\end{reference}
Show that $\text{Pic}(\mathbb{P}^n) \simeq
\mathbb{Z}$.
\end{exercise}

\noindent
The next step is to check
that most results we've developed so far for
complex vector spaces and complex manifolds
descend to cohomology in Kähler manifolds.


\medskip\noindent
Let's put here the part on the cones until
I find a better place.

Recall that $\text{NS}(X)$ is the quotient
group
$\text{Pic}(X)/H^1(\mathcal{O}_X)\subset
H^2(X,\mathbb{Z})
\subset H^2(X,\mathbb{R})$
coming from the exponential exact sequence.
This has an intersection form.
In the case of a K3 we readily see that
$h^1=0$ so $\text{NS}=\text{Pic}$.
In fact, the intersection form coincides with
the usual cup product on 2-cocycles of
singular cohomology.

For a Kähler manifold we have
\begin{itemize}
\item The {\it Kähler cone}. Classes
in $H^{1,1}(X,\mathbb{R})$ that are
represented by a
Kähler form.

\item The {\it positive cone}. Classes
$\text{NS}(X)_\mathbb{R}
:=\text{NS}(X) \otimes_\mathbb{Z}\mathbb{R}$
with $\alpha^2>0$.

\item The {\it ample cone}. Classes
$\text{NS}(X)_\mathbb{R}$
with $\alpha^2>0$ and $\alpha.C>0$ for all
algebraic curves 0.
(See
\ref{algebraic-geometry-theorem-%
Nakai-Moishezon-criterion}.)
This is true, but hard to prove, see
\cite[p. 150]{huk}.
Another description is that the ample cone
is the set of classes in
$\text{NS}(X)_\mathbb{R}$
that are finite sums $\sum a_iL_i$
for $L_i \in \text{NS}(X)$ ample
and $a_i \in \mathbb{R}_{>0}$.

\item The {\it nef cone}. Classes (in
$\text{NS}(X)?)$
with $\alpha.C=\int_C\alpha\geq 0$ for any
algebraic curve.
There is also a notion of nef for analytic
cycles
which basically says that the form is
approximated
by Kähler forms.
In \cite[Theorem 4.7(ii)]{Demailly-Paun} we
have that
a $(1,1)$-class is nef (in the analytic
sense) if and only if
$\int_Y \alpha\wedge\omega^{p-1}\geq 0$ for
every
irreducible analytic set $Y \subset X$,
$p=\dim Y$
and $\bar{\omega}$ Kähler class.

\end{itemize}

\noindent
The upshot about Demailly-Pa\u un theorem
is that Kähler classes are numerically
positive
(recall that ``numerically'' means w.r.t.
intersection form)
on {\it analytic} cycles. This distinguishes
complex analytic geometry from algebraic
geometry.
This is why Demailly-Pa\u un say that
their theorem is a generalization of
Nakai-Moishezon:
because the latter is for {\it algebraic}
cycles,
indeed, found in Hartshorne's book.

The proof goes roughly as follows.
It's clear that $\mathcal{K} \subset
\mathcal{P}$
($\mathcal{P}$ is the analytically
numerically positive cone)
because, since the Kähler form
comes from the Riemannian metric
and the curves are complex,
which means that they are invariant under
the complex structure,
we have integrals
$\int_Y\omega$ for curves $Y \subset X$
but $\omega(v,Jv)>0$ for a basis $v,Jv$
of $T_pY$.
That is to say, there are no complex
Lagrangian curves.

But that's elementary.
The point of the theorem is proving that
a form in $\overline{\mathcal{K}}\cap
\mathcal{P}$,
i.e. in the closure of the Kähler cone
and which is also numerically analytically
positive,
is in fact Kähler.

So, the theorem says that the boundary of the
Kähler cone does
not intersect the numerically analytically
positive cone.

\begin{theorem}[Demailly-Pa\u un]
\label{theorem-demailly-paun}
Let $X$ be a compact Kähler manifold.
The Kähler cone is a connected component
of the cone of 1,1-classes numerically
positive on
analytic cycles.
\end{theorem}

\begin{exercise}
\label{exercise-demailly-paun}
Let $M$ be a K3 surface such that
$\text{rk}\text{NS}(M)<20$.
Using Demailly-Paun theorem,
prove that there exists a non-zero vector
$\eta \in H^{1,1}(M)$ which satisfies
$\eta^2=0$
and belongs to the boundary of the Kähler
cone.
\end{exercise}

\begin{proof}
Let $\alpha \in
\text{NS}(X)^{\perp}\setminus\{0\}$,
which is nonempty since $\dim
\text{NS}(X)<20=h^{1,1}$.
By the Hodge index theorem
\ref{theorem-hodge-index}
the intersection form on
$H^{1,1}(X,\mathbb{R})$
has index $(1,19)$.
If we assume that $X$ is projective
then we know that the restriction of
the Fubini-Study form $\omega$ is integral.
Since $\omega$ is positive (by being a Kähler
form),
we see by the index of the intersection form
that $\alpha$ must be negative.

Define $\eta_t:=t\alpha+(1-t)\omega$ for
$t\in[0,1]$.
Then $\eta_t^2=t^2\alpha^2+(1-t)^2\omega^2$.
Since at $t=0$ we have a positive number
and at  $t=1$ a negative number,
we know that there exists  $t_0$
such that $\eta_{t_0}^2=0$
by continuity of the intersection form.

Now we need to show that $\eta_{t_0}$ is
in the boundary of the Kähler cone.
For this it's enough to show that for every
$\varepsilon>0$
$\eta_{t_0+\varepsilon}$ is in the Kähler,
which
by Demailly-P\u aun
means that we have to show
that it is positive on analytic cycles.

By our choice of $\omega$ and $\alpha$
we see that
$$
\int_Y \eta_{t_0+\varepsilon}
=\int_Y
(t_0+\varepsilon)\alpha+(1-(t_0+\varepsilon))
\omega
$$
and the integral of $\alpha$ vanishes
since $\alpha$ is in the orthogonal
complement
of $\text{NS}(X)$, which contains
all analytic curves by $X$ being projective
(Chow theorem).
\end{proof}






\section{Lefschetz theorems}
\label{section-lefschetz-theorems}

\begin{theorem}[Lefschetz theorem on
(1,1)-classes]
\label{theorem-11-classes}
All the classes in $H^{1,1}(X,\mathbb{Z})
:=H^{1,1}(X) \cap H^2(X,\mathbb{Z})$ are
algebraic.
\end{theorem}

\noindent
This implies that
$\text{NS}(X):=\text{Pic}(X)/\ker c_1$
is isomorphic to $H^{1,1}(X,\mathbb{Z})$.
The ``realification'' of the integer classes,
$H^{1,1}(X,\mathbb{Z}) \otimes \mathbb{R}
=\text{NS}(X) \otimes \mathbb{R}$
is a vector subspace
$H^{1,1}(X) \cap H^2(X,\mathbb{R})$.

\begin{theorem}[Hodge index]
\label{theorem-hodge-index}
The index of the intersection pairing on
$H^{1,1}(X,\mathbb{R})$
is $(1,\rho(X)-1)$.
\end{theorem}

\begin{theorem}[Hard Lefschetz Theorem]
\label{theorem-hard-lefschetz}
\begin{reference}
\cite[p. 122]{gri}
\end{reference}
The map
$$
L^k:H^{n-k}(M) \to H^{n+k}(M)
$$
is an isomorphism.
If we define the {\it primitive} cohomology
as
$$
P^{n-k}(M) = Ker L^{k+1}
$$
we have
$$
H^m(M)=\bigoplus_k L^kP^{m-2k}(M),
$$
called the Lefschetz decomposition.
\end{theorem}

\section{Vector bundles, connections,
curvature and Chern classes}
\label{section-vector-bundles}

\begin{lemma}[Vector bundle construction from
fibers]
\label{lemma-vector-bundle-from-fibers}
\begin{reference}
\cite[Lemma 3.4]{lec}
\end{reference}
Suppose $M$ is a complex manifold,
and for each $p \in M$ we are given a
$k$-dimensional vector space $E_p$.
Let $E=\coprod_{p \in M}E_p$ and let $\pi:E
\to M$ be the obvious projection.
Suppose further that we are given
\begin{enumerate}
\item an indexed open cover $U_i$ of $M$
\item for each $i$, a bijection $
\psi_i:\pi^{-1}(U_i) \to U_i
\times\mathbb{C}^k$,
\item for each $i,j$ with $U_i\cap U_j \neq
\emptyset$
a holomorphic map $g_{ij}:U_i \cap U_J \to
\text{GL}(k,\mathbb{C})$
such that $\psi_i\circ
\psi_j^{-1}(p,v)=(p,g_{ij}(p)v$.
\end{enumerate}
\noindent
Then $\pi:E \to M$ has a unique structure as
a holomorphic vector bundle
with $(U_i,\psi_i)$ as holomorphic local
trivializations,
and with $g_{ij}$ as transition functions.
\end{lemma}

\begin{proposition}[Vector bundle
construction from cocycle]
\label{proposition-vecto-bundl-cocyc}
Let $M$ be a complex manifold and let $U_i$
be an open cover of $M$.
Suppose that for each $i,j$ such that $U_i
\cap U_j \neq \emptyset$
we are given a holomorphic map $g_{ij}:U_i
\cap U_j \to \text{GL}(k,\mathbb{C})$
such that
$$
g_{ij}(p)g_{jk}(p)g_{ki}(p)=\text{Id}
$$
for any $p \in U_i\cap U_j\cap U_k$.

Then there is a rank-$k$ holomorphic vector
bundle $E \to M$
that admits a holomorphic trivialization over
each set $U_i$
with transition functions given by $g_{ij}$.
\end{proposition}

\begin{proposition}[Holomorphic Vector Bundle
Isomorphism Criterion]
\label{proposition-isomo-vecto-bundl}
\begin{reference}
\cite[Proposition 3.7]{lec}
\end{reference}
Two vector bundles $E$ and $F$
with transition functions $g_{ij}$ and
$h_{ij}$
are isomorphic if and only if
$g_{ij}=A^{-1}_ih_{ij}A_j$
for some $A_i:U_i \to
\text{GL}_n(\mathbb{C})$.
\end{proposition}

\begin{proof}
The idea is to define a map in local
trivializations
and confirm that this glues to a global
bundle map.
This means: invent
$\varphi_i:E(U_i):=\Gamma_E(U_i) \to F(U_i)$,
i.e. a local map of sections
(yes, a map of sections, since a bundle map
is constructed as map of the total spaces,
which are points and vectors,
but we need to preserve the projection,
so all we care is about the vectors).
Then check that on $U_i\cap U_j:= U_{ij}$
we have $\varphi_i=\varphi_j$.
This literally creates a bundle map.

Try to do it for line bundles,
so $A_i$ are just functions.
Take a local frame $e_i$ for $E$
and a local frame $f_j$ for $F$.
Change coordinates and apply the map,
apply the map and change coordinates
- this should give both sides of the cocycle
condition
of the proposition.
\end{proof}

\begin{exercise}[Tensor product bundle]
\label{exercise-tensor-product-bundle}
\begin{reference}
\cite[Exercise 2.2.1]{huc}
\end{reference}
Let $E$ and $F$ be vector bundles determined
by the cocycles
$g_{ij}: U_i \cap U_j \to
\text{GL}(r,\mathbb{C})$
and $\tilde{g}_{ij}:U_i \cap U_j \to
\text{GL}(\tilde{r},\mathbb{C})$.
Verify the cocycle description for
\begin{itemize}
\item The direct sum $E \oplus F$,
\item the tensor product $E \otimes F$ and
\item the dual bundle $E^\vee$.
\end{itemize}
\end{exercise}

\begin{proof}
Let's do the tensor product.
I think that, slightly more generally than in
the
formulation of the question, if $U_i$
is a trivializing cover for $E$ and
$\tilde{U}_j$ is a trivializing cover for
$F$,
we should consider the cover $U_i\cap
\tilde{U}_j$.
Let us denote such a cover simply by $U_i$.

Define a cocycle by
$$
\begin{aligned}
h_{ij}: U_i\cap U_j &\longrightarrow
\text{GL}(r\tilde{r},\mathbb{C}) \\
p &\longmapsto g_{ij}(p)\otimes
\tilde{g}_{ij}(p).
\end{aligned}
$$

\noindent
The key observation is that while formally
the image of $h_{ij}$
lies in $\text{GL}(\mathbb{C}^r \otimes
\mathbb{C}^{\tilde{r}})$,
we can regard such a space as
$\text{GL}(r\tilde{r},\mathbb{C})$
by defining, for $A \in
\text{GL}(r,\mathbb{C})$
and $A' \in \text{GL}(\tilde{r},\mathbb{C})$,
the map $A(v^ie_i) \otimes
\tilde{A}(\tilde{v}^j\tilde{e}_j)
=v^i\tilde{v}^j Ae_i \otimes
\tilde{A}\tilde{e}_j$
where $e_i$ and $\tilde{e}_j$ are bases
of $\mathbb{C}^r$ and
$\mathbb{C}^{\tilde{r}}$.
This is indeed a linear transformation of the
$r\tilde{r}$-dimensional $\mathbb{C}$-vector
space
generated by the basis $e_i \otimes
\tilde{e}_j$.

Thus the cocycle condition follows simply by
the cocycles of $E$ and $F$, and
commutativity of the
coordinate functions.
\end{proof}

\begin{example}[Determinant bundle]
\label{example-determinant-bundle}
Let $\pi:E \to M$ be a vector bundle of rank
$r$.
The {\it determinant bundle} of $E$ is given
by the following data.
A space $\tilde{E}:=\coprod_{p \in
M}\Lambda^rE_p$
with the obvious projection map
and topology induced from such map.
Since we have a cover $U_i$ of $M$
along with local trivializations
$U_i\overset{\varphi_i}{\simeq}
U_i\times\mathbb{C}^r$
we may just postcompose $\varphi_i$ with the
functor $\Lambda^r$
on the second entry, giving trivializations
$$
\xymatrix{
\pi(U_i)\ar[r]^-{\varphi_i}_-{\simeq}
& U_i \times
\mathbb{C}^r\ar[r]^-{\text{Id}\times
\Lambda^r}
& U_i \times \Lambda^r\mathbb{C}^r=U_i \times
\mathbb{C}
}
$$
We can also apply $\Lambda^r$ to the
transition functions
$g_{ij}$ on the intersections
$U_{ij}:=U_i\cap U_j$
$$
\xymatrix{
U_{ij}\ar[r]^-{\varphi_i}&U_{ij} \times
\mathbb{C}^r
\ar[d]^{g_{ij}\in\text{GL}_r(\mathbb{C})}\\
U_{ij}\ar[r]^-{\varphi_j}&U_{ij}
\times\mathbb{C}^r
}
$$
giving maps in $\text{GL}_1(\mathbb{C})$.
Recall that $\Lambda^r$ is functorial on
linear maps
$g_{ij}\in \text{GL}_n(\mathbb{C})$ because
it is given by the determinant of $g$,
and the determinant is multiplicative.
We thus obtain the cocycle condition on the
maps $g_{ij}$
after applying $\Lambda^r$.
\end{example}

\begin{definition}
\label{definition-connection-vector-bundle}
Let $E$ be a vector bundle.
A {\it connection} on $E$ is an
$\mathbb{R}$-linear map
$\nabla:H^0(E)\to H^0(E) \otimes \Omega^1(M)$
satisfying the identity
$\nabla(fs)=dfs + f\nabla s$.
\end{definition}

\begin{definition}
\label{definition-hermitian-connection}
\begin{reference}
\cite[Definition 4.2.9]{huc}
\end{reference}
Let $(E,h)$ be a hermitian vector bundle.
A connection $\nabla$ on $E$ is called
{\it hermitian} if
$d(h(s,s'))=h(\nabla s,s')+h(s,\nabla s')$.
\end{definition}

\noindent
Perhaps this notion defined on a seminar
talk by Beatrice Brienza is the same:
let $(M,I,g)$ be a hermitian manifold. A
linear connection $\nabla$ on $TM$ is
{\it hermitian} if
$\nabla g=0$ and $\nabla I =0$.

\begin{definition}
\label{definition-holomorphic-connection}
A connection $\nabla$ on a vector bundle $E$
is said to be {\it compatible with the
complex structure $I$} if
$\pi^{0,1}\nabla s=\bar\partial s$
where $\pi^{0,1}:\mathcal{A}^1(M)
\to \mathcal{A}^{0,1}(M)$ is the projection.
\end{definition}

\begin{proposition}
\label{proposition-chern-connection}
On a holomorphic vector bundle
equipped with a hermitian metric there exists
a unique hermitian connection
compatible with the complex structure
called the {\it Chern connection}.
\end{proposition}

\medskip\noindent
Now we discuss line bundles on
$\mathbb{P}^n$.

\begin{definition}
\label{definition-line-bundles-in-pn}
$$
\mathcal{O}(1):=\mathcal{O}(-1)^\vee,\qquad
\mathcal{O}(d):=\bigotimes_d \mathcal{O}(1).
$$
\end{definition}

\noindent
To show that the sections of $\mathcal{O}(d)$
are degree-$d$ homogeneous polynomials we
first
show that the fibers are degree-$d$
homogeneous functions
(not necessarily polynomials)
and then pass to the global situation.

\begin{lemma}
\label{lemma-fiber-od-homog-degre-d-funct}
The fiber $\mathcal{O}(d)_p$ is in bijection
with
degree-$d$ homogeneous functions on $p$.
\end{lemma}

\begin{proof}
So far we can only say that the fibers of
$\mathcal{O}(-1)\otimes \mathcal{O}(1)$ are
the fibes of $\mathcal{O}$.
This means that for a fixed $f \in
\mathcal{O}(1)_p$
we have a linear functional $z \mapsto
f\otimes z$
on $\mathcal{O}(-1)_p$. Any linear functional
on $\mathcal{O}(-1)_p$
can be realised this way.

To generalize this for degree $d$ we do as
follows.
First note that
$$
\Hom(\mathcal{O}(-1)_p,\mathbb{C})
\otimes\Hom(\mathcal{O}(-1)_p,\mathbb{C})
\cong
\Hom(\mathcal{O}(-1)_p\otimes\mathcal{O}(-1)
_p,\mathbb{C}).
$$
(This is not true for general modules,
but in this case it is using a basis.)

By the universal property of tensor product
(see Commutative Algebra Theorem
\cite{theorem-unive-prope-tenso-produ})
we know that
$\Hom(\mathcal{O}(-1)_p \otimes … \otimes
\mathcal{O}(-1)_p,\mathbb{C})$
is the set of multilinear maps
$\text{Lin}(\mathcal{O}(-1)_p
\times…\times\mathcal{O}(-1)_p,\mathbb{C})$.
Any such map $L$ gives a degree-$d$ function
(a polynomial) $P$ by
$P(x):=L(x,…,x)$.
In fact, since the fiber $\mathcal{O}(-1)_p$
is one-dimensional,
we can reconstruct $L$ from $P$ by defining
the multilinear map on a basis (which is
composed by a single vector $x$)
by $L(x,…,x):=P(x)$. This independs on the
choice
of basis since if $x'=\lambda x$ and we
define $L'(x',…,x'):=P(x')$,
We see that $\lambda^d L'(x,…,x)=L'(x',…,x')=
=P(x')=P(\lambda x)=\lambda^d P(x)=\lambda^d
L(x,…,x)$.
\end{proof}

\noindent
To pass from the local situation in Lemma
\ref{lemma-fibers-of-od}
to a section we do as follows.

\begin{exercise}
\label{exercise-polynomials-are-sections}
\begin{reference}
\cite[Exercise 2.2.8]{huc}, and
\cite[Lemma 3.30]{lec} for the answer.
\end{reference}
Show that any non-trivial homogeneous
degree-$d$ polynomial $0 \neq f \in
\mathbb{C}[z_0,…,z_n]_d$ can be considered as
a non-trivial section of $\mathcal{O}(d)$ on
$\mathbb{P}^n$.
\end{exercise}

\begin{proof}
Let $P$ be a homogeneous degree-$d$
polynomial. We define a section fiberwise and
then show it is holomorphic. For $p \in
\mathbb{P}^n$ we restrict $P$ to the
1-dimensional vector space that $p$ is in
$\mathbb{C}^{n+1}$. To show this is
holomorphic we can evaluate our section on a
local frame for $\mathcal{O}(d)$, which gives
a holomorphic function since $P$ is a
polynoimal (see \cite[Proposition 3.35]{lec}.
\end{proof}

\begin{theorem}
\label{theorem-secti-od-degre-d-homog-poly}
Sections of $\mathcal{O}(d)$ are in
correspondence with homogeneous degree-$d$
polynomials.
\end{theorem}

\begin{proof}
It only remains to prove the direct
implication.
Let $s \in H^0(\mathcal{O}(d))$.
Then $s(p) \in \mathcal{O}(d)_p$ is a
homogeneous degree-$d$ function $f_p$
by Lemma
\ref{lemma-fiber-od-homog-degre-d-funct}
defined on the vector space $p$.

Now define a holomorphic homogeneous
degree-$d$ function $f$ on
$\mathbb{C}^{n+1}\setminus\{0\}$ as follows.
For any $z \in
\mathbb{C}^{n+1}\setminus\{0\}$ we can take
$p$
to be the line passing through $z$
and assign the number $f_p(z)$
(since $z \in p$).

Finally we show that $f$ must be a polynomial
by considerations on its Taylor series
after extending to $\mathbb{C}^{n+1}$ via
Hartog's theorem
(see \cite[Theorem 3.36]{lec}).
\end{proof}


\begin{exercise}
\label{exercise-line-bundle-trivial-iff}
\begin{reference}
\cite[Exercise 2.2.5]{huc}
\end{reference}
Let $L$ be a holomorphic line bundle on a
compact complex manifold $X$.
Show that $L$ is trivial if and only if $L$
and its dual $L^\vee$
admit non-trivial global sections.
\end{exercise}

\begin{proof}
Direct implication is easy: trivial bundle
has 1-dimensional space of sections.
For the converse choose nonvtrivial sections
$s$ of $L$
and $s^\vee$ of $L^\vee$.
Then $s \otimes s^\vee$ which is a section
of the trivial bundle.
If such a section has a zero, then it must be
zero throughout $X$.
\end{proof}

\medskip\noindent
Now let's discuss Serre's duality.
It is a consequence of Hodge theorem,
which says that Hodge star operator descends
to cohomology.
In fact, the general result is

\begin{theorem}[Hodge]
\label{theorem-hodge-star-descends}$$
*: H^{p,q}(X) \longrightarrow H^{n-p,n-q}(X)
$$
is an isomorphism.
\end{theorem}

\noindent
However, we shall not use this isomorphism
directly for Serre duality, see \cite[Remark
4.1.17]{huc}. What we do instead is show that
the pairing of forms given by integration is
nondegenerate. The upshot is that the proof
of nondegeneracy is given by the Hodge star
operator in cohomology.

\begin{proposition}[Serre duality]
\label{proposition-serre-duality}
\begin{reference}
\cite[Proposition 4.1.15]{huc}
\end{reference}
Let $X$ be a compact complex manifold.
For any holomorphic vector bundle $E$ on $X$,
the natural pairing
$$
H^{p,q}(X,E) \times H^{n-p,n-q}(X,E) \to
\mathbb{C}
$$
is nondegenerate.
\end{proposition}

\begin{proof}
We must show that for any $\alpha \in
H^{p,q}(X,E)$ there exists $\beta \in
H^{n-p,n-q}(X,E)$ such that
$$
\int_X\alpha\wedge \beta\neq 0.
$$
Choose
$$
\beta=\overline{*}\beta
$$
where $\overline{*}$ is the conjugate of the
Hodge star operator in cohomology.
\end{proof}

\noindent
Of course, this proposition lacks context in
these notes; we should properly extend Hodge
theory to sections of arbitrary sections of
vector bundles; but it looks as if the theory
for sections of the cotangent algebra
generalizes without the need of introducing
any new ideas.

We put $p=0$ to obtain
$$
H^{0,q}(X)=H^q(X,\Omega^0)=H^q(X,\mathcal{O})
$$
and
$$
H^{n,n-q}(X)=H^{n-q}(X,\Omega^n).
$$

\section{Line bundles and divisors}
\label{section-line-bundles}

\noindent
A simple way of associating a line bundle to
an irreducible
hypersurface $Y\subset X$ is by noting that
by Weierstrass Division Theorem
\ref{complex-analysis-theorem-%
weierstrass-division-theorem}
the ideal sheaf of $Y$ is locally generated,
i.e. there is an open cover $U_i$ of $X$
such that $\mathcal{I}(U_i)=(f_i)$.
This means that at intersections we have
$f_j \in (f_i)$ and $f_i\in (f_j)$,
which says that $f_i$ and $f_j$ are
associates,
i.e. there is a unit $u_{ij} \in
\mathcal{O}^*_X(U)$
(the units in the ring of holomorphic
functions
are nowhere vanishing functions)
such that $f_i=u_{ij}f_j$.
Thus $u_{ij}:U_i\cap U_j \to
\text{GL}(1,\mathbb{C})$
define a cocycle in $H^1(X,\mathcal{O}^*_X) =
\text{Pic}(X)$.

\begin{definition}
\label{definition-divisor}
A {\it Weil divisor} is a formal sum
$$
D=\sum a_i Y_i
$$
where $a_i\in \mathbb{Z}$ and $Y_i$ are
irreducible hypersurfaces.
\end{definition}

\noindent
An irreducible hypersurface $Y$ clearly
defines a Weil divisor,
and so does a reducible reduced hypersurface
with irreducible
components $Y_i$ putting $a_i=1$.
To allow for coefficients $>1$ we may admit
nonreduced hypersurfaces,
so that we have larger multiplicites.

This can also be presented as follows.
Given a holomorphic function $g$, we define
{\it order of vanishing of $g$ at $P \in Y$}
to be the largest number $\text{ord}_{P,Y}g$
such that
$$
g=f_i^{\text{ord}_{P,Y}g}h,\qquad h \in
\mathcal{O}^*_{X,P},
$$
where $f_i$ is an irreducible function
vanishing along $Y$ near $P$,
i.e. a generator of the local ring of $Y$.
Further, by coherence/preservation relative
primes in nearby stalks
(see \cite[p. 696]{gri})
the order of $g$ is constant for any $P \in
Y$.

Thus for any holomorphic function $g$
we can define the divisor
$$
\text{div}(g):=\sum_{Y}\text{ord}_Y(g)Y
$$
where the sum ranges in all irreducible
hypersurfaces of $X$.

Similarly, for a meromorphic function
$g=g_1/g_2$
we define the order of vanishing of $g$ at
$Y$ to be
$$
\text{ord}_Y(g)=\text{ord}_Y(g_1)-\text{ord}
_Y(g_2).
$$
The associated divisor to a meromorphic
function
is defined identically as that of a
holomorphic function,
and we observe that there might appear
negative coefficients.
Divisors obtained in this way from
meromorphic functions
are called {\it principal divisors}.

This shows how to assign a divisor from any
meromorphic function,
i.e. it's a map
$$
H^0(X,\mathcal{M}^*) \to \text{Div}(X).
$$
Note that if $g$ is in fact a nowhere
vanishing meromorphic or
holomorphic function,
the corresponding divisor is trivial,
i.e. it is in the kernel of our map.
In fact,
$$
H^0(X,\mathcal{M}^*/\mathcal{O}^*) \simeq
\text{Div}(X).
$$
Indeed, given a divisor we can assign a
section of
$\mathcal{M}^*/\mathcal{O}_X^*$ as follows.
If the divisor is an irreducible hypersurface
locally given by $V(f_i)$
we assign the meromorphic function given
locally as $f_i$.
If the divisor has several irreducible
connected components
we consider the product of the irreducible
generators
at every point.

\medskip\noindent
It remains to connect $\text{Div}(X)$ and
$\text{Pic}(X)$.
The idea is that a section $f$ of
$\mathcal{M}^*/\mathcal{O}_X^*$
defines a cocycle in $H^1(\mathcal{O}_X^*)$.
Indeed, $f$ is given, by definition of
quotient sheaf (?),
by the data of an open cover $U_i$ and
meromorphic functions
$$
f_i=\frac{g_i}{h_i} \in \mathcal{M}^*
$$
which coincide at the intersections $U_i \cap
U_j$,
that is,
$$
f_i|_{U_i\cap U_j} = f_j|_{U_i\cap U_j}.
$$
But since we are in the quotient ring,
equality means that
$$
f_i f_j^{-1} \in \mathcal{O}_X^*.
$$
This defines a cocycle, and thus a map
$$
\text{Div}(X) \to \text{Pic}(X).
$$
This map can also be though of as induced
from
$$
\xymatrix{
0\ar[r]
&\mathcal{O}^*_X\ar[r]
&\mathcal{M}^*\ar[r]
&\mathcal{M}^*/\mathcal{O}_X^*\ar[r]
&0
}
$$
via the cohomology long exact sequence giving
$$
\xymatrix{
\cdots\ar[r]
& H^0(\mathcal{M}^*_X)\ar[r]
&H^0(\mathcal{M}^*_X/\mathcal{O}^*_X)\ar[r]
&H^1(\mathcal{O}_X^*)\ar[r]
&\cdots
}
$$
Indeed, the kernel of the map $\text{Div}(X)
\to \text{Pic}(X)$
is given by the principal divisors,
i.e. those coming from meromorphic sections.

In general there is no way to ensure this map
is an isomorphism.

\medskip\noindent
Now we discuss how to obtain submanifolds
from sections.
We will see that $\text{Div}(X)$
is isomorphic to the group of line bundles
equipped with a meromorphic section.


Bear in mind the following observation for a
rank-$r$
vector bundle $E \to X$.
Suppose $(\varphi_i,U_i)$ are local
trivializations of $E$,
i.e. $\varphi_i:\pi^{-1}(U_i)
\xrightarrow{\simeq}U_i \times \mathbb{C}^k$.
Take a local section $s :U_i \to
\pi^{-1}(U_i)$.
Take a local basis $e_i$ of $\mathbb{C}^k$.
Then
$$
s_i(x)=f_i^k(x)\varphi_i^{-1}(x,e_k)
$$
for some functions $f_i^k:U_i \to
\mathbb{C}$.

At an intersection $U_i \cap U_j$,
we have
$$
s_i(x)=f_i^k(x)\varphi_i^{-1}(x,e_k)
=
s_j(x)=f_j^k(x)\varphi_j(x,e_k).
$$
Applying $\varphi_i$ we see that
$$
f_i^k e_k = f_j^k g_{ij}(e_k).
$$

In the case of a line bundle this just says
that sections
are given by local functions $f_i:U_i \to
\mathbb{C}$ satisfying
$$
f_i=g_{ij}f_j
$$
where $g_{ij}$ is the cocycle defining the
line bundle.

Let $L$ be a line bundle.
A {\it meromorphic section} of $L$ is a
section
of $\mathcal{O}(L)
\otimes_\mathcal{O}\mathcal{M}$,
which is given by a collection of meromorphic
functions
$s_i \in \mathcal{M}(U_i)$
such that
$$
s_i=g_{ij}s_j.
$$
That is, sections of the line bundle viewed
as a module over $\mathcal{O}$
with coefficients that can be meromorphic
functions.

The orders of the coefficient meromorphic
functions
is well defined along hypersurfaces of $X$,
and thus
we can define the
{\it divisor $\text{div}(s)$ associated to
the section $s$ of $L$}
by
$$
\text{div}(s)
:= \sum_Y\text{ord}_Y(s)Y.
$$

Conversely, given any divisor $D$
we can assign a line bundle $L$ as before,
namely if the divisor is given by the
meromorphic functions $f_i$
we put $g_{ij}=f_i/f_j$.
But now we observe
the existence of a distinguished section $s$
such that $\text{div}(s)=L$.
Indeed, we note that
$$
f_i=f_if_j^{-1}f_j=g_{ij}f_j,
$$
which is a section of $L$,
and in fact (check…) $\text{div}(s)=L$.

We conclude that $L$ is the line bundle
associated to a divisor
if and only if it has a meromorphic section
not identically zero.
It is the line bundle associated to an
effective divisor
if and only if it has a holomorphic section
not identically zero
--- whose zero set is the divisor.

\medskip\noindent
Upshot. We have more than the submanifold
associated to a section:
we have the divisor associated to a
meromorphic section,
which gains to becoming less geometrically
intuitive
the virtue of being an element of a group.

\medskip\noindent
Now we turn to line systems.
Take a divisor
$$
D=\sum a_iV_i.
$$
The divisor
$$
D+(f)
$$
is linearly equivalent to $D$ --- their
difference is principal.
If
$$
D+(f)\geq 0,
$$
then $f$ must be holomorphic away from $V_i$,
and on $V_i$ it can have poles with the
condition
$$
\text{ord}_{V_i}(f) \geq  -a_i.
$$
Denote $\mathcal{L}(D)$ the space of such
functions.

Note that if $s_0$ is a section such that
$(s_0)=D$, then multiplication by $s_0$ gives
an identification
$$
\mathcal{L}(D) \xrightarrow{\otimes s_0}
H^0(X,L).
$$

\begin{definition}
\label{definition-line-system}
We denote by
$$
|D|\subset \text{Div}(X)
$$
the set of all effective divisors linearly
equivalent to $D$.
\end{definition}

\noindent
$\mathcal{L}(D)=H^0(X,L)$ is almost the same
thing as $|D|$.
At least in the compact case,
two functions in $f,f' \in \mathcal{L}(D)$
give the same divisor when $f=\lambda f'$
for some constant function $\lambda \in
\mathbb{C}$.

Thus
$$
|D|=\mathbb{P}H^0(X,L).
$$
This is called a {\it complete linear
system}.
More generally, a {\it linear system} is a
linear
subspace of $\mathbb{P}H^0(X,L)$.
The {\it base locus} of $|D|$
is the set of points where all the global
sections
generating $H^0(X,L)$ vanish.
We will see in Section
\ref{section-kodaira-embedding}
the relationship between line systems and
embeddings to projective space.

\begin{theorem}[Bertini]
\label{theorem-bertini}
The generic divisor of a linear system
is smooth away from the base locus.
\end{theorem}

\medskip\noindent
Now let's introduce the first adjunction
formula,
which says that when a divisor is a smooth
hypersurface,
it's associated line bundle is its
holomorphic normal bundle.

\begin{theorem}[Adjunction formula I]
\label{theorem-adjunction-formula-i}
\begin{reference}
\cite[p.146]{GH}
\end{reference}
If $X$ is a smooth hypersurface of a smooth
complex manifold $Y$, then
\begin{equation}
\label{equation-adjunction-i}
\mathcal{O}_Y(X)|_{X} \simeq
\mathcal{N}_{X/Y},
\end{equation}
where $\mathcal{N}_{X/Y}:=T_Y'/T_X'$.
\end{theorem}

\begin{proof}
If $X$ is locally defined
by the vanishing of $f_i$ on an open cover
$U_i$ of $Y$,
then $\mathcal{O}_Y(X)$ is given by the
cocycle $f_{ij}:=f_if_j$.

To prove Equation \ref{equation-adjunction-i}
it's enough to show that
$\mathcal{O}_Y(X)|_{X}\otimes
\mathcal{N}^\vee_{X/Y}$
is trivial, which can be done by constructing
a nonvanishing section.
Such a section is given by
$$
df_i=d(f_{ij}f_j)=df_{ij}\cdot
f_j+f_{ij}df_j,
$$
since the first term vanishes along $X$
(because $f_j$ does)
and the remaining terms define precisely a
section
of ${T'}^\vee_Y/{T'}^\vee_X$,
the conormal bundle ---
indeed, sections of such a bundle
are holomorphic differencials vanishing along
$X$.
\end{proof}

\noindent
Now let's find the canonical bundle of
$\mathbb{P}^n$.
By the above discussion, we need the inverse
map of
$$
\text{Div}(\mathbb{P}^n) \to \{L \in
\text{Pic}(X):
\text{$L$ has a nonzero meromorphic
section}\},
$$
that is, we need to find a nonzero
meromorphic section $s$
of the canonical bundle $\Omega^{n,0}=K$.
Then the divisor will be given by
$$
\sum_Y \text{ord}_Ys [Y].
$$
A meromorphic section is given by local
sections $s_i$
so that
$$
s_i=g_{ij}s_j
$$
where $g_{ij}$ is the cocycle of $K.$

Consider in the canonical chart $z_0\neq 0$
of $\mathbb{P}^n$
the local meromorphic section of $K$
$$
s_0=
\frac{dw_1}{w_1}\wedge\frac{dw_2}{w_2}\wedge…
\wedge\frac{dw_n}{w_n}
$$
where $w_i=z_i/z_0$ are the coordinate
functions in $z_0\neq 0$.

In any chart $z_j\neq 0$ we obtain the local
expression
$$
(-1)^j\frac{dw_0'}{w_0'}
\wedge…
\wedge \widehat{\frac{dw_j'}{w_j'}}
\wedge…
\wedge \frac{dw_n'}{w_n'}
$$
according to the canonical change of
coordinates, which says
$$
w_i=\frac{w_i'}{w'0}, i\neq j,\qquad
w_j=\frac{1}{w_0'}.
$$

The point is that we constructed a
meromorphic section which has
a single pole along each hyperplane $z_i=0$,
which means
$$
K=-(n+1)H.
$$
\medskip\noindent
The next adjunction formula is the key to
computing
the canonical bundle of submanifolds.
It is a consequence of this linear algebra
fact:

\begin{lemma}
\label{lemma-detrminant-of-ses}
If
$$
\xymatrix{
0\ar[r]
&E'\ar[r]
&E\ar[r]
&E''\ar[r]
&0
}
$$
is a short exact sequence of vector bundles
(of possibly different ranks),
then
$$
\det(E)\cong \det(E')\otimes \det(E'').
$$
\end{lemma}

\begin{proof}
First notice that each of the three
determinant bundles gives fiberwise the
top-exterior
power of the corresponding vector spaces,
thus varying the rank of the exterior power
taken.
Then define a map
$$
\begin{aligned}
x_1 \wedge…\wedge x_{\text{rk}E'}
&\otimes x_{\text{rk}E'+1}\wedge…\wedge
x_{\text{rk}E'+\text{rk}E''}\\
&\qquad \qquad \mapsto x_1 \wedge…\wedge
x_{\text{rk}E'}
\wedge x'_{\text{rk}E'+1}\wedge…\wedge
x'_{\text{rk}E'+\text{rk}E''}
\end{aligned}
$$
where $x'_{\text{rk}E'+i}$ maps to
$x_{\text{rk}E'+i}$
under the surjection given by the exact
sequence.

It's clear that such a map is injective:
if any monomial in the image is zero, then
one of the
factors must be zero, and then the element
we began with also vanishes.
Being two vector spaces of dimension 1,
the map must also be surjective.
Check this map is well-defined!
This gives an isomorphism in every fiber.

In fact the isomorphism above acts like 1.
To show that this glues to a bundle
isomorphism
we have to check some cocycle condition
as in Proposition
\ref{proposition-isomo-vecto-bundl}.
In this case the transformation functions
$A_i$
are just 1 by definition of the morphism.
The cocycle condition follows from the fact
that
any short exact sequence splits,
so on $U_{ij}$ we have the isomorphism
$E(U_{ij}) \simeq E'(U_{ij})\oplus
E''(U_{ij})$.
Recall that the cocycle of the direct sum of
bundles is given by
$\begin{pmatrix}h'_{ij}&0\\
0&h''_{ij}\end{pmatrix}$
so that when taking determinants
we obtain that
$g_{ij}=h'_{ij}h''_{ij}:=h_{ij}$
(since taking determinant bundle is taking
determinant of
the transition functions).
Thus the transition functions of the two
bundles in question
actually coincide,
and the cocycle condition in Proposition
\ref{proposition-isomo-vecto-bundl}
is the trivial equality $g_{ij}=h_{ij}$.
\end{proof}

\noindent
Here's a cleaner version of this proof:

\begin{proof}
Choose an open cover $\{U_i\}$ on which the
sequence splits, so that
\[
E|_{U_i}\simeq E'|_{U_i}\oplus E''|_{U_i}.
\]
Pick local frames $(e'_{i,1},\dots,e'_{i,r})$
for $E'|_{U_i}$ and
$(e''_{i,1},\dots,e''_{i,s})$ for
$E''|_{U_i}$, and use the splitting to obtain
a frame $(e'_{i,\bullet},e''_{i,\bullet})$ of
$E|_{U_i}$.

On overlaps $U_{ij}$ the transition matrix of
$E$ in these adapted frames has
block diagonal form
\[
g^E_{ij}=
\begin{pmatrix}
g'_{ij} & 0\\
0 & g''_{ij}
\end{pmatrix},
\]
hence
\[
\det(g^E_{ij})=\det(g'_{ij})\,\det(g''_{ij}).
\]
Thus the transition functions of $\det(E)$
coincide with those of
$\det(E')\otimes\det(E'')$, giving an
isomorphism of line bundles. In these
frames it is given by
\[
(e'_{i,1}\wedge\cdots\wedge e'_{i,r})\otimes
(e''_{i,1}\wedge\cdots\wedge e''_{i,s})
\longmapsto
e'_{i,1}\wedge\cdots\wedge e'_{i,r}\wedge
e''_{i,1}\wedge\cdots\wedge e''_{i,s}.
\]

If a different local splitting is chosen, the
adapted frames differ by a
matrix of the form
\[
\begin{pmatrix}
I & B\\
0 & I
\end{pmatrix},
\]
whose determinant is $1$. Therefore the
induced map on determinant lines is
unchanged, showing that the above isomorphism
is independent of the splitting
and hence canonical.
\end{proof}

\begin{theorem}[Adjunction formula II]
\label{theorem-adjunction-ii}
Let $X$ be a complex submanifold of a complex
manifold $Y$.
Then
\begin{equation}
\label{equation-adjunction-ii}
(K_Y\otimes \mathcal{O}_Y(X))|_{X}\cong K_X.
\end{equation}
\end{theorem}

\begin{proof}
From the holomorphic tangent bundle exact
sequence we get
$$
\xymatrix{
0\ar[r]
&T'_X\ar[r]
&T'_Y|_{X}\ar[r]
&N_{X/Y}:=T'_Y/T'_X\ar[r]
&0
}
$$
Taking determinant of this exact sequence we
obtain the result.
\end{proof}


\medskip\noindent
Now we turn to Chern classes of line bundles.

Slogan: the first Chern class of a line
bundle
is the Poincaré dual of the divisor's
fundamental class
and the cohomology class of the curvature
form
of any connection of $L$.
It is an integral $(1,1)$-curvature form.

The {\it curvature form} is the differential
2-form
that represents the curvature operator with
respect to a given connection
on the line bundle.
This can also be defined for affine vector
bundles of any rank.

\begin{proposition}
\label{proposition-first-chern-class}
\begin{reference}
\cite[p. 141]{gri}
\end{reference}
\begin{enumerate}
\item For any line bundle $L$ with curvature
form $\Theta$,
$$
c_1(L)=\left[\frac{\sqrt{-1}}{2\pi}
\Theta\right]
\in H^2_{dR}(X).
$$

\item If $L=[D]$ for some $D \in
\text{Div}(M)$,
$c_1(L)$ is the Poincaré dual of the
fundamental class of $D$.
\end{enumerate}
\end{proposition}

\medskip\noindent
Finally let's just put some more exercises
here.

\begin{exercise}[Curves on surfaces and the
Hodge index theorem]
\label{exercise-curve-surfa-hodge-index}
\begin{reference}
\cite[Exercise 2, Chapter 11]{voi}
\end{reference}
Let $S$ be a smooth projective surface and
let $\phi:S \to \mathbb{P}^n$
be a holomorphic map satisfying the following
properties:
\begin{itemize}
\item There exists a curve $C \subset S$
such that $\phi(C)$ is a finite set of
points.
(Here $C$ is a not necessarily connected
finite union
of irreducible hypersurfaces $C_i$ which are
not necessarily smooth.)

\item There exists an open set of $S$ where
$\phi$ is an embedding.
\end{itemize}
\begin{enumerate}
\item Show that the class
$h=\phi^*c_1(\mathcal{O}_{\mathbb{P}^n}(1))$
satisfies $\left\langle{}h,h\right\rangle{} >0$ where
$\left\langle{}\cdot,\cdot\right\rangle{}$ is the
intersection form on $H^2(S,\mathbb{Z})$.

\item Show that $\left\langle{}[C_i],h\right\rangle{}=0$ for
any $i$.
Deduce from the Hodge index theorem that the
intersection form
$\left\langle{}\cdot,\cdot\right\rangle{}$ is negative
definite on the subgroup
of $H^2(S,\mathbb{Z})$ generated by the
$[C_i]$'s.

\item Show that $\left\langle{}C_i,C_j\right\rangle{}\geq 0$
for $i\neq j$.
Deduce then from the above that the $[C_i]$'s
are independent in $H^2(S,\mathbb{Z})$.
\end{enumerate}
\end{exercise}

\begin{proof}
\begin{enumerate}
\item Here's an outline of how to prove this
result.
First notice that $h$ is the class of a line
bundle
since $\varphi^*c_1=c_1\varphi^*$.
Thus $\left\langle{}h,h\right\rangle{}$ is the intersection
number
of two divisors in the class $h$.

To compute their intersection number
we use that $\varphi$ is an embedding on an
open set of $S$,
say $U \simeq \varphi(U) \subset
\mathbb{P}^n$.
To compute the intersection number we take
two generic hyperplanes $H_1$ and $H_2$
in $\mathbb{P}^n$ and intersect them with
$\phi(U)$.
The dimension of the resulting submanifold is
zero.
Thus the intersection is a finite number of
points.

\item Since $C_i$ is mapped to a set of
points in $\mathbb{P}^n$,
we compute the intersection number as above,
which gives zero
for generic choice of $H_1,H_2$. Hodge index
theorem tells us that
the signature of $\left\langle{}\cdot,\cdot\right\rangle{}$
is $(1,\rho-1)$ on
$H^2(S,\mathbb{Z})$. Since $h$ is positive
and $C_i$ is orthogonal to $h$,
it's necessary that $C_i$ are negative
vectors.

\item This is by definition/construction of
the intersection form on a surface:
for two irreducible divisors of the surface,
the intersection number is the number of
intersection points
of the divisors counted with multiplicity.
Such a number cannot be negative.
The fact that the $C_i$'s are independent
follows from a linear-algebra
argument.
\end{enumerate}
\end{proof}

\section{Kodaira Vanishing theorem}
\label{section-kodaira-vanishing}

\noindent
In \cite[p. 150]{gri}
we see that the curvature form of the
tautological bundle
is given by $\partial\bar\partial \text{log}
|z|^2$
where $z$ is a nonzero section.
Not giving further details for now,
this implies that the Chern class of the
hyperplane bundle
is the Fubini-Study form in $\mathbb{P}^n$.
We already discussed that this form is
positive
in the sense of Definition
\ref{definition-positive-11form}.

\begin{definition}
\label{definition-positive-line-bundle}
\begin{reference}
\cite[p. 148]{gri}
\end{reference}
Let $M$ be a compact Kähler manifold.
A line bundle $L \to M$ is {\it positive}
if there exists a metric on $L$ with
curvature form $\Theta$
such that $(\sqrt{-1}/2\pi)\Theta$ is a
positive $(1,1)$-form.
\end{definition}

\begin{theorem}[Kodaira-Nakano Vanishing
Theorem]
\label{theorem-kodaira-nakano-vanishing}
\begin{reference}
\cite[p. 154]{gri}
\end{reference}
[Let $M$ be a compact Kähler manifold of
complex dimension $n$.]
If $L \to M$ is a positive line bundle,
then
$$
H^q(M,\Omega^p(L)) = 0,\qquad
p+q>n=\dim_\mathbb{C} M.
$$
\end{theorem}

\noindent
Here's the same result as it appears in
\cite[sea], where it is not proved.

\begin{theorem}[Kodaira Vanishing]
\label{theorem-Kodaira-vanishing-sea}
\begin{reference}
\cite[21.5.8]{sea}
\end{reference}
Suppose $k$ is a field of characteristic 0,
and $X$ is a
smooth projective $k$-variety. Then for any
ample invertible sheaf $L$,
$H^{i}(X,K_X \otimes L)=0$ for $i>0$.
\end{theorem}

\noindent
For completeness we put another vanishing
theorem
concerning higher-rank vector bundles
(probably due to Serre though Griffith-Harris
don't say so).
The idea is that twisting an ample line
bundle with itself
enough times and tensoring with any vector
bundle $E$
will kill positive cohomology.

\begin{theorem}
\label{theorem-vanis-highe-rank-verct-bund}
\begin{reference}
\cite[p. 159]{gri}
\end{reference}
Let $M$ be a compact, complex manifold and $L
\to M$
a positive line bundle.
Then for any holomorphic vector bundle $E$,
there exists $\mu$ such that
$$
H^q(M,L^\mu \otimes E)=0\qquad \text{for
}q>0, \mu\geq \mu_0.
$$

\end{theorem}

\section{Lefschetz hyperplane sections
theorem}
\label{section-lefschetz-hyperplane}

Lefschetz hyperplane theorem holds for
any $n$-dimensional compact complex manifold
$M$
and $V \subset M$ a smooth hypersurface with
$L=[V]$ positive.
An example of this is when $M$ is projective
and $V$ is the pullback of the hyperplane
bundle
(which, as explained in Section
\ref{section-kodaira-vanishing},
is positive since its curvature form is
Fubini-Study),
that is, $V=M \cap H$ for a generic
hyperplane $H \subset \mathbb{P}^N$.

\begin{theorem}[Lefschetz hyperplane
sections]
\label{theorem-lefschetz-hyperplane}
The map
$$
H^q(M,\mathbb{Q}) \to H^q(V,\mathbb{Q})
$$
induced by the inclusion $V
\xymatrix{\ar@{^{(}->}[r]&} M$
is an isomorphism for $q \leq  n-2$
and injective for $q=n-1$.
\end{theorem}

\begin{proof}
Looks accessible! ``Purely topological''.
\end{proof}

\section{Complex varieties}
\label{section-complex-varieties}

For a definition of sheaf
see Algebraic Geometry Definition
\ref{algebraic-geometry-definition-sheaf}.

\begin{lemma}
\label{lemma-O_X-is-a-sheaf}
Let $X$ be a complex manifold of complex
dimension $n$.

The functor
$$
\begin{aligned}
\mathcal{O}_X: \mathit{Open}_X^{\text{op}}
&\longrightarrow \mathit{Set}
\\
U &\longmapsto \mathcal{O}_X(U):=\{f:U\to
\mathbb{C}^n:
f \text{ is holomorphic}\}\\
i:V \hookrightarrow V &\longmapsto
\mathcal{O}_X(i):\mathcal{O}_X(U) \to
\mathcal{O}_X(V)
\end{aligned}
$$
where $\mathcal{O}_X(i)$ is given by
restriction,
is a sheaf.
\end{lemma}

\begin{proof}
$\mathcal{O}_X$ is a presheaf,
i.e. a functor,
by definition: restriction
maps the identity to the identity
and preserves compositions.

To check the condition of being a sheaf let
$U_i$ be an open cover of some open set $U$
of $X$
and $f_i:U_i \to X$ sections of
$\mathcal{O}_X(U_i)$,
that is, holomorphic functions, which
coincide pairwise in intersections.
Then it is obvious that the function $f:U \to
\mathbb{C}$,
$x \mapsto f_i(x)$ for any $U_i \ni x$
is well defined. It is also holomorphic
since holomorphicity is a condition defined
on
open sets.
\end{proof}

\medskip\noindent
A meromorphic function is a function that
can be written locally as a quotient of two
holomorphic functions:

\begin{definition}
\label{definition-meromorphic-functions}
\begin{reference}
\cite[Definition 2.1.8]{huc}
\end{reference}
A {\it meromorphic function} on a complex
manifold $X$
is a map
$$
f:X \to \coprod_{x \in
X}\text{Frac}(\mathcal{O}_{X,x})
$$
which associates to any $x\in X$
an element $f_x \in
\text{Frac}(\mathcal{O}_{X,x})$
such that for any $x_0 \in X$
there exists $U\ni x$ open and $f,g
\in\mathcal{O}(U)$
such that $f_x=g/h$ for all $x \in U$.
The {\it sheaf of meromorphic functions} is
$\mathcal{K}_X$.
\end{definition}


\medskip\noindent
Now we pass to the theory of analytic
subvarieties.

\begin{definition}
\label{definition-analytic-subvariety}
An {\it analytic set} or {\it subvariety} $Y$
is a subset of a complex manifold $X$
given locally as the zeros of a finite
collection of holomorphic functions.
\end{definition}

\begin{lemma}
\label{lemma-structure-sheaf-of-subvariety}
Let $Y \subset X$ be an analytic subvariety
of a complex manifold $X$
equipped with the subspace topology.
Then
$$
\begin{aligned}
\mathcal{O}_Y: \mathit{Open}_Y^{\text{op}}
&\longrightarrow \mathit{Set} \\
U\cap Y &\longmapsto \{f|_{U \cap Y}:f \in
\mathcal{O}_X(U)\}
\end{aligned}
$$

\noindent
is a sheaf.
\end{lemma}

\begin{proof}
This is a presheaf for the same reasons as
the sheaf $\mathcal{O}_X$ is a presheaf.
The fact that it is a sheaf follows
from the fact that $\mathcal{O}_X$ is a
sheaf and by definition of subspace topology;
namely for any family of functions
defined over an open cover $\{U_i \cap Y\}$
of some open set  $U\cap Y$ we can
reconstruct a function
on $U$ that restricted to $U \cap Y$ is as
desired.
\end{proof}

Recall that any map of sheaves of groups
has well-defined kernel, see Lemma
\ref{algebraic-geometry-lemma-sheaves-%
valued-on-groups-have-kernels}.

\begin{definition}
\label{definition-ideal-sheaf}
Let $Y \subset X$ be an analytic subvariety
of a complex manifold $X$.
The {\it ideal sheaf} of $Y$, denoted
$\mathcal{I}_Y$,
is the kernel
of the restriction map $\mathcal{O}_X \to
\mathcal{O}_Y$.
\end{definition}

\noindent
This just says that $\mathcal{I}_X(U)$ is the
set of functions $f \in \mathcal{O}_X$ that
vanish under the restriction map, i.e.
that vanish at points of $Y$.

Since the restriction map is surjective,
we have the following short exact sequence:
$$
\xymatrix{
0\ar[r]&\mathcal{I}_Y\ar[r]&\mathcal{O}
_X\ar[r]&\mathcal{O}_Y\ar[r]&0
}
$$

\medskip\noindent
Notice that for any open set $U\subset X$,
$\mathcal{I}_Y(U)$ is an ideal of
$\mathcal{O}_X(U)$. (Indeed,
if a function vanishes at $Y$ then its
product
with any other function will also vanish.)

\begin{definition}
\label{definition-irreducible-variety}
An analytic variety $V\subset U\subset
\mathbb{C}^n$ is {\it irreducible} if
$V$ cannot be written as the union of two
distinct analytic varieties
$V_1,V_2\subset U$, both distinct to $V$.
\end{definition}

\medskip\noindent
We will show that for the case of an
irreducible hypersurface
the ideal sheaf is locally generated by a
single element.
This will allow us to construct a line bundle
from the defining functions
of the hypersurface, which is in general not
possible
for an arbitrary smooth hypersurface,
--- although it suffices to ask for
orientability in the
smooth case.

We need two ingredients:
(1) the fact that the stalks $\mathcal{O}_p$
of the
sheaf of holomorphic functions are UFDs,
and (2) the following lemma.

\begin{lemma}
\label{lemma-ideal-irred-hyper-heigh-1}
The ideal
(correction: this should be the stalk)
of an irreducible analytic hypersurface is a
height-1 prime ideal.
\end{lemma}

\begin{proof}
We already know that irreducible varieties
correspond with prime ideals.

To show the ideal $I$ defining the variety
$Y$ has height 1
suppose that $0 \subset \mathfrak{p}
\subseteq I$ is prime.
Then $V(\mathfrak{p})$ is an analytic variety
that contains $Y$.
At regular points of both $V(\mathfrak{p})$
and $Y$,
both are smooth manifolds, but since $Y$ is
of codimension 1 and
$V(\mathfrak{p})$ does not equal all of  $Y$,
we conclude that they coincide.
Thus, in a neighbourhood
of a regular point we have
$\mathfrak{p}=I(V(\mathfrak{p}))=I$ by the
Nullstellensatz.
\end{proof}

\begin{lemma}
\label{lemma-stalk-ideal-sheaf-princ-hyper}
The stalks of the ideal sheaf of
an irreducible hypersurface
are principal rings.
\end{lemma}

\begin{proof}
Since $\mathcal{O}_p$
is a UFD, we pick any element
$f \in \mathcal{I}_p$, and factor it as a
product of irreducibles.
Since $\mathcal{I}$ is prime, some of the
irreducible factors,
say $f_i$, is in $\mathcal{I}_p$.
The ideal generated by $f_i$ is prime since
$\mathcal{O}_p$
js a UFD (see
\href{https://stacks.math.columbia.edu/tag/%
034T}{034T}).
But by Lemma
\ref{lemma-ideal-irred-hyper-heigh-1},
$\mathcal{I}_p$ is a height-1 prime ideal, so
the prime ideal $(f_i)$
must equal all of $\mathcal{I}_p$.
\end{proof}

\begin{lemma}
\label{lemma-ideal-sheaf-irred-hyper-line}
The ideal sheaf $\mathcal{I}$
of the irreducible hypersurface $Y \subset X$
corresponds to a unique line bundle.
\end{lemma}

\begin{proof}
Let $p \in U$ and consider for any stalk
$\mathcal{I}_p$
its generator $[f_p]$,
which has a representative $(f_p,U_p)$.
These form an open cover of $X$,
coincide up to a unit in the intersections as
follows.
These units shall give us a cocycle.

For any $U_i\cap U_j \neq 0$
we have
[This is slightly imprecise.
We may extend $f_j|_{U_i\cap U_j}$ to $U_i$
and $f_i|_{U_i\cap U_j}$ to $U_j$,
and then restrict back to $U_i\cap U_j$.
This works because we can extend holomorphic
functions by their local power series
expansion!],
$f_j\in (f_i)$, so there exist functions
$g_{ij}$ such that
$$
f_i=g_{ij}f_j.
$$
These satisfy the cocycle condition since
$$
f_if_jf_k=g_{ij}f_jg_{jk}f_kg_{ki}f_i.
$$
By Proposition
\ref{proposition-vecto-bundl-cocyc} we have a
line
bundle.
\end{proof}

\noindent
Finally we have concluded that the ideal
sheaf
of a hypersurface is a line bundle. Now we
can take its dual to define:

\begin{definition}
\label{definition-sheaf-assoc-effec-carti}
Let $D:=S$ be an effective Cartier divisor,
that is,
an analytic hypersurface of a complex
manifold.
The line bundle $\mathcal{O}(D)$ is the dual
line bundle
of the ideal sheaf of $D$.
\end{definition}

\section{Chow's theorem}
\label{section-chow-theorem}

\begin{theorem}[Chow]
\label{theorem-chow}
Any analytic subvariety of $\mathbb{C}P^{n}$
is algebraic.
\end{theorem}

\begin{proof}
First we prove it for hypersurfaces.
A hypersurface in $\mathbb{P}^n$ defines a
line bundle
along with a section vanishing on the
hypersurface.
Such a line bundle must be $\mathcal{O}(d)$
for some $d$
by Exercise \ref{exercise-pic-of-pn-is-z}.
Sections of $\mathcal{O}(d)$ are homogeneous
degree-$d$ polynomials
by Theorem
\ref{theorem-secti-od-degre-d-homog-poly}.

For higher codimension the challenge is find
adequate planes
in $\mathbb{P}^n$, so that we can project
$V$, $\dim V=k$
to $\mathbb{P}^{k+1}$ and apply the
hypersurface version.
Pick any point $p\in \mathbb{P}^n$ not in
$V$.
We can choose a $(n-k-1)$-plane not
intersecting $V$ and containing $p$.
Inside such $\mathbb{P}^{n-k-1}$ we can pick
a $\mathbb{P}^{n-k-2}$
not touching $p$.
Now pick a complementary $\mathbb{P}^{k+1}$
in $\mathbb{P}^n$,
i.e. $\mathbb{P}^n=\mathbb{P}^{k+1}\oplus
\mathbb{P}^{n-k-2}$.

Now project $V$ to $\mathbb{P}^{k+1}$ and
apply Chow for hypersurfaces.
We obtain a polynomial $F$ defining the
projection of $V$.
Extend $F$ to $\tilde{F}$ on all
$\mathbb{P}^n$ by the rule
$\tilde{F}(x_0,…,x_n)=F(x_0,…,x_{k+1})$.
Note that since $p$ is not in $V$,
$\tilde{F}$ does not vanish in $p$.
Thus $V$ is the intersection of all such
polynomials $F$,
one for each $p$, which ensures that no
points other than those
of $V$ are in the common zeroes.
\end{proof}

\section{Kodaira Embedding theorem}
\label{section-kodaira-embedding}

\noindent
We introduced line systems in Section
\ref{section-line-bundles}.
Let $L$ be a holomorphic bundle on a complex
manifold $X$.
A point $x \in X$ is
a {\it base point} of $L$ if $s(x)=0$ for all
$s\in H^{0}(X,L)$
and the {\it base locus} $\text{Bs}(L)$
is the set of all base points of $L$.

Let $s_0,\ldots,s_N\in H^{0}(X,L)$ are a
basis
for some subspace $E \subset H^0(X,L)$.
Define a map
$$
\begin{aligned}
\varphi_L: X\setminus\text{Bs}(L)
&\longrightarrow \mathbb{P}^N \\
x &\longmapsto [s_0(x):\ldots:s_N(x)]
\end{aligned}
$$

\noindent
which is well-defined since we excluded
the base locus,
and holomorphic since we chose holomorphic
sections of $L$ for the coordinates.
Further, the pullback of the hyperplane
bundle of $\mathbb{P}^N$ is $L$,
i.e.
$$
\varphi^*_L \mathcal{O}_{\mathbb{P}^N}(1)
\simeq L|_{X\setminus\text{Bs}(L)}
$$
since $s_i=0$ is the equation of a hyperplane
in $\mathbb{P}^N$,
so, it is a section of $H$.
Its pullback under $\varphi$ gives the
section $s_i$
whose associated divisor $(s_i)$
corresponds to the line bundle $L$ ---
the line bundle associated to a divisor
associated to a section of the line bundle
is the line bundle itself.

\medskip\noindent
The question is: when is $\varphi$ an
embedding?

\begin{definition}
\label{definition-ample}
\begin{reference}
\cite[Definition 2.3.28]{huc}
\end{reference}
A line bundle $L$ on a complex manifold is
called {\it ample} if for some $k>0$
and some linear system in $H^{0}(X,L^k)$ the
associated map $\varphi$ is an
embedding.
\end{definition}

\begin{theorem}[Kodaira Embedding]
\label{theorem-Kodaira-embedding}
\begin{reference}
\cite[Theorem 10.12]{lec}
\end{reference}
Suppose $M$ is a compact complex manifold.
A holomorphic line bundle $L\to M$ is
ample if and only if it is positive.
Thus $M$ is projective if and only if it
admits a positive holomorphic line bundle.
\end{theorem}

\begin{proof}
We must show that $\varphi$ separates points
and has
surjective differential everywhere…
\end{proof}

\noindent
Compare this theorem with Nakai-Moishezon
Criterion
\ref{algebraic-geometry-theorem-%
Nakai-Moishezon-criterion}.

\begin{exercise}
\label{exercise-l-ample-impli-lotim-globa}
Let $L$ be an ample bundle on a K3 surface
$M$.
Prove that $L^{\otimes 2}$ is
globally generated
(that is, for each $x \in M$ there exists a
section $h \in
H^{0}(L^{\otimes 2})$ which does not vanish
in $x$).
\end{exercise}

\begin{proof}
\begin{reference}
\href{https://math.stackexchange.com/%
questions/5081481/
ample-line-bundle-on-a-k3-surface-implies-l-
otimes-2-is-base-point-free/
5081492?noredirect=1#comment10928708_5081492}
{StackExchange}
\end{reference}
As an outline for exposition:
\begin{enumerate}
\item Show that $L$ has sections
(by ampleness+Kodaira Vanishing+Riemann-Roch
for surfaces),
pick a section $C$ and note that if
$L^{\otimes2}$ had base points they would be
in $C$.

\item  Using Kodaira Vanishing again we show
that every section of the line bundle
restricted to $C$
comes from a section of the line bundle on
$M$.
Thus it's enough to find a nonvanishing
section on the restricted bundle.

\item Then we aim to apply Riemann-Roch on
$C$.
If we can show that $h^0((L^{\otimes 2})^\vee
\otimes \omega_C)$
and $h^0((L^{\otimes 2}(-p))^\vee \otimes
\omega_C)$
both vanish, we apply Riemann-Roch and see
that
$h^0(L^{\otimes 2})> h^0(L^{\otimes 2}(-p)$,
which proves the result since that says
that not every section is a section vanishing
at $p \in C$.

\item To compute the vanishings of the $h^0$
in the previous step we show they have
negative degrees
(genus formula + degree of canonical bundle
of a curve).
\end{enumerate}

First we need to show that $L$ has sections.
To see this I will show that since $L$ is
ample, $L^2>0$. Indeed, if $mL$ is very ample
we apply by Riemann-Roch, Kodaira Vanishing
and the fact that $\chi(\mathcal{O}_X)=2$ on
a K3 to obtain $h^0(mL)=2+\frac{1}{2}m^2L^2$
which must be positive since $mL$ has
sections by being very ample. Notice $L^2$
cannot be zero as $h^0(mL)=2$ would imply we
have an embedding of a K3 surface into
$\mathbb{P}^1$.

Applying again Riemann-Roch and Kodaira
vanishing we see that $h^0(L)>0$. Let $C$ be
a section of $L$.

Notice that if $L^{\otimes 2}$ had any base
points, they would have to be on $C$. Indeed,
since $C$ is the vanishing set of a section
$s\in H^{0}(L)$, we must have $s \otimes s
\in H^{0}(L^{\otimes 2})$ vanishing in any
base point of $L^{\otimes 2}$.

Then we consider the ideal exact sequence for
$C$ and tensor by $L^{\otimes 2}$ (which
preserves left-exactness because line bundles
are locally free of rank 1, that is, we are
tensoring by the underlying ring) to obtain
$$ \xymatrix{ 0\ar[r]&L^{\otimes
2}(-C)\ar[r]&L^{\otimes 2}\ar[r]& L^{\otimes
2}\otimes\mathcal{O}_C\ar[r]&0 } $$ Notice
that the term on the left is actually $L$
since the ideal sheaf
$\mathcal{O}_M(-C)=\mathcal{I}_C$ is dual to
$L$ because $C$ is defined as the vanishing
set of a section of $L$. Passing to
cohomology we obtain $$ \xymatrix{
H^{0}(L^{\otimes 2})\ar[r]&H^{0}(L^{\otimes
2}\otimes \mathcal{O}_C)\ar[r]& H^{1}(L) } $$
where the latter term vanishes by Kodaira
Vanishing because $L$ is ample. This says
that every section of $L^{\otimes 2}|_C$ is
the restriction to $C$ of a section of
$L^{\otimes 2}$. In turn this shows that it's
enough to show that the bundle $L^{\otimes
2}$ has no base points along $C$.

Now pick a point $p\in C$. Saying
that $p$ is not a base point of $L^{\otimes
2}|_{C}$ is the same as saying that not every
section of $L^{\otimes 2}|_{C}$ vanishes at
$p$, that is, that $h^0(L^{\otimes
2}|_{C}(-p))<h^0(L^{\otimes 2}|_{C})$,
which concludes the exercise.

The key step to prove
the last inequality is showing that
$h^0(\omega_C\otimes(L^{\otimes2}|_{C})^\vee)
=0=h^0(\omega_C\otimes
(L^{\otimes2}|_{C}(-p))^\vee)$. For this we
need to see these line bundles as Weil
divisors, so that their degree is the sum of
their multiplicities. First note that by the
genus formula for (possibly singular) curves
on smooth surfaces and by the fact that $M$
is smooth, we see that $2p_a(C)-2=(L.L)$ so
that
$$
4p_a(C)-4=(2L.L)=\text{deg}_C(L^{\otimes 2})
$$
Notice that since $L^2>0$ we exclude the
cases that $p_a(C)=0,1$.

By Riemann-Roch on the curve
$C$, we know that
$\text{deg}\omega_C=2p_a-2$. Then
$$
\text{deg}_C
(\omega_C\otimes (L^{\otimes2})^\vee)
=\text{deg}_C(\omega_C)-\text{deg}
_C(L^{\otimes
2})
=2p_a-2-(4p_a-4)<0
$$
and a similar computation works for
$L^{\otimes 2}(-p)$ as it has the same
degree of $L^{\otimes 2}$ minus 1. This
implies that neither of these bundles
can have sections since any section would
provide a linearly equivalent
effective divisor
--- degree is preserved under linear
equivalence,
and divisors given by sections are effective.
Therefore
$h^0(\omega_C\otimes
(L^{\otimes2}|_{C})^\vee)=0
=h^0(\omega_C\otimes (L^{\otimes
2}|_{C}(-p))^\vee)$ as claimed.

Finally we apply Riemann-Roch to $L^{\otimes
2}$ and $L^{\otimes 2}(-p)$ to
obtain that $h^0(L^{\otimes
2})>h^0(L^{\otimes 2}(-p))$.
\end{proof}

\section{Riemann-Roch formulas}
\label{section-riemann-roch}

\noindent
We want to know the dimension of
$H^0(X,\mathcal{O}(D))$,
where $X$ is a curve and $D$ is a divisor on
$X$.
To refresh out memory recall that if $D$ is a
smooth hypersurface
(in this case this just means that it's a set
of points
without multiplicities)
then that $\mathcal{O}(D)$ is given by the
cocycle
$g_{ij}=f_j/f_i$ where $D=\text{div}(f_i)$ on
$U_i$.
The $g_{ij}$ give a cocycle, which in turn is
a line bundle on $X$.
If $e_i$ is a local basis then we have
$e_i=g_{ij}e_j$
by definition of transition function.
A section is given locally by $h_ie_i$ and it
has to glue,
that is, $h_je_j=h_ie_i=e_if_j/f_ie_j$
so $h_j/f_j=h_i/f_i$.
That's literally the data of a meromorphic
function,
namely some local meromorphic functions which
coincide in intersections.

What's the vanishing of this meromorphic
function along $D$?
Now suppose $D$ is an arbitrary Weil divisor,
so the $f_i$ have some multiplicities.
If $h_i$ has no zeroes, then the multiplicity
of the new
meromorphic function remains unaltered, but
$h_i$ can also vanish
along $D$ and then the vanishing becomes as
positive as we want it.
But it cannot have poles, the least positive
it can be is
the prescribed negativity of $D$.

This is why looking for $h^0(D)$ is just as
much as looking for
meromorphic functions $h$ on $X$ such that
$$
\text{div}h + D \geq 0\qquad \text{(is
effective).}
$$



\begin{theorem}[Riemann-Roch for curves]
\label{theorem-Riemann-Roch-for-curves}
Let $D$ be a divisor on a compact Riemann
surface $X$, that is, $D$ is a
collection of points with multiplicities on
$X$.
\begin{equation}
\label{equation-riema-roch-riema-surfa}
h^0(D)-h^0(K-D)=d-g+1
\end{equation}
\end{theorem}

\begin{exercise}
\label{exercise-no-holom-vecto-field-g}
Prove that if $X$ is a complex smooth curve
with $g\geq 2$ then $h^0(T_X)=0$.
\end{exercise}

\begin{proof}
First you notice that the dimension of the
holomorphic 1-forms is the genus.
This is just because the genus $p_a$ is
defined as $h^1(\mathcal{O}_X)$, which
is just $h^0(K_X)$ by Serre duality. Then you
say well if you have a holomorphic
vector field, pair it with the nonzero
holomorphic 1-form. This gives a nonzero
function, but it must be constant because $X$
is compact. Then it is actually
nowhere vanishing. This says the vector field
cannot vanish anywhere, which
means the tangent bundle of the curve is
trivial. Apparently elliptic curves,
i.e. $g=1$ are the
\end{proof}

\begin{theorem}[Riemann-Roch for line bundles
on surfaces]
\label{theorem-riema-roch-line-bundl-surfa}
Let $L$ be a line bundle on a complex surface
$X$, and $K_X=\Omega^2X$ the
canonical bundle of $X$. Then
\begin{equation}
\label{equation-riema-roch-line-bundl-surf}
\chi(L)=\chi(\mathcal{O}_X)+\frac{1}{2}(L,L-
K_X)
\end{equation}
\end{theorem}

\begin{lemma}[Genus formula]
\label{lemma-genus-formula}
Let $C$ be any curve on a surface $S$. Then
\begin{equation}
\label{equation-genus-formula}
2p_a(C)-2=(C.C-K_S).
\end{equation}
\end{lemma}

\begin{proof}
If $C$ is nonsingular we may use the normal
bundle short exact sequence. If $C$ is
singular we need to use the ideal sheaf short
exact sequence, which means we think of $C$
either as a closed embedded subscheme of $S$
or as an analytic subvariety. Taking Euler
characteristic we obtain
$$
\chi(\mathcal{O}(-C))+\chi(\mathcal{O}_C)=
\chi(\mathcal{O}_S)
$$
Since $p_a(C)=h^0(\Omega^1)
=1-\chi(\mathcal{O}_C)$, we
obtain that
$p_a(C)=1-\chi(\mathcal{O}_S)-
\chi(\mathcal{O}(-C))$.
To compute $\chi(\mathcal{O}(-C))$ we use
Riemann-Roch formula for curves on surfaces
to get
$$
\chi(\mathcal{O}(-C))=\chi(\mathcal{O}_S)
+\frac{(\mathcal{O}(-C).
\mathcal{O}(-C)-\omega_S^\vee)}{2}
$$
Taking duals on both entries of the
intersection form and substituting in our
previous expression for $p_a(C)$ we obtain
the result.
\end{proof}

\begin{remark}
\label{remark-euler-characteristic-additive}
To prove that Euler characteristic is
additive, i.e. that for an exact sequence
$$
\xymatrix{
0\ar[r]
&E\ar[r]
&F\ar[r]
&G\ar[r]
&0
}
$$
we must have
$$
\chi(F)=\chi(E)+\chi(G),
$$
we take the long exact sequence in cohomology
and use the general fact that for
a finite exact sequence of vector spaces
$$
\xymatrix{
0\ar[r]
&V_0\ar[r]
&V_1\ar[r]
&\cdots\ar[r]
&V_n\ar[r]
&0
}
$$
we must have
$$
\sum_i(-1)^i\dim V_i,
$$
which basically follows by dimension theorem.
\end{remark}

\begin{exercise}
\label{exercise-singular-curve-on-K3}
Let $S$ be a singular, irreducible complex
curve on a K3 surface. Prove that
$(S.S)\geq 0$.
\end{exercise}

\begin{proof}
This is just an application of Eq.
\ref{equation-genus-formula}. Since
$K_X=\mathcal{O}_X$ it suffices to show that
$p_a(S)$ is not zero. This follows
fact that $S$ is singular, since its
arithmetic genus is defined as the genus of
its normalization, which must be strictly
positive since there exists at least
one singularity.
\end{proof}


\section{Symplectic isogenies and
Shafarevich}
\label{section-symplectic-isogenies-shaf}

These are notes from Misha Verbitsky's talk
{\it Symplectic isogenies for hyperkahler
manifolds and Shafarevich conjecture},
given at IMPA on May 7, 2026.

\medskip
Abstract.
A correspondence between hyperkahler
manifolds is a symplectic isogeny if it is
generically finite over both factors and its
push-pull operation sends the holomorphic
symplectic form to a nonzero holomorphic
symplectic form.
The conjecture, attributed to Shafarevich by
Buskin, says that hyperkahler manifolds in
the same deformation family with isomorphic
rational period lattices should be related
by such an isogeny.
The talk explains degenerate twistor
families and shows that projective
hyperkahler manifolds in the same such
family are symplectically isogenous.

\medskip
Absolute upshot.
Same degenerate twistor family is a much
more structured condition than being merely
deformation-equivalent.
Misha proves the Shafarevich-type isogeny
statement first in this structured setting:
the common Lagrangian fibration gives torus
fibers and torsors, which can be trivialized
after finite base change.
The resulting graph produces the desired
symplectic correspondence.

\medskip
Upshot.
Find the simplest Hodge-theoretic piece
forced by the holomorphic symplectic form:
the transcendental lattice $V^{tr}$.
Then prove that this piece behaves well:
it is usually irreducible, and in the
degenerate case has only one controlled
defect.
See Theorem
\ref{theorem-transcendental-k3-type}.

\medskip
{\bf Motivation: ADE contractions.}

Let $X$ be a smooth K3 surface and let
$C \subset X$ be a smooth rational curve.
Since $K_X=\mathcal{O}_X$, adjunction gives
$$
2g(C)-2=C^2+K_X\cdot C=C^2.
$$
Thus a smooth rational curve on a K3 surface
has self-intersection $C^2=-2$.
Conversely, a smooth curve with $C^2=-2$
has genus $0$.
Such a curve is called a $(-2)$-curve.

\medskip
A $(-2)$-curve is not obtained by blowing up
another smooth K3 surface.
Rather, it can be contracted:
$$
X \longrightarrow Y.
$$
The target $Y$ is usually singular.
For a single $(-2)$-curve, the image point is
a Du Val singularity of type $A_1$.
Locally this is the rational double point
$$
xy=z^2,
$$
or equivalently $\mathbb{C}^2/\{\pm 1\}$.

\medskip
More generally, suppose
$C_1,\ldots,C_r \subset X$ are $(-2)$-curves
whose intersection matrix is negative
definite.
Then the whole configuration can be
contracted.
The dual graph has one vertex for each
$C_i$ and an edge when two curves meet
transversely once.
For rational double points this graph is one
of the ADE Dynkin diagrams:
$$
A_n,\qquad D_n,\qquad E_6,E_7,E_8.
$$

\medskip
Conversely, the minimal resolution of a
Du Val singularity replaces the singular
point by a configuration of $(-2)$-curves.
Their intersection graph is the corresponding
Dynkin diagram.
Thus the slogan is
$$
\text{Du Val singularity}
\longleftrightarrow
\text{ADE configuration of }
(-2)\text{-curves}.
$$
The lattice generated by these curves is the
corresponding root lattice.

\medskip
In the motivating picture, one contracts a
Dynkin configuration
$$
f:X\longrightarrow Y
$$
and then resolves the singular surface $Y$:
$$
\widetilde{Y}\longrightarrow Y.
$$
If $X$ is not already the chosen resolution,
the common singular surface $Y$ gives a
relation between $X$ and $\widetilde{Y}$.
One studies this relation by the
correspondence
$$
X \xleftarrow{\ p_1\ }
X\times_Y\widetilde{Y}
\xrightarrow{\ p_2\ }\widetilde{Y}.
$$
This is the prototype for replacing a
birational map by a cycle in a product.

\medskip
This is the end of the motivating example.
The real subject of the talk is to replace
maps between hyperkahler manifolds by
correspondences which transport holomorphic
symplectic forms.

\medskip
{\bf Correspondences and symplectic
isogenies.}

\begin{definition}
\label{definition-correspondence-verbitsky}
Let $M$ be a complex variety.
A {\it cycle} in $M$ is a subvariety
together with a positive integer
multiplicity on each irreducible component.
Let $A,B$ be irreducible compact complex
varieties of dimension $n$.
A {\it correspondence} from $A$ to $B$ is an
$n$-dimensional cycle
$$
Z\subset A\times B
$$
such that every irreducible component of
$Z$ projects surjectively to both $A$ and
$B$.
\end{definition}

\noindent
Thus a correspondence is a multivalued map.
When using cohomology, we often think of the
cycle $Z$ as its Poincare dual differential
form. The caveat is that pushforward of a
differential form is generally a current.
For holomorphic forms in this setting, the
pushforward is holomorphic away from the
exceptional locus and extends by Hartogs.

\begin{definition}
\label{definition-symplectic-isogeny}
Let $A,B$ be holomorphically symplectic
manifolds and let
$Z\subset A\times B$ be a correspondence.
Let $\pi_1,\pi_2$ be the projections from
$Z$ to $A$ and $B$.
The correspondence $Z$ is a
{\it symplectic isogeny} if, for a
holomorphic symplectic form $\Omega_A$ on
$A$,
$$
(\pi_2)_*(\pi_1^*\Omega_A)\neq 0.
$$
\end{definition}

\begin{lemma}
\label{lemma-sympl-isog-transcendental}
Let $M_1,M_2$ be holomorphically symplectic
manifolds of K3 type, and let
$Z\subset M_1\times M_2$ be a symplectic
isogeny.
Then the push-pull map
$$
\Phi_Z(\eta):=(\pi_2)_*(\pi_1^*\eta)
$$
induces an isomorphism of rational
transcendental Hodge lattices
$$
\Phi_Z:T(M_1)\longrightarrow T(M_2).
$$
\end{lemma}

\noindent
Here $T(M)$ means the transcendental Hodge
lattice $V^{tr}\subset H^2(M,\mathbb{Q})$.
This is not literally the complex torus
attached to a weight-one Hodge structure.
The torus story is the model case: isogenies
of tori are the same as isomorphisms of
rational weight-one Hodge structures.
For K3-type Hodge structures, the analogous
object is $T(M)$.

\medskip
Recall also the dictionary between complex
tori and Hodge structures; see
{\it mumford-tate-groups-in-hodge-theory.tex}
for the Hodge-theoretic background.
If $T$ is a compact complex torus, then
$H_1(T,\mathbb{Z})$ carries an integral
Hodge structure of weight $1$.
Conversely, such a Hodge structure
$$
V_{\mathbb{C}}=V^{1,0}\oplus V^{0,1}
$$
with a lattice $V_{\mathbb{Z}}\subset V$,
gives the complex torus
$$
T=V^{1,0}/V_{\mathbb{Z}}.
$$
After tensoring the lattice with
$\mathbb{Q}$, this is the category of
complex tori up to isogeny.

\begin{definition}
\label{definition-isogeny-tori}
Let $T_1,T_2$ be compact complex tori of the
same dimension.
An {\it isogeny} from $T_1$ to $T_2$ is a
subtorus
$$
T\subset T_1\times T_2
$$
such that the two projections
$T\to T_1$ and $T\to T_2$ are coverings.
We say that $T_1$ and $T_2$ are isogenous
if such a subtorus exists.
\end{definition}

\medskip
Briefly, the Mumford-Tate group records the
rational algebraic symmetries forced by a
Hodge structure.
See Section
\ref{mumford-tate-section-mumford-tate-group}
and Definition
\ref{mumford-tate-definition-mumford-tate-%
group}.
If $V$ is a rational Hodge structure and
$$
\rho:U(1)\longrightarrow
GL(V_{\mathbb{R}})
$$
is the Hodge circle action, then $MT(V)$ is
the smallest algebraic subgroup of $GL(V)$
defined over $\mathbb{Q}$ whose real points
contain $\rho(U(1))$.

\begin{proposition}
\label{proposition-mt-generated-galois}
Let $V$ be a rational Hodge structure.
For each
$\sigma\in\operatorname{Aut}_{\mathbb{Q}}
(\mathbb{C})$,
let $U(1)^\sigma$ denote the Galois conjugate
of the Hodge circle in
$GL(V_{\mathbb{C}})$.
Then $MT(V)_{\mathbb{C}}$ is the Zariski
closure of the subgroup generated by all
$U(1)^\sigma$.
\end{proposition}

\begin{theorem}
\label{theorem-transcendental-hodge-lattice}
Let $V_{\mathbb{Q}}$ be a rational Hodge
structure of weight $d$.
Let $V^{tr}\subset V_{\mathbb{Q}}$ be the
smallest rational subspace such that
$$
V^{d,0}\subset V^{tr}_{\mathbb{C}}.
$$
Equivalently, $V^{tr}_{\mathbb{C}}$ is
generated by all Galois conjugates of
$V^{d,0}$.
Then $V^{tr}$ is a Hodge substructure.
\end{theorem}

\begin{definition}
\label{definition-polarization-hodge}
A {\it polarization} on a rational Hodge
structure $V$ of weight $k$ is a
$U(1)$-invariant nondegenerate bilinear form
$$
h\in V_{\mathbb{Q}}^*\otimes
V_{\mathbb{Q}}^*
$$
which is symmetric or skew-symmetric
depending on the parity of $k$, and which
satisfies the Hodge-Riemann positivity
condition.
\end{definition}

\noindent
This is Misha's formulation.
In
Definition
\ref{mumford-tate-definition-polarization},
Andrey packages the same positivity using a
morphism to $\mathbb{Q}(-k)$ and the Weil
operator.

\medskip
The basic source of such polarizations is the
Hodge-Riemann pairing.
For a compact Kahler manifold $X$, a Kahler
class $\omega$ gives, on primitive
$H^k(X)$, the pairing
$$
(\alpha,\beta)\longmapsto
\int_X\alpha\wedge\beta\wedge\omega^{n-k}.
$$
After the usual sign and Weil-operator
normalization, the Hodge-Riemann bilinear
relations say this pairing is positive.
Thus the abstract polarization in Hodge
theory is the linear-algebra shadow of the
intersection pairing from complex geometry.

\begin{definition}
\label{definition-simple-semisimple-object}
Let $\mathcal{A}$ be an abelian category.
An object $V$ is {\it simple} if its only
subobjects are $0$ and $V$.
It is {\it semisimple} if it is a finite
direct sum of simple objects.
\end{definition}

\noindent
The upshot is that polarizations prevent
extensions from hiding inside a Hodge
structure.
If $W\subset V$ is a Hodge substructure, the
orthogonal complement $W^\perp$ is again a
Hodge substructure, so $V$ splits as
$W\oplus W^\perp$.

\begin{proposition}
\label{proposition-polarizable-hs-semisimple}
Every polarizable rational Hodge structure
is semisimple.
\end{proposition}

\noindent
In particular, the rational Hodge structures
$H^k(X,\mathbb{Q})$ of a compact Kahler
manifold are semisimple.
Indeed, the Lefschetz decomposition reduces
the statement to primitive cohomology, where
the Hodge-Riemann bilinear relations give
the required polarization.

\begin{definition}
\label{definition-k3-type-hodge-structure}
A rational Hodge structure of weight $2$
is of {\it K3 type} if
$$
V_{\mathbb{C}}=
V^{2,0}\oplus V^{1,1}\oplus V^{0,2}
$$
and $\dim V^{2,0}=1$.
\end{definition}

\begin{example}
\label{example-bbf-k3-type}
Let $M$ be an irreducible holomorphic
symplectic manifold.
Then $H^2(M,\mathbb{Q})$ is a Hodge
structure of K3 type.
It carries the Beauville-Bogomolov-Fujiki
quadratic form
$$
q:H^2(M,\mathbb{Q})\times H^2(M,\mathbb{Q})
\longrightarrow \mathbb{Q}.
$$
Fujiki's formula says that, if
$\dim_{\mathbb{C}}M=2n$, then
$$
\int_M\alpha^{2n}=c\,q(\alpha,\alpha)^n
$$
for a positive rational constant $c$.
\end{example}

\noindent
Thus, for hyperkahler manifolds, the BBF
form plays the role of the natural quadratic
form on the K3-type Hodge structure
$H^2(M)$.
It is not literally the ordinary
Hodge-Riemann pairing on all cohomology, but
it is the replacement on $H^2$ which controls
the period lattice and the Torelli theory.

\begin{theorem}
\label{theorem-transcendental-k3-type}
Let $(V,q)$ be a pseudo-polarized rational
Hodge structure of K3 type.
Let $V^{tr}\subset V$ be the transcendental
Hodge sublattice, i.e. the smallest rational
Hodge substructure containing $V^{2,0}$.
Then exactly one of the following happens.
\begin{enumerate}
\item If $V^{1,1}_{\mathbb{Q}}$ contains a
vector $x$ with $q(x,x)>0$, then
$q|_{V^{tr}}$ is a polarization and
$V^{tr}$ is irreducible.
\item If $q|_{V^{1,1}_{\mathbb{Q}}}$ is
negative definite, then $V^{tr}$ is
irreducible and
$q|_{V^{tr}}$ is a pseudo-polarization.
\item If $q|_{V^{1,1}_{\mathbb{Q}}}$ is
negative semidefinite with one-dimensional
kernel, then
$$
V^{1,1}_{\mathbb{Q}}\cap
(V^{1,1}_{\mathbb{Q}})^\perp
=\langle\eta\rangle,
\qquad q(\eta,\eta)=0.
$$
In this case $\langle\eta\rangle\subset
V^{tr}$ is the only nonzero proper Hodge
substructure of $V^{tr}$, and
$V^{tr}/\langle\eta\rangle$ is irreducible.
\end{enumerate}
\end{theorem}

\noindent
The third case is the strange one.
The transcendental lattice is no longer
irreducible, but it fails in the smallest
possible way: it has a unique isotropic
rational line, and the quotient by this line
is irreducible.

\medskip
Upshot: for K3-type Hodge structures the
transcendental part is almost simple.
Usually it is irreducible.
In the degenerate case it has exactly one
visible defect, the isotropic line
$\langle\eta\rangle$.
Thus a correspondence which is nonzero on
$V^{2,0}$ is forced to control essentially
all of $V^{tr}$.
This is why the condition
$(\pi_2)_*\pi_1^*\Omega\neq 0$ is so strong.

\begin{lemma}
\label{lemma-sympl-isog-controls-tm}
The point of the symplectic isogeny
condition is that the correspondence is
nonzero on the line $H^{2,0}$.
Since $T(M)$ is the smallest rational Hodge
substructure generated by this line, the
correspondence is forced to identify the
transcendental lattices:
$$
T(M_1)\cong T(M_2).
$$
\end{lemma}

\begin{theorem}[Witt]
\label{theorem-witt-quadratic-spaces}
Let $(V,q)$ be a nondegenerate rational
quadratic space.
Let $W_1,W_2\subset V$ be subspaces such that
$q|_{W_i}$ is nondegenerate.
If $W_1$ and $W_2$ are isomorphic as
quadratic spaces, then
$W_1^\perp$ and $W_2^\perp$ are isomorphic
as quadratic spaces.
\end{theorem}

\begin{lemma}
\label{lemma-isogeny-gives-h2-hodge-isom}
Let $M_1,M_2$ be projective hyperkahler
manifolds in the same deformation class.
If $M_1$ and $M_2$ are symplectically
isogenous, then
$$
H^2(M_1,\mathbb{Q})\cong
H^2(M_2,\mathbb{Q})
$$
as rational Hodge structures.
\end{lemma}

\noindent
Indeed, the isogeny identifies the
transcendental Hodge lattices $T(M_i)$.
The algebraic parts are the orthogonal
complements with respect to the BBF form.
Witt's theorem then gives the corresponding
isomorphism on the complements, hence on all
of $H^2$.

\begin{theorem}[Torelli for K3 surfaces]
\label{theorem-torelli-k3}
Let $S_1,S_2$ be K3 surfaces.
Then $S_1\cong S_2$ if and only if their
integral Hodge structures on $H^2$ are
isomorphic, with the Kahler cone condition.
\end{theorem}

\begin{theorem}[Buskin]
\label{theorem-buskin-k3-isogeny}
Let $S_1,S_2$ be projective K3 surfaces.
If
$$
H^2(S_1,\mathbb{Q})\cong
H^2(S_2,\mathbb{Q})
$$
as rational Hodge structures, then $S_1$ and
$S_2$ are symplectically isogenous.
\end{theorem}

\noindent
This is essentially a Hodge-conjecture
statement for $S_1\times S_2$.
An isomorphism of rational Hodge structures
on $H^2$ gives a rational Hodge class in
$$
H^2(S_1,\mathbb{Q})^*
\otimes H^2(S_2,\mathbb{Q})
\subset H^4(S_1\times S_2,\mathbb{Q}).
$$
To produce a symplectic isogeny, one must
realize this Hodge class by an algebraic
cycle, i.e. by a correspondence.
That is exactly the kind of assertion the
Hodge conjecture predicts.

The Shafarevich conjecture says:
if $M_1,M_2$ are hyperkahler manifolds in
the same deformation class and their rational
Hodge structures on $H^2$ are isomorphic,
then $M_1$ and $M_2$ are symplectically
isogenous.

\medskip
Very briefly: Torelli says integral periods
recover a K3 surface.
Buskin says rational periods for projective
K3 surfaces recover a symplectic isogeny.
Shafarevich asks for the analogous rational
period statement for hyperkahler manifolds.
The earlier lemma gives only the reverse
direction: an isogeny forces an isomorphism
of rational $H^2$ Hodge structures.

\begin{definition}
\label{definition-holomorphic-lagrangian-fib}
Let $(M,\Omega)$ be a holomorphically
symplectic manifold of complex dimension
$2n$.
A holomorphic map
$$
\pi:M\longrightarrow B
$$
is a {\it holomorphic Lagrangian fibration}
if its fibers have dimension $n$ and
$\Omega$ restricts to zero on the smooth
part of every fiber.
\end{definition}

\noindent
By the holomorphic Liouville theorem, the
smooth compact fibers are complex tori.
Thus a holomorphic Lagrangian fibration is a
family of complex tori over its smooth base
locus.

\medskip
{\bf Degenerate twistor deformations.}

Now let $\pi:M\to X$ be a holomorphic
Lagrangian fibration, with $M$ irreducible
holomorphic symplectic and holomorphic
symplectic form $\Omega$.
Usually $X=\mathbb{P}^n$, but this is not
essential here.
Let $\omega_{FS}$ be a Kahler form on $X$.
For $\lambda\in\mathbb{C}$, put
$$
\Omega_\lambda=
\Omega+\lambda\pi^*\omega_{FS}.
$$

\begin{lemma}
\label{claim-degenerate-twistor-form}
The closed complex form $\Omega_\lambda$
defines a complex structure on the smooth
manifold $M$ for which $\Omega_\lambda$ is a
holomorphic symplectic form.
\end{lemma}

\noindent
This is the basic degenerate twistor
deformation.
We deform the complex structure by adding a
form pulled back from the base of the
Lagrangian fibration.

\begin{definition}
\label{definition-degenerate-twistor-def}
The family of holomorphic symplectic
manifolds obtained from the forms
$$
\Omega_\lambda=
\Omega+\lambda\pi^*\omega_{FS},
\qquad \lambda\in\mathbb{C},
$$
is called the {\it degenerate twistor
deformation} associated to the Lagrangian
fibration $\pi:M\to X$.
\end{definition}

\noindent
Basic properties:
\begin{enumerate}
\item If the added form on $X$ is exact, the
deformation is trivial by Moser's theorem.
\item When $M$ is compact and of Kahler type,
these deformations are limits of ordinary
twistor deformations.
\item Each member still carries a
holomorphic Lagrangian fibration over the
same base $X$.
\item The fibers and the base are unchanged;
only the complex structure on the total
space is varied.
\item Algebro-geometrically, this is a
Tate-Shafarevich deformation: it is described
by a class in
$$
H^1(X,\operatorname{Aut}_\pi(M)),
$$
where $\operatorname{Aut}_\pi(M)$ is the
sheaf of fiberwise automorphisms.
\end{enumerate}

\begin{definition}
\label{definition-torsor}
Let $G\to X$ be a group object over $X$.
A {\it torsor} over $G$ is a space
$P\to X$ with a fiberwise action of $G$ such
that, locally on $X$, $P$ is isomorphic to
$G$ with its translation action.
Equivalently, $G$ acts freely and
transitively on the fibers of $P\to X$.
\end{definition}

\noindent
In the Lagrangian fibration setting, the
smooth fibers are tori.
Different fibrations in the same degenerate
twistor family can be viewed as torsors over
the same relative torus fibration.

\medskip
Upshot.
The culminating goal is to show that any two
projective hyperkahler manifolds in the same
degenerate twistor family are symplectically
isogenous.
To construct this isogeny, first make the two
torus fibrations literally the same after a
finite cover of the base.
Then the graph of the resulting isomorphism,
after compactification, gives the desired
correspondence.

\begin{theorem}
\label{theorem-degenerate-twistor-cover}
Let $M_1,M_2\to X$ be projective
hyperkahler manifolds in the same degenerate
twistor family.
Then there is an open subset $U\subset X$
and a finite cover $\widetilde{U}\to U$ such
that
$$
M_1\times_X\widetilde{U}
\cong
M_2\times_X\widetilde{U}
$$
over $\widetilde{U}$.
\end{theorem}

\noindent
This is the theorem where torsors enter.
After finite base change, the torsors acquire
sections, hence become trivial.
The graph of the resulting isomorphism over
an open set is then compactified to produce
the correspondence, giving the symplectic
isogeny.

\section{Hypercomplex manifolds}
\label{section-hypercomplex-manifolds}

\begin{exercise}
\label{exercise-SU2-invariant-forms}
Let $M$ be a compact hypercomplex manifold of
real dimension 4, equipped with a
quaternionic Hermitian structure, and $V$ the
space of closed
$\text{SU}(2)$-invariant 2-forms. Prove that
$V$ is finite-dimensional.
\end{exercise}

\begin{proof}[Proof by Arpan.]
Consider the Hodge star operator with respect
to the Riemannian metric on $M$.
\begin{enumerate}
\item A closed anti-self-dual form is
harmonic.
\item The space of harmonic forms is
finite-dimensional.
\item An $\text{SU}(2)$-invariant closed form
is
anti self-dual.
\begin{enumerate}
\item
$\Lambda^+=\text{span}(\omega_I,\omega_J,
\omega_K)$.
$M$ is a manifold of real dimension 4,
which means that $\dim\Lambda^2(M)=6$ and
thus
$\dim\Lambda^+(M)=3$
(here we can check using a basis of
$\mathbb{R}^6$
that $*$ permutes the basic covectors with no
fixed points,
so that it has zero trace;
in turn this means that $\dim \Lambda^+ +
\dim\Lambda^-=0$).
Further, each of $\omega_I$, $\omega_J$ and
$\omega_K$
squared and multiplied by $\frac{1}{2}$ give
the Riemannian
volume form, and therefore the three are in
$\Lambda^+(M)$.
Clearly they are linearly independent and
thus span $\Lambda^+(M)$.

\item $\omega_I,\omega_J$ and $\omega_K$ are
not fixed by the
$\text{SU}(2)$ action
Indeed, $J\omega_I=-\omega_I$ and similarly
for other cases. Therefore there are no
$\text{SU}(2)$-fixed
points in $\Lambda^+$.

\item
$\text{SU}(2)\Lambda^+(M)=\Lambda^+(M)$.
This is clear since $J\omega_I = -\omega_I$
and so on.
Since $*$ is an isomorphism of $\Lambda^+(M)$
it follows that
$\text{SU}(2)\Lambda^-(M)=\Lambda^-(M)$.

\medskip\noindent
We conclude that if $\alpha$ is
$\text{SU}(2)$-invariant (i.e. a fixed point
of
the action of $\text{SU}(2)$ on $\Lambda^2$)
its positive part will vanish, so that
$\alpha$ is anti-self-dual. More precisely,
if
$\alpha = \alpha^+ +\alpha^-$ is
$\text{SU}(2)$-invariant and $A \in
\text{SU}(2)$,
$$
\alpha^+ +\alpha^- = \alpha =
A\alpha=A\alpha^+ + A\alpha^-,
$$
so $A\alpha^+=\alpha^+$, but $\text{SU}(2)$
has no fixed points in $\Lambda^+$, so
$\alpha^+=0$.
Therefore $\alpha \in \Lambda^-$ as required.
\end{enumerate}
\end{enumerate}
\end{proof}

\begin{exercise}
\label{exercise-close-funda-forms-almos}
Let $(M,K,J,K,g)$ be an almost hypercomplex
Hermitian structure
on a K3 surface, and $\omega_I, \omega_J,
\omega_K$
its fundamental forms.

Suppse that $\omega_I$ is closed.
Prove that $\omega_J,\omega_K$ are closed, or
find a counterexample.
\end{exercise}

\noindent
The idea of the following solution
(provided entirely by ChatGPT)
is that we can ``rotate''
(in the $J,K$-plane within $\text{End}(TM)$)
two of the complex structures while
preserving the
quaternionic relations.
Any orthonormal triple of such endomorphisms
should be expected to satisfy the
quaternionic relations (algebra),
but the closedness condition on the
associated fundamental
forms will be modified (differential).

\begin{proof}
Let $(M,I,J,K,g)$ be a hyperkähler structure
on a K3 surface.
Define
$$
\begin{aligned}
\tilde{J}&=\cos \theta J+ \sin \theta K\\
\tilde{K}&=-\sin \theta J + \cos \theta K
\end{aligned}
$$

\noindent
Then for any $v \in TM$,
$$
\begin{aligned}
\tilde{J}^2v&=\cos J(\cos \theta Jv+ \sin
\theta Kv)
+\sin \theta K (\cos J v + \sin Kv)\\
&=\cos^2J^2v+\cos \theta \sin \theta J K v
+ \sin \theta \cos \theta K J v + \sin^2
\theta K^2 v\\
&=-\text{Id} v
\end{aligned}
$$

\noindent
And likewise $\tilde{K}^2=-\text{Id}$.
Also,
$$
\begin{aligned}
I \tilde{J} \tilde{K}v&=
I(\cos\theta J(-\sin \theta Jv+\cos \theta
Kv)
+\sin \theta K (- \sin Jv+ \cos \theta Kv))\\
&=I (-\cos \theta \sin \theta J^2 v
+ \cos^2 \theta J K v
- \sin^2 \theta K J v
+\sin \theta \cos \theta K^2 v)\\
&=-\cos \theta \sin \theta I J^2 v
+ \cos^2 \theta I J K
- \sin^2 \theta I K J v
+ \sin \theta \cos \theta I K^2 v\\
&=-\cos^2v- \sin^2v=-\text{Id}v.
\end{aligned}
$$

\noindent
Thus $(I,\tilde{J},\tilde{K})$ is an almost
hypercomplex
structure on $M$.
However,
$$
\begin{aligned}
d
\omega_{\tilde{J}}&=d(g(\cdot,\tilde{J}\cdot)
)\\
&=d(g(\cdot,\cos \theta J\cdot + \sin \theta
K\cdot))\\
&=d(\cos \theta \omega_J + \sin \theta
\omega_K)\\
&=d \cos \theta \wedge \omega_J + d \sin
\theta \wedge \omega_K\\
&=-\sin \theta
d\theta\wedge\omega_J+\cos\theta
d\theta\wedge\omega_K\\
&=d \theta \wedge \omega_{\tilde{K}},
\end{aligned}
$$

\noindent
which is not always zero if $\theta$ is not
constant.
\end{proof}

\section{Holonomy}
\label{section-holonomy}

\noindent
Let $M$ be a smooth manifold
and $\nabla$ a linear connection.
Any loop $\gamma$ gives a map
$P_\gamma:T_pM \to T_pM$.

Define
$$
\text{Hol}_p(\nabla)
=\{P_\gamma:\text{loops starting at }p\}
\subseteq \text{GL}(n,\mathbb{R}),
$$
which indeed is a subgroup of
$\text{GL}(n,\mathbb{R})$
since parallel transport is invertible
by changing the direction of the
parametrization.
If $M$ is connected, $\text{Hol}_p(\nabla)$
is independent of $p$.

\begin{example}
\label{example-holonomy}
\begin{enumerate}
\item For $\nabla^{\text{LC}}$
we have $\text{Hol}(\nabla^{\text{LC}})
\subset O(n,\mathbb{R})$.

\end{enumerate}
\end{example}

\begin{theorem}[Berger, 1955]
\label{theorem-berger}
If $(M,g)$ a complete Riemannian
simply connected
manifold,
not locally a product
and not a symmetric space,
then $\text{Hol}(\nabla^{\text{LC}})$
can only be one of the following:
\begin{enumerate}
\item $\text{SO}(n)$ oriented Riemannian
manifold.
\item $\text{SU}(n/2)$ Calabi-Yau.
\item $Sp(n/4)$ Hyperkähler.
\item $Sp(n/4)Sp(1)$ quaternionic Kähler.
\item $U(n/2)$ Kähler.
\item $G_2$ exceptional geometry -+.
\item $\text{Spin}(7)$ exceptional
geometry 8.
\end{enumerate}
\end{theorem}

\section{Bismut connection}
\label{section-bistmuth}

\noindent
These notes document Beatrice Brienza's
work based on a few talks by her.

\begin{theorem}[Gauduchon]
\label{theorem-gauduchon-connections}
On $TM$ there exists an affine line
of Hermitian connections known as
{\it canonical} or
{\it Gauduchon connections}
$$
g(\nabla_X^tY,Z)
=g(\nabla_X^{\text{LC}}Y,Z)
+\frac{t-1}{4}
d^c\omega(X,Y,Z)
+\frac{t+1}{4}
d^c(X,JY,JZ).
$$
\end{theorem}

\noindent
Note that:
\begin{enumerate}
\item $\nabla^{\text{LC}}$ is Hermitian
if and only if $d \omega=0$.
\item $d \omega \neq 0 \implies
\forall t \in \mathbb{R}$,
$\nabla^t$ has torsion. This leads to
``Hermitian geometry with torsion''.
\item $t=1$ gives the Chern connection,
i.e. the unique Hermitian connection
such that $\nabla^{0,1}=\bar\partial$
(see Definition
\ref{definition-chern-connection}).
\item $t=-1$ gives the Bismut connection,
defined next, satisfying
$$
g(\nabla_X^BY,Z)
=g(\nabla^{\text{LC}}_XY,Z)
-\frac{1}{2}d^c\omega(X,Y,Z),
$$
where $\omega$ is the {\it Bismut torsion}.
\end{enumerate}

\noindent
Upshot: Chern connection is adapted
for the complex structure,
Bistmut connection is adapted
for skew torsion,
both equal Levi-Civita connection in
Kähler case.

\begin{definition}[Stve, Bis, Fl, 2002]
\label{definition-bismut-connection}
The {\it Bismut connection} is the
unique Hermitian connection with
totally skew symmetric torsion, i.e.
$g(T^B(X,Y),Z) \in \Omega^3(M)$,
where of course
$T^B(X,Y)=\nabla^B_XY-\nabla^B_YX-[X,Y]$
is the usual torsion.
\end{definition}

After lowering an index of the torsion,
we obtain a tensor $T$:
$$
T \in \Lambda^2(T^*M) \otimes T^*M
\simeq \underbrace{T^*M}_{\text{vectorial
type}} \oplus\underbrace{\Lambda^3T^*M}_{
\text{skew type}}
\oplus \underbrace{\tau}_{\text{
terra incognita}}
$$
and we are interested in the case
where $T$ is in $\Lambda^3T^*M$.

\medskip\noindent
It turns out that
$$
\rho^B=0 \iff
\text{Hol}(\nabla^B) \subseteq
\text{SU}(n)
$$
where $\rho^B$ is the
Bismut Ricci curvature,
$$
\rho^B(X,Y)
=\frac{1}{2}\sum_i^{2n}
R^B(X,Y,e_i,J e_i)
\in \Omega^2(M),
$$
in fact $\rho^B \in H^2$ is
the Chern class of $M$.

\begin{definition}
\label{definition-bismut-hermite-einstein}
A {\it Bismut Hermite Einstein} manifold
is when
$$
dH=d d^c\omega=0,\qquad
\rho^B=0.
$$
\end{definition}

\noindent
And $\rho^B=0 \iff \text{Hol}(\nabla^B)
\subseteq \text{SU}(n)$.

\medskip\noindent
In dimension 4, a BHE
can only be a Kähler Ricci flat manifold
or $S^1 \times \text{SU}(2)$. On
dimension 6, ABSL.

\begin{definition}
\label{definition-almost-contact}
An {\it almost contact manifold}
is a triple $(\varphi,\eta,\xi)$
where $\varphi \in \text{End}(TM)$,
$\eta \in \Omega^1(M)$,
$\xi \in \mathfrak{X}(M)$,
$\varphi^2 = -\text{Id} + \eta \otimes \xi$,
$\eta(\xi)=1$.
\end{definition}

\noindent
This means that
$$
TM = \mathcal{F}_{\xi}\oplus Ker\xi.
$$

We then introduce a metric
compatible with $\varphi$.




\section{Instantons}
\label{section-instantons}

\noindent
Let $X$ be a compact 4-dimensional compact
oriented Riemannian manifold, $G$ a compact
Lie group (e.g. $U(n), \text{SO}(n),Sp(n)$).

\begin{definition}
\label{definition-g-instanton}
A {\it $G$-instanton} on $X$
is a $G$-vector bundle $E$
with a connection $\nabla$ such that its
curvature $F_{\nabla}$ is self-dual, i.e.
$*F_{\nabla}=F_{\nabla}$
where $*$ is Hodge operator.
\end{definition}

ADHM studied the case when $X=S^4$,
and $X=\mathbb{C}P^2$ was also studied in a
paper by N. Buchdahl.

\section{Twistor space}
\label{section-twistor-space}

\noindent

Let $X$ be a compact oriented 4-dimensional
Riemannian manifold. Recall that since the
Hodge star operator $*:\Lambda^2\to\Lambda^2$
has eigenvalues $\pm 1$, we have a
decomposition
$$
\Lambda^2=\Lambda^+ \oplus \Lambda^- \cong
\text{SO}(T_X)
$$
Consider the spherization of $\Lambda^+$:
$$
Z:=S\Lambda^+=\{S \in \Lambda^+:\text{unit
vectors}\}
$$
$(m,s) \in Z$, $m \in X$, $s_m \in
\Lambda^+_m$,

$$
T_{(m,s)}Z=
T_mX \oplus_{LC}T_sS \Lambda^+_m
\cong T\mathbb{C}P^1
$$
which has a complex structure
$\mathcal{I}_{\mathbb{C}P^1}$
(metric orientation).

\begin{theorem}
\label{theorem-twistor-space}
$(Z,\mathcal{I}_Z=(s,\mathcal{I}_{\mathbb{C}
P^1}))$
is a complex manifold of complex dimension 3.
\end{theorem}

\noindent
We also have a projection $p:Z \to X $
whose fibers are $\mathbb{C}P^1$,
called {\it real lines}.

\begin{example}
\label{example-twistor-space}
\begin{itemize}
\item $X=S^4$ then $Z=\mathbb{C}P^3$
(projective).
\item $X=\mathbb{C}P^2$ then
$Z=\mathbb{F}_3(\mathbb{C})$ (projective).
\end{itemize}
\end{example}

Let $Z$ be a 3-dimension complex manifold.
Let $L \xymatrix{\ar@{^{(}->}[r]&} Z$
be a rational curve with a normal bundle
$N_L \cong \mathcal{O}(1) \oplus
\mathcal{O}(1)$.

\begin{remark}
\label{remark-small-deformations}
If $Z \supset W$ where $W$ is a complex
submanifold,
the small deformations of $W$ are encoded by
$H^0(N_W)$ and $H^1(N_W)$.
\end{remark}

Let $\mathcal{M}$ be the space of
deformations of $L$.
Since $H^1(N_L)=0$ then
$TM \simeq H^0(N_L) = H^0(\mathcal{O}(1)
\oplus \mathcal{O}(1))$,
so $\dim_\mathbb{C}TM=4$.

Let $\sigma:Z \to Z$ be an antiholomorphic
involution with no fixed points but $L$ ($L$
is fixed). There is an induced
antiholomorphic involution $\tilde{\sigma}$
on $\mathcal{M}$.

\begin{remark}
\label{remark-antiholomorphic-involution}
\begin{itemize}
\item We can recover $X$ as a connected
component
of the fixed point set of $\tilde{\sigma}$ in
$\mathcal{M}$.
\item We can recover the
conformal structure.
\end{itemize}
\end{remark}

\medskip\noindent
Now we discuss Ward-Penrose transform.
Let $E$ be a $G$-instanton
(i.e. a vector bundle with a connection such
that
$*F_\nabla=F_\nabla$) on $X$.
Let $\phi:V \to \overline{V}^*$ be an
Hermitian form.

There is a correspondence (the {\it Ward
correspondence})
which allows to construct a holomorphic
vector bundle on $Z$
such that $\tilde{E}$ restricted to the fiber
of $p:Z \to X$
(the fibers of $p$ are $\mathbb{C}P^1$)
trivially.

Conversely, starting from a hermitian
holomorphic vector bundle $\tilde{E}$
on $Z$ which restricts trivially to the
fibers of $p:Z \to X$.

\medskip\noindent
{\bf Problem (ADHM).} Describe all
$G$-instanton bundles on $S^4$
with $r \geq 2$, $c_2=k$ (fixed second Chern
class).

This makes us want to study the vector
bundles on $\mathbb{C}P^3=Z$.

See Horrseks for the original construction of
monads.
Upshot: ``All bundles can be represented by
the cohomology of some monad''.

\begin{definition}[Monad]
\label{definition-monad}
Let $U,V,W$ be complex vector spaces over
$\mathbb{C}P^1$.
Consider the sequence $h \circ a=0$, where
$a,b$are bundle homomorphisms:
$$
\xymatrix{
0\ar[r]&U(-1):=U \otimes \mathcal{O}(1)
\ar[r]^a
& V \otimes \mathcal{O}
\ar[r]^b
& W(1):=W \otimes \mathcal{O}(1)\ar[r]&0
}
$$
Let $E:=\frac{\Ker b}{\text{Im}a}$.
\end{definition}

\medskip\noindent
Let's study the case of $\mathbb{C}P^3$.
If $Z_i$ are homogeneous coordinates on
$\mathbb{C}P^3$,
$$
a=\sum_i A_iZ_i,\qquad b=\sum_jB_jZ_j
$$
the {\it monad condition} is
$$
B_jA_i+A_jB_j=0
$$
[which gives the conditions for a hyperkähler
metric on the moduli space].

\begin{theorem}[Barth, Barth-Hulek]
\label{theorem-bbh}
Let $E$ be a holomorphic vector bundle of
rank $r \geq 2$
on $\mathbb{C}P^3$ with a $c_2=k$ and
$E|_L=$trivially.
Then $E$ is the cohomology of a monad.
Also, it satisfies some vanishing conditions:
$$
H^0(E)=H^0(E^*)=H^1(E(-2))=H^1(E^*(-2))=0.
$$
\end{theorem}






\section{Stability}
\label{section-stability}

\noindent
This section outlines some concepts
from a paper by Verbitsky from approx. 2002.

A {\it Gauduchon metric} is such that
$\partial\bar\partial \omega^{n-1}=0$. The
degree of a holomorphic vector bundle is
given by $$ \text{deg}_g(\mathcal{F}) =
\int_M c_1(\det \mathcal{F}) \wedge
\omega^{n-1}, $$ which is well defined when
$g$ is Gauduchon.

\begin{definition}
\label{definition-slope-and-semistability}
The {\it slope} of a vector bundle is
$$
\mu_g(\mathcal{F}) = \frac{\text{deg}_g
\mathcal{F}}{\text{rk}\mathcal{F}}.
$$
A vector bundle $\mathcal{F}$ is {\it stable}
if for any coherent subsheaf $\mathcal{E}
\subset \mathbb{F}$
$$
\mu_g(\mathcal{E}) < \mu_g(\mathcal{F}).
$$
It is called {\it semistable} if we put $\le
$
instead.
\end{definition}

\begin{example}
\label{example-line-bundles-are-stable}
Line bundles are stable.
\end{example}

\begin{theorem}[Hitchin-Kobayashi
correspondence]
\label{theorem-hitchin-kobayashi}
Let $(M,g)$ be a compact complex manifold
with $g$ Gauduchon.
Then an hermitian vector bundle $(F,h)$
is Hermitian-Einstein,
i.e. it satisfies
$$
\sqrt{-1}\Lambda\Theta=\gamma\text{id}
$$
if and only if
it is polystable, i.e.
it is the sum of stable bundles
with the same slope.
\end{theorem}

\noindent
The aim of Misha's paper is to
classify stable bundles in positive elliptic
fibrations using Hitchin-Kobayashi
correspondence. Positive elliptic fibrations
are $\pi:M \to X$ with an action of a torus
on the fibers, and such that
$\pi^*\omega_X=d\eta$ where $\omega_X$ is a
Kähler form on $X$. Such a metric is called
preferred if it is invariant on the fibres
and is a Riemannian submersion.

The main theorem is the following:

\begin{theorem}
\label{theorem-stability-misha}
For a positive elliptic fibration $\pi:M \to
X$,
if $(B,h)$ is a stable bundle on $M$,
$$
B \simeq L \otimes \pi^*B_0
$$
where $L$ is a line bundle and
$B_0$ is stable bundle on $X$.
\end{theorem}

\noindent
The upshot is that a locally conformal
hyperkähler manifold
has an associated elliptic fibration…


\section{Deformation theory (introduction)}
\label{section-deformations}

\noindent
Here I discuss deformations in the category
of complex spaces.

Let's just do a summary of \cite[Chapter
9]{Voisin}.

Let $\mathcal{X}$ be a complex manifold, $B$
a complex manifold
and $\phi:\mathcal{X} \to B$ a holomorphic
map.
We call $\phi$ a {\it family} if it is a
proper holomorphic submersion.

A theorem by Ehresmann says that any
``deformation'' of
smooth real manifolds is trivial:

\begin{theorem}[Ehresmann]
\label{theorem-ehresmann}
\begin{reference}
\cite[Theorem 9.3]{voi}
\end{reference}
If $\phi:\mathcal{X} \to B$ is a
proper submersion of differentiable manifolds
(no complex structure), with $B$ contractible
and $0 \in B$
some distinguished point,
then the total space $\mathcal{X}$ is trivial
as a
manifold over $B$, i.e.
there exists a diffeomorphism
 $$
T:\mathcal{X}\xrightarrow{\cong}X_0 \times B
$$
of manifolds over $B$, i.e. commuting with
projection.
\end{theorem}

\begin{proof}
The main tool is the theorem of tubular
neighbourhoods,
which says that there is a diffeomorphism of
a nieghbourhood $W$ of $X_0$ in $\mathcal{X}$
and a neighbourhood of $X_0$ in its normal
bundle.
(We have a normal bundle because $X_0$ is a
submanifold of $\mathcal{X}$.)

In particular this gives a differentiable
retraction
$T_0:W \to X_0$, which allows us to consider
the map
$$
(T_0,\phi):W \to X_0 \times B.
$$
Since $T_0$ is a retraction
(it's the identity when restricted to $X_0$)
and $\phi$ is a submersion,
the map $(T_0,\phi)$ has nonsingular
differential along $X_0$.
Then we can invert it locally along $X_0$,
and since $X_0$ is compact we obtain a whole
neighbourhood $W'\subset W \subset
\mathcal{X}$
where the differential is nonzero and thus
$(T_0,\phi)|_{W'}$ is an embedding
(using the inverse function theorem).

Finally we restrict the image of
$(T_0,\phi)|_{W'}$
to a neighbourhood $U\subset B$ of $0$,
using properness of $\phi$ to ensure
that $\phi^{-1}(U) \subset W'$ --- though I'm
not exactly sure how.
We have constructed a diffeomorphism
$\phi^{-1}(U) \to X_0 \times U$
commuting with the projection.

This proof was only local; the global version
is in \cite{voi} though it's not used in
the rest of the text.
\end{proof}

In the complex situation, if
$\phi:\mathcal{X} \to B$ is a family
of complex manifolds and $0 \in B$, then
``up to replacing $B$ by a neighbourhood of
0'',
there exists a smooth trivialization
$T=(T_0,\phi):\mathcal{X} \to X_0 \times B$
such that the fibres of $T_0$ are complex
submanifolds of $\mathcal{X}$.

Now let $\phi:\mathcal{X} \to B$ be a family
of complex manifolds.
Recall the pullback bundle $\phi^*B$ and the
fact
that the differential $\phi_*$ can be thought
of as a gadget
that maps tangent vectors to $\mathcal{X}$
not to vectors tangent to $B$ but
to vectors in the pullback bundle $\phi^*B$.
Restricting everything to the fibre
$X=\phi^{-1}(0)$ we get
\begin{equation}
\label{equation-short-exact-seque-defor}
\xymatrix{
0\ar[r]&T_X\ar[r]&T_{\mathcal{X}|X}\ar[r]&
\phi^*T_{B|X}\ar[r]&0.
}
\end{equation}
There's what seems to be a really nice result
saying that
$H^1(\mathcal{F})$ parametrizes isomorphisms
classes
of extensions of $\mathcal{F}$ by the trivial
bundle,
but for now let me just define it roughly as
follows.
The category of $\mathcal{O}_X$ is most
surely abelian
with enough injectives. This means that
objects
have associated complexes, and we can take
cohomologies of these complexes.
This says, roughly, that short exact
sequences of vector bundles have associated
cohomology
long exact sequences. Then the
Kodaira-Spencer map
is the map
$$
\rho:T_{B,0}=H^0(X,\phi^* T_{B|X})\to
H^1(X,T_X)
$$
induced by the long exact sequence associated
to
\ref{equation-short-exact-seque-defor}.
Here the equality $T_{B,0}=H^1(X,\phi^*
T_{B|X})$
is just because the pullback of the tangent
bundle of $B$
to any fiber is trivial.

\medskip\noindent
Now I will explain first order deformations.
Consider a family $\mathcal{X} \to B$.
We introduce a complex analytic subscheme
$B_\varepsilon$ of $B$
giving an ideal sheaf $\mathcal{M}^2_0$
where $\mathcal{M}_0$ is the maximal ideal
of $\mathcal{O}_{B,0}$ consisting of the
holomorphic
functions vanishing at 0.
This means that a function in
$\mathcal{M}^2_0$
is (a finite sum of terms) of the form $f=gh$
for $g,h \in \mathcal{M}_0$,
which means that $df=d(gh)=hdg+gdh=0$ at $0$
since both $g$ and $h$ vanish at $0$.
That is, this sheaf contains functions that
vanish
at $0$ and whose first derivative vanishes at
$0$.

This means that $B_\varepsilon$ is a point:
it's the vanishing locus of the set of
functions
that vanish at a single point.
However, it is endowed with an unusual
structure sheaf.
Let $v$ be a tangent vector to $B$ at $0$ and
$f$ a function in the strange sheaf
$\mathcal{M}_0/\mathcal{M}_0^2$.
Then by the computation above, since $v$ is a
differential
operator, $v$ annihilates $f$.
This means that this sheaf detects not only
the point $0$
but also any tangent vector at $0$.
Let's denote
$B_\varepsilon=\Spec
\mathbb{C}[\varepsilon]/(\varepsilon^2)$
just because its structure sheaf is
$\mathbb{C}[\varepsilon]/(\varepsilon^2)$.

In fact, this is so throughout deformation
theory:
$\Spec k[\varepsilon]/(\varepsilon^2)$ is
topologically a point!
But it's a {\it fat point}: the support of
its
structure sheaf is a point, but as a scheme,
it's not a point.
And the same happens with the family
$\mathcal{X}$ over
$\Spec k[\varepsilon]/(\varepsilon^2)$,
topologically it's just the fibre $X$!
But its structure sheaf is different:
``it contains the directions in $\mathcal{X}$
which
are normal to $X$''.

\medskip\noindent
Perhaps the most important part of this
construction
comes when considering a fat point
$B_\varepsilon \subset B$.
This gives a first-order deformation
$\phi^{-1}(B_\varepsilon)) \subset
\mathcal{X}$
$$
\xymatrix{
X\ar[r]\ar[d]
&\phi^{-1}(B_\varepsilon) \subset
\mathcal{X}\ar[d]\\
\ast\ar[r]
&B_\varepsilon \subset B
}
$$
Using the constructions above, one shows that
to any such deformation we can associate a
cocycle
in $H^1(X,T_X)$.
Thus, first order deformations are ``exactly
parametrised''
by the group $H^1(X,T_X)$.
This is the main idea behind Kodaira-Spencer
map,
which is well-explained in
\href{https://en.wikipedia.org/wiki/Kodaira
% E2%80%93Spencer_map}{Wikipedia}.

I will present a similar construction
to Kodaira-Spencer map
that works for embedded deformations
instead of just deformations.
(That's with the intention of understanding
the
forthcoming exercise Exercise
\ref{exercise-deformations-curve-in-k3}.)

Since we are in complex space category
I follow \cite[p. 224]{voi} for the main
setup,
but this is also a mixture
of \cite[p. 147, 150]{Sernesi}.

The result is in fact quite reasonable
geometrically:
the dimension of directions in which we can
deform a
submanifold should be non other than
the dimension of the normal bundle.

\section{Deformation theory}
\label{section-deformation-theory}

\begin{definition}
\label{definition-deformation}
Let $X$ be a complex manifold.
A {\it deformation} of $X$ is a holomorphic
submersion
$\pi:\mathcal{X} \to B$ of complex manifolds
such that $\pi^{-1}(0)\cong X$ for some $0
\in B$.

If $\mathcal{X}$ and $\mathcal{Y}$ are two
deformations
of $X$ over the same base,
a {\it morphism of deformations} is a map
$f:\mathcal{X} \to \mathcal{Y}$ such that the
diagram
$$
\xymatrix{
& X\ar[dl]\ar[dr]\\
\mathcal{X}\ar[dr]\ar[rr]^f&&\mathcal{Y}
\ar[dl]\\
&B
}
$$
commutes. The {\it space of deformations} of
$X$
is the set of deformations modulo
the equivalence relation of two deformations
being isomorphic. A priori we don't know
if this set has any structure.

Let $Y$ be another complex manifold.
A {\it deformation of $X$ in $Y$}
is a deformation of $X$ such that
$\mathcal{X}\subset Y\times B$.

The {\it fat point} is a ringed space
$(B_\varepsilon,\mathbb{C}[\varepsilon]/
(\varepsilon^2))$
whose underlying topological space
$B_\varepsilon$ is a point
and structure sheaf is the ring of dual
numbers
$\mathbb{C}[\varepsilon]/(\varepsilon^2)$.

A {\it first order deformation} of $X$ is a
deformation such that
the base $B$ is a fat point.
\end{definition}

The following proposition is a
mixture of ideas from \cite[p. 224]{voi}
and \cite[Examples
3.2.4(ii)]{Sernesi-deformations}.

\begin{proposition}
\label{proposition-kodai-spenc-embed-defor}
Let $X$ be a smooth hypersurface in a complex
manifold $Y$.
For every deformation of $X$ in $Y$
with base of dimension 1 we can construct a
section of the bundle
$N_{X/Y}$. Different deformations give
different sections up to rescaling.
\end{proposition}

\begin{proof}
Let
$$
\xymatrix{
X\ar[r]\ar[d]
&\mathcal{X}\subset Y\times B\ar[d]^\pi\\
0\ar[r]
&B
}
$$
be a deformation of $X$ in $Y$
and suppose that $\dim B=1$.
This means that $\mathcal{X}$ is a complex
submanifold of
the product $Y \times B$
and $\pi:\mathcal{X} \to B_\varepsilon$ is
a holomorphic submersion.

Since $X$ is a hypersurface,
we can find an open cover $V_i$
of $M$ such that $V_i \cap X:= U_i$
is given locally by the vanishing of
a function $f_i \in \mathcal{O}_Y(V_i)$.
In terms of coordinates we find that
that $V_i \cap X$ is given by
$(z_1,…,z_{n-1},0)$.

Since $\pi$ is a submersion,
by the holomorphic inverse function theorem
we obtain for every $V_i \subset \mathcal{X}$
(such that $V_i \cap X \neq  \emptyset$)
a biholomorphism from some $B_i \subset B$
containing $0$ to $V_i$.
This gives coordinates $(0,…,b)$ to $V_i$.
We conclude that $V_i \cong U_i\times B_i$
along $X$.

Now we ``restrict'' to a first order
deformation.
Consider the fat point subscheme
$B_\varepsilon$ of $B$
(I think this is the reason why $B$ must have
dimension 1,
in Voisin notation, that a first order
neighbourhood
of $0$ in $B$ is indeed a fat point)
and let
$X_\varepsilon=\pi^{-1}(B_\varepsilon)$,
which is a ringed space whose topological
space
is the same as $X$, but its
structure sheaf is induced from
$B_\varepsilon$.

(I think that to prove this part I have
to look carefully at the induced sheaf
$\pi^*(\mathbb{C}[\varepsilon]/
(\varepsilon^2)$.)
Since we are on a local trivialization,
the structure sheaf of the product $U_i
\times B_i$
is the tensor product of the structure
sheaves
(reference?),
so that the structure sheaf of
$X_\varepsilon$ at
$U_i\times B_\varepsilon$ is given by
$$
\mathcal{O}_{X_\varepsilon}(U_i)=
 \mathcal{O}_Y(U_i) \otimes
 \mathbb{C}[\varepsilon]/(\varepsilon^2).
$$
Functions here look like finite sums of
functions of the form $f+ \varepsilon g$ for
$f,g \in \Gamma(U_i,\mathcal{O}_{U_i})$.

The next key step is to argue that
$X_\varepsilon$ is locally cut by
functions of the kind $f_i + \varepsilon g_i$
for some $g_i \in \mathcal{O}_M(U_i)$.
Here the $f_i$ are the same as those which
define $X$.
I think a proof is that from the ideal sheaf
exact sequence of $X$ in $Y$
$$
\xymatrix{
0\ar[r]
&\mathcal{I}_X\ar[r]
&\mathcal{O}_Y\ar[r]
&\mathcal{O}_X\ar[r]
&0
}
$$
we obtain after tensoring by
$\mathbb{C}[\varepsilon]/(\varepsilon^2)$
$$
\xymatrix{
0\ar[r]
&\mathcal{I}_X \otimes
\mathbb{C}[\varepsilon]/(\varepsilon^2)\ar[r]
&\mathcal{O}_Y \otimes
\mathbb{C}[\varepsilon]/(\varepsilon^2)\ar[r]
&\mathcal{O}_X \otimes
\mathbb{C}[\varepsilon]/(\varepsilon^2)\ar[r]
&0.
}
$$
(But is the functor $-\otimes
\mathbb{C}[\varepsilon]/(\varepsilon^2)$
exact? Why? So far we didn't impose a
flatness condition on $\pi$.)
The sheaf in the right
coincides with the structure sheaf of
$X_\varepsilon$ when
evaluated on a trivial open set of the form
$U_i \otimes B_\varepsilon$.
This shows that the ideal sheaf of
$X_\varepsilon$
as a ringed space is
$\mathcal{I}_X
\otimes\mathbb{C}[\varepsilon]/
(\varepsilon^3)$.
Since $\mathcal{I}_X$ is locally principal
generated by $f_i$,
we conclude that indeed $X_\varepsilon$
must be cut locally by functions of the form
$f_i + \varepsilon g_i$
for some $g_i \in \mathcal{O}_Y(U_i)$

Recall that the line bundle associated to $X$
is $\mathcal{O}_Y(X)$, given by the cocycle
$f_i/f_j:=f_{ij}$.
A version of adjunction formula
(more precisely, \cite[Adjunction formula I,
p. 146]{GH})
establishes that this line bundle,
when restricted to $C$, is in fact
the normal bundle of $X$ in $Y$, defined as
que quotient
bundle $T_Y'/T_X':=N_{X/Y}$.
In symbols, $\mathcal{O}_Y(X)|_{X}\cong
N_{X/Y}$.

Similarly, the line bundle associated to
$\mathcal{X}$
is given by the cocycle
$$
F_{ij}=f_{ij}+\varepsilon g_{ij}
$$
for $g_{ij}=\frac{g_i-f_{ij}g_j}{f_j}$.
Indeed, note that
$\frac{1}{f_j+\varepsilon
g_j}=\frac{1}{f_j}-\varepsilon
\frac{g_j}{f^2_j}$
and compute $F_{ij}=(f_i+\varepsilon
g_j)/(f_j+\varepsilon g_j)$.

Then
$$
f_i+\varepsilon g_i=(f_{ij}+\varepsilon
g_{ij})(f_j+\varepsilon g_j)
$$
gives
$$
g_i=f_{ij}g_j+g_{ij}f_j
$$
and since $f_j$ vanishes along $X$,
we are left with $g_i=f_{ij}g_j$
which is how a section of the sheaf $N_{X/Y}$
is obtained.
\end{proof}

\begin{exercise}
\label{exercise-deformations-curve-in-k3}
Let $C$ be a smooth genus $g$ curve which
can be embedded in a K3 surface $M$, and
$\mathcal{X}$
``the family of all deformations'' of $C$ in
$M$.
Prove that $\dim \mathcal{X} \leq g$.
\end{exercise}

\begin{proof}
By Adjunction formula II
(Theorem \ref{theorem-adjunction-ii})
and $M$ being a K3 surface
we get that $h^0(C,N_{C,M})=g$.

If the deformation space $X$ exists and is a
complex manifold,
the map constructed in
Proposition
\ref{proposition-kodai-spenc-embed-defor}
would be a holomorphic (why holomorphic!?)
injective map $X \to H^0(C,N_{C/Y})$.
This obliges $\dim X \leq  h^0(C,N_{C/Y})=g$.
Indeed, if not, the differential must have a
kernel,
and integrating along a direction in the
kernel
we obtain the germ of a complex curve along
which
the map is constant, which is impossible
since it is injective.
\end{proof}







% --- End complex-geometry.tex ---

% --- Begin physics.tex ---
\chapter{Physics}
\label{chapter-physics}





\section{History}
\label{section-history}

The first force discovered was gravity, by
Newton. Then there was
electromagnetism by Coulomb. The formulas are
very similar.

Faraday introduced notion of field as a
``convenient" way to study
electromagnetism. Then Maxwell's equations.
These were about 1870's.

Special relativity in 1905. General
relativity in 1915: Einstein equation,
Hilbert
ation.

QM. Schroedinger equation 1920. Heisenberg
equation 1929. Bohr Rule 1925.

QFT: particles created/annihilated. Dirac,
1929, spinors. Feynman, Tomonaga,
Dyson. Wilson in 1970s with
``Renormalization".

In 1953 Yang-Mills introduced non-abelian
gauge theory.

Higgs interaction field, 1964.

\section{Electromagnetism}
\label{section-electromagnetism}

Probably the one Gauge theory you are all
familiar with is electromagnetism.
It's was the first \textit{field theory} in
which there were dynamical objects
which are functions $\phi(\vec{x},t)$ which
could be measured, etc. It is based
in what we call an abelian Gauge theory (for
historical reasons) with group
$\mathsf{U}(1)$. Its \textit{dynamical field}
is a connection $A$, which is
(locally?) a 1-form on Minkowsky space
$\mathbb{R}^{3,1}$, over a
$\mathsf{U}(1)$-bundle. And the
\textit{physical observables} are the
\textit{field strengths} (i.e. curvature)
$F=dA=dt\wedge \vec{x}
\vec{E}\ldots$, whose components are the
familiar electric and
magnetics fields.

Maxwell unified electricity and magnetism
into these theory, whose equations are
simple second order PDEs:
\[dF=0,\qquad d\wedge F=0\]
which follow from an action functional
\[\delta=\frac{1}{4e^2}\int_{M}F\wedge * F.\]
These equations were the first \textit{field
theory}. Having dynamics described
by a function $\phi(\vec{x},t)$ of space and
time. The action is invariant under
Lorentz transformations i.e. in the group
$\mathsf{SO}(3,1)$. The solutions of
Maxwell equations are light $\sim$ waves of
ripples in this field strength.

\section{Quantum Electrodynamics (QED)}
\label{section-QED}

In some diagram from the book Lightness of
the being, electrons are attracted to
each other by the presence of photons
(probably this is why we say that Maxwell
equations ``are light''). In my
interpretation: electric charges are
attracted
by light. Anyway, then we have that more than
one photon may appear in the
trajectories of two electrons. And more
interestingly: an excited electron can
emit a photon: ``let there be light''. To
quote the book ``that's what we call
electromagnetic radiation, one form of which
is light''.

\section{Yang-Mills theory}
\label{section-Yang-Mills-theory}

A generalization of gauge theory to non

YM theory has a mass gap, which produces a
scale that ED in the vaccum does not
have. We will be discussing non-abelian
electromagnetism just like
electromagnetism has no scale as a
dimensionless coupling. And we will produce a
scale, a mass gap, from quantum effects.

What does it mean to define or prove the
existence of a QFT?

We need a Hilbert space $\mathcal{H}$, states
in a Hilbert space $|\psi\rangle$,
positive norm states (we identify probability
with the norm of states, and
probability should be positive). We want {\it
relativistic invariant},
symmetries, Poincaré invariants: translation
in space and time $x \mapsto x+a$ +
Lorentz transformations $\mathsf{SO}(3,1)$.
In quantum mechanics, such
symmetries of the theories are represented by
unitary operators under which
states, the fields, transform covariantly.
And we want a {\it vacuum
state} $|\Omega\rangle$, a ground state of
the system, that is invariant under
translations and Lorentz transformations.
We're looking for all this.

We also want that the generators of these
symmetry transofrmations, the momenta
$P^n$, have definite spectral properties. In
particular, the energy $P^0$ should
have spectrum bounded from below $P^0\geq 0$.
We call
$\mathcal{P}^2=E^2-P^2=m^2$ the {\it mass}.

In the case of ENM, the spectrum of the
energy  $P^0$ is continuous.

The spectrum of the mass operator has a base
state $|\Omega\rangle$, and then
there is a discrete state $m_G$, $G$ standing
for \textit{gap}. An the pairs of
these particles of twice the mass, $2m_G$,
and some extra kinetic energy, are
continuum of these multiparticle sates. So,
in a QFT where we have massive
particles and no massless particles like
lightrays, \textit{there is a gap
between the vacuum, the ground state, and the
first state that can be created
with a minimal energy equal to the mass of
that particle}, and then
multiparticle states, and perhaps some
isolated states.

\section{Tensions in QM and GR}
\label{section-tensions}

Gravity and QM don't agree at a very basic
observation: vacuum energy.

Quantization of a very simple variation
problem, for a scalar field, says that
energy of the vacum is infinite. But there's
an experiment that says that the
cosmological constant (which is energy of the
vacuum) is
 $\rho_{\operatorname{cc}}\approx 5 \times
 10^{-10}\frac{J}{m^3}$.

One there was a variation problem for the
functional
$$
S=\int dt \int d^3x \frac{1}{2}\left\{
\frac{1}{2}\dot\varphi^2-(\nabla
\varphi)^2-m^2\varphi^2\right\}
$$
whose solution, upon quantization, gives an
infinite integral. But it should be
the cosmological constant! A tiny number.

We want to experimentally find a system that
is governed both by quantum
mechanics and relativity. An experiment
called ``Micro g" by Aspa-Meyer shows
that gravitational laws in a 1-millimetre,
i.e. $10^{-6}$ kg particles. The size
of the proton is $10^{-27}$ kg, where we see
quantum mechanics. Experimentalists
at PUC-Rio have work with sizes of the orders
$10^{-18}$ kg.

\section{General Relativity}
\label{section-general-relativity}
\begin{slogan}
The mass determines the metric.
\end{slogan}
Consider the variational problem:

\section{Standard Model}
\label{section-standard-model}

\begin{slogan}
Standard model = symmetries + representations
\end{slogan}

Recall: in electromagnetism we have
$$
\psi\rightsquigarrow
e^{i\alpha(x)}\psi,\qquad
A_\mu\rightsquigarrow
A_\mu+\partial_\mu\alpha(x)
$$
so that the gauge group is $\text{U}(1)$.

In the standard model we have gauge group
$\text{SU}(3)\times\text{SU}(2)\times\text{U}
(1)$.
So for a gauge transformation
we do
$$
\psi\rightsquigarrow U(x)\psi,\qquad
A_\mu\rightsquigarrow
U(x)(A_\mu+i\partial_\mu)U^\dag(x)
$$
We have the following table as a result of
several theoretic and experimental
resereach:
\begin{table}
\centering
\caption{Standard Model (incomplete)}
\label{tab:standard-model}
\begin{tabular}{c c c c c}
Field &	$\text{SU}(3)$ & $\text{SU}(2)$ &
$\text{U}(1)$ &
$\text{SL}(2,\mathbb{C})$ \\
gluons, $G_\mu$ & 8 & 1 & 0 (charge) & \\
$W_\mu$ & & & &\\
$B_\mu$ &&&&\\
\hline
$q^i_L$, $i=1,2,3$ &&&&\\
$\ell^i_L$ &&&&\\
$u^i_R$ &&&&\\
$d^i_R$ &&&&\\
$e^i_R$ &&&&\\
\hline
$H$ Higgs &&&&
\end{tabular}
\end{table}

Confinement: as two quarks spread apart, the
potential energy between them grows
linearly, i.e. $V(r)=\sigma r$ where $r$ is
the distance. (Due to strong
coupling.)

\section{Neutrinos}
\label{section-neutrinos}


Protons and neutrons have very similar mass
of $m_p\approx940$ MeV, while
electrons have mass of $m_e \approx 0.5$ MeV.
MeV is $10^{6}$ electronvolts,
where one eV is approximately $1.6\times
10^{-19}$ J. This is standard in high
energy physics, they use electronvolts
instead of Joules. Recall that
2J=1N$\times$1m.

So most of the things we see are protons
since they are so much larger than
electrons. But protons (and neutrons) are not
elementary particles.

Here's the standard model:
\begin{tabular}{c c c c}
Quarks & $\begin{bmatrix} u\\d
\end{bmatrix}_L $ & $\begin{bmatrix} c\\s
\end{bmatrix}_L$ & $\begin{bmatrix} t\\b
\end{bmatrix}_L$\\
Leptons & $\begin{bmatrix} \nu_e\\e^-
\end{bmatrix}_L $ & $\begin{bmatrix}
\nu_\mu\\ \mu^-
\end{bmatrix}_L$ & $\begin{bmatrix}
\nu_\tau\\ \tau^- \end{bmatrix}_L$\\
Generation & 1st & 2nd & 3rd\\
Bosons & $g$, $\gamma$ , $\omega^\pm$, $z$ \\
Higgs Bosson & $H$
\end{tabular}
It is very particular that nature repeats
itself three times. The $L$ in those
matrix actually means left-handed, and
accounts for chirality. Only left-handed
fermions have weak interaction. Right-handed
have electromagnetic interaction,
gravitational interaction, but not weak
interaction.

\medskip

And then there's neutrinos. They have
negative helicity (chirality). Being
left-handed, mathematically, means to have
helicity $-1$. I think this means
that the spin is left-handed. But chirality
and helicity are not the same:
helicity is observer-dependent, and chirality
is not. Almost all neutrinos we
can see (\% 99.99999…) have negative
helicity, but not all of them.

Consider the following:
$$
n+\nu_e\leftrightarrow p+e^-
$$
But it's not completely correct: we'd better
put $d$ instead of $n$, and $u$
instead of $p$: the $d$ and $u$ quarks,
instead of the neutrons and protons.

Now consider the following reaction: a
neutron decays into a proton, an electron
and an antineutrino:
$$
n\to p+e^+\overline{\nu}_e
$$
Protons is very stable, that's why we are
here. But neutron decays in only 15
minutes.

By experimental data, we can conclude that
neutrinos' mass is consistent with
zero. But if they have mass, it should be
much smaller than the electron's $m_v
\leq 0.5$ eV. And the electron is already the
lightest fermion!

If the mass of the neutrino was zero, i.e.
$m_\nu=0$, then $v_\nu=c$ in vacuum,
which would imply that
$$
\xymatrix{
\nu_e\ar[r]^{L}&\nu_e\ar[r]\nu_e\\
0:00&0:00&0:00
}
$$
meaning: time doesn't pass! And this means
the state of the particle cannot
change.

\section{Stochastic thermodynamics}
\label{section-stochastic-thermodynamics}

Just to recall that there is an analogue of
length functional given by some sort
of information function $\iota$. It looks
something like
$$
\ell(\tau)=\int_0^\tau \iota_x
$$
As I understand, the notion of information is
somehow inversely proportional to
entropy: zero entropy means you have all the
information.


% --- End physics.tex ---

\bibliographystyle{alpha}
\bibliography{my}

\end{document}
